\documentclass[11pt]{article}
\usepackage[T1]{fontenc}
\usepackage{lmodern,amsmath,amssymb,amsthm,microtype}
\usepackage[margin=1in]{geometry}
\usepackage[hidelinks]{hyperref}
\newtheorem{theorem}{Theorem}[section]
\newtheorem{proposition}[theorem]{Proposition}
\newtheorem{lemma}[theorem]{Lemma}
\newtheorem{corollary}[theorem]{Corollary}
\theoremstyle{definition}
\newtheorem{assumption}[theorem]{Assumption}
\newcommand{\one}{\mathbf 1}
\newcommand{\supp}{\operatorname{supp}}
\newcommand{\dist}{\operatorname{dist}}
\newcommand{\E}{\mathbb E}
\newcommand{\Var}{\operatorname{Var}}
\allowdisplaybreaks
\setlength\emergencystretch{2em}
\title{Pseudocritical point of the spread-out Ising model with free boundaries at the upper critical dimension in \(1\le d\le4\)}
\author{Jiwoon Park}
\date{\today}

\begin{document}
\maketitle

\begin{abstract}
We study the susceptibility of the spread-out Ising model on a cube
of volume \(V=n^d\) with free boundary conditions. The interaction
\((-\Delta_R)^{d/4}\) places the model at its upper critical dimension
for each \(d\in\{1,2,3,4\}\). For sufficiently large fixed spread \(R\),
we prove that there is an inverse temperature at which the effective
potential of the principal mode has zero curvature. At this
pseudocritical point the susceptibility has order \(\sqrt{V\log V}\).
At the bulk critical point, Camia, Jiang and Newman obtain order \(n^2\)
for nearest-neighbour Ising boxes with free boundaries in \(d\ge4\)
under a bulk two-point hypothesis \cite{CamiaJiangNewman2021}.
The logarithmic factor at our pseudocritical point also occurs for the critical weakly coupled
\(\lvert\varphi\rvert^4\) model on the four-dimensional torus
\cite{TorusFSS} and in the free-boundary pseudocritical window of the
four-dimensional hierarchical model \cite{HierarchicalFBC}.
\end{abstract}

\subsubsection*{Version note: you can find earlier versions on DOI:10.5281/zenodo.24167641}

v1. November 20, 2026

\section{Model and main result}\label{sec:model}

Fix \(d\in\{1,2,3,4\}\), and write \(\alpha=d/4\), \(s=2\alpha=d/2\).
On \(\mathbb Z^d\) let
\begin{align}
 D_R(x)=\frac{\one_{\{0<|x|_\infty\le R\}}}{(2R+1)^d-1},
 \qquad -\Delta_R=L_R=I-D_R,\qquad
 \Lambda=\{1,\ldots,n\}^d,\qquad V=n^d .
\end{align}
Here \(D_R\) also denotes the convolution operator. 
We also let $P_{\Lambda}$ is the simple restriction $\ell^2 (\mathbb{Z}) \rightarrow \ell^2 (\Lambda)$
and $P_\Lambda^*$ be the extension of $\ell^2 (\Lambda)$ to $\ell^2 (\mathbb{Z})$ by zero.
On the full
lattice, \(L_R^\alpha\) is the Fourier multiplier
\((1-\widehat D_R(k))^\alpha\), with value zero at \(k=0\).
The physical interaction is
\(P_\Lambda L_R^\alpha P_\Lambda^*\).
Thus the power is taken before restriction, and free boundaries
mean deleting interaction bonds to the exterior. The original diagonal
is retained and contributes a spin-independent factor.
We consider the spread-out Ising model with free boundary condition at inverse temperature \(\tau\) 
defined by the measure
\begin{align}
      \mathbb{E}_\tau^F [F(\sigma)]
      = \frac{1}{Z_\Lambda^{F}} \sum_{ \sigma\in\{-1,1\}^\Lambda } F(\sigma)  \exp\bigl[-\tfrac{\tau}{2}
      \langle\sigma, P_\Lambda L_R^\alpha P_\Lambda^*\sigma\rangle\bigr] 
      \label{eq:measure}
\end{align}
where $Z_\Lambda^F$ is the normalisation and the susceptibility
\begin{align}
 \chi_{\Lambda}^F(\tau)=V^{-1}\E_\tau^F
                         (\one^{\mathsf T}\sigma)^2 .
\end{align}
For \(d<4\) the interaction is long-range. At \(d=4\) it is the local
operator \(L_R\), with \(R\) large as required by the Ising initialization.
Quartic coupling is marginal in all dimensions because $d = 4\alpha$.

On a cube with free boundaries the principal mode of the Gaussian
reference gives the leading contribution to the susceptibility.
We fix the reference and its decomposition into scales, then vary the
physical temperature until the terminal effective potential of this mode
has zero curvature at the origin.
This is soon explained in Section~\ref{sec:TheMatchingCondition}.

\begin{theorem}[Pseudocritical susceptibility]
\label{thm:susceptibility}
Fix \(d\in\{1,2,3,4\}\). For all sufficiently large fixed \(R\)
and all \(n\ge n_0(R)\), there is an inverse temperature
\(\tau_{\Lambda}^{\mathrm{pc}}\) at which the terminal effective
potential has zero curvature at the origin, and
\begin{align}
 c_{d,R}\sqrt{V\log V}\ \le\
 \chi_{\Lambda}^F(\tau_{\Lambda}^{\mathrm{pc}})
 \ \le\ C_{d,R}\sqrt{V\log V}.
\end{align}
Here \(0<c_{d,R}\le C_{d,R}<\infty\) are independent of \(n\).
\end{theorem}


In terms of side length the four orders are
\(\sqrt{n\log n}\), \(n\sqrt{\log n}\),
\(n^{3/2}\sqrt{\log n}\), and \(n^2\sqrt{\log n}\), respectively.
In $d=4$, this is comparable to that with periodic boundary condition at the bulk critical point proved in \cite{TorusFSS},
justifying the `criticality' of $\tau_{\Lambda}^{\mathrm{pc}}$.
Also, note that the susceptibility at the bulk critical point lacks the logarithmic correction as show in \cite{CamiaJiangNewman2021}, 
showing a trait of a near-critical Gaussian scaling \cite[Theorem~1.9(2)]{TorusFSS}.

\subsubsection{The matching condition for $\tau^{\mathrm{pc}}_{\Lambda}$} \label{sec:TheMatchingCondition}

Let us first introduce some essential terminologies.
\begin{itemize}
\item Let \(\ell\) be the dual principal vector of the full reference. 
This is defined in Section~\ref{sec:gaussian}, and as $n \rightarrow \infty$, 
equivalent to $v_g$, the eigenvector of $P_{\Lambda} L_\Lambda^\alpha P_\Lambda^*$ 
with the smallest eigenvalue.
Thus, everything below can be interpreted with $v_g$ in place of $\ell$.
\item $C_T$ is the terminal covariance, defined in Section~\ref{sec:FinitePropagationTerminalConditioning}.
It contains the Gaussian fluctuation of the field in the direction of $\ell$.
The remaining covariance is denoted $C_{<}$, containing the bulk Gaussian fluctuation of the field, 
so that $C_{<} + C_T$ represents the covariance of the Ising model, when we view the Ising model as a Gaussian field
restricted to values $\pm 1$. This view is particularly useful in RG implementations.
\end{itemize}
\noindent
Let \(F_T\) denote the remaining interaction factor after integrating
the fluctuation $C_{<}$
and \(\phi\) for the terminal Gaussian field with covariance \(C_T\).
With \(a=\ell^{\mathsf T}\phi\), we can decompose $\phi$ as
\begin{align}
 \sigma_T^2=\ell^{\mathsf T}C_T\ell,\qquad
 u_T=C_T\ell/\sigma_T^2,\qquad
 \phi=a u_T+Y
\end{align}
with \(Y\) independent of \(a\). The density of \(a\) is proportional to
\begin{align}
  e^{-a^2/(2\sigma_T^2)}\,\E^Y F_T(a u_T+Y)
\end{align}
where the expectation is over \(Y\). Noting that this is now defined a distribution on $\mathbb{R}$,
the second derivative at zero of the negative logarithm of this density
defines the curvature in the theorem. The condition is invariant under a
nonzero linear rescaling of \(a\). 
In dimension 4, the actual implementation differs, only because
we first need a quadratic counterterm to tune the reference field.

We also prove in the companion paper \cite{obs-extension} that 
$\Delta_n=\tau_{\Lambda_n}^{\rm pc}-\tau_c$ is a non-trivial shift from the 
bulk critical point $\tau_c=\tau_c(d,R)=\inf\{\tau>0:\langle\sigma_0\rangle_\tau^+>0\}$,
where $\langle\cdot\rangle_\tau^+$ is the infinite-volume plus state for
the Hamiltonian used in \eqref{eq:measure}.
Write $(u)_+=\max\{u,0\}$.

\begin{theorem}[Shift of the pseudocritical point] \cite{obs-extension} 
\label{thm:pc-shift}
For $d\in\{1,2,3,4\}$ and sufficiently large fixed $R$,
$\tau_{\Lambda_n}^{\rm pc}\to\tau_c$ as $n\to\infty$.
For $n\ge n_1(d,R)$ one has $(\Delta_n)_+\le C_dR^{-d}$.
For $d=1,2,3$ one also has
\begin{equation}\label{eq:pc-shift-upper}
 (\Delta_n)_+\le C_{d,R}
 \begin{cases}
   n^{-1/10}(\log n)^{1/10}, & d=1,3,\\
   n^{-1/5}(\log n)^{1/2}, & d=2 .
 \end{cases}
\end{equation}
There is $c_{d,R}>0$ such that, for all sufficiently large $n$,
\begin{equation}\label{eq:pc-shift}
 \Delta_n\ge c_{d,R}\,n^{-s}=c_{d,R}\,n^{-d/2}.
\end{equation}
\end{theorem}

% Its dependence on the reference
% matters in dimension four, where the proof first constructs a matched
% point for a tunable reference.


% Let \(F_T\) denote the remaining interaction factor after integrating
% the fluctuation $C_{<}$. 
% After a change of variable
% \begin{align}
%       \phi \mapsto \phi=a u_T+Y , 
%       \qquad \text{where} \qquad 
%       \sigma_T^2=\ell^{\mathsf T}C_T\ell
% \end{align}
% for the integral with covariance $C_{T}$

% with \(Y\) independent of \(a\). After the earlier fluctuation scales have
% been integrated, let \(F_T\) denote the remaining interaction factor.
% The density of \(a\) is proportional to
% \[
%   e^{-a^2/(2\sigma_T^2)}\,\E F_T(a u_T+Y)
% \]
% where the expectation is over \(Y\). The second derivative at zero of the negative logarithm of this density
% defines the curvature in the theorem. The condition is invariant under a
% nonzero linear rescaling of \(a\). Its dependence on the reference
% matters in dimension four, where the proof first constructs a matched
% point for a tunable reference.


\subsection{Historical context and related results}

\paragraph{Long-range Ising criticality.}
Dyson proved the existence of a phase transition for the ferromagnetic
chain with interaction proportional to $|x-y|^{-1-s}$,
$0<s<1$ \cite{Dyson1969}. For an interaction of order
$|x-y|^{-d-s}$, Fisher, Ma and Nickel identified $d=2s$ as the upper
critical dimension in their renormalisation group analysis and derived
bulk logarithmic corrections at that dimension \cite{FisherMaNickel1972}.
Aizenman and Fern\'andez proved mean-field values of bulk critical
exponents for a class of reflection-positive long-range Ising interactions
when $d<4$ and $0<s\le d/2$ \cite{AizenmanFernandez1988}.
For the periodic system at $d=2s$, Luijten and Bl\"ote derived the
finite-size susceptibility prefactor $L^{d/2}(\log L)^{1/2}$ and tested
it numerically, including for the chain with $s=1/2$
\cite[Eq.~(27) and Table~V]{LuijtenBlote1997}.

\paragraph{Shift and rounding.}
Fisher and Barber distinguished the shift of the effective finite-volume
critical point from the rounding of the transition
\cite{FisherBarber1972}. Br\'ezin and Zinn-Justin computed an effective
Hamiltonian for the homogeneous mode by integrating out the other modes
\cite{BrezinZinnJustin1985}. Above the upper critical dimension the
quartic coupling is dangerously irrelevant for the homogeneous mode:
although it decreases under rescaling, it determines the scale of the
spatial average at the effective critical point. Rudnick, Gaspari and
Privman analysed the free-boundary analogue using standing-wave modes
and found a boundary-induced shift of the effective critical temperature
\cite{RudnickGaspariPrivman1985}. Wittmann and Young later found numerical
evidence in the five-dimensional Ising model for susceptibility of
order $L^2$ at the bulk critical temperature and $L^{d/2}$ at a
free-boundary pseudocritical temperature \cite{WittmannYoung2014}.
In this historical discussion, $L$ denotes the side length of a
$d$-dimensional box and $V=L^d$ its volume.

\paragraph{Bulk scaling and finite-volume observables.}
Bulk Gaussian scaling and the finite-volume law of the spatial average
use different observables and normalizations. Aizenman and Fr\"ohlich proved
Gaussianity of continuum limits for high-dimensional spin or scalar
field models \cite{Aizenman1982,Frohlich1982}. Aizenman and
Duminil-Copin established marginal triviality at $d=4$ for the
nearest-neighbour Ising model and related scalar models, and Panis
extended Gaussianity to critical reflection-positive Ising and scalar
models of effective dimension at least four
\cite{AizenmanDuminilCopin2021,Panis2026}. On periodic lattices, Park's
results for weakly coupled $|\varphi|^4$ models give a Gaussian field
limit when tested against mean-zero functions and a quartic law for the
spatial average at its own normalization \cite[Theorem~1.9]{TorusFSS}.
These observables probe different parts of the fluctuation spectrum at
criticality.

\paragraph{Self-avoiding walk and percolation.}
Self-avoiding walk provides a path analogue. Above its
short-range upper critical dimension four, the lace expansion gives a
Brownian scaling limit on the infinite lattice, first for weak
self-avoidance and then for nearest-neighbour strict self-avoidance
\cite{BrydgesSpencer1985,HaraSlade1991}. For weakly self-avoiding walk
on a torus in \(d>4\) with sufficiently small self-repulsion, Michta
obtains Brownian motion on the continuum torus when the step count
satisfies \(m=o(V^{1/2})\) and \(m/L^2\) tends to a positive limit.
The diffusion constant agrees with that on the infinite lattice
\cite{MichtaTorWalk}. At the subcritical edge of the finite-size window,
where the critical fugacity is approached on the scale $V^{-1/2}$,
the torus susceptibility is of order $V^{1/2}$. The two-point function
has a bulk power-law term and a plateau of order $V^{-1/2}$.
Slade proved this for weakly self-avoiding walk, and Liu established
the corresponding bounds for nearest-neighbour strict self-avoiding
walk \cite{SladeWSAWPlateau,LiuSAWPlateau}. The unwrapping picture associates the plateau with the sum over
windings around the torus, as illustrated in numerical studies of
self-avoiding walk and the Ising model \cite{GrimmEtAl2017}.

For short-range percolation above dimension six, the lace expansion
gives mean-field critical behaviour and random-walk power-law decay
of the bulk connectivity for the models covered by the expansion
\cite{HaraSladePercolation,HaraHofstadSlade2003}. Aizenman used
mean-field bounds to study the proliferation of incipient spanning
clusters in critical boxes \cite{AizenmanSpanning1997}. The finite-graph
triangle condition of Borgs, Chayes, van der Hofstad, Slade and Spencer
yields the Erd\H{o}s--R\'enyi scales $V^{-1/3}$ for
the critical window and $V^{2/3}$ for the largest clusters
\cite{BorgsEtAlI,BorgsEtAlII}. In the high-dimensional torus window,
Hutchcroft, Michta and Slade prove a two-point plateau of order
$V^{-2/3}$ and susceptibility of order $V^{1/3}$. At the bulk threshold
the free-box susceptibility has order $L^2$ in the setting of
Chatterjee and Hanson \cite{HutchcroftMichtaSlade2023,ChatterjeeHanson2020}.
For sufficiently high-dimensional nearest-neighbour and sufficiently
spread-out models with $d>6$, Blanc-Renaudie and Nachmias identify the
limiting ranked torus cluster volumes, after model-dependent
rescaling, with the ranked excursion lengths
of Brownian motion with parabolic drift above its past minimum
\cite{BlancRenaudieNachmias}. This is the random-graph law for cluster
sizes.

\paragraph{The pseudocritical window.}
Liu, Park and Slade put the torus plateau estimates into a common
framework by comparing the periodic two-point function with the sum of
the infinite-lattice two-point function over periodic images. Their
comparison theorem gives the scale of the susceptibility and the
plateau, and they verify its hypotheses for spread-out long-range
self-avoiding walk above its upper critical dimension. They formulate conjectures for the profiles throughout
the scaling window and for the corresponding free-boundary
pseudocritical behaviour \cite{LiuParkSlade2026}. For spins, the
distinction between the bulk critical temperature and the centre of the
finite-size window already has rigorous counterparts. Camia, Jiang
and Newman obtain the $L^2$ critical susceptibility for free-boundary
nearest-neighbour Ising boxes in \(d\ge4\), assuming that the bulk
critical two-point function has order \(\lvert x-y\rvert^{2-d}\)
\cite{CamiaJiangNewman2021}. In the hierarchical $|\varphi|^4$ model,
Michta, Park and Slade prove $\sqrt{V}$ susceptibility in the periodic and free-boundary windows, each centred at 
its own effective critical point, with an additional logarithmic correction in $d=4$.
For $d>4$ the free-boundary shift is of order $L^{-2}$, larger
than the window width $L^{-d/2}$ \cite{HierarchicalFBC}.
This clearly distinguishe the finite-size scalings with different boundary conditions.
This distinction and the presence of pseudocritical point also predicted by \cite{flores2016}.

Theorem~\ref{thm:susceptibility} addresses the marginal Euclidean box.
It gives susceptibility bounds at a point specified by the exact
principal-mode curvature. For the same spread-out interaction on
\(\mathbb Z^d\), Park proves critical two-point decay of order
\(\lvert x\rvert^{-d/2}\) \cite[Theorem~1]{IsingCorrelation}.
Here the principal profile of the finite cube enters both the quartic
energy and the coupling to a uniform magnetic source.

\subsection{Comparison with other RG formulations}\label{sec:status}

The perturbative flow and the estimates for a single renormalisation
group step have counterparts in the lattice constructions of
Bauerschmidt, Brydges and Slade \cite{BauerschmidtBrydgesSlade2014}.
Park's map treats weakly coupled \(\lvert\varphi\rvert^4\)
interactions with periodic boundary conditions \cite{RGMap}.
Slade's long-range analysis treats weak coupling at a positive distance
below the upper critical dimension \cite{SladeLR}. Here the interaction
starts from the Ising spin measure, and the estimates must hold at
\(d=2s\) on blocks meeting a physical boundary. The pressure expansion
of Section~\ref{sec:pressure-singleton} and the boundary estimates of
Section~\ref{sec:geometry} address these two features.

In the hierarchical model the renormalisation group identifies both the
free-boundary critical window and its displacement from the bulk
critical point \cite{HierarchicalFBC}. On the Euclidean cube
Theorem~\ref{thm:susceptibility} estimates the susceptibility at a
matched temperature. Determining the displacement of this temperature
from the bulk critical point is a separate question.

Two comparisons clarify the choice of reference Gaussian measure 
(equivalently, choice of the counterterm) in dimension four.
Appendix~\ref{app:fixed-kinetic} treats the Gaussian measure 
with a fixed kinetic coefficient.
Lemma~\ref{lem:dressed-connected-expansion} gives the normalized
fluctuation expansion in an independent background, and
Proposition~\ref{prop:kinetic-quartic-pairs} exhibits second-order
terms with support. Using this expansion in the
original coordinates would also require a local representation of the
transformed background. 
% Appendix~\ref{app:direct-comparison} estimates
% the hard covariance when the original field is retained. Positivity
% throughout the full matching window remains open for this direct
% comparison. 
% Corollary~\ref{cor:d4-prescribed} uses the covariance
% column \(T\ell\) obtained after changing the microscopic field.

\subsection{Outline of the proof}
\numberwithin{equation}{section}

We use the renormalisation group flow to control the flow of the coupling constants.
The running quartic copuling constant appears as the logarithmic correction in Theorem~\ref{thm:susceptibility}---at \(d=2s\), 
the quartic interaction is marginal both on the full
lattice and in the interior of a finite cube. 

In \cite[Theorem~1 and Sections~7, 10 and~22]{IsingCorrelation}
the RG is performed on periodic boxes. Translation invariance gives
one set of bulk coefficients, and the thermodynamic limit at positive
auxiliary mass before computing the limit to the critical two-point function
on \(\mathbb Z^d\). Here, with free boundary condition,
we keep the cube finite and tune the temperature
to the curvature of its principal mode. Deleting bonds at a face
changes the short-range spin pressure already at the initial scale.
Integrating a fluctuation then produces a position-dependent quadratic
interaction even from a quartic with constant coefficient.
The central challenge is to control these terms.

\paragraph{From the spins to a local RG step.}
We use the positive Gaussian split and the conditional Gaussian
representation of \cite[Propositions~2.4 and~17.1]{IsingCorrelation},
with the local split of Lemma~\ref{lem:local-split} at \(d=4\).
After subtraction of the bulk mass, the difference between the cube
and bulk interaction Hessians decays exponentially in the distance to
the boundary in units of \(R\) (Lemma~\ref{lem:initial-hessian-comparison}).
At the initial block size
\(b_0\asymp R\log R\), Proposition~\ref{prop:initial-cube-heads}
separates the quadratic terms of range \(O(b_0)\) from the tails.
Proposition~\ref{prop:cube-joint-initialization} assigns the tails to
a polymer activity so that its value on every polymer avoiding the
grouped boundary labels agrees exactly with its bulk value. In
\cite[Sections~6--7 and~17--18]{IsingCorrelation} one can absorb the
translation-invariant quadratic correction into the Gaussian covariance.
Here we keep the boundary correction in the interaction.

For the Gaussian integration we apply the positive polynomial
decomposition of \cite[Lemma~2.1]{IsingCorrelation} to the finite-range
auxiliary precision. A polynomial of degree \(m\) in this precision
has range \(O(mR)\). Thus a fluctuation covariance row agrees with the
bulk row whenever its propagation neighbourhood lies inside the cube.
This property holds also for the fractional interaction, whose
original spin couplings have infinite range.

\paragraph{Reblocking at a face.}
Let \(b\) be a fine block size and \(B=Lb\) the next size.
A coarse block at a face contains \(O(L^{d-1})\) fine blocks within
distance \(O(b)\) of the boundary. The field radius changes from
\(\ell_b=b^{-s/2}\) to \(\ell_B=L^{-s/2}\ell_b\). For an unrestricted
quadratic the resulting factor is \(L^{d-1-s}=L^{s-1}\), attained
on constant fields (Lemma~\ref{lem:face-count}). It gives no
contraction in dimensions two through four.

The covariance itself supplies the missing power. Set
\(\rho_b(x)=\min\{b,R+\dist(x,\Lambda^c)\}\). For an ordinary
fluctuation slice \(\Gamma_j\) at scale \(b\), the boundary resolvent
estimate of Proposition~\ref{prop:boundary-resolvent-diagonal} gives,
with \(\lambda\) the covariance parameter of Section~\ref{sec:gaussian},
\begin{equation}
 \Gamma_j(x,x)\le C_L\lambda b^{-s}
                  \bigl(\rho_b(x)/b\bigr)^{2p},
 \qquad 0<p<s/2 .
 \label{eq:outline-face-variance}
\end{equation}
A positive supersolution for each face bounds the killed resolvent near the boundary,
and there is a freedom of choosing $p \in (0,s/2)$ in the construction of $\Gamma_j$.
Taking the minimum of the supersolutions covers points near several
faces, and applying the estimate on both observation legs gives the
square of the weight in \eqref{eq:outline-face-variance}. The annular
difference estimates of Section~\ref{sec:ext-cov} give Gaussian moments
in the corresponding weighted field norm. On the old boundary layer
\(\rho_B/\rho_b\) is bounded. Restricting a coarse test field to a
fine block in this layer therefore gains \(L^{-p}\) in each argument.
For a quadratic, the combined factor is
\begin{equation}
 L^{d-1}\bigl(L^{-s/2}L^{-p}\bigr)^2=L^{s-1-2p}.
 \label{eq:outline-face-power}
\end{equation}
We choose \(\max\{0,(s-1)/2\}<p<s/2\), with \(p\ne1/2\).
Lemma~\ref{lem:cube-boundary-heat} then gives a strict contraction
for large \(L\), including the contributions from edges and corners.

There is a separate estimate for the quadratic created at the
current step. Contracting two legs of the quartic gives a coefficient
proportional to the diagonal of the filtered slice. Its difference
from the bulk value can extend to distance \(O(B)\) from a face,
so all \(O(B^d)\) sites in a coarse block can contribute.
Lemma~\ref{lem:new-boundary-quadratic} bounds this sum using the
slice diagonal and the coarse field radius. We retain the resulting
quadratic on that block. At the next step its support is already
in the old boundary layer, and \eqref{eq:outline-face-power} applies.
Corollary~\ref{cor:accumulated-boundary-source} sums these contributions
over the scales. The bound is uniform in the side length, with the surface
interaction fixed by the microscopic spin measure.

\paragraph{Localisation and the quartic flow.}
On a boundary block we keep the full quadratic, quartic and sextic
Taylor coefficients. In the interior the bulk symmetries allow a
smaller set of coefficients. Proposition~\ref{prop:cube-jet-projection}
matches the quadratic Taylor coefficient on pairs of spatial
polynomials of total degree at most two, and the quartic on constant
fields. The resulting Taylor remainder contracts after summation
over the fine blocks. When a covariance contraction joins different
blocks, Lemma~\ref{lem:cube-oriented} assigns its terms to the input
labels before reblocking. This keeps the second-order subtraction
exact across block interfaces. The bulk quartic consequently has
the same leading decrement \(\mathfrak b_j g_j^2\) on the cube
interior. Proposition~\ref{prop:filtered-bulk-bubble} bounds the
bubble increment \(\mathfrak b_j\) above and below by multiples of
\(\lambda^2\) on the ordinary scales away from the terminal cut.
The reciprocal increments along the trajectory of
Theorem~\ref{thm:path} give \(g_N\asymp_{d,R}(\log V)^{-1}\).

\paragraph{Matching the mass and the kinetic coefficient.}
For \(d<4\) a bulk quadratic with two spatial differences scales
by \(L^{s-2}<1\). Corollary~\ref{cor:fractional-cube-path} controls
this coefficient together with the boundary tensors and the polymer
activity. We solve the mass recurrence backwards from the equation of
zero terminal curvature. The microscopic mass has nonzero thermal
derivative, so varying the temperature crosses the required initial
mass (Corollary~\ref{cor:physical-mass-crossing}).

In \(d=4\) the kinetic factor \(L^{s-2}\) equals one.
Proposition~\ref{prop:bulk-kinetic-source} estimates the increment
of this coefficient, and Lemma~\ref{lem:kinetic-backward-sum} sums
the increments in the backward recurrence. We vary an auxiliary
Gaussian reference to match the initial kinetic coefficient as well
as the mass. The terminal curvature depends on this choice of
reference. To impose the curvature prescribed in
Section~\ref{sec:TheMatchingCondition}, we change the microscopic
field so that its covariance with the original coordinate is the
original column \(C_T\ell\). Stopping one scale earlier adds a
variance margin of order \(\lambda n^2\) in the Gaussian regression
(Proposition~\ref{prop:prescribed-positive}). The conditioned
covariance therefore stays positive throughout the matching rectangle.
Corollary~\ref{cor:d4-prescribed} uses it to impose the original
curvature condition at the physical temperature.

\paragraph{The magnetic moment.}
In the natural terminal field write \(\phi=a u_T+Y\), as in
Section~\ref{sec:TheMatchingCondition}. Its profile satisfies
\(\sum_x u_T(x)^4\asymp V^{-1}\) and
\(\lvert\one^{\mathsf T}u_T\rvert\asymp V^{1/2}\).
The quartic along this direction has size \(g_N a^4/V\), which puts
\(a\) on the scale \((V/g_N)^{1/4}\). After integration over \(Y\),
Proposition~\ref{prop:terminal} gives a magnetic second moment of
order \(V^{3/2}g_N^{-1/2}\). Restoring the uniform source in the
earlier slices costs \(O_R(V^{3/2})\). The positive-row comparison of
Lemma~\ref{lem:spin} transfers the moment to the spin variables.
In dimension four this moment is computed in the transformed field
at the earlier cut. The uniform source there gives a spin weight
bounded above and away from zero at every site, so the same comparison
applies (Corollary~\ref{cor:d4-prescribed}). Division by \(V\) gives
the factor \(\sqrt{V\log V}\) in Theorem~\ref{thm:susceptibility}.

\setcounter{tocdepth}{1}
\tableofcontents

\section{Microscopic and Gaussian estimates}\label{sec:gaussian}

The microscopic split has two roles. Its short-range part gives a
pressure whose quadratic coefficient can be varied with the temperature.
Its positive part gives a Gaussian reference whose principal variance
can be prescribed before the temperature is chosen. We first construct
the split and control the short-range pressure. We then estimate the
principal mode, the remaining covariance, and the bulk bubble that
drives the quartic flow.

For \(d<4\), we use the expansion
\begin{equation}
 I-L_R^\alpha=\sum_{m\ge1} c_m D_R^{*m},\qquad
 c_m=\frac{\alpha\,\Gamma(m-\alpha)}
            {\Gamma(1-\alpha)\Gamma(m+1)}>0.
\end{equation}
Use the positive decomposition of \cite[Proposition~17.1]{IsingCorrelation}
and split the series at a sufficiently large fixed even
\(N_{\rm sm}\). The short part contains the terms below
\(N_{\rm sm}\), and the long part contains the remaining terms.
After diagonal regularization and scalar normalization, the long
part gives \(A>0\), the short part gives \(Q\ge0\) entrywise, and
\(A+Q\) has bulk row sum \(1\). We write \(B=I-Q\).
Positivity of the long part follows from the beta integral for
\(\widehat D_R(k)^{N_{\rm sm}}/(1-u\widehat D_R(k))\) with a
positive measure on \(u\in(0,1)\). For finite matrices we suppress
\(\Lambda\) when there is no ambiguity. The short part has finite
range \(O(R)\), off-diagonal entries \(O(R^{-d})\) and row sum
bounded away from one. Retain the diagonal of \(Q\) in the definitions of \(A\) and \(B\).
It contributes only a constant to the spin energy and drops out of
the normalized short-range spin pressure. The temperature convention
for this split is \(\beta=(1+e^{-R})\tau\).

To obtain a finite-propagation decomposition for the Gaussian
reference, we realize the long part on an auxiliary space. On the
continuous channel take
\(\mathcal D_u=I-uD_R\) and
\(\mathcal L_u=(\zeta w_\alpha(u))^{1/2}D_R^{N_{\rm sm}/2}\), where
\begin{equation}
 \zeta=(1+e^{-R})^{-1},\qquad
 w_\alpha(u)=
 \frac{\alpha\,u^{N_{\rm sm}-\alpha-1}(1-u)^\alpha}
      {\Gamma(1-\alpha)\Gamma(1+\alpha)} .
\end{equation}
The white channel has diagonal \(I\) and coupling
\((\zeta e^{-R})^{1/2}I\).
With \(\mathcal L_\Lambda=\mathcal L P_\Lambda^*\), the
two-block precision has diagonal blocks \(\mathcal D,B_\Lambda/c\)
and off-diagonal blocks \(-\mathcal L_\Lambda,-\mathcal L_\Lambda^*\).
The observed covariance is its inverse tested through
\(\mathcal O=(\mathcal L_\Lambda^*,0)\), which gives \(V_c\) by the
Schur complement. The continuous auxiliary inverse is interpreted
as a quadratic form on the observation space.

We use the same field normalization for both splits. Write
\(g_{\mathrm{mic}}=R^{-d}\),
\(A_{\mathrm{sc}}=g_{\mathrm{mic}}^{-a_*}\), with
\(1/12<a_*<1/10\), and \(\lambda=A_{\mathrm{sc}}^{-2}\).
The canonical field is defined by
\(h_{\mathrm{phys}}=A_{\mathrm{sc}} R^{-d/4}\phi\).
It gives canonical covariance \(a_R V_c\), where
\begin{equation}
 a_R=\lambda R^s,\qquad
 V_c=(A^{-1}-cB^{-1})^{-1},\qquad
 E=B^{-1/2}AB^{-1/2},\qquad c_{*,\Lambda}=\eta_1(E)^{-1}.
\end{equation}
The split is chosen so that \(B-A\) is a positive scalar multiple
of the compressed \(L_R^\alpha\). Write \(e\) for the unit top
eigenvector of \(E\) with \(B^{-1/2}e>0\), and
\(\ell=B^{-1/2}e\). For
\(c=c_{*,\Lambda}-\theta_0\lambda^2\bar g_N(R/n)^s\)
the full principal variance is
\(n^s/(\theta_0\lambda\bar g_N)\).
The parameter \(\bar g_N=(G_{\rm ref}^{-1}+\lambda^2N)^{-1}\),
with \(G_{\rm ref}=A_{\rm sc}^4g_{\mathrm{mic}}\), depends only on
the reference scales, so this construction precedes the temperature
variation.

\subsection{The microscopic split and the thermal crossing}
\label{sec:micro-extension}

For a thermal crossing, the short-range reference needs a row sum
bounded away from one uniformly in \(R\). In dimension four the
binomial series for \(I-(I-D_R)^\alpha\)
terminates at degree one. Taking the original even-tail split with
\(N_{\mathrm{sm}}\ge 2\) would leave only the white regularization in
\(A\). The short-range row sum would then tend to one and the
uniform high-temperature gap would be lost. For the spread-out
operator we use the convolution-square construction of
\cite[Section~3]{IsingCorrelation}, with the explicit kernel below.

\begin{lemma}[Positive split for the local power]
\label{lem:local-split}
Take \(d=4\) and let \(D_R\) be uniform on the nonzero sites of
the spread-out cube. Fix a positive integer \(m\) and let
\(E_R\) be uniform, including zero, on the cube of radius
\(\lfloor R/(4m)\rfloor\). Write \(F_R=E_R^{*2m}\). For a fixed
sufficiently small \(\kappa>0\) and \(\epsilon_R=e^{-R}\), put
\begin{equation}
 \begin{gathered}
 t_R=\kappa F_R(0)+\epsilon_R,\qquad
 A=\frac{\kappa F_R+\epsilon_R I}{1+t_R},\\
 Q=\frac{D_R-\kappa F_R+\kappa F_R(0)I}{1+t_R},\qquad B=I-Q .
 \end{gathered}
\end{equation}
For sufficiently large \(R\), \(A>0\), \(Q\) is nonnegative
entrywise and has zero diagonal, and \(A+Q\) has row sum one.
The bulk row sums \(p=\sum_xA(x)\), \(q=\sum_xQ(x)\) satisfy
\(p+q=1\) and \(p\ge\kappa/2\). Furthermore,
\(\sup_x |Q(x)|\le C_m R^{-4}\) and
\begin{equation}
 B-A=(I-D_R)/(1+t_R).
\end{equation}
The physical inverse temperatures are related by
\(\beta=(1+t_R)\tau\), and the replacement does not change the
finite-volume spin measure.
\end{lemma}

\begin{proof}
The convolution \(F_R\) has support in the cube of radius \(R/2\),
row sum one and supremum at most \(\|E_R\|_\infty\le C_m R^{-4}\).
Every nonzero site in its support has \(D_R(x)\ge cR^{-4}\).
Taking \(\kappa\) small gives the off-diagonal positivity of \(Q\).
The Fourier multiplier of \(F_R\) is \(\widehat E_R^{\,2m}\ge0\).
The white term makes \(A\) strictly positive, including on
principal restrictions. Direct addition gives
\(A+Q=(D_R+t_RI)/(1+t_R)\). Its diagonal correction contributes a
constant in Ising variables, since \(\sigma_x^2=1\).
\end{proof}

The proof also works for a uniform Euclidean spread-out ball by
using a smaller inscribed smoothing cube. The smallness
\(R^{-4}\) is a large-coordination property. It is absent for a
fixed nearest-neighbour kernel. The identity \(s=2\) below concerns
the local power of the spread-out Laplacian.

For \(d<4\), the original positive binomial coefficients give the
same short-range hypotheses. Indeed, the finite head has a row
gap equal to the positive tail mass, uniformly at fixed
\(N_{\mathrm{sm}}\), and its maximum entry is \(O(R^{-d})\) because
\(\|D_R^{*i}\|_\infty\le\|D_R\|_\infty\). Write \(Q^\circ=Q-Q(0)I\) for the spin graph. Its normalized
pressure equals that of \(Q\). The two walk resolvents differ by
\(\beta Q(0)(I-\beta Q)^{-1}(I-\beta Q^\circ)^{-1}\), whose absolute
exponentially weighted row sum is \(O(R^{-d})\). On restricting to
the cube, a term in the difference from the bulk must also visit
the exterior. The unused walk length weight supplies the boundary
factor in Lemma~\ref{lem:extension-cavity}. Thus that lemma applies
with the frozen \(B=I-Q\) used in the Gaussian reference.

In dimension four the auxiliary realization can use two copies
of the lattice and a finite-range observation coupling
\begin{equation}
 \mathcal L=\frac{1}{\sqrt{1+t_R}}
 \begin{pmatrix}\sqrt\kappa\, E_R^{*m}\\ \sqrt{\epsilon_R}\,I\end{pmatrix},
 \qquad
 \mathcal Q_c=\begin{pmatrix}I&-\mathcal L_\Lambda\\
                     -\mathcal L_\Lambda^*&B_\Lambda/c
       \end{pmatrix}.
\end{equation}
Here \(\mathcal L_\Lambda=\mathcal L P_\Lambda^*\). At \(r\ge0\)
the observed bare resolvent is \(A_\Lambda/(1+r)\), and the
physical Schur complement of \(\mathcal Q_c+rI\) is
\(B_\Lambda/c+rI-A_\Lambda/(1+r)\). In particular
\(\mathcal O\mathcal Q_c^{-1}\mathcal O^*=V_c\) for
\(\mathcal O=(\mathcal L_\Lambda^*,0)\). Polynomials in the block precision therefore give finite-range
operators in the observed coordinates. Section~\ref{sec:ext-cov}
estimates their covariance up to the boundary.

The Hessian of the short-range pressure fixes the initial mass.
The next lemma compares it with the walk resolvent in the bulk and
estimates the error caused by deleting exterior bonds. We use the
first bound to differentiate the bulk mass with respect to temperature
and the second to localize the initial quadratic interaction near the
boundary.

\begin{lemma}[A weighted short-range Hessian estimate]
\label{lem:extension-cavity}
Suppose the symmetric kernel \(Q\) is nonnegative, has zero
diagonal and range at most \(C_0R\), and satisfies
\(\sup_x\sum_y Q_{xy}\le q_0<1\) and \(\sup_{x,y}Q_{xy}\le C_0 g\).
For \(g\) sufficiently small there are \(a,r>0\) such that for
\(|\beta-1|<r\) and every finite induced set \(U\), with
\(S_{\beta,U}=D^2\log Z_{\beta Q,U}(0)\), the error
\(E_U=S_{\beta,U}-(I-\beta Q_U)^{-1}\) obeys
\begin{equation}
 \sup_{x\in U}\sum_{y\in U}
      e^{a|x-y|/R}|E_U(x,y)|\le Cg .
\end{equation}
If \(U=\widetilde U\cap\Lambda\), then, with
\(\delta(x)=\operatorname{dist}(x,\Lambda^c)\),
\begin{equation}
 \sup_{x\in U}\sum_{y\in U}
 e^{a(\delta(x)+|x-y|)/R}
 |E_{\widetilde U}(x,y)-E_U(x,y)|\le Cg .
\end{equation}
The bounds pass to the bulk limit and give thermal derivative
bounds on a smaller complex disc.
\end{lemma}

\begin{proof}
Write \(t_{xy}=\tanh(\beta Q_{xy})\). By the support bound, the
gap, and \(\tanh z=z+O(z^3)\), there are \(\rho<1\) and \(a_0>0\)
such that the matrix
\begin{equation}
 T_{xy}=|t_{xy}|\exp\{a_0(1+|x-y|/R)\}
\end{equation}
has row sum at most \(\rho\) and maximum at most \(Cg\).
The constants are uniform on a fixed small complex \(\beta\)-disc.
A simple cycle \(C\) through \(x\) has a rooted oriented walk of
the same length, so its \(T\)-weight summed over all \(C\ni x\)
is at most \(\sum_{k\ge3}(T^k)_{xx}\le
Cg\sum_{k\ge3}\rho^{k-1}\).

Let \(Z_{\mathrm{ev}}(U)\) be the high-temperature partition sum
over even edge sets on \(U\), normalized by stripping the
\(2^{|U|}\prod_{xy}\cosh(\beta Q_{xy})\) factor. Its polymers are
nonempty connected even edge sets, compatible when their vertex
sets are disjoint. Each such set decomposes into edge-disjoint
simple cycles. Fix one decomposition and a spanning tree of the
intersecting cycles, rooted at a cycle through a specified anchor
\(x\). Attach each child at one shared site and order the children
at each site. The union of edge sets is determined by these
data. Thus the sum over ordered rooted cycle trees majorizes
the anchored sum. If \(S\) bounds the total descendant-tree
weight rooted at any site, ordering the children costs at most
\((1-S)^{-k}\) at a cycle of length \(k\). For \(S\le M g\),
increasing \(T\) by this factor per edge still leaves row sum
less than one for small \(g\). Choosing \(M\) large and
inducting in tree height gives \(S\le M g\).
The same argument with part of \(a_0\) reserved proves an
anchored \(O(g)\) bound with spare factors exponential in the
number of edges and in their total length divided by \(R\).

For any \(v\in U\) let
\(r_U(v)=Z_{\mathrm{ev}}(U)/Z_{\mathrm{ev}}(U\setminus\{v\})\).
Separate the unique polymer containing \(v\), if present. Then
\begin{equation}
 r_U(v)=1+
   \sum_{\substack{F\ {\rm connected,\ even}\\v\in V(F)\subset U}}
      t(F)\frac{Z_{\mathrm{ev}}(U\setminus V(F))}
                 {Z_{\mathrm{ev}}(U\setminus\{v\})},
 \qquad t(F)=\prod_{e\in F}t_e .
\end{equation}
The quotient on the right is a product of inverse deletion
ratios on smaller induced sets. For ratios within \(Cg\) of
one it costs \(\exp(C'g|V(F)|)\), absorbed by the reserve.
Induction in the set size therefore gives nonvanishing and
\(r_U(v)=1+O(g)\). Comparing the same recurrence on \(U\) and
\(\widetilde U\), the missing polymers reach outside \(\Lambda\),
so their span weight supplies \(\exp(-a\delta(v)/R)\).
For common polymers replace the inverse deletion factors one
at a time. A site \(w\) in the polymer costs
\(\exp(a|w-v|/R)\) to move the distance weight from \(w\) to \(v\).
The spare length weight pays this cost, and the spare size weight
pays \(|V(F)|-1\) choices of the changed factor. The resulting
anchored bound is \(Cg\) times the induction bound on the
weighted difference of deletion ratios. Absorb this factor
for small \(g\) to obtain
\begin{equation}
 |r_{\widetilde U}(v)-r_U(v)|\le Cg e^{-a\delta(v)/R}.
\end{equation}
Removing the same common vertices from both domains preserves
the comparison hypothesis, so this induction applies to all
products in the recurrence.

For \(x\ne y\), the spin two-point formula sums the connected
edge sets \(F\) with odd set \(\{x,y\}\), times
\(Z_{\mathrm{ev}}(U\setminus V(F))/Z_{\mathrm{ev}}(U)\). The two
odd vertices must belong to the same component. Fix a
self-avoiding \(x\)-to-\(y\) path in \(F\) and decompose the
remaining even set into cycles. For a path of length \(k\),
the product of deletion ratios differs from one by at most
\(Cg(k+1)e^{Cg(k+1)}\). The contribution with at least one
cycle has the same bound times the path weight. Indeed,
the ordered cycle forests attached to the \(k+1\) path
vertices sum to at most
\((1-Mg)^{-(k+1)}-1\), with the deletion product absorbed by
a small part of the size weight. Hence subtracting the
self-avoiding path sum leaves a weighted row error \(O(g)\).
For domain comparison, any lost path with a nonzero
deletion correction or any lost cycle-bearing graph has
total length at least \(\delta(x)\) and still has its
\(g\) factor. For common graphs telescope the deletion
products and use the weighted ratio estimate just proved,
moving \(\delta(w)\) to \(\delta(x)\) along the connected
graph. This costs its reserved length weight.

To replace the self-avoiding path sum by the walk resolvent,
a repeated site at times \(i<j\) contributes at most
\(Cg\rho^{k-1}\) at length \(k\) after the weighted row sum.
This follows by bounding the return between \(i\) and \(j\)
by \(\sup_z (T^{j-i})_{zz}\le Cg\rho^{j-i-1}\) and
summing the two remaining path sections. The sum over \(k\)
and \(i,j\) is \(O(g)\). Missing repeated walks have spare
length at least \(\delta(x)\). Finally, the expansion of
\((I-t)^{-1}-(I-\beta Q)^{-1}\) has one marked edge
\(t_{xy}-\beta Q_{xy}=O(g^2 Q_{xy})\) on each walk,
so its comparison obeys the same (in fact stronger) bounds.
These arguments include \(x=y\) with the trivial path, since
\(S_{\beta,U}(x,x)=1\). Choose the final \(a\) smaller than
the length reserve both for \(|x-y|\) and for \(\delta(x)\).
Exhaustion and Cauchy's formula give the last assertion.
\end{proof}

To identify the Hessian in the initial field, we apply the Gaussian
reconstruction of \cite[Proposition~2.4 and Section~17]{IsingCorrelation}
to the compressed interaction. Let \(P_{\beta,\Lambda}(v)=\log\langle
 e^{v^{\mathsf T}\sigma}\rangle_{\beta Q_\Lambda}\), and write
\(h=h_{\rm phys}\). Completing the square in the joint weight
\begin{equation}
 \exp\!\left\{-\tfrac12h^{\mathsf T}A_\Lambda^{-1}h
       +\sqrt\beta\,h^{\mathsf T}\sigma
       +\tfrac\beta2\sigma^{\mathsf T}Q_\Lambda\sigma\right\}\,dh
\end{equation}
gives the spin interaction \(\beta(A_\Lambda+Q_\Lambda)\) and the
conditional law \(h\mid\sigma\sim N(\sqrt\beta A_\Lambda\sigma,A_\Lambda)\).
Relative to the centered Gaussian of covariance \(V_c\), the negative
log density of \(h\), up to an additive constant, is
\begin{equation}
 \tfrac c2h^{\mathsf T}B_\Lambda^{-1}h-P_{\beta,\Lambda}(\sqrt\beta\,h).
\end{equation}
In the pressure coordinate \(u=B_\Lambda^{-1}h\), its Hessian at zero
is \(cB_\Lambda-\beta B_\Lambda S_{\beta,\Lambda}B_\Lambda\).

Write \(T_{\beta,\Lambda}=cB_\Lambda-\beta B_\Lambda S_{\beta,\Lambda}B_\Lambda\)
for this Hessian, \(T_\beta^\infty(z)\) for its bulk kernel, and
\(B_\beta=I-\beta Q\). Separate the bulk row mass from its zero-row
remainder by setting \(r(\beta)=\sum_z T_\beta^\infty(z)\) and
\(Z_\beta(z)=T_\beta^\infty(z)-r(\beta)\one_{\{z=0\}}\).

\begin{lemma}[Bulk mass and localization of the initial Hessian]
\label{lem:initial-hessian-comparison}
Let \(g=R^{-d}\) and \(|c-1|+|\beta-1|\le C_0g\) with \(C_0\)
fixed. For some \(a>0\),
\begin{equation}
 \begin{aligned}
 \sum_z e^{a|z|/R}|T_\beta^\infty(z)|&\le Cg, &
 \sum_z |z|^2 |Z_\beta(z)|&\le Cg R^2,\\
 \sup_{x\in\Lambda}\sum_{y\in\Lambda}
 e^{a(\delta(x)+|x-y|)/R}
 |T_{\beta,\Lambda}(x,y)-T_\beta^\infty(x-y)|&\le Cg .
 \end{aligned}
\end{equation}
At fixed \(c\), the bulk and cube Hessians satisfy
\(\partial_\beta T_\beta=-I+O(g)\) in the exponentially weighted row norm. The derivative of the
finite-volume difference obeys the same boundary estimate.
\end{lemma}

\begin{proof}
On either domain the identity
\begin{equation}
 T_\beta=(c-1)B-(\beta-1)I-(\beta-1)^2 QB_\beta^{-1}
             -\beta B E_\beta B
\end{equation}
follows by multiplying
\(\beta B B_\beta^{-1}B=B+(\beta-1)I+
(\beta-1)^2QB_\beta^{-1}\) by \(B_\beta\).
Apply Lemma~\ref{lem:extension-cavity} to the last term and use
exponentially weighted convolution.
The walk expansion of \(QB_\beta^{-1}\) converges in the
weighted row norm. A walk present only in the bulk has length at
least \(\delta(x)\). The same argument applies to the frozen primary
filter \(B_\Lambda^{-1}\) on both sides of the Hessian.
The second moment follows by using a smaller exponential weight.
Differentiation at fixed \(c\) gives \(\partial_\beta T_\beta=-I+O(g)\)
in the weighted row norm on the stated window. For the pressure error
use Cauchy's formula on the fixed disc of Lemma~\ref{lem:extension-cavity}.
The same differentiation and walk comparison give the boundary
weight for the derivative of the finite-volume difference.
\end{proof}

\begin{corollary}[Thermal crossing of the bare mass]
\label{cor:bare-thermal-crossing}
Let \(p=\sum_z B(z)>0\) and \(h_0=\sum_z B^{-1}(z)=p^{-1}\). Let \(c\) be
fixed with \(|c-1|\le C_0 g\). Given \(M<\infty\), take \(A\)
sufficiently large in terms of \(M,C_0\) and the split.
For all sufficiently large \(R\), the bare mass \(h_0^2r(\beta)\)
crosses every continuous real target of absolute value at most
\(Mg\) on \(|\beta-1|\le Ag\).
\end{corollary}

\begin{proof}
With \(q=\sum_z Q(z)=1-p\), the bulk row sum is
\begin{equation}
 r(\beta)=cp-\beta p^2\{(1-\beta q)^{-1}+e_\beta\},
 \qquad |e_\beta|+|e_\beta'|\le Cg
\end{equation}
by Lemma~\ref{lem:extension-cavity}. At \(\beta=1\) this gives
\(r(1)=O(g)\) with a constant independent of \(A\).
Lemma~\ref{lem:initial-hessian-comparison} gives \(r'(\beta)=-1+O_A(g)\)
throughout the interval. Thus the two endpoint values of the
primary-coordinate mass \(h_0^2r\) bracket the target values
for \(A\) sufficiently large, and the conclusion follows by
continuity. The prescribed \(c\) satisfies \(|c-1|=O(g)\) for
all sufficiently large \(n\) at fixed \(R\), as required for the
shooting argument in Section~\ref{sec:path}.
\end{proof}

\subsection{All-field pressure damping and the singleton seed}
\label{sec:pressure-singleton}

The polymer expansion requires control of the pressure at every
real background field. The all-field forest expansion for the periodic
short-range reference in \cite[Proposition~5.1]{IsingCorrelation}
motivates the construction here. Deleted bonds and the frozen filter
require estimates uniform over the induced graphs.
For one cell we compare the Hessian-subtracted pressure with the
onsite quartic and keep the damping in both weights. For a connected
set of cells we also have to differentiate in the bonds between cells.
Lemma~\ref{lem:pressure-singleton} gives the one-cell estimate on
every induced graph, uniformly in the real interpolated couplings.

\begin{lemma}[Short-reference pressure on an induced graph]
\label{lem:pressure-singleton}
Let \(U\) be a finite set of sites, \(J\) a real symmetric
matrix with zero diagonal and absolute row sum at most \(q<1\),
and \(\sup_{x,y}|J_{xy}|\le Cg\). Write
\(B=I-J\),
\begin{equation}
 P(h)=\log\frac{\sum_\sigma e^{\sigma^{\mathsf T} J\sigma/2+h^{\mathsf T}\sigma}}
                         {\sum_\sigma e^{\sigma^{\mathsf T}J\sigma/2}},
 \qquad S=D^2P(0).
\end{equation}
If a real matrix \(L\) satisfies
\(\|L-B\|_{1\to1}+\|L-B\|_{\infty\to\infty}\le Cg\), set
\begin{equation}
 {\cal U}(u)=\tfrac12(Lu)^{\mathsf T}S(Lu)-P(Lu),\qquad
 V_4(u)=\tfrac1{12}\sum_xu_x^4,\qquad
 \Psi(t)=\frac{t^4}{1+t^2}.
\end{equation}
For \(g\) sufficiently small and real \(u\),
\begin{equation}
 {\cal U}(u)\ge c\sum_x\Psi(u_x),\qquad
 |{\cal U}(u)-V_4(u)|\le C\sum_x(|u_x|^6+g|u_x|^4).
\end{equation}
For each fixed \(k\) and fixed field-radius multiple, if all
\(\|v_i\|_\infty\le1\), then at sufficiently small \(r>0\),
\begin{equation}
 \begin{aligned}
 r^k\big|D^k(e^{-{\cal U}}-e^{-V_4})(u)[v_1,\ldots,v_k]\big|
 &\le C_k e^{-c\sum_x\Psi(u_x)+C_k|U|r^4}\\
 &\quad{}\times
     \sum_x\big((|u_x|+r)^6+g(|u_x|+r)^4\big).
 \end{aligned}
\end{equation}
The bound for \(e^{-{\cal U}}\) holds with the last sum replaced
by \(1\). The constants are independent of \(|U|\) and of all
coordinate values of the real background \(u\).
\end{lemma}

\begin{proof}
We first compare the exact magnetization with its mean-field value,
including derivatives in the field. The mean-field variational
formula will then give the damping and identify the quartic
coefficient. To control the derivatives, we use a Taylor norm. For a power series in field increments
\(\boldsymbol\xi\), use the norm
\(\|F\|_a=\sum_{\boldsymbol j}|[\boldsymbol\xi^{\boldsymbol j}]F|
a^{|\boldsymbol j|}\), with \(a>0\) fixed and small. This is a
submultiplicative norm and each coordinate increment has norm
\(a\). At a real center \(z_0\), the Taylor coefficients of
\(\tanh(z_0+\cdot)\) and of its fixed derivatives have uniform
Cauchy bounds in a fixed disc, independent of \(z_0\).
In particular \(\|(\tanh)'(z)\|_a\le1+C_k a\) if
\(\|z-z(0)\|_a\le k a\). The exact cavity message for an edge
\(b\in\mathbb R\) is
\begin{equation}
 f_b(z)=\tfrac12\bigl(\log\cosh(z+b)-\log\cosh(z-b)\bigr).
\end{equation}
Integrating the derivatives of \(\tanh\) along the real segment
from \(-b\) to \(b\) bounds its derivative by
\(|b|(1+C_k a)\) in this algebra. The same segment formula,
with its centered second-order remainder, gives
\(\|f_b(z)-b\tanh z\|_a\le C|b|^3\).

For a site \(i\), the exact field giving its magnetization
as \(\tanh z_i\) is obtained by removing \(i\) from the graph
and taking the ratio of partition functions at the fields
shifted by \(\pm J_{i,\cdot}\). Telescope the ratio by shifting
one neighbor field at a time. Each term is
\(f_{J_{ij}}(z_j^{(i,j)})\), where the field in the smaller graph
has shifts by \(\pm J_{ik}\) at other neighbors \(k\).
Iteration deletes one ancestor at each depth. In particular
the direct field shift and the sum of omitted couplings in a
row are both bounded by \(C g l\) at depth \(l\).
The mean-field field vector is the unique solution
\(z^{\rm mf}=h+J\tanh z^{\rm mf}\) on the real axis.
For the increment series, choose a fixed \(k\) with
\(1+qk<k\) and then \(a\) so small that
\(1+q(1+C_k a)k\le k\). Center the cavity fields and the
mean-field fields at their respective real values.
The exact recursion and the contraction defining the mean-field
solution give the same bound \(ka\) for every increment series.
The cavity ratio is used first at real fields; its analytic
formula gives the germ on the field polydisc. The increment
majorant gives a convergent extension on a disc of uniform
radius. In particular it gives that extension for the
gradient of \(P\) at every real center.

Compare an exact cavity field with \(z^{\rm mf}\) at the same
root. The segment joining their series has real center and
increment norm at most \(ka\), so the preceding Lipschitz bound
applies even when the centers differ. With
\(q'=q(1+C_k a)<1\), the maximum discrepancy at depth \(l\)
is at most \(Cg(l+1)\) plus \(q'\) times the corresponding maximum
at depth \(l+1\). In this estimate omitted couplings and
ancestor shifts cost \(Cgl\), and replacing each \(f_b\)
by \(b\tanh\) costs \(\sum_j C|J_{ij}|^3=O(g^2)\).
Backward summation gives the root discrepancy \(O(g)\).
Applying \(\tanh\) proves a Taylor-algebra bound \(O(g)\)
on the difference of each magnetization component,
uniformly around every real background.

Define the mean-field pressure and its entropy function by
\begin{equation}
 P_{\rm mf}(h)=\sup_{m\in[-1,1]^U}
  \bigl(h^{\mathsf T}m+\tfrac12m^{\mathsf T}J m-\textstyle\sum_x I(m_x)\bigr),
 \qquad
 I(t)=\tfrac12\bigl((1+t)\log(1+t)+(1-t)\log(1-t)\bigr).
\end{equation}
Its maximizing vector is \(m=\tanh z^{\rm mf}\), since the
entropy Hessian on the box is bounded below by the identity
and the maximization is strictly concave.
In particular \(D^2P_{\rm mf}(0)=B^{-1}\).
The Taylor-algebra bound just proved gives, for every fixed
\(l\ge1\), the estimates
\begin{equation}
 \sup_{x_1}\sum_{x_2,\ldots,x_l}
    |\partial_{x_1}\cdots\partial_{x_l}(P-P_{\rm mf})(h)|
       \le C_l g .
\end{equation}
Indeed ordered index tuples at fixed degree multiply the
absolute sum of power-series coefficients by \((l-1)!\).
For \(P\) alone the same estimate holds with \(C_l\)
instead of \(C_lg\).
The bounds hold on a smaller complex field strip by
translating the Taylor series and using its spare radius.
They also hold after composing with \(L\), whose absolute
row and column sums are bounded. The pressure itself
extends by integrating its gradient on the strip,
with the real normalization above. Its mixed derivatives
are symmetric, so any of their indices can serve as the root.

Write \(\Delta=(P-P_{\rm mf})\circ L\) and
\(\Delta_{\ge4}(u)=\Delta(u)-D^2\Delta(0)[u,u]/2\).
By spin flip \(\Delta\) is even, and \(\Delta(0)=0\).
For \(u^s\) equal to \(u\) on \(|u_x|\le1\) and zero
elsewhere, Taylor's formula and the fourth rooted bound
give \(|\Delta_{\ge4}(u^s)|\le Cg\sum_x|u_x^s|^4\).
Here and below, bounds on a symmetric derivative tensor
are converted to estimates of sums of powers by the
arithmetic--geometric mean inequality in the indices.
To compare \(u\) with \(u^s\), the gradient of \(\Delta\)
is bounded by \(Cg\) in supremum norm, and \(D^2\Delta(0)\)
has absolute row sum at most \(Cg\). The change of \(\Delta\)
is at most \(Cg\sum_{|u_x|>1}|u_x|\).
For the change of the quadratic use the same bound
on the cross term, since \(\|u^s\|_\infty\le1\), and
\(Cg\sum_{|u_x|>1}|u_x|^2\) for the pure large term.
Therefore
\begin{equation}
 |\Delta_{\ge4}(u)|\le Cg\sum_x\Psi(u_x).
\end{equation}

For \(v=B^{-1}Lu\), the exact variational identity gives
\begin{equation}
 \tfrac12 v^{\mathsf T} Bv-P_{\rm mf}(Bv)
   =\inf_{m\in[-1,1]^U}
      \left\{\tfrac12(v-m)^{\mathsf T}B(v-m)
                  +\sum_x [I(m_x)-m_x^2/2]\right\}.
\end{equation}
The right side is bounded below by \(c\sum_x\Psi(v_x)\):
\(B\ge (1-q)I\), and \(I(m)-m^2/2\ge c m^4\) for \(|m|\le1\).
If \(|v_x|\le2\), separate \(|m_x|\ge |v_x|/2\) from its
complement; if \(|v_x|>2\), use \(|m_x|\le1\).
The convex function \(\Psi\) has
\(\Psi(a t)\le \max\{a^2,a^4\}\Psi(t)\) for \(a>0\).
Jensen's inequality shows that a matrix with bounded
absolute row and column sums is bounded for the modular
\(\sum_x\Psi(v_x)\), up to fixed constants.
Since \(u\mapsto v\) and its inverse have that property
by the row gap and the bound on \(L-B\), this proves the
real damping after using \(\Delta_{\ge4}=O(g\sum\Psi)\).

The mean-field quartic coefficient at \(v=0\) is
\(\sum_x v_x^4/12\). In fact \(m=v+O(\|v\|_\infty^3)\)
by the stationarity equation and the invertibility of \(B\),
so this is the quartic coefficient of the variational
identity. Since \(B^{-1}L=I+O(g)\) in absolute row and
column norms, its quartic tensor in \(u\) differs from
that of \(V_4\) by \(O(g)\) in rooted norm.
The fourth rooted bound on \(\Delta\) adds \(O(g)\).
The analytic bounds for \(P\) give bounded sixth rooted
derivatives of \({\cal U}\) on a fixed strip. Together
with evenness and Hessian subtraction, Taylor's formula
along the segment from zero to \(\zeta\) gives
\begin{equation}
 |{\cal U}(\zeta)-V_4(\zeta)|
     \le C\sum_x (|\zeta_x|^6+g|\zeta_x|^4)
\end{equation}
for all \(\zeta\) on a smaller strip.

For the derivative estimate we also need complex damping.
For real \(u,w\) with \(\|w\|_\infty\) small, the linear
imaginary shift contributes zero to the real part.
To bound its second-order remainder use
\(D^2{\cal U}(0)=D^3{\cal U}(0)=0\) and the integral
of the fourth derivative along the segment from zero to
\(u+itw\). By the fourth rooted bound, after separating
\(|u_x|\le1\) from the other sites, one has for any
\(\varepsilon>0\) and \(0\le t\le1\),
\begin{equation}
 |D^2{\cal U}(u+itw)[w,w]|
 \le C\|w\|_\infty^2\sum_{|u_x|>1}|u_x|^2
       +C\varepsilon\sum_{|u_x|\le1}|u_x|^4
       +C_\varepsilon\sum_x|w_x|^4.
\end{equation}
Indeed the tensor contracts with two \(w\) factors and two
\(u+itw\) factors. In a term containing a large \(u_x\)
bound both \(w\) factors by their supremum and use the
quadratic inequality on the large factors while bounding
the small factors by a constant. For terms with small
factors only, use Young's inequality, giving an arbitrarily
small coefficient to the fourth powers of the small \(u\)
factors. Sum the resulting one-index powers using the
rooted bounds. Take \(\varepsilon\) and then
\(\|w\|_\infty\) small to absorb the first two terms into
real damping. This proves
\begin{equation}
 \operatorname{Re}{\cal U}(u+iw)\ge
                   c\sum_x\Psi(u_x)-C\sum_x|w_x|^4 .
\end{equation}
The same lower bound holds for \(V_4\) by direct expansion
and Young's inequality. A real displacement bounded by
\(C_k r\) is handled by the convexity and scaling inequality
for \(\Psi\), at an additional cost \(C_k |U|r^4\) in the exponent.
On the polydisc of radii a fixed multiple of \(r\) in the
\(k\) prescribed field directions, the difference of the
two exponential weights is bounded using its exact
interpolation integral and the complex bound on
\(|{\cal U}-V_4|\). Cauchy's formula on this finite polydisc
gives the stated directional jet bounds.
\end{proof}

At the two initial radii the factor \(\exp(C_k|U|r^4)\)
is bounded on a single cell of volume of order \(b_0^d\).
The last sum in the singleton estimate is at most
\begin{equation}
 C|U|(r^6+g r^4)(1+\|u\|_\infty/r)^6 .
\end{equation}
Thus the lemma supplies the \(b_0^d(r^6+g r^4)\) singleton
factor with stable damping and a fixed polynomial loss.
It also applies to a fixed real row-filter substitution,
with the field-direction radii adjusted by its absolute
row bound. For several cells, the derivatives in the forest formula
connect the cells of an atom. The next proposition retains the same
small factor after summing over its spanning trees.

\begin{proposition}[The multicell pressure forest]\label{prop:pressure-trees}
Let \(J^0\) be the real symmetric, zero-diagonal matrix of microscopic
bonds, with range at most \(C_0R\) and \(\sup_{i,j}|J^0_{ij}|\le Cg\).
Assume the absolute row gap of Lemma~\ref{lem:pressure-singleton}
for all real forest interpolants, including \(J^0\). Let the cells have side
comparable to \(b_0=M_RR\), with \(M_R\) a sufficiently large multiple
of \(\log R\), and use the frozen thermal chart
\(L_J=\sqrt\beta(I-J/\beta)\), \(|\beta-1|\le Cg\).
A fixed diagonal correction to \(L_J\) of size \(O(g)\) is
allowed. This accommodates removing a constant short diagonal
from the spin Hamiltonian while keeping the original field filter.
For either initial radius \(r\) suppose \(b_0^d r^4\le C\) and
\(\nu_r=b_0^d(r^6+g r^4)=o(1)\) with \(\nu_r\ge R^{-C_0}\).
For each atom \(C\) of at least two cells, its sum of BKAR tree
derivatives of \(e^{-\mathcal U_J(u)}\) satisfies, through any fixed
number of field jets tested on directions of supremum norm at most \(r\),
the bound
\begin{equation}
 C_{A,\epsilon}\nu_r
  \exp\left[-A|C|-c\sum_{x\in \cup C}\Psi(u_x)
       +\epsilon\sum_{B\in C}\frac{\sup_{x\in B}|u_x|^2}{r^2}\right].
\end{equation}
Here \(A<\infty\) and \(\epsilon>0\) are prescribed, \(c>0\) is
fixed, and \(R\) is sufficiently large in terms of these constants.
The bound is uniform in the real interpolation parameters.
\end{proposition}

\begin{proof}
Fix a tree \(T\) on the cells. We first estimate the field and bond
coefficients of the pressure. Their spatial decay will pay for the
sum over the trees. Write \(U=\cup C\). For each microscopic bond
\(\{i,j\}\) over a cell edge \(e\in T\), introduce an independent
variable \(z_{ij}\), so the interpolated coupling varies by
\(J^0_{ij}z_{ij}\). Identifying all such variables above \(e\)
with \(\theta_e\) at the end gives the desired cell derivative.
For a small fixed \(\eta>0\), use the field envelopes
\begin{equation}
 W_B=1+\sup_{y\in U}e^{-\eta\,\operatorname{dist}(y,B)/R}|u_y|/r,
 \qquad W_e=\max(W_B,W_D)\quad(e=\{B,D\}).
\end{equation}
The analytic radius of a microscopic variable above \(e\) will
be \(a/W_e\) for a small fixed \(a>0\).

To estimate the joint coefficients at a root \(x\), use the
absolute field and bond Taylor norm with field radii \(a\) and
the above bond radii, multiplying each monomial by
\(e^{\mu\mathcal L/R}\). Here \(\mathcal L\) is the minimum length
of a connecting network for the root, the field indices, and both
endpoints of each microscopic bond in the monomial. The norm is
submultiplicative at a fixed root by taking the union of two
networks. Moving the root from \(i\) to \(j\) costs at most
\(e^{\mu|i-j|/R}\). We take \(\mu>0\) small depending on the row
gap and then \(\eta\) small enough that \(6\eta<\mu/2\).

In this norm the cavity recursion in Lemma~\ref{lem:pressure-singleton}
gives uniform bounds on each gradient component
\(\partial_{u_x}(P_J\circ L_J)\), and a bound \(Cg\) on its
difference from \(\partial_{u_x}(P_{{\rm mf},J}\circ L_J)\).
For clarity, the expansion can be centered at any \(t u\),
\(0\le t\le1\), and at the real interpolated bond values.
The change in \(L_J t u\) per row has norm \(Ca r\), since
\(|u_j|\le r W_e\) at the endpoints of the varied bond.
Independent field increments cost \(Ca\). The nonconstant part
of the sum of ancestor shifts also costs \(Ca\), by the row sum
bound and bounded length of each shifted bond. Thus the increments
of all cavity messages and of the mean-field solution are bounded
by \(K a\) as in the previous proof.
At a real forest interpolant the bond is \(s_{ij}J^0_{ij}\), with
\(0\le s_{ij}\le1\). Its increment \(J^0_{ij}z_{ij}\) has radius
at most \(a|J^0_{ij}|\) since \(W_e\ge1\). The integral formula
\(f_b(z)=\tfrac b2\int_{-1}^1\tanh(z+t b)\,dt\) and the field
derivative bound in Lemma~\ref{lem:pressure-singleton} give the
rooted Lipschitz majorant
\begin{equation}
 |J^0_{ij}|(1+C_K a)e^{\mu |i-j|/R}.
\end{equation}
For \(q=\sup_i\sum_j|J^0_{ij}|<1\) its row sum is at most
\(q(1+C_K a)e^{C_0\mu}<1\) when \(\mu,a\) are sufficiently small.
Expanding the integrand through first order in \(t b\), the odd
term integrates to zero. The bounded second derivative of \(\tanh\)
on the smaller field disc gives a weighted error \(C|J^0_{ij}|^3\)
in replacing \(f_b(z)\) by \(b\tanh z\). Its row sum is
\(O(g^2)\), since \(\sum_j|J^0_{ij}|^3\le q\sup_{k,l}|J^0_{kl}|^2\).
Comparing the depleted message to the
full mean-field message, the omitted ancestors and their field
shifts still have weighted norm at most \(Cg(l+1)\) at depth
\(l\); each shift uses an actual bond of bounded length.
Summation of the same geometric recursion gives the \(Cg\) error.
Composition by \(L_J\) and the derivative chain rule cost a
bounded weighted row sum. These arguments prove the asserted
joint field/bond coefficient bounds. Fixed field derivatives
use only a fixed part of the field radius, so they obey the
same bounds with constants depending on their orders. The bounds
also apply at \(t(u+\xi)\) if \(\|\xi\|_\infty\le C_k r\) is a
complex shift: translation of the field series by this constant
shift uses a spare part of \(a\) and only decreases the network
length when an index is removed.

Set \(\mathcal R_J=\mathcal U_J-V_4\), and decompose it as
\(\sum_x \zeta_x\int_0^1\partial_x\mathcal R_J(t\zeta)\,dt\),
where \(\zeta=u+\xi\). In the integral use Taylor expansion of
\(\partial_x\mathcal R_J\) through degree four at zero, with
the fifth derivative integral remainder. Evenness and exact
Hessian subtraction hold as identities in the bond variables.
The cubic part has weighted coefficient bound \(Cg\): for
the exact-minus-mean-field part this follows from the joint
estimate above. For the mean-field part its quartic before
subtracting \(V_4\) is
\(\sum_x ((I-J)^{-1} L_J u)_x^4/12\).
The Neumann series gives
\((I-J)^{-1}L_J-I=O(g)\) in weighted row and column sums of
absolute coefficients. Indeed \(L_J-(I-J)\) is a linear
combination of \(I,J\) with coefficients \(O(g)\) and an optional
fixed \(O(g)\) diagonal, and
each additional step of the Neumann series has weighted row
sum less than one. Replacing one of the four corresponding
matrix factors at a time gives the claimed cubic gradient
bound. The fifth gradient derivative has bounded weighted
coefficients by the joint estimate for the exact pressure.
Consequently a term in the representation of \(\mathcal R_J\)
has either four field insertions and coefficient bound \(Cg\)
or six field insertions and coefficient bound \(C\), summed
over the rooted tensors. All inserted field indices, including
\(x\), occur in the weighted network.

We next sum square-free coupling coefficients. For
\(\varnothing\ne F\subset T\), a microscopic monomial contributing
to \(\theta_F=\prod_{e\in F}\theta_e\) uses one bond above each
edge in \(F\). Choose any edge \(e_0=\{B,D\}\in F\). Every inserted
field site \(y\) has distance at most \(\mathcal L\) from \(B\).
Thus \(|\zeta_y|\le C_k r W_{e_0}e^{\eta\mathcal L/R}\).
For a square-free coefficient, take the reduced radius
\(\rho_e=aW_e^{-7}\). Relative to the analytic radii this
pays all four or six factors \(W_{e_0}\), leaving at least half
the network weight available for a bound on \(|F|\).

A network of length \(\mathcal L\) intersects at most
\(C_d(1+\mathcal L/b_0)\) cells containing its marked sites.
In fact, a traversal of a spanning tree of the network,
subdivided into segments of length at most \(b_0\), covers the
marked sites by \(O(1+\mathcal L/b_0)\) balls of radius \(b_0\),
each meeting a bounded number of cells. Since \(F\) is a forest,
its edges meet at least \(|F|+1\) distinct cells. There are
fixed \(m_0,c_0>0\) such that \(\mathcal L\ge c_0 b_0|F|\)
when \(|F|\ge m_0\).

Let \(t_r=\nu_r^{1/(2m_0)}\). The last half of the network
weight pays \(t_r^{-|F|}\) for \(|F|\ge m_0\) on taking
\(M_R\) large, by the assumed lower power bound on \(\nu_r\).
For \(|F|<m_0\) this cost is at most \(\nu_r^{-1/2}\).
Summing the weighted rooted coefficient bounds over \(x\in U\)
therefore gives
\begin{equation}
 \sum_{\varnothing\ne F\subset T}
  \left(\prod_{e\in F}\rho_e\right)t_r^{-|F|}
          |[\theta_F]\mathcal R_J(\zeta)|
       \le C |C|\sqrt{\nu_r}.
\end{equation}
Without \(t_r^{-|F|}\) the bound is \(C|C|\nu_r\).
Both bounds are for absolute sums of coefficients, so they
include the sum over different microscopic representatives
of each cell edge.

In calculating \([\theta_T]\exp(-\mathcal U_J)\) one can work
in the algebra \(\theta_e^2=0\). The quartic \(V_4\) is independent
of these variables. The coefficient norm of the nonconstant
exponential is bounded by the exponential of the norm of the
nonconstant log series minus one. The displayed estimate and
the damping of the constant term, which is the pressure at
\(\theta=0\), bound the desired coefficient by
\begin{equation}
 e^{-c\sum_U\Psi(u_x)+C|C|}
    \left(\prod_{e\in T}\rho_e^{-1}\right)
    t_r^{|T|}\bigl(e^{C|C|\sqrt{\nu_r}}-1\bigr).
\end{equation}
The complex field shift costs only \(C|C|\), since
\(b_0^d r^4\le C\). When \(|T|\ge m_0\), the first \(m_0\)
powers of \(t_r\) supply another factor \(\sqrt{\nu_r}\).
For \(|T|<m_0\) use the unweighted estimate instead.
The remaining powers of \(t_r\) give any prescribed exponential
margin by taking \(R\) large. This pays both the fixed exponential
costs and the number of spanning trees of the bounded-degree
cell graph, at most a fixed exponential in \(|C|\).

Finally the penalty after removing \(\prod\rho_e\) is at most
\(a^{-|T|}\prod_e W_e^7\). Each cell has bounded degree and
\begin{equation}
 W_B\le 1+\sum_D e^{-\eta\operatorname{dist}(B,D)/R}
                       \sup_{y\in D}|u_y|/r .
\end{equation}
The scalar matrix on the right has uniformly bounded row and
column sums. Thus \(\log(1+t)\le \delta t^2+C_\delta\) and
the \(\ell^2\) bound on that matrix absorb this penalty in
\(\epsilon\sum_B\sup_B|u|^2/r^2+C_\epsilon|C|\).
The constant per cell is paid from the exponential size margin.
All estimates were uniform for \(\|\xi\|_\infty\le C_k r\);
Cauchy's formula in the fixed tested field directions proves the
jet version. Integration over the BKAR parameters preserves the
bound.
\end{proof}

The coefficient \(\epsilon\) in Proposition~\ref{prop:pressure-trees}
can be chosen small enough to absorb the quadratic loss into the
regulator for a packing of pressure atoms.
Indeed the absolute rows of the filters give a bounded
\(\ell^2\) matrix on cell suprema, including their fixed collars.
The normalized square of a cell supremum is bounded by the sum
of its normalized \(L^2\) square and its quotient square in the
large-field recovery estimate. At the small radius one uses the
spare growing regulator. Hence a product of disjoint pressure
atoms consumes only a fixed spare fraction of those two reserves.
Together with Lemma~\ref{lem:pressure-singleton} this supplies
the short-reference pressure input to the packing argument.
Proposition~\ref{prop:cube-joint-initialization} combines these
pressure estimates with the full filter and the quadratic tails.

For thermal matching, fix the width \(A_\beta\) of the window
\(|\beta-1|\le A_\beta g\). The part of the quartic defect coming
from \(L_J-(I-J)\) is \(O(A_\beta g)\). At both initial radii
\(g/r^2\to0\), so for sufficiently large \(R\) depending on the
fixed \(A_\beta\) its contribution is absorbed by the sextic
factor \(b_0^d r^6\) with a constant independent of \(A_\beta\).
The row gap and constants in the cavity and damping estimates
can be fixed in a small temperature neighbourhood before
choosing this width.

\subsection{Spectral and covariance estimates}\label{sec:ext-cov}
We turn to the Gaussian reference on the cube. The principal
eigenvalue determines how close the reference can be taken to its
pole. After conditioning on the principal coordinate we need estimates
on the remaining covariance that retain its decay near each face.
For the cube walk of Section~\ref{sec:model}, write \(s=2\alpha=d/2\),
\(\varrho=R/n\), and \(v=n^d\). The physical observation \(\mathcal O\)
acts on the full auxiliary lattice, while the physical block of
\(\mathcal Q_c\) is restricted to the cube. Below, \(a_R=\lambda R^s\).
For \(d<4\) use the beta-integral split and for \(d=4\) use the
finite-range split of Lemma~\ref{lem:local-split}, with auxiliary diagonal
\(\mathcal D=I\). Denote the observed bulk auxiliary resolvent by
\(A_r=\mathcal L^*(\mathcal D+r)^{-1}\mathcal L\) and put
\(H_{r,\Lambda}=(B_\Lambda/c)+rI-A_{r,\Lambda}\).

\subsubsection{Generalized principal eigenvalue}
The first two estimates locate the pole and separate it from the
next mode. The last two control the principal vector and its overlap
with a uniform magnetic field.

\begin{proposition}[Principal spectral gap and magnetic overlap]
\label{prop:cube-principal-mode}
Let \(\eta_1\ge\eta_2\) be the two largest eigenvalues of
\(E=B_\Lambda^{-1/2}A_\Lambda B_\Lambda^{-1/2}\) and choose
the sign of its unit principal eigenvector \(e\) so that
\(B_\Lambda^{-1/2}e>0\). For \(n/R\) sufficiently large,
\begin{equation}
 \begin{aligned}
 c_{*,\Lambda}-1&\asymp\varrho^s, &
 \eta_2^{-1}-\eta_1^{-1}&\asymp\varrho^s,\\
 \|B_\Lambda^{-1/2}e\|_\infty&\lesssim v^{-1/2}, &
 \sum_{x\in\Lambda}(B_\Lambda^{1/2}e)_x&\asymp v^{1/2}.
 \end{aligned}
\end{equation}
Moreover, for \(1\le c<c_{*,\Lambda}=\eta_1^{-1}\),
the covariance conditioned on \(\ell=B_\Lambda^{-1/2}e\) satisfies
\(\mathcal S_\ell(V_c)\le C V_1\) and
\(\|V_1\|_{2\to2}\le C\varrho^{-s}\).
\end{proposition}

\begin{proof}
We have \(B-A=\zeta (I-D_R)^\alpha\) for a uniformly bounded positive
normalization \(\zeta\). For the cube walk
\(1-\widehat D_R(k)\asymp\min\{R^2|k|^2,1\}\), by the Taylor expansion
on \(|Rk|\le\kappa\), the product Dirichlet-sum formula away from
zero, and the strict bound \(|\widehat D_R(k)|<1\) on the intermediate
compact scaled annulus. The latter bound is uniform for large \(R\)
by the Riemann-sum limit, and for the finitely many smaller \(R\ge2\)
by the aperiodicity of the support. For a unit vector supported in the
cube, \(|\widehat u|\le v^{1/2}\) shows that at least half its Fourier
mass lies outside \(|k|\le a/n\) for a small fixed \(a\). The resulting
energy lower bound on zero extensions is \(c\varrho^s\). For an
upper bound on the span of two fixed smooth test bumps at scale \(n\),
\(\|\tau_z u-u\|_2\le C |z|\|u\|_2/n\) for \(|z|\le C R\),
so Jensen's inequality for the spectral measure and \(0<\alpha\le1\)
give \(C\varrho^s\). Since \(B_\Lambda\) has uniformly bounded
condition number, min--max gives \(1-\eta_i\asymp\varrho^s\) for
\(i=1,2\).

For \(\alpha<1\), the jump kernel also satisfies
\begin{equation}
 A(x)\ge c R^s(R+|x|)^{-d-s}.
\end{equation}
To verify this, Fourier inversion gives
\(D_R^{*m}(x)\ge c R^{-d}m^{-d/2}\) for \(m\ge m_0\) and
\(|x|\le a_0R\sqrt m\). On \(|k|\le M/(R\sqrt m)\), take \(a_0\)
small enough that the phase has cosine at least \(1/2\); the
ball of radius \(1/(R\sqrt m)\) gives the stated lower bound.
The region up to \(\kappa/R\) outside the first ball of radius
\(M/(R\sqrt m)\) costs at most
\(C R^{-d}m^{-d/2}\int_{|y|>M}e^{-c|y|^2}\,dy\).
The remaining integral is at most
\(C R^{-d}q_\kappa^{m-2}\), using
\(\int|\widehat D_R(k)|^2\,dk=O(R^{-d})\).
Sum this lower bound over \(m\asymp 1+(|x|/R)^2\) with
binomial weights comparable to \(m^{-1-\alpha}\) above the
smoothing order.

Set \(P=B_\Lambda^{-1}A_\Lambda\). The Neumann series dominates
\(P\) entrywise by the restriction of the stochastic bulk convolution
\(P_\infty=B^{-1}A\). In the fractional case \(P\) has all
entries positive; in the local case it is irreducible because
\(Q\) includes the nearest-neighbour walk edges and \(A\) includes
a strictly positive diagonal. Thus its positive eigenvector
\(u=B_\Lambda^{-1/2}e\) satisfies \(\|u\|_2\asymp1\).
For \(m=\lceil\varrho^{-s}\rceil\), Fourier inversion gives
\begin{equation}
 \|P_\infty^m(0,\cdot)\|_2^2\le C R^{-d}
                   (m^{-d/s}+e^{-a m}).
\end{equation}
On \(|Rk|\le\kappa\) the multiplier is bounded by
\(\exp[-c(R|k|)^s]\). On the complement it is bounded strictly
below one, with square integral \(O(R^{-d})\) by the fixed smoothing
power of the cube walk (or of the smaller cube walk in the local
split) and the exponentially small white term. Therefore
\(\|u\|_\infty\le C\|P_\infty^m(0,\cdot)\|_2\le C v^{-1/2}\),
since \(\eta_1^{-m}\le C\). Positivity gives
\(\sum u\ge \|u\|_2^2/\|u\|_\infty\gtrsim v^{1/2}\); Cauchy--Schwarz
gives the upper bound. The finite-volume rows of \(B\) have sum in
\([p,1]\) with \(p>0\), proving the magnetic overlap for \(Bu\).
For \(\alpha<1\), \(P(x,y)\ge c\varrho^s/v\); hence the stochastic
ground-state transform has rows bounded below by the common mass
\(c\varrho^s u(y)/\sum u\). Oscillation contraction proves
\(\eta_1-\eta_2\gtrsim\varrho^s\).

For the local case the gap follows from the first two Dirichlet
min--max levels of the rescaled walk form. Set
\(\sigma_R^2=\sum_z z_1^2 D_R(z)\asymp R^2\) and divide the walk
form by \(\sigma_R^2/(2n^2)\). Here is a compactness argument
valid along every sequence \(R/n\to0\). Unit vectors of bounded
rescaled energy satisfy
\(\|\tau_z u-u\|_2^2\le C[(|z|/n)^2+(R/n)^2]\) by splitting Fourier
space at \(|k|=\kappa/R\). Embed them as piecewise-constant functions
with mesh \(1/n\) on the zero-extended unit cube. A fractional-cell
shift is bounded by the sum of the adjacent integer-shift squared
norms, so the translation estimate gives \(L^2\) compactness.
On every fixed ball of rescaled Fourier frequencies the
normalized symbol converges uniformly to \(|\xi|^2\), because the
jumps have length \(O(R)\). The Fourier liminf gives the Dirichlet
energy of the limit, which is in \(H^1_0((0,1)^d)\) by zero extension.
The denominator satisfies
\(\langle u,B_\Lambda u\rangle=p\|u\|_2^2+o(1)\); each shift in the
finite-range kernel contributes \(O(R/n)\) in norm by the same
translation estimate. Conversely, samples of smooth compactly
supported functions give the continuum energy by Taylor expansion,
uniformly for \(R/n\to0\). Applying these bounds to the two min--max
spaces (whose orthogonality passes to the limit by the denominator
estimate) gives
\begin{equation}
 \frac{1-\eta_i}{\zeta\sigma_R^2/(2pn^2)}
     \longrightarrow \lambda_i(-\Delta_{(0,1)^d}^{\rm Dir}),
     \qquad i=1,2.
\end{equation}
These Dirichlet levels are distinct. The resulting gap is uniform
when \(n/R\) is sufficiently large.

The identity
\(V_c=B_\Lambda^{1/2}(E^{-1}-cI)^{-1}B_\Lambda^{1/2}\) shows
that conditioning on \(\ell=B_\Lambda^{-1/2}e\) deletes the
principal summand in \(V_c\). Writing \(\mathcal S_\ell\) for
the covariance after conditioning on this coordinate, we obtain
\(\mathcal S_\ell(V_c)\le C V_1\) and
\(\|V_1\|_{2\to2}\le C\varrho^{-s}\) for \(1\le c<c_{*,\Lambda}\).
\end{proof}

\subsubsection{Finite propagation and terminal conditioning}
\label{sec:FinitePropagationTerminalConditioning}
Let \(q_i\) be the positive scalar polynomials of
\cite[Lemma~2.1]{IsingCorrelation}, rescaled to a common spectral
interval \([0,S]\) for the auxiliary precisions. Put \(r_i=L^{-2i}\).
Choose the terminal cut \(r=L^{-2J}\asymp M_*(R/n)^2\) with
\(M_*\) a sufficiently large fixed constant. Put
\begin{equation}
 C_i=a_R\mathcal O q_i(\mathcal Q_c)\mathcal O^*,\qquad
 C_<=\sum_{i<J} C_i,\qquad C_T=a_R V_c-C_< .
\end{equation}
Let \(X_<,X_T\) be independent centred Gaussian fields with covariances \(C_<,C_T\),
respectively, and write \(a=\ell^{\mathsf T}(X_<+X_T)\).

\begin{lemma}[Terminal Gaussian comparison]
\label{lem:terminal-gaussian-comparison}
For \(1\le c<c_{*,\Lambda}\), the slices sum to \(a_RV_c\) and
\(C_i\) has range \(O_L(RL^i)\). For each fixed integer \(m\ge1\),
\begin{equation}
 \mathcal S_\ell(C_T)\le\operatorname{Cov}(X_T\mid a=0)
 \le C_m\mathcal C_r^{(m)},\qquad
 \mathcal C_r^{(m)}=
 a_R r^{2m-1}\mathcal O(\mathcal Q_c+rI)^{-2m}\mathcal O^* .
\end{equation}
The preterminal variance satisfies
\(\ell^{\mathsf T}C_<\ell\asymp\lambda n^s\). The comparison
constants are uniform as \(c\uparrow c_{*,\Lambda}\).
\end{lemma}

\begin{proof}
Lemma~\ref{lem:scalar-polynomial-slices} gives the sum of the slices.
Each application of \(\mathcal Q_c\) propagates at most \(CR\) and
\(\deg q_i\le C_LL^i\), which gives the range bound. The covariance
comparison is Corollary~\ref{cor:terminal-resolvent-majorant},
applied to the observed field in the canonical normalization.

By \eqref{eq:gaussian-partial-scalar}, \(C_<\asymp a_R V_r\), with
\(V_r=\mathcal O(\mathcal Q_c+rI)^{-1}\mathcal O^*\).
Write \(Y_r=B_\Lambda/c+rI\). The Schur formulas give
\begin{equation}
 V_r=A_{r,\Lambda}+A_{r,\Lambda}H_{r,\Lambda}^{-1}A_{r,\Lambda},
 \qquad V_r^{-1}=A_{r,\Lambda}^{-1}-Y_r^{-1}.
\end{equation}
The bulk lower bound \(H_r(k)\ge c_1(r+\omega(k))^\alpha\), derived
below, survives compression when \(M_*\) is sufficiently large. Since
\(A_r\) is bounded, \(\|V_r\|\le C r^{-\alpha}\). For the reverse
estimate, the beta integral split at \(1-u=r\) gives
\(0\le A-A_r\le Cr^\alpha A\) as bulk forms and after compression.
Indeed \(1-u\widehat D_R\ge1-u\), and the integrable density in the
difference is bounded by
\(C(1-u)^{\alpha-1}\min\{1,r/(1-u)\}\) times
\(\widehat D_R^{N_{\rm sm}}\); this multiplier also occurs with a
uniformly positive coefficient in \(A\). The white channel has
relative change \(O(r)\). In the local model \(A_r=A/(1+r)\).
Inverting the comparison and using the uniform gap of \(Y_0\) gives
\begin{equation}
 0\le V_r^{-1}-V_c^{-1}\le
 Cr^\alpha A_\Lambda^{-1}+Cr I .
\end{equation}
Put \(h=B_\Lambda^{1/2}e\). Then \(\ell^{\mathsf T}h=1\),
\(h^{\mathsf T}A_\Lambda^{-1}h=\eta_1^{-1}\), and
\(h^{\mathsf T}V_c^{-1}h=\eta_1^{-1}-c=O(r^\alpha)\).
Cauchy--Schwarz for \(V_r^{1/2}\ell\) and \(V_r^{-1/2}h\)
yields \(\ell^{\mathsf T}V_r\ell\ge c r^{-\alpha}\).
Since \(a_R r^{-\alpha}\asymp\lambda n^s\), we obtain
\(\ell^{\mathsf T}C_<\ell\asymp\lambda n^s\).
\end{proof}

\subsubsection{Boundary diagonal of the higher resolvent}
The preceding comparison reduces the boundary estimate for the
conditioned Gaussian field to a diagonal bound on the higher
resolvent. For a point of depth \(\delta(x)=\operatorname{dist}(x,\Lambda^c)\)
put \(b=Rr^{-1/2}\), \(\rho_b=\min\{b,R+\delta\}\), and
\(w(x)=(\rho_b(x)/b)^{p_0}\).

\begin{proposition}[Boundary diagonal of the higher resolvent]
\label{prop:boundary-resolvent-diagonal}
Fix \(0<p_0<\alpha\). For \(0<r\le1\),
\(|c-1|\le c_0r^\alpha\) with \(c_0\) sufficiently small, and
all sufficiently large fixed \(m\), we have, uniformly in \(R,n\),
\begin{equation}
 \mathcal C_r^{(m)}(x,x)\le C_m\lambda b^{-s}w(x)^2.
\end{equation}
The estimate holds with a constant depending on \(C\) if
\(|c-1|\le C r^\alpha\) and the finite-volume precision is positive.
In the microscopic collar the variance of a fixed \(k\)-difference
stencil has the additional factor \(C_k R^{-2k}\), for sufficiently
large fixed smoothing order.
\end{proposition}

\begin{proof}
For \(|c-1|\le C r^\alpha\), take fixed \(\kappa\ge1\) such that
\(C\le c_0\kappa^\alpha\). On the positive finite-volume spectrum,
\((t+r)^{-2m}\le \kappa^{2m}(t+\kappa r)^{-2m}\), and the scales and
weights for \(r\) and \(\kappa r\) differ by fixed factors.

We first bound the row sum of the physical Schur inverse, then use
this bound on each of the two outer resolvent factors. For the bulk Schur multiplier,
\begin{equation}
 H_r(k)=\zeta\omega(k)^\alpha +(c^{-1}-1)B(k)+r+A(k)-A_r(k).
\end{equation}
If \(\omega\ge r\), the first term absorbs the displacement
\(|c-1|\le c_0 r^\alpha\). When \(\omega\le r\ll1\), restrict the
positive beta integral for \(A-A_r\) to \(1-u\in[r,2r]\).
There \(1-u\widehat D_R\asymp r\), the difference of the two inverse
denominators has order \(r^{-1}\), and the density has order
\(r^\alpha\). Thus \(A-A_r\ge c_1 r^\alpha\) in the fractional
case. For \(\alpha=1\) the term \(r\) suffices. When \(r\) is bounded
below the same term absorbs the displacement by reducing \(c_0\).
Together with the upper bound on \(A-A_r\) above this proves
\(H_r(k)\asymp(r+\omega(k))^\alpha\).

We shall also use the inverse bound for the full bulk precision.
Put \(h=r+\omega(k)\) and \(D_r(k)=\mathcal D(k)+rI\). For the
fractional split, \(D_r(k)\ge c h I\), and the squared norm of
\(D_r(k)^{-1}\mathcal L(k)\) is bounded by
\begin{equation}
 C\left(1+\int_0^1\frac{t^\alpha}{(h+t)^2}\,dt\right)
 \le C' h^{\alpha-1}.
\end{equation}
Indeed \(t=1-u\), \(1-u\widehat D_R(k)+r\ge c(t+h)\), and the
remaining factors in the density and observation are bounded.
For the local split the same squared norm is bounded uniformly.
Block inversion, using the Schur bound just proved, now gives
\begin{equation}
 \|(\mathcal Q_c(k)+rI)^{-1}\|
 \le \|D_r(k)^{-1}\|+
       H_r(k)^{-1}(1+\|D_r(k)^{-1}\mathcal L(k)\|^2)
 \le C h^{-1}.
\end{equation}

The entries of
\(H_{r,\Lambda}^{-1}\) are nonnegative. Moreover
\begin{equation}
 (H_{r,\Lambda}^{-1}\mathbf1)(x)\le C r^{-\alpha}w(x).
\end{equation}
To see this, consider each face separately. In the fractional
case any fixed \(0<p_0<\alpha\) is available. Use
\(\min\{1,((C_1 R+x_1)/b)^{p_0}\}\) for \(x_1>0\),
zero outside, where \(C_1\) is fixed and sufficiently large.
We verify the required generator sign. For the one-coordinate
projected walk write \(S_m\) for its \(m\)-step sum and
\(\sigma_R^2=\operatorname{Var}S_1\asymp R^2\).
For \(y/R\to\infty\), the binomial coefficients \(c_m\asymp
m^{-1-\alpha}\), with \(c_m m^{1+\alpha}\) converging to a
positive limit, give
\begin{equation}
 \left(\frac y{\sigma_R}\right)^{2\alpha}
 \sum_{m\ge1} c_m\,\mathbb E[1-(1+S_m/y)_+^{p_0}]
 \ \longrightarrow\
 c_\alpha\int_0^\infty t^{-1-\alpha}
                \mathbb E[1-(1+\sqrt t Z)_+^{p_0}]\,dt,
 \qquad Z\sim N(0,1).
\end{equation}
The convergence is uniform along such sequences of \(y,R\).
In fact pair opposite steps in the expectation. The paired test
is \(O(|S_m/y|^2)\) at small displacement and
\(O(1+|S_m/y|^{p_0})\) at large displacement. Thus the part
\(m\le\epsilon(y/\sigma_R)^2\) contributes \(O(\epsilon^{1-\alpha})\)
after scaling, and the part \(m\ge M(y/\sigma_R)^2\) contributes
\(O(M^{-\alpha+p_0/2})\). In the middle range the bounded
standardized third moments give a uniform central limit theorem,
uniform integrability follows from the second moments, and the
rescaled sum is a Riemann sum. Integrating the Gaussian density
in \(t\) gives a positive constant times
\begin{equation}
 \int_0^\infty
      \bigl(2-(1+u)^{p_0}-(1-u)_+^{p_0}\bigr)u^{-1-2\alpha}\,du .
\end{equation}
This integral vanishes at \(p_0=\alpha\): integration by
parts leaves a constant times the difference of the integrals
of \((1-u)^{\alpha-1}u^{-2\alpha}\) on \((0,1)\) and
\((1+u)^{\alpha-1}u^{-2\alpha}\) on \((0,\infty)\).
The substitution \(v=u/(1-u)\) identifies them; when
\(\alpha\ge1/2\), the difference of lower cutoffs contributes
\(O(\epsilon^{2-2\alpha})\) and vanishes in the paired integral.
For \(p_0<\alpha\), strict concavity of the power \(p_0/\alpha\)
bounds its generator from below by the tangent-line multiple
of the generator at exponent \(\alpha\), with strictly
positive integral. Taking \(C_1\) large therefore gives
\((I-D_R)^\alpha y_+^{p_0}\ge c(R/y)^s y^{p_0}\) for \(y\ge C_1 R\).
The cutoff at height one and killing on \(x_1\le0\) only
increase this bound at points where the candidate is less
than one. By the identity
\begin{equation}
 H_r=c^{-1}\zeta(I-D_R)^\alpha+(c^{-1}-1)A+rI+(A-A_r)
\end{equation}
in bulk, the possibly adverse correction is at most
\(C|c-1|\) times the candidate, using the projected jump
\(p_0\)-moment bound for \(A\).
For \(C_1R+x_1\le b\), this is absorbed by the positive
term by choosing \(c_0\) small. At \(C_1R+x_1>b\) the
candidate is one. This proves the required bound by
\(cr^\alpha\); the harmless shift \(C_1R\) changes the
claimed weight by a fixed factor.

For \(d=4\), use \(R+x_1\) in the barrier and any fixed
\(0<p_0<1\). There is a uniform positive component of the
projected step kernel comparable to the projected \(D_R\)
away from its diagonal, with projected step range at most
\(R\). Near \(x_1=0\), say \(x_1<c_1R\) with fixed small
\(c_1\), a fixed step slab exits and gives the lower bound
by a loss of one term. All crossing pairs of opposite
jumps have nonnegative contributions, since the ratio of
the largest forward barrier to its value at \(x_1\) is at
most \(2^{p_0}<2\). At \(x_1\ge c_1 R\) and
\(y=R+x_1<b\), use a smaller fixed slab of projected
increments of size comparable to \(R\) that does not exit.
The concave second difference of the barrier on this slab
is at least \(c(p_0)(R/y)^2(y/b)^{p_0}\) in magnitude.
Truncation at height one only improves the lower bound.
This again yields \(cr^\alpha\) with \(\alpha=1\).
The bulk mass gives the
lower bound when the barrier is one. Minimum over all faces
preserves supersolutions by positivity of the jumps. This
proves the row estimate. The weight is preserved by convolution
with \(A_r\) by its uniform moment in \((1+|z|/R)^{p_0}\).

To reach the full diagonal bound we use a high resolvent power
rather than the first derivative of the observed resolvent.
Set \(v_x=\mathcal O^*\delta_x\) and
\(F_x=(\mathcal Q_c+r)^{-1}v_x\). These are nonnegative
in the auxiliary and physical variables. The physical component
of \(F_x\) is \(H_{r,\Lambda}^{-1}A_{r,\Lambda}\delta_x\) and
has spatial mass \(C r^{-\alpha}w(x)\). In the auxiliary
channel the component is \((\mathcal D+r)^{-1}(v_x+
\mathcal L F_{x,\mathrm{phys}})\). The \(L^2\) norm over the
channel variable of its spatial sum is at most
\(C r^{-(1+\alpha)/2}w(x)\), using
\(\int_0^1(1-u)^\alpha(r+1-u)^{-2}\,du\le Cr^{\alpha-1}\).
The direct \(v_x\) term fits since \(1\le r^{-\alpha}w(x)\).
For the local split the same bound follows from \(\mathcal D=I\)
with \(\alpha=1\). After extending the physical component by
zero, we thus have
\(\sup_k\|\widehat F_x(k)\|\le Cr^{-(1+\alpha)/2}w(x)\).

The resolvent powers in the finite domain are bounded
entrywise by the bulk powers. This follows by expanding the
positive-path semigroup of a sufficiently large constant minus
\(\mathcal Q_c\) and integrating in time with \(e^{-rt}\).
Consequently \(\langle F_x,(\mathcal Q_c+r)^{-(2m-2)}F_x\rangle\)
is bounded by the quadratic form with the bulk multiplier.
The low-frequency operator norm is
\(C_m(r+R^2|k|^2)^{-(2m-2)}\) by the Schur spectral bound.
Integration there, and the two factors of \(\widehat F_x\),
give \(C R^{-d}r^{d/2-2m+2}r^{-1-\alpha}w(x)^2\).
On high frequencies use the bounded multiplier and
\(\|F_x\|^2\le C R^{-d}r^{-1-\alpha}\).
The last estimate follows by entrywise bulk comparison on
the two resolvent factors. At low frequency use
\(\langle v(k),(\mathcal Q_c(k)+r)^{-2}v(k)\rangle
 \le Cr^{-1-\alpha}\), by the spectral and Schur bounds,
and integrate over a ball of volume \(O(R^{-d})\).
At high frequency the spectral gap is bounded below and
\(\|v(k)\|^2\) has integral \(O(R^{-d})\) by smoothing.
The factor \(r^{2m-1}\) absorbs this crude high-frequency
estimate for large \(m\), since \(w^2\ge r^{p_0}\).
The resulting scale is \(a_R R^{-d}r^{d/2-\alpha}
=\lambda b^{-s}\), as asserted.

Taking the absolute value at each differentiated observation leg gives
\(C_k\lambda b^{-s}R^{-2k}w(x)^2\) for the variance of a
\(k\)-difference stencil in the microscopic collar.
For the walk part, move each difference onto a separate step
in \(D_R^{*N/2}\). Here \(N\) denotes the fixed smoothing order \(N_{\rm sm}\); in the local
split the observation leg has the corresponding convolution of \(E_R\).
The absolute marked step has mass
\(O(R^{-1})\) and range \(O(R)\), and the unmarked steps retain
an arbitrarily large fixed smoothing power. The absolute white
term has mass at most \(C_k\sqrt{\varepsilon_R}\) and support
in a fixed stencil. After taking \(R^{-k}\) out of each leg
the spatial sum, moment, and high-frequency integral bounds
in the proof are unchanged.
\end{proof}

\subsubsection{Intermediate annuli for fractional powers}
For fractional powers the field norm also tests differences on
regions where \(R\ll\delta(x)\lesssim b\). The next estimate
supplies the derivatives used in that norm.

\begin{lemma}[Annular stencils for fractional powers]
\label{lem:fractional-annular-stencils}
Let \(d<4\), \(0<r\le1\), \(|c-1|\le c_0r^\alpha\), and
\(X\sim N(0,\mathcal C_r^{(m)})\). Given a fixed derivative order
\(P\) and \(0<p<\alpha\), choose the smoothing and resolvent orders
sufficiently large. For stencils in the cube,
\begin{equation}
 \|\rho_b(x)^{|\gamma|}\nabla^\gamma X(x)\|_{L^2}
 \le C(\lambda b^{-s})^{1/2}(\rho_b(x)/b)^p,
 \qquad |\gamma|\le P+1 .
\end{equation}
For each scale-\(b\) cell \(B\), the corresponding supremum has
a uniform Gaussian tail. More precisely, for \(t\ge1\),
\begin{equation}
 \mathbb P\!\left(\sup_{\substack{x\in B\,;\ |\gamma|\le P}}
 \frac{\rho_b(x)^{|\gamma|}|\nabla^\gamma X(x)|}
 {(\lambda b^{-s})^{1/2}(\rho_b(x)/b)^p}>t\right)
 \le C e^{-ct^2}.
\end{equation}
\end{lemma}

\begin{proof}
We first prove
\begin{equation}
 \operatorname{Var}(\nabla^\gamma X(x))
 \le C_\gamma\lambda b^{-s}\rho_b(x)^{-2|\gamma|}
\end{equation}
for \(|\gamma|>s\). We use interior bounds for bulk Green functions
and then apply them to the difference between bulk and cube solutions.
On a fixed complex disc \(|z-r|\le\epsilon r\)
write the bulk Green kernel \(G_z=H_z^{-1}\). The estimates
from the beta integral and the Fourier symbol bound give
for \(|y|\ge C_k R\)
\begin{equation}
 |\nabla^k G_z(y)|\le C_k R^{-s}|y|^{s-d-k},
 \qquad |\nabla^k A_z(y)|\le C_k R^s |y|^{-d-s-k}.
\end{equation}
Indeed the derivatives in \(\theta=Rk\) of the low-frequency
symbols on a window \(T\ge\sqrt r\) are bounded by
\(C_jT^{-s-j}\) and \(C_j T^{s-j}\), respectively, for
\(G_z\) and \(A_z-A_z(0)\) (the value subtracted here is
at zero Fourier frequency). This follows by differentiating
the beta integral in \(\widehat D_R\). Low-frequency window
integrals are \(C R^{-d-k} T^{d+k\mp s}
(1+T|y|/R)^{-J_1}\) for any fixed \(J_1\), by integration
by parts. At high frequencies subtract a finite-range
partial Neumann sum of the short kernel from \(G_z\).
The remainders and their finitely many derivatives have
rapid fixed decay by the cube Dirichlet-sum bounds and
arbitrarily large fixed smoothing order. The zero-exclusion term is \(O(R^{-d})\), and hence
\(O((1+|\theta|)^{-1})\) on the scaled torus; the same
single inverse power bounds all fixed derivatives by
the product Dirichlet-sum formula. Taking a sufficiently
large fixed smoothing power supplies the required
integrability of the differentiated remainders. Summing
the windows gives the two kernel bounds.

Put \(d_x=\min\{b,\delta(x)\}\gg R\). For a bounded
\(H_z\)-harmonic function \(u\) on the cube, take a cutoff
\(\chi=1\) near \(x\) with support in a ball of radius
\(d_x/2\). Set \(h=G_z*\nabla^{\gamma,*}\delta_x\).
The identity
\(\nabla^\gamma u(x)=\langle H_z((1-\chi)h),u\rangle\)
uses the bilinear pairing. The tail bounds give
\begin{equation}
 \|(1-\chi)h\|_1\le C R^{-s}d_x^{s-|\gamma|},\qquad
 \|\nabla((1-\chi)h)\|_1\le C R^{-s}d_x^{s-|\gamma|-1}
\end{equation}
when \(|\gamma|>s\). The kernel tail gives a jump bound
\(\sum_y |K_z(y)|\min(1,|y|/d_x)\le C(R/d_x)^s\)
for \(s<1\); for \(s\ge1\) use the centered second difference
instead. The same tail computation gives
\(\sum_y |K_z(y)|\min(1,|y|^2/d_x^2)\le C(R/d_x)^s\)
and the cutoff kernel has second difference norm
\(C R^{-s}d_x^{s-|\gamma|-2}\).
The mass contribution is bounded by using
\(|H_z(0)|\lesssim r^\alpha\), with zero denoting Fourier
frequency. Thus
\(|\nabla^\gamma u(x)|\le C d_x^{-|\gamma|}\|u\|_\infty\).
The same estimate holds for \(A_z*u\) by using this bound
on the near part of the convolution and the differentiated
\(A_z\) tail on the far part.

For a unit source in the auxiliary-plus-physical Hilbert
space compare the bulk and cube solutions, extending the
physical source by zero. Their physical difference is
\(H_z\)-harmonic with exterior data from the bulk solution.
Its supremum is bounded by the exterior supremum by the
substochastic exit kernel for the real part of \(z\).
The exterior physical-source row omits the onsite term and
has squared \(L^2\) norm \(C R^{-d}(1+\log(1/r))\).
The auxiliary-source row has squared norm
\(C R^{-d}r^{\alpha-1}\), by integrating the low-frequency
majorant \((r+R^2|k|^2)^{-1-\alpha}\). On the high frequencies
both row bounds follow from the smoothing estimates.
Specifically the physical-source bound uses the Neumann
difference \(G_z-(c^{-1}+z)^{-1}I\), whose high-frequency
square is integrable at order \(R^{-d}\) by Plancherel for
the short kernel. The auxiliary-source bound uses the
squared norm of \(\mathcal L(k)\) on high frequencies
and the auxiliary beta density on low frequencies.
Since \(\alpha<1\), the logarithm is dominated by
\(Cr^{\alpha-1}\). The observed difference of the solutions
is \(A_z*u\). Cauchy differentiation on the disc for the
\((m-1)\)-st derivative multiplies its bound by
\(C_m r^{1-m}\). Squaring the resulting dual norm and
multiplying by \(a_R r^{2m-1}\) gives
\(C\lambda b^{-s}d_x^{-2|\gamma|}\).
The bulk contribution has the same bound with \(d_x=b\)
by Fourier integration of the bulk higher resolvent, so
the triangle inequality completes the estimate.

To interpolate and take suprema, on a cube of diameter
\(d_x\) use the Hilbert space of Gaussian random variables,
the improved diagonal bound, and stencils through order
\(k_0>\max\{P+1,s\}\) (as large as needed). In the tensor-product discrete Taylor
formula use mesh \(\max\{1,\lfloor c d_x (d_x/b)^{p_0/k_0}\rfloor\}\),
with stencils directed into the cube. The fixed-order
Vandermonde inverse and Taylor remainder bound give
\begin{equation}
 \|\rho_b(x)^{|\gamma|}\nabla^\gamma X(x)\|_{L^2}
 \le C(\lambda b^{-s})^{1/2}(\rho_b(x)/b)^{p_1},
 \qquad |\gamma|\le P+1,\quad
 p_1=p_0(1-(P+1)/k_0).
\end{equation}
Here we choose \(p<p_0<\alpha\) and then \(k_0\) large enough that
\(p<p_1\).
If the real-valued optimal mesh is less than a fixed constant,
use the immediate finite differences of the diagonal estimate.
For each derivative and each cube of diameter comparable to
the shell depth, the extra derivative gives a Lipschitz
\(L^2\) metric on the normalized stencil field, with scale
\((\lambda b^{-s})^{1/2}(\rho_b/b)^{p_1}\).
Dyadic spatial grids prove a subgaussian supremum bound
with this scale: at level \(j\) there are \(O(2^{dj})\)
increments whose standard deviations are \(O(2^{-j})\)
times the scale; thresholds proportional to \(t2^{-j/2}\)
give a summed tail \(C e^{-c t^2}\) for \(t\ge1\).
A shell of depth \(2^{-i}b\) in a scale-\(b\) cell needs
at most \(C2^{i(d-1)}\) such cubes. For a field weight
\(p<p_1\), the shell tail after normalization is bounded
by \(C2^{i(d-1)}\exp[-ct^2 2^{2i(p_1-p)}]\).
Its sum gives the exponential field-norm moment uniform
in the cube and the scale, including edges and corners.
\end{proof}

\subsubsection{The local annular estimate in dimension four}
\label{subsec:local-annulus}

For the local split of Lemma~\ref{lem:local-split}, the auxiliary
diagonal is the identity and every block has range \(O(R)\).
Consequently the interior equation can be iterated on a ball
contained in the cube. This yields the annular derivative estimate
needed for the four-dimensional field norm.

\begin{lemma}[Local resolvent derivatives]\label{lem:local-annulus}
Take \(d=4\), \(0<r\le1\), \(|c-1|\le c_0r\) for a sufficiently
small fixed \(c_0\), \(b=Rr^{-1/2}\), and
\(\rho_b(x)=\min\{b,R+\operatorname{dist}(x,\Lambda^c)\}\).
For each fixed derivative multiindex \(\gamma\) and all sufficiently
large fixed integers \(m\), the covariance
\begin{equation}
 \mathcal C_r^{(m)}
 =a_Rr^{2m-1}\mathcal O(\mathcal Q_c+rI)^{-2m}\mathcal O^*
\end{equation}
satisfies
\begin{equation}
 \operatorname{Var}(\nabla^\gamma X(x))
 \le C_{m,\gamma}\lambda b^{-2}\rho_b(x)^{-2|\gamma|},
 \qquad X\sim N(0,\mathcal C_r^{(m)}).
\end{equation}
Near the boundary the differences are taken on stencils lying in
the cube. The short smoothing split can be chosen once for any
finite family of derivative orders.
\end{lemma}

The proof is given in Appendix~\ref{app:local-resolvent}.

Proposition~\ref{prop:boundary-resolvent-diagonal} and
Lemma~\ref{lem:local-annulus} give the two inputs to the discrete
Taylor interpolation and shell sum in the proof of
Lemma~\ref{lem:fractional-annular-stencils}. Thus the same weighted
stencil and cell-supremum estimates hold in dimension four for
any fixed \(0<p<1\).

For the merged slice set \(M_R=b_0/R\asymp\log R\). Measured at
scale \(b_0\), a stencil of order at most \(P\) at a constituent
scale \(b\le b_0\) gains at most \((b_0/b)^{P+p}\) from the change
of weights. A merged cell contains \(O((b_0/b)^d)\) cells at that
scale. The subgaussian tail on each cell therefore bounds their
maximum at an additional factor \(C\sqrt{1+\log(b_0/b)}\).
Summing the geometric scales, the squared fluctuation parameter in
the \(b_0^{-s/2}\) field units is at most
\(C\lambda M_R^{2P+2p+s}(1+\log M_R)\), which tends to zero.
For the diagonal alone, summing the variances gives
\(C\lambda R^{-s}=C\lambda M_R^s b_0^{-s}\).

\begin{corollary}[Profile of the terminal principal coordinate]
\label{cor:terminal-principal-profile}
At the prescribed enhancement let \(\sigma_T^2=\ell^{\mathsf T}C_T\ell\)
and \(u_T=C_T\ell/\sigma_T^2\). For \(1\le d\le4\), any fixed
\(0<p<\alpha\), and each of the finitely many tested stencils,
\begin{equation}
 |\nabla^\gamma u_T(x)|\le C v^{-1/2}\rho_n(x)^{-|\gamma|}
                         (\rho_n(x)/n)^p .
\end{equation}
Moreover, as \(n\to\infty\) at fixed \(R\),
\(\|u_T-B_\Lambda^{1/2}e\|_2=o(1)\),
\(\one^{\mathsf T}u_T\asymp v^{1/2}\), and
\(\sum_x u_T(x)^4\asymp v^{-1}\).
\end{corollary}

\begin{proof}
Put \(\sigma^2=\sigma_T^2+\sigma_L^2\). By regression, the covariance
of \(X_T\) conditioned on \(a=0\) contains the positive rank-one term
\(\sigma_T^2\sigma_L^2 u_Tu_T^{\mathsf T}/\sigma^2\),
where \(\sigma_L^2=\ell^{\mathsf T} C_<\ell\).
By Lemma~\ref{lem:terminal-gaussian-comparison},
\(\sigma_L^2\asymp\lambda n^s\) and the conditioned covariance is
bounded by \(C_m\mathcal C_r^{(m)}\). Since
\(\sigma_T^2/\sigma^2\to1\) under enhancement, testing the rank-one
term against a stencil and applying the weighted stencil estimate
at \(b\asymp n\) gives the stated pointwise bound. With
\(u_\star=B_\Lambda^{1/2}e\) the exact identity
\begin{equation}
 u_T-u_\star=\frac{\sigma_L^2u_\star-C_<\ell}{\sigma_T^2}
\end{equation}
has right side \(o(1)\) in \(\ell^2\) by the regularized inverse
bound. Proposition~\ref{prop:cube-principal-mode} and
Cauchy--Schwarz now give \(\one^{\mathsf T}u_T\asymp v^{1/2}\).
Finally \(\|u_T\|_2\asymp1\), so the fourth-moment bounds follow
from \(v^{-1}\|u_T\|_2^4\le\sum_xu_T(x)^4\le
\|u_T\|_\infty^2\|u_T\|_2^2\).
\end{proof}

\subsubsection{The integrated uniform source}
A uniform magnetic source induces the deterministic field
\(C_<\one\) after the earlier slices are integrated. We estimate it
in the raw weighted test norm \(F_{n,p}\) of Section~\ref{sec:geometry},
taken at scale \(n\).

\begin{corollary}[Integrated uniform source]
\label{cor:integrated-uniform-source}
For an ordinary
slice \(C_j\) of side length \(b\), for any fixed derivative order
and \(0<p<p'<\alpha\) one has
\begin{equation}
 \|C_j\one\|_{F_{n,p}}\le C\lambda n^p b^{s-p}.
\end{equation}
Consequently \(\|C_<\one\|_{F_{n,p}}\le C_R n^s\) at fixed \(R\).
\end{corollary}

\begin{proof}
The variance bounds at weight \(p'\) and covariance
Cauchy--Schwarz give
\begin{equation}
 |\nabla^\gamma(C_j\one)(x)|
 \le C \lambda b^{d-s}\rho_b(x)^{-|\gamma|}
       (\rho_b(x)/b)^{p'} .
\end{equation}
Indeed a row and its fixed stencil have support of volume \(O(b^d)\),
and the other covariance leg has diagonal variance at most
\(C\lambda b^{-s}\). For positive derivative order the left side
vanishes more than \(Cb\) from the boundary by finite propagation
and the constant bulk row sum. In the collar
\(\rho_n(x)\le C\rho_b(x)\).
Everywhere \((\rho_b/b)^{p'}/(\rho_n/n)^p\le (n/b)^p\).
Together with \(d=2s\) these inequalities give the bound.
Choose \(p<s\). Summation over the preterminal scales gives
\(\|C_<\one\|_{F_{n,p}}\le C_{R} n^s\), with a scale-independent
constant at fixed \(R\) also for the merged slice. The same bound
holds in any test norm controlled by this raw norm.
\end{proof}

\subsection{The filtered bulk bubble}\label{sec:bubble}

The filtered bubble gives the leading decrease of the bulk quartic
coupling at each scale. The logarithm in the marginal flow requires a
positive contribution of order \(\lambda^2\) on the ordinary scales
outside a fixed number of terminal scales.
Write \(C=\Gamma_j^{\rm bulk}\), \(W=\sum_{i<j}\Gamma_i^{\rm bulk}\),
\(D_j=\overline T C\overline T^*\) and
\(w_j^T=\overline T W\overline T^*\).

\begin{proposition}[The filtered bulk bubble]
\label{prop:filtered-bulk-bubble}
For the prescribed reference and the normalized finite filter
\(\overline T\), the subtraction in the
compensated quartic coordinate is \(\mathfrak b_j g^2\), where
\begin{equation}
 \mathfrak b_j=9\sum_x\bigl((w_j^T+D_j)(x)^2-w_j^T(x)^2\bigr).
\end{equation}
On ordinary scales \(0\le\mathfrak b_j\le C\lambda^2\), with
\(\mathfrak b_j\ge c\lambda^2\) outside a fixed number of
terminal scales. The merged initial slice satisfies
\(0\le\mathfrak b_0\le C\lambda^2\log(2+b_0/R)\).
\end{proposition}

\begin{proof}
To verify the coefficient, periodize the kernels on a torus larger
than the contraction ranges and test on a constant field. For a
symmetric kernel \(S\), write
\(P_S=\exp(\tfrac12\sum_{x,y} S(x,y)\partial_x\partial_y)\) and
\(X_S(U)=\tfrac12[P_S((P_{-S}U)^2)-U^2]\).
The heat semigroup identity
\begin{equation}
 X_{W+C}(P_CU)=P_C X_W(U)+\tfrac12[P_C(U^2)-(P_CU)^2]
\end{equation}
expresses the second-order subtraction as the outgoing localized
\(X_{W+C}\) minus the incoming localized \(X_W\) evolved by \(P_C\).
Localization preserves the quartic evaluation at a constant field
since \(\overline T\one=\one\). Write \(G\) for the local quartic
coefficient before mass compensation. Two cross contractions between filtered
quartic vertices give
\(\tfrac12(G/4)^2\binom42^2 2!=9G^2/4\) times the square of the
filtered covariance per ordered pair of sites. With the quartic
convention \(G/4\), the difference gives the displayed coefficient.
For the single cross contraction with the quadratic vertex, write
\(u\) for its dimensionful row mass, \(u_+\) for the row mass of
\(P_CU\), and \(s_j=\sum_x w_j^T(x)\). The resulting quartic
increment is \(-4G(u_+s_{j+1}-us_j)\), because both the mixed and
filtered covariances have the unfiltered row sum. It cancels to
second order in the coordinate \(g=G(1+4u s_j)\) used in
Section~\ref{sec:cube-retained-chart}.
A zero-row quadratic vertex has zero contribution on a constant field.
In an interior cube cell all contraction stencils agree with their
bulk counterparts. We estimate the
coefficient in Fourier space, first from below on a frequency band and
then from above by the pointwise slice bound.

For the observation vector \(v(k)=\mathcal O(k)^*1\) write
\(V_\rho(k)=\langle v(k),(\mathcal Q_c(k)+\rho)^{-1}v(k)\rangle\).
Let \(\omega(k)=1-\widehat D_R(k)\). Near zero,
\(\omega(k)\asymp R^2|k|^2\), with isotropic quadratic leading
term. With a regularization \(\rho\ge0\), the physical Schur
complement is
\begin{equation}
 S_\rho(k)=B(k)/c+\rho-A_\rho(k),
 \qquad
 V_\rho(k)=A_\rho(k)+A_\rho(k)^2/S_\rho(k)
\end{equation}
for scalar bulk filters. In the fractional case \(A_\rho\) is
obtained by replacing \(1-u\widehat D_R\) in the beta integral
by \(1-u\widehat D_R+\rho\). It is bounded below near
zero frequency, uniformly for \(0\le\rho\le1\). In the local
split of Lemma~\ref{lem:local-split} use
\(A_\rho=A/(1+\rho)\).

Let \(r\) be the polynomial scale. At \(\omega(k)\asymp\delta r\)
and \(|c-1|\le \eta r^\alpha\), with \(\eta\) sufficiently small
depending on \(\delta\), \(S_0(k)\) is positive. The estimates
\begin{equation}
   V_{\delta r}(k)\ge c\delta^{-\alpha}r^{-\alpha},
           \qquad V_r(k)\le C r^{-\alpha}
\end{equation}
hold with constants independent of small \(\delta\). To see this
for \(\alpha<1\), use \(B-A\asymp\omega^\alpha\) and the upper
bound \(A_0-A_\rho\le C\rho^\alpha\), obtained by splitting the
beta integral at \(1-u=\rho\). For \(\omega\le\delta r\),
integration over \(1-u\asymp r\) gives
\(A_0-A_r\ge c r^\alpha\). In dimension four the corresponding
bounds follow from \(A_0-A_\rho=A\rho/(1+\rho)\) and \(s=2\).
The added positive \(\rho\) term in the Schur complement gives
the same estimates. This is a calculation on the full-lattice
bulk fibers.

Let \(\nu_k\) be the spectral measure of the observation vector
on this band. For the positive polynomial \(q_i\) of scale
\(r=L^{-2i}\), choose the fixed cutoff \(c_*>0\) in
\(q_i(t)\ge c/r\) for \(0\le t\le c_*r\).
The part of \(V_{\delta r}\) from \(t>c_*r\) is at most
\((1+c_*^{-1})V_r\). Choose \(\delta\) small enough that this
is at most half the lower bound on \(V_{\delta r}\). Hence
\(\nu_k([0,c_*r])\ge c_\delta r^{1-\alpha}\) and
\begin{equation}
 \langle v(k),q_i(\mathcal Q_c(k))v(k)\rangle
       \ge c_\delta r^{-\alpha}.
\end{equation}
The normalized finite filter is bounded below there by its row
sum one and second moment \(O(R^2)\).
The band has measure \(\asymp_{\delta,d}R^{-d}r^{d/2}\);
multiplication by \(a_R\) and Plancherel give
\begin{equation}
 \|D_i\|_2^2\ge c\lambda^2 R^{2s-d} r^{d/2-2\alpha}
             =c\lambda^2 .
\end{equation}
Here \(D_i\) denotes an individual filtered polynomial slice in
this estimate. The nonnegative cosine coefficients of the scalar construction
in Appendix~\ref{app:gaussian-scalar} give positivity also on the
negative bulk spectrum.
Thus the two contributions to \(\mathfrak b_j\) are nonnegative
by Plancherel, and a merged slice inherits the lower bound of
any one component. The estimate \(|c-1|=O((R/n)^s)\) puts
all but a fixed number of terminal scales in the required
displacement range.

For the upper bound the ordinary pointwise estimate is
\(\|\Gamma_j^{\rm bulk}\|_\infty\le C\lambda b_j^{-s}\).
Here is its direct bulk proof, which also covers preterminal scales
with a fixed enlarged displacement constant
\(|c-1|\le C r^\alpha\). Shift the auxiliary precision by
\(\kappa' r\) to make it nonnegative, using the Schur lower bound
at that regularization, and choose \(\kappa>\kappa'\) fixed.
The spectral majorant of \(q_i(t)\) is
\(C_m r^{2m-1}(t+\kappa r)^{-2m}\). On \(t\ge0\) this follows
from the positive-spectrum polynomial bound. On
\(-\kappa' r\le t<0\), the cosine coefficients, \(q_i(0)=O(r^{-1})\)
and the degree \(O(r^{-1/2})\) bound the hyperbolic cosines by a
constant. At \(|Rk|\le c_1\), the Schur spectral lower bound for
the shifted precision reduces the observed inverse power to
an observed resolvent times \(C(r+R^2|k|^2)^{-(2m-1)}\).
The observed resolvent is bounded by
\(C(r+R^2|k|^2)^{-\alpha}\).
Integration yields \(C R^{-d} r^{d/2-\alpha}\) before the
factor \(a_R\). On the remaining frequencies the gap is
bounded below and the square integral of the observation
vector is \(O(R^{-d})\), by the smoothing in the split.
The factor \(r^{2m-1}\) thus controls this range as well.
If \(\kappa r\) is bounded below, the same inequalities
hold with constants depending on \(\kappa\) and fixed spectral bounds.
Positivity of the bulk multiplier controls every entry by its
diagonal, and \(a_R R^{-d}r^{d/2-\alpha}=\lambda b_j^{-s}\).

Finite range and bounded absolute filter rows now give
\(\|D_j\|_1\le C\lambda b_j^s\) and
\(\|w_j^T\|_1\le C\lambda b_j^s\).
It follows that \(\mathfrak b_j\le C\lambda^2\).
For a merged slice, the same ordinary bound and finite range
give the envelope \(C\lambda(R+|x|)^{-s}\) on its support of
radius \(O(b_0)\) before filtering. Squaring and summing over
radial layers costs \(C\lambda^2\log(2+b_0/R)\) because \(2s=d\).
The filter is bounded in \(\ell^2\), giving the asserted
first-step upper bound. This calculation isolates the bulk
covariance input needed for the logarithmic flow.
\end{proof}

\subsection{Relative form bound on the boundary Hessian}

Subtracting the bulk row mass leaves a zero-row convolution and
a correction concentrated near the boundary. We first bound both
terms relative to the kinetic energy. The next subsection estimates
them in the weighted field norm required by the polymer expansion.

\begin{lemma}[Form bound on the initial boundary Hessian]
\label{lem:boundary-form-extension}
Write the zero-extension kinetic form as
\begin{equation}
 \mathcal E_\Lambda(u)=
 \langle P_\Lambda^*u,L_R^{d/4}P_\Lambda^*u\rangle .
\end{equation}
Suppose \(D_\Lambda\) is symmetric and
\(\sum_y|D_\Lambda(x,y)|\le Cg e^{-a\delta(x)/R}\).
For \(d=1,2,3,4\), uniformly in the cube,
\begin{equation}
 |\langle u,D_\Lambda u\rangle|\le C_a g\mathcal E_\Lambda(u).
\end{equation}
The same relative form bound holds for a compressed even
zero-row convolution \(Z\) with
\(\sum |Z(z)|\le Cg\) and \(\sum|z|^2|Z(z)|\le CgR^2\).
\end{lemma}

\begin{proof}
Symmetry and \(2|u_xu_y|\le |u_x|^2+|u_y|^2\) reduce the
boundary bound to \(\sum_xe^{-a\delta(x)/R}|u_x|^2\).
For \(d<4\), the zero-extension jump form includes the killing
row \(r_\Lambda(x)|u_x|^2\). The jump lower bound gives
\(r_\Lambda(x)\ge c[R/(R+\delta(x))]^s\) by summing an exterior
box across a closest face at distance comparable to
\(R+\delta(x)\). This dominates the required exponential.

For \(d=4\), the Fourier lower bound
\(1-\widehat D_R(k)\ge c\min\{R^2|k|^2,1\}\)
controls the zero-extension energy of the one-dimensional
uniform walk with jumps \(\pm1,\ldots,\pm R\) in any chosen
coordinate. Consider the lower face of that coordinate and
a fixed step \(\ell\in[R/4,R/2]\). If \(x_i\) is the normal
distance in lattice units, telescope from \(x\) to the first
point outside the face along the negative \(\ell\)-steps.
Cauchy--Schwarz gives
\begin{equation}
 |u_x|^2\le C(1+x_i/R)
 \sum_{\substack{j\ge0\\x_i-j\ell\ge 1}}
        |u_{x-j\ell e_i}-u_{x-(j+1)\ell e_i}|^2 .
\end{equation}
With weight \(e^{-a x_i/R}\), summing the right side costs
only a fixed multiple of each edge with interior upper
endpoint \(t\), whose coefficient is bounded by
\begin{equation}
 \sum_{j\ge0}e^{-a(t_i+j\ell)/R}(1+(t_i+j\ell)/R).
\end{equation}
This bound is uniform in \(\ell\) in the specified range
and in \(t_i\ge1\). Average over those \(\ell\), whose
number is comparable to \(R\), to obtain the normalized
axis energy. Repeat for the upper faces. The exponential
of the distance to the complement is bounded by the sum
of the face exponentials, up to a harmless lattice convention.

For the bulk term the zero row and evenness give
\begin{equation}
 |\widehat Z(k)|
 \le \sum_z|Z(z)|\,|1-\cos(k\cdot z)|
 \le Cg\min\{R^2|k|^2,1\}
 \le C'g(1-\widehat D_R(k))^{d/4}.
\end{equation}
Use Plancherel on the zero extension.
\end{proof}

The affine-completed field norm defined in Section~\ref{sec:geometry} admits constant
test directions at the boundary. A shrinking boundary coefficient in
that norm therefore requires a separate estimate. Absorbing the
perturbation into the Gaussian reference would also require a new
principal pole and regulator estimates for the modified slices.

\subsection{The initial bulk and boundary quadratic heads}
\label{sec:initial-heads}

We now split each initial quadratic interaction into a head, whose
two endpoints lie within a fixed multiple of the initial cell size,
and a tail that enters the polymer activity. Writing the zero-row
bulk form as a sum of squared differences gains two powers of the
jump length. For rows lost at the physical boundary, the gain comes
from the boundary weight in the field norm.

\begin{proposition}[Initial quadratic heads in the cube chart]
\label{prop:initial-cube-heads}
Let \(d\in\{1,2,3,4\}\), \(b=b_0=M_R R\), and choose
\(\max\{0,(s-1)/2\}<p<s/2\), with \(p\ne 1/2\). Subtract the bulk row
mass from the canonical Hessian of Section~\ref{sec:micro-extension}.
On a fixed enlarged cube of side comparable to \(b\), the resulting
bulk difference-square form and the boundary form have normalized
bilinear \(F_b\)-norms, at radius \(\ell_b\), bounded respectively by
\begin{equation}
 C\lambda G_0\big(M_R^{s-2}+M_R^{s-1-2p}\big),
 \qquad C\lambda G_0 M_R^{s-1-2p}.
\end{equation}
The interior bulk estimate is \(C\lambda G_0 M_R^{s-2}\).
The forms admit exact decompositions into heads of span \(O(b)\)
and exponentially small tails. Anchors of the boundary head are
within \(O(b)\) of the physical boundary. The tail row bounds and
the corresponding local tensor bounds gain a factor
\(C\exp(-c b/R)\). All bounds hold on the complex thermal window
\(|\beta-1|\le A_0g_{\mathrm{mic}}\) at fixed \(A_0\).
\end{proposition}

\begin{proof}
In the primary field, let \(Z\) denote the translation-invariant
zero-row Hessian after the scalar subtraction and let \(D\)
denote the finite-volume comparison. Put
\(\theta=\lambda G_0\). Lemma~\ref{lem:initial-hessian-comparison}
and the exponentially weighted
Neumann norm of the primary pullback give
\begin{equation}
 \sum_z |Z(z)|e^{a|z|/R}\le C\theta R^{-s},\qquad
 \sup_{x\in\Lambda}\sum_{y\in\Lambda}
   |D(x,y)|e^{a(\delta(x)+|x-y|)/R}\le C\theta R^{-s},
 \qquad \delta(x)=\operatorname{dist}(x,\Lambda^c).
\end{equation}
Here the canonical conversion multiplies the physical Hessian
by \(\lambda^{-1}R^{-s}\), and
\(\lambda^{-1}g_{\mathrm{mic}}\asymp\theta\).
All norms in these two inequalities use absolute values, so the
argument applies also to complex coefficients. For the compressed
bulk form use the exact identity
\begin{equation}
 \frac12\sum_{x,y\in\Lambda}\phi_x Z(x-y)\phi_y
 =-\frac14\sum_{x,y\in\Lambda}Z(x-y)(\phi_x-\phi_y)^2
   +\frac12\sum_{x\in\Lambda}m_x\phi_x^2,
 \qquad
 m_x=-\sum_{y\notin\Lambda}Z(x-y).
\end{equation}
The missing row \(m_x\) obeys
\(|m_x|\le C\theta R^{-s}e^{-a\delta(x)/R}\).
We anchor the two-site summands at a site of the pair and use
the same assignment for the head and tail. For a local test the
chart is taken on a fixed enlarged tile containing the anchor
and the endpoints of the head summands. When testing long
summands on several tiles we use directions satisfying the
stated field bound at each endpoint.

Write \(w_b(x)=(\rho_b(x)/b)^p\). For \(F_b(v)\le\ell_b\),
one has \(|v_x|\le\ell_b w_b(x)\) and
\(w_b(y)\le Cw_b(x)(1+|x-y|/R)^p\). For the boundary forms,
absorb the latter polynomial in a fraction of the exponential
row weight. Each physical normal layer intersects an enlarged
tile in at most \(Cb^{d-1}\) sites, including tiles with several
faces. Consequently
\begin{equation}
 \sum_{x\in U}e^{-a\delta(x)/(2R)}w_b(x)^2
       \le C b^{d-1}R(R/b)^{2p}.
\end{equation}
Multiplying by \(\theta R^{-s}\ell_b^2\) proves the boundary
estimate.

For span \(|x-y|\le c\rho_b(x)\), summing first differences in
the enlarged tile gives
\begin{equation}
 |v_x-v_y|\le C\ell_b b^{-p}\rho_b(x)^{p-1}|x-y|.
\end{equation}
For longer spans one uses the endpoint bound and
\(\sum_{|z|>c\rho}|Z(z)|(1+|z|/R)^{2p}
 \le C\theta R^{-s}e^{-c'\rho/R}\).
These two estimates bound the difference-square sum on the
tile by
\begin{equation}
 C\theta \ell_b^2 b^{-2p}R^{2-s}
              \sum_{x\in U}\rho_b(x)^{2p-2}.
\end{equation}
For the long spans the same bound follows from
\(t^2e^{-c't}\le C\) for \(t=\rho_b(x)/R\).
On interior tiles \(\rho_b\asymp b\). In general, summing the
layers with \(\delta(x)<b\) and then the interior gives
\begin{equation}
 \sum_{x\in U}\rho_b(x)^{2p-2}
       \le C b^{d-1}(b^{2p-1}+R^{2p-1}).
\end{equation}
Since \(\ell_b^2=b^{-s}\) and \(d=2s\), these are the desired
powers of \(M_R\). The same estimates bound the bilinear form
by applying the pointwise bounds to each direction.

Take the bulk head to have \(|x-y|\le\vartheta b\), the
\(D\)-head to have \(|x-y|+\min(\delta(x),\delta(y))\le\vartheta b\),
and the missing-row head to have \(\delta(x)\le\vartheta b\),
for any fixed \(\vartheta>0\). Cutting the bulk difference
square preserves its zero row in the infinite lattice. On a
discarded term, the unused exponential row weight supplies the
factor \(C e^{-c b/R}\); for a long bulk span one can use the
long-span estimate without taking a derivative along that
span. One can put back every discarded term using the same pair
anchor. This proves the tail bounds and the exact decomposition.
\end{proof}

The leading thermal term before the pullback is
\(- (\beta-1)I\). Its pullback has a bulk constant part,
a bulk zero-row part, and a boundary comparison with the same
weights as in Proposition~\ref{prop:initial-cube-heads}.
After division by the initial mass radius \(\Lambda_R G_0\),
the two nonscalar terms change by
\(O(A_0(M_R^{-2}+M_R^{-1-2p}))\) across the window.
All other thermal changes in the weighted microscopic
Hessian are \(O((A_0+A_0^2)g_{\mathrm{mic}}^2)\) before
canonical conversion, by the exact formula for \(T_\beta\)
and its weighted derivative estimates. Thus the nonscalar
initial coordinates can use constants fixed before choosing
the thermal window, by taking \(R\) large after fixing
\(A_0\). The boundary estimate decreases with \(M_R\) in all four dimensions. The bulk estimate leaves an initial
kinetic term of size \(O(\lambda G_0)\) in dimension four.

\section{Boundary transport and localization}\label{sec:geometry}

Passing from an interval to a cube changes the boundary contribution
to a single renormalization-group step. A coarse face contains
many fine tiles. To sum their contributions we test fields with a
weight that decays towards the physical faces and retain the local
quadratic forms supported in a collar of those faces. An
\emph{anchor} specifies the site or tile to which a polynomial is
assigned. Several polynomials may depend on fields at the same site
while having different anchors. Keeping track of the anchors makes the
covariance subtraction compatible with the reblocking of polymers.

Two results connect this construction to the shooting argument of
Section~\ref{sec:path}. Corollary~\ref{cor:fractional-cube-path}
gives the one-step estimates in dimensions one through three. In
dimension four the bulk kinetic coefficient is marginal.
Proposition~\ref{prop:residual-coupled-chart} gives the corresponding
estimates with the Gaussian reference fixed and a separate equation
for the residual kinetic coefficient.

\subsection{Face geometry and the change in the RG chart}

Set \(s=d/2\), \(\ell_b=b^{-s/2}\), and \(B=Lb\), where \(b\) is
the side length of a fine tile. Fixed enlargements of the supports
are absorbed in the constants independent of \(L\). On a cube
touching a physical face, let
\begin{equation}
 \rho_b(x)=\min(b,R+\operatorname{dist}(x,\Lambda^c)),\qquad
 F_b(v)=\max_{|\gamma|\le P}\sup_x
 \frac{\rho_b(x)^{|\gamma|}|\nabla^\gamma v(x)|}{(\rho_b(x)/b)^p}.
\end{equation}
Write \(E_b(v)=\inf_{a\ {\rm affine}}\{F_b(v-a)+b^{-d/2}\|a\|_2\}\)
for the norm obtained by adjoining affine fields. Two estimates for
unrestricted boundary tensors explain our choice of test space.

\begin{lemma}[Unweighted sum of face tensors]
\label{lem:face-count}
Let a coarse face tile of side \(B=Lb\) meet a full layer of fine
face tiles of side \(b\). On each fine tile let a homogeneous
\(k\)-linear form have absolute value at most one whenever all
its arguments satisfy \(E_b(v)\le\ell_b\).
The norm of their sum on directions with \(E_B(v)\le\ell_B\) is at
most \(C L^{d-1-ks/2}\). The power of \(L\) is attained up to a constant.
\end{lemma}

\begin{proof}
A new face tile of side \(B\) meets order \(L^{d-1}\) old face
tiles of side \(b\). Restriction of a coarse \(E_B\)-unit direction
to a fine tile has bounded fine \(E_b\)-norm. This follows by
restricting the affine decomposition, since affine coefficients
on the coarse cube have size \(O(1)\) in supremum and slope
\(O(B^{-1})\), while \(\rho_b\le\rho_B\) and
\((\rho_B/B)^p\le(\rho_b/b)^p\) for stencils in the tile.
Thus a homogeneous degree-\(k\) form of normalized fine norm one
has coarse norm at most \(C L^{-ks/2}\) on one fine tile. The sum of
the old face forms has norm at most \(C L^{d-1-ks/2}\).
This power is sharp on unrestricted tensors. Take a site \(x_i\)
in each tile and forms \(\ell_b^{-k}v(x_i)^k\). Their fine norms
are uniformly bounded since \(E_b\) controls the supremum. Evaluate
their sum on the constant coarse direction \(\ell_B\mathbf 1\),
whose \(E_B\)-norm is at most \(\ell_B\). The value is of order
\(L^{d-1-ks/2}\).
\end{proof}

In particular, the quadratic power
\begin{equation}
 L^{d-1-s}=L^{s-1}
\end{equation}
equals one in \(d=2\) and grows with \(L\) in \(d=3,4\).
Gaussian integration of a quadratic changes only its constant
term. Thus the same power occurs in the linearized RG map.
In dimension one there are only two endpoints and the factor is \(L^{-1/2}\). In higher
dimensions our boundary test space gives an additional gain for each
field direction.

\paragraph{Loss of the affine product bound.}
In dimensions \(d\ge2\), on a cube of side \(b\) touching \(x_1=1\), take
\(v(x)=x_2/b\). As \(v\) is affine, \(E_b(v)=O(1)\).
For any affine \(a\), evaluate \(v^2-a\) on three equally spaced
tangential sites of spacing \(\lfloor (b-1)/3\rfloor\) in the layer
\(x_1=1\). Its second difference across these sites has absolute
value bounded below by a positive constant, since the affine part
has zero second difference there. Therefore some one of these
values is bounded below in modulus, whereas its boundary weight
is \(O((R/b)^p)\). Taking the infimum over \(a\) gives
\begin{equation}
 E_b(v^2)\ge c(b/R)^p .
\end{equation}
Thus in \(d\ge2\) the product estimate for \(E_b\) fails for every \(p>0\). This affects the
large-field regulator and the matching of local spatial jets.
At the level of the quotient restriction estimate, the direct
summing argument also encounters order \(L^{d-1}\) tiles in each
boundary layer. Indeed a raw coarse direction on the layer at
distance \(i b\) has local affine quotient bounded by
\(C L^{-p}(1+i)^{p-2}\) after subtracting its first-order Taylor
polynomial. Summing the second differences along coordinate
segments of length \(O(b)\) gives the stated bound. For the first layer use direct raw
restriction and the factor \(L^{-p}\). Summing the squares gives
\(C L^{d-1-2p}\) if \(0<p<1\), to which the large-field radii
contribute the factor \(L^{-s}\) at comparable couplings.
The resulting power \(L^{s-1-2p}\) requires \(2p>s-1\) to
yield a small factor.

\paragraph{Ordinary spatial localisation.}
For mass localisation, cancelling evaluation on constant fields
alone gives a factor \(O(b/\rho_B(x_i))\) from the first-order
variation of the field across a fine placement \(x_i\). Even
on a coarse cube at distance at least \(B\) from the physical
boundary, summing over order \(L^d\) placements gives the power
\(L^d L^{-s}L^{-1}=L^{s-1}\). The order in this estimate cannot be
improved for arbitrary jets from this cancellation: the diagonal
quadratic difference \(v(x_i+te)^2-v(x_i)^2\), with
\(t\asymp b\), has zero constant-field evaluation and zero
relative second moment, yet is of order \(b/B\) on affine
coarse fields in a region of fixed rescaled width. An additional
first-spatial-moment cancellation gives \(O((b/B)^2)\) in the
interior, hence \(O(L^{s-2})\), which suffices there for \(d=1,2,3\).
Proposition~\ref{prop:cube-jet-projection} implements the moment
cancellations also for marks crossing coarse faces. A fixed coarse
face can meet order \(L^{d-1}\) translates of a mark crossing the
interface. Even with \(O(L^d)\) placements, the nine-direction Taylor
remainder gains \(L^{d-9s/2}=L^{-5s/2}=L^{-5d/4}\), since \(d=2s\).

\paragraph{A weighted sunset bound for \(d=1,2,3\).}
If a cubic sunset difference-square kernel obeys
\begin{equation}
 |S(x,y)|\le Cg^2\lambda^3(R+|x-y|)^{-d-s},
 \qquad S(x,y)=0\quad (|x-y|>Cb),
\end{equation}
then its quadratic form on an enlarged face tile has bounded
normalized \(E_b\)-norm whenever \(2p>s-1\). To see this, write a
direction of \(E_b\)-norm at most one as affine plus a raw
component of bounded \(F_b\)-norm. For the affine component use
\(|v_x-v_y|\le C |x-y|/b\); the summed kernel row is \(O(b^{-s})\)
since \(s<2\). For a raw component \(f\), for
\(|x-y|\le c\rho_b(x)\) its difference is at most
\(C(\rho_b(x)/b)^p |x-y|/\rho_b(x)\), by summing first
differences on the tile. For longer spans it is bounded by
\(C(\min(b,\rho_b(x)+|x-y|)/b)^p\).
For \(0<2p<s\), shell summation at fixed \(x\) gives an energy
row bounded by \(C g^2\lambda^3 b^{-2p}\rho_b(x)^{2p-s}\).
The assumption \(2p>s-1\) then yields
\(\sum_x \rho_b(x)^{2p-s}\le C b^{d+2p-s}\), by summing
normal layers each having at most \(Cb^{d-1}\) sites.
After multiplying by the two small-field radii
\(\ell_b^2=b^{-s}\) the estimate is \(O(g^2\lambda^3)\), since
\(d=2s\). For \(2p\ge s\) reduce \(p\) to any value with
\(s-1<2p<s\), which weakens the weight. Polarization supplies
the symmetric tensor norm. The estimate uses only the kernel
envelope and the weight. In dimension three it requires
\(p>1/4\). The product estimate for the field regulator will use
the raw norm \(F_b\).

\paragraph{Kinetic coordinates at \(d=4\).}
A normalized coefficient of a bulk quadratic with two spatial
differences transports with factor \(L^{s-2}\), which is one
at \(s=2\). Thus it cannot be placed in the contracting stable
block of the shooting theorem. Also, for a cumulative
Green envelope of order \((R+|x|)^{-s}\), the absolute second
moment of its cube is bounded by a constant times
\begin{equation}
 \sum_{m=1}^{Cb} m^{d+1}(R+m)^{-3s}.
\end{equation}
This is \(O(b^{2-s})\) for \(s<2\) and logarithmic at \(s=2\).
At this endpoint the kinetic coefficient and its derivatives must
therefore be followed over the full trajectory. A source
of order \(O(g_j^2)\) is summable when
\(g_j\asymp (g_0^{-1}+c j)^{-1}\).
Proposition~\ref{prop:bulk-kinetic-source} bounds the successive
increments, and Proposition~\ref{prop:residual-coupled-chart}
treats their coupling to the collar.

\paragraph{The boundary field norm and its affine quotient.}
The raw norm \(F_b\) is a uniform algebra by the discrete
Leibniz rule, since the weights are at most one and comparable
on fixed stencils. For a tile \(U\) of side comparable to \(b\),
let \(\mathcal A_U\) be the affine fields vanishing on the
planes of all physical faces within distance \(2b\) of \(U\).
Other fixed dilation factors work as well, using nested
enlargements for fine and coarse tiles. If there is one such
face, an element of \(\mathcal A_U\) is linear in the distance
from that plane. If there are two distinct face planes, it is
zero. If there are none, the raw weights on \(U\) are one.
These observations and scaling of affine polynomials on a
cube give
\begin{equation}
 F_b(a)\le Cb^{-d/2}\|a\|_{2,U}\quad (a\in\mathcal A_U,\ 0<p<1).
\end{equation}
For one face, \(|a(x)|\le Cb^{-d/2}\|a\|_2
\operatorname{dist}(x,\text{face})/b\); when \(\rho_b(x)<b\) the
nearest physical face must be among the included faces.
The derivative estimates follow from the bounded normalized
slope. In particular, with
\(q_U(v)=\inf_{a\in\mathcal A_U}F_b(v-a)\) one has
\begin{equation}
 F_b(v)\le C\{q_U(v)+b^{-d/2}\|v\|_{2,U}\}.
\end{equation}
Indeed an approximating \(a\) has normalized \(L^2\)-norm at
most \(b^{-d/2}\|v\|_2+C F_b(v-a)\). This quotient therefore
suits quadratic damping by the local \(L^2\)-norm and does not
require an affine-completed algebra. Coarse affines in this
definition restrict to fine affines of the allowed type.

In the raw norm, restriction from a coarse unit direction to
fine tiles in the outer \(O(b)\) physical face layer gains
\(C L^{-p}\) in each direction, since the two distance functions
are comparable on that layer. Therefore full degree-two face
forms with bounded fine norms give coarse norm
\(O(L^{s-1-2p})\), a power strictly below one when \(2p>s-1\).
The bound contracts for any \(p>0\) in dimensions one and two, and for \(p>1/4\) in dimension three.
The same computation gives the summed quotient estimate, since
the fine quotients annihilate the restricted coarse affine.
On the remaining layer \(i\ge1\) at distance comparable to
\(ib\) use fine affine Taylor subtraction, giving the local
quotient factor \(C L^{-p}i^{p-2}\).

For ordinary quadratic residuals, suppose the local projection
cancels the constant evaluation and the evaluation with one
constant and one linear coordinate. Taylor expansion of the
coarse raw directions on a layer at depth \(ib\) then costs
\(C((i+1)/L)^{2p}(i+1)^{-2}\) per placement, before multiplication
by the ratio \(\ell_B^2/\ell_b^2=L^{-s}\). For this estimate,
terms with two linear increments cost the square of the
first-order variation, and a single Taylor remainder has the
same bound. There are at most \(C L^{d-1}\) placements at any
fixed layer. Summation gives \(O(L^{s-1-2p})\) for \(p<1/2\),
including the interior contribution \(O(L^{s-2})\). For the
quartic, constant evaluation cancellation alone gives
\(C((i+1)/L)^{4p}(i+1)^{-1}\) per placement before the factor
\(L^{-d}\), and hence \(O(L^{-1})\) for \(p>0\).
These powers give contraction for \(d=1,2\) and, with the stated
boundary weight, for \(d=3\). To use the estimates in a polymer
activity we must allocate the moments on the actual cell labels and
control the Gaussian translation in the regulator. The allocation
is given in Proposition~\ref{prop:cube-jet-projection}. The
fluctuation moments of Section~\ref{sec:ext-cov} enter the
product estimate of Lemma~\ref{lem:cube-joint-reblocking}.

\paragraph{Activity norm and circle product.}
At scale \(b\) let \(X\) be a set of disjoint cell labels, each with a
fixed enlarged field support \(P^*\), clipped to the physical cube.
A field direction on \(X\) is tested in
\(\|v\|_X=\max_{P\in X}F_{b,P^*}(v)\).
For a field-dependent activity \(K(X)\) use the finite jet seminorm
\begin{equation}
 |K(X)|_{r,\phi}
 =\sum_{k=0}^{M}\frac{r^k}{k!}
   \sup_{\|v_i\|_X\le1}
   |D^k K(X,\phi)[v_1,\ldots,v_k]|,\qquad M\ge9 .
\end{equation}
Here \(g_b>0\) is the reference coupling used to set the radii
\(\ell_b=b^{-s/2}\) and \(h_b=g_b^{-1/4}\ell_b\). The two regulators are
\begin{equation}
 \begin{aligned}
 G_b(X,\phi)&=\exp\!\left(c_\ell\sum_{P\in X}
           F_{b,P^*}(\phi)^2/\ell_b^2\right),\\
 \mathcal R_b(X,\phi)&=\exp\!\left(\sum_{P\in X}
     \left[-\kappa\frac{\|\phi\|_{2,P}^2}{b^d h_b^2}
           +c_h\frac{q_{P^*}(\phi)^2}{h_b^2}\right]\right).
 \end{aligned}
\end{equation}
The mixed norm on one polymer is
\begin{equation}
 \|K(X)\|_{\mathrm{mix},b}
 =\max\left\{\sup_\phi\frac{|K(X)|_{\ell_b,\phi}}{G_b(X,\phi)},
       \;g_b^{9/4}\sup_\phi\frac{|K(X)|_{h_b,\phi}}{\mathcal R_b(X,\phi)}\right\}.
\end{equation}
The activity norm takes the supremum over connected \(X\), with the
size weight \(\varrho^{(|X|-N_s)_+}\), \(\varrho>1\), of
Section~\ref{sec:cube-linked-size}. The anchored sum used at the
initial scale bounds this supremum for any fixed size weight.
A direction on a union restricts to each factor with norm at most
one. The finite-jet product rule and the product over disjoint
cell labels therefore apply even when their stencils overlap.

Write \(I(X)=\prod_{P\in X}I(P)\), extend \(K\) by products over the
connected components of its argument, and set \(K(\varnothing)=1\).
The weight on \(X\) is the circle product
\begin{equation}
 (I\circ K)(X)=\sum_{Y\subseteq X}I(X\setminus Y)K(Y).
\end{equation}
A step integrates this weight against the fluctuation covariance,
regroups the fine labels by their coarse supports, and allocates the
retained Taylor coefficients by Proposition~\ref{prop:cube-jet-projection}.

\paragraph{The corrected cubic row.}
The three-line covariance term contains a logarithmic row sum
when \(2s=d\). Subtracting the local diagonal from the current
slice gives a bound uniform in the scale under the following
hypotheses. Suppose the cumulative kernel \(W\)
and the current Gaussian covariance \(\Gamma\) have range
\(Cb\), \(|W(x,y)|\le C\lambda(R+|x-y|)^{-s}\), and the
diagonal and first-difference standard deviations of \(\Gamma\)
are bounded by \(C(\lambda b^{-s})^{1/2}w(x)\) and
\(C(\lambda b^{-s})^{1/2}w(x)/\rho_b(x)\), respectively,
with \(w(x)=(\rho_b(x)/b)^p\) and \(p>0\).
The Cauchy--Schwarz inequality for \(\Gamma\), applied to a
difference stencil and then telescoped, gives
\begin{equation}
 |\Gamma(x,y)-\Gamma(x,x)|
 \le C\lambda b^{-s}w(x)^2 |x-y|/\rho_b(x)
 \qquad (|x-y|\le c\rho_b(x)).
\end{equation}
The sum of \((R+|x-y|)^{-2s}|x-y|/\rho_b(x)\)
over this near region is bounded by a constant, since \(2s=d\).
On the far region use
\(|\Gamma(x,y)|\le C\lambda b^{-s}w(x)w(y)\) and
\(w(y)\le C(\min(b,\rho_b(x)+|x-y|)/b)^p\).
There the integral of \(r^{-1}(r/b)^p\) on \([c\rho_b(x),Cb]\)
is bounded. The subtracted diagonal contributes at most
\(C\lambda b^{-s}w(x)^2\log(2+b/\rho_b(x))\), also bounded by
\(C_p\lambda b^{-s}\). Thus
\begin{equation}
 \sup_x\sum_y |W(x,y)|^2
          |\Gamma(x,y)-\Gamma(x,x)|
 \le C_p\lambda^3 b^{-s}.
\end{equation}
The same argument treats the symmetric diagonal average, with the
far \(y\)-diagonal integral using \(r^{-1}(r/b)^{2p}\).
The other two increment terms \(W\Gamma^2\) and \(\Gamma^3\)
give the same row bound by finite range and the pointwise estimates.
This proves the corrected row inequality in all four dimensions
under the stated slice bounds.

\paragraph{Connected covers in the initialization.}
In higher dimension a rectangular hull can have cardinality
of order the \(d\)-th power of its diameter, so its exponential
polymer weight cannot generally be paid by a linearly exponential
filter-tail bound. For a tail edge between two fine tiles choose
a grid path of length at most \(C(1+\operatorname{dist}/b_0)\).
A boundary mark can likewise be joined to a closest physical
face by a normal path. Include the actual occupied labels and these
paths, rather than filling a rectangle. For fixed root tile \(D\),
the number of pairs of separation at most \(k\) in grid units
whose path contains \(D\) is at most \(C(1+k)^{2d}\), since
both endpoints must be within \(C(1+k)\) of \(D\).
For a boundary promotion use the corresponding count with \(k\)
equal to the pair span plus the distance to the selected face,
in grid units. The boundary-defect weight contains the
exponential in this additional distance. Consequently these
weights pay \(e^{A|\mathrm{path}|}\) for any fixed \(A\) when
\(b_0/R\) is large, and leave a summable edge weight at each root.
Averages over
coordinate-order paths and over tied closest faces preserve the cube
symmetries and change the rooted bounds by at most fixed constants.
Proposition~\ref{prop:cube-joint-initialization} uses this count to
regroup the filter tails together with the microscopic pressure
terms in the activity norm.

\paragraph{Large-field reserve for the adapted quotient.}
For \(d=1,2,3\) and \(2p>s-1\), the adapted quotient controls
the field on the complement of the occupied labels. Assume a finite primary filter \(T\) with uniformly
bounded absolute rows and columns and a finite forward filter
\(B_{\mathrm{rec}}\) of range \(O(R)\) such that
\(\|I-B_{\mathrm{rec}}T\|_{2\to2}\le\epsilon\). Use the local decaying regulator with tile volume \(b^d\).
In a coarse tile, normalize its squared background norm \(Z^2\),
the squared norm \(A^2\) on occupied sites, and the squared
filtered norm \(S^2\) on unoccupied rows by \(B^d h_B^2\).
Write \(Q=q_{C^*}(\phi)/h_B\), where \(C^*\) is its enlarged tile.

Outside strips of width \(O(R)\) along the fine occupied
interfaces and the coarse faces, forward recovery on an
unoccupied site uses only unoccupied rows of \(T\phi\).
The total strip volume is \(O((R/b)B^d)\). The norm bound
\(F_B(\phi)/h_B\le C(Z+Q)\) follows as above with the affine
norm on the smaller tile in the right-hand side; normalized
\(L^2\) on a tile controls the affine norm on its fixed
enlargement. The recovery identity, squared on the
unoccupied sites, yields
\begin{equation}
 (1-\delta)Z^2\le A^2+C S^2+\delta Q^2,\qquad
 \delta=C(R/b+\epsilon^2).
\end{equation}
If \(S\) uses a rowwise convex combination of \(T\phi\) and
\(T(\phi+\zeta)\), add
\(C\|T\zeta\|_{2,C}^2/(B^d h_B^2)\) on the right.
The quartic bound on the unoccupied labels supplies any
fixed negative multiple of \(S^2\), since \(g_b B^d h_B^4\)
is bounded above and below. The occupied decay supplies
a multiple of \(A^2\) of order
\(\kappa B^d h_B^2/(b^d h_b^2)\asymp \kappa L^s\).
Using the summed quotient gain and then taking \(R/b+\epsilon^2\)
small yields a strict negative \(Z^2\) margin and positive
\(Q^2\) coefficient strictly below the output allowance.
With uniform Gaussian exponential moments for \(F_b\),
finite-range coloring bounds the additional fluctuation cost
by a constant per fine label. Fixed field jets and small
quadratic perturbations are absorbed by the strict margins in the
two squared norms, since any fixed power is bounded by a Gaussian
exponential with arbitrarily small positive coefficient. Lemma~\ref{lem:cube-joint-reblocking}
applies this estimate to products whose occupied labels and
field stencils can differ.

\subsection{An exact cube projection for the low field jets}
\label{sec:cube-jet-projection}

For a quadratic Taylor coefficient, the spatial moments of order
two test a bilinear form in two ways. Both directions may be linear,
or one may be constant and the other quadratic. Matching both,
as well as the lower moments, gives a small remainder on each tile
when \(2p>s-1\), including at \(d=4\). The bulk kinetic coefficient
retained by this projection is marginal at \(s=2\).

\begin{proposition}[Local moment allocation on cube tiles]
\label{prop:cube-jet-projection}
Use an integer scale ratio \(L\), fine tiles of side \(b\), and
coarse tiles of side \(B=Lb\). A small connected set \(X\) of fine
tiles has at most \(N_s\) tiles and its field support lies in an
enlargement of diameter \(O(b)\). Fix a boundary collar of a
sufficiently large constant number of fine layers. Declare \(X\)
interior if it meets no tile of this collar. Assume the ordinary
filtered quartic uses a finite-range filter of row sum one and
bounded absolute row sum on interior rows. There is a linear
projection on even Taylor jets through field order eight, with
an allocation to the labels in \(X\), such that:
\begin{enumerate}
\item The allocation is bounded in the radius-\(\ell_b\) jet norm
using the raw field space \(F_b\), uniformly in \(b\) and the cube.
It preserves the vacuum. At boundary placements it extracts the
full tensors of degrees two, four and six.
\item At interior placements, the quadratic residual has zero
evaluation on any pair of polynomial directions whose spatial
degrees sum to at most two. The quartic residual vanishes on
constant fields.
\item Suppose \(0<p<1\), \(2p>s-1\), and \(P\ge3\) in the field norm.
For the sum of the residuals of small marks with a fixed coarse
closure, the normalized zero-field jet norm is bounded by the
sum's maximal fine input norm times
\begin{equation}
 C\bigl(L^{s-1-2p}+L^{s-3}+L^{-1}+L^{-s}\bigr).
\end{equation}
Here and in this estimate the small-set threshold and all stencil
enlargements are fixed before \(L\).
\end{enumerate}
The allocation is covariant under translations and cube
symmetries respecting the collar. In the uniform bulk quadratic
potential its additional spatial-moment densities reduce to
finite-range kinetic densities and lattice divergences.
\end{proposition}

\begin{proof}
Write an even jet as \(f_0+f_2+f_4+f_6+f_8\), and write
\(f_2(\phi)=A(\phi,\phi)\) with \(A\) symmetric bilinear. Use
\(z_i(x)=(x_i-o_i)/b\) for coordinates on the fixed support
enlargement of a small set. On interior sets all these coordinates
and their products of degree at most two have \(F_b\)-norm bounded
by a constant. The quadratic moments to reproduce are
\begin{equation}
 A(1,1),\qquad A(1,z_i),\qquad
 A(z_i,z_j),\qquad A(1,z_i z_j)\quad (i\le j)
\end{equation}
with all \(i\) in the first-moment family. Moment matching in this
list does not depend on the choice of \(o\).

Let \(e_i\) be the lattice basis vectors, set
\(D_i\phi(x)=\tfrac b2(\phi(x+e_i)-\phi(x-e_i))\), and write
\(\langle\cdot\rangle_P\) for averaging over the sites of an
ordinary tile \(P\). On this tile use
\begin{equation}
 \begin{aligned}
 M_P(\phi)&=\langle\phi^2\rangle_P,&
 J_{P,i}(\phi)&=\langle \phi D_i\phi\rangle_P,\\
 Z_{P,ij}(\phi)&=\langle D_i\phi D_j\phi\rangle_P,&
 H_{P,ij}(\phi)&=\langle\phi D_iD_j\phi\rangle_P .
 \end{aligned}
\end{equation}
The inverse of their moment matrix is uniformly bounded. In
fact \(M_P\) has mass moment one. Each \(J_{P,i}\) has mass
moment zero and first-moment vector \(\tfrac12 e_i\). The
\(Z_{P,ij}\) and \(H_{P,ij}\) have both lower moments zero.
In the two remaining moment families, \(Z_{P,ij}\) gives the
elementary symmetric \((i,j)\)-matrix in \(A(z_k,z_l)\) and
zero in \(A(1,z_kz_l)\). The \(H_{P,ij}\) gives zero in the first
of these families and the elementary symmetric matrix in the
second. The nonzero diagonal entries in these last two blocks
are \(1\) or \(1/2\). All lower-order cross entries are bounded
averages of products of the \(z_i\). This proves the inverse
bound and gives, for each \(P\), a unique quadratic combination
reproducing the moments of \(f_2\). Its coefficient norm is
bounded by that of \(f_2\), because \(D_i\) and \(D_iD_j\)
are bounded evaluation operators on the interior stencils
in the fine \(F_b\) units.

Multiply this combination by \(|P|/\sum_{Q\in X}|Q|\)
and assign it to \(P\). With the same weights assign the
vacuum \(f_0\) and the quartic
\(f_4(1)\langle (T\phi)^4\rangle_P\), where \(T\) is the fixed
ordinary filter. The latter density equals one on the unit
constant field. Summing the allocations gives a projection
because it preserves all the input moments and depends only on
them. On a boundary set use the full homogeneous degrees
zero, two, four and six, dividing the tensor equally among
the tiles of \(X\) in the collar and retaining its actual
field stencils. Fine field restriction to the selected stencils
has uniformly bounded norm in a fixed enlargement. The
construction is symmetric in cases of ties between boundary
tiles. In the interior the moments and representing span
transform covariantly under coordinate reflections and
permutations, so uniqueness of the per-tile combinations
gives covariance of the allocation.

For the spatial bound, all ordinary supports assigned to a
fixed coarse closure have diameter \(O(b)\) each, and the
closure uses a bounded number of coarse tiles. Group the
supports in layers indexed by \(i\ge1\) with
\(\rho_B\asymp ib\) for \(i<L\), and put the region
\(\rho_B=B\) in the last layer. A layer below \(B\) contains
\(O(L^{d-1})\) fine tiles, including near edges and corners.
The interior part with \(\rho_B=B\) contains \(O(L^d)\)
tiles. Each fine tile belongs to only a bounded number of
small sets. For a coarse unit direction, its spatial Taylor
polynomial of degree two and the remainder on an ordinary
support have fine bounds, respectively,
\begin{equation}
 C(i/L)^p i^{-a}\quad(a=0,1,2),\qquad
 C(i/L)^p i^{-3}\quad\text{for the remainder}.
\end{equation}
To obtain these bounds take the discrete Newton polynomial
using \(\prod_j\binom{x_j-o_j}{\gamma_j}\) through total degree
two. Repeated telescoping on coordinate segments within the
enlargement bounds its remainder after \(|\beta|\le2\) differences
by \(C b^{3-|\beta|}\max_{|\gamma|=3}|\nabla^\gamma\phi|\).
Higher differences are those of \(\phi\) itself. Because the
fine weights here are one and the coarse depth on the segments
is comparable to \(ib\), these bounds give the asserted
\(F_b\) estimates. Polynomial terms of total degree at most
two cancel in testing the residual bilinear form. Every
remaining term costs \(C(i/L)^{2p} i^{-3}\).
The ratio of quadratic radii is \(L^{-s}\). Summing the
layer estimate gives \(C L^{d-1-s-2p}\), since
\(\sum_{i\ge1}i^{2p-3}<\infty\). The last layer costs
\(C L^{d-s-3}\). Both exponents have the stated form because
\(d=2s\). The quartic uses constant cancellation and costs
\(C(i/L)^{4p} i^{-1}\) before the radius factor \(L^{-d}\).
The resulting layer and interior sums are \(O(L^{-1})\)
for \(p>0\). Boundary placements have no residual degree
below eight. Everywhere, all residual degrees at least
six gain \(L^{-s}\) after summing \(O(L^d)\) terms.
This proves the zero-field bound.

Finally, the necessary extra densities can be retained literally
during the allocation. At common coefficients on identical tiles,
\(J_{P,i}\) and \(Z_{P,ij}+H_{P,ij}\) sum to zero in the bulk
by skew symmetry of \(D_i\). With the summation restricted
to the ordinary sites \(S\) the two identities are
\begin{equation}
 \begin{aligned}
 \sum_{x\in S}\phi(x)D_i\phi(x)
   &=\tfrac b2\sum_x \phi(x)\phi(x+e_i)
                       \bigl(\mathbf1_S(x)-\mathbf1_S(x+e_i)\bigr),\\[2pt]
 \sum_{x\in S}\bigl(D_i\phi(x)D_j\phi(x)
                      +\phi(x)D_iD_j\phi(x)\bigr)
   &=-\sum_x [D_i,\mathbf1_S]\phi(x)D_j\phi(x).
 \end{aligned}
\end{equation}
Their stencils meet the interface of the ordinary region
and the collar. For a fine tile on this interface, the
number of crossing bonds is \(O(b^{d-1})\), and the
commutator costs \(O(b)\). Before the tile normalization,
\(\phi\) and \(D_j\phi\) have bounded evaluation norms
at this positive distance from the physical boundary.
Consequently the corresponding shifted collar tensor has
bounded fine norm after dividing by the tile volume.
An allocation which agrees with the bulk outside the
fixed halo has its nonuniform coefficients supported in
the halo in addition to these interface terms. Thus the
reassignment can be made without identifying distinct
block densities pointwise.
\end{proof}

In particular, on a small set the allocation can be applied to
the finite jet quotient \(K(X)/I^X\). Suppose the normalized
jet of \(I(P)\) has vacuum bounded away from zero and, after
division by its vacuum, equals \(1+O(g)\). The finite jet
inverse has bounded norm and is analytic on the corresponding
complex reference box. If the reference is analytic there
with bounded first derivatives in the normalized potential
directions and \(\tau=\|K\|\) in the mixed norm, the allocation
satisfies
\begin{equation}
 \|q_P\|_{\ell_b,0}\le C\tau,\qquad
 \|D_Kq_P[\dot K]\|_{\ell_b,0}\le C\|\dot K\|,\qquad
 \|D_Vq_P[\dot V]\|_{\ell_b,0}\le C\tau\|\dot V\| .
\end{equation}
Indeed the first two bounds follow by composing the finite
jet inverse with the bounded allocations and summing the
bounded list of small sets through \(P\). The last follows
from \(D(I^{-1})=-I^{-1}(D I)I^{-1}\) in the same algebra.
Writing \(q_P=\sum_{X\ni P}q_{X,P}\), the usual reapportioning
\(J_P(X)=q_{X,P}\) for \(X\ne\{P\}\), adjusted at the singleton
to have zero sum through \(P\), is therefore exact. Changing
a block factor \(I\) to \(\widehat I\) in a circle product
uses the multiplicative difference with factors
\(I(P)-\widehat I(P)\). This identity follows by the binomial
expansion on the uncovered labels. It requires no pointwise
inverse of the reference.

\paragraph{Size of the coarse closure.}
For any connected fine set \(X\), with fixed bounded graph range
in fine units, its coarse closure satisfies
\begin{equation}
 |\overline X|\le C_d(1+|X|/L).
\end{equation}
Here the closure may include a fixed fine-stencil enlargement.
Color the coarse lattice modulo three in every coordinate
and choose a fine site of \(X\) within fixed fine distance
of each coarse tile counted. Two chosen sites of the same
color have distance at least \(cL\) for \(L\) large.
A spanning-tree walk in the fine graph visits every site in
at most \(2(|X|-1)\) steps, each of bounded length.
Consecutive first visits to the chosen sites of that color
therefore cost at least \(cL\) in fine distance. Summing over
colors proves the estimate. Choose \(N_s\) and then \(L\)
large. This gives \(|\overline X|\le\vartheta|X|\) with
\(\vartheta<1\) for \(|X|>N_s\), and hence the large-set gap
in the size exponent \((|X|-N_s)_+\). It does not use a
rectangular hull, which can expand a connected sparse fine
set by an arbitrarily large factor.

Proposition~\ref{prop:cube-jet-projection} also covers small marks
meeting several coarse tiles, since their coarse closure has bounded
size. Section~\ref{sec:ext-cov} supplies the Gaussian weights,
including \(1/2<p<1\) when \(d=4\). We now use the projection to
bound the term linear in one marked activity. To pass to the full
derivative we will also compare the reference factors on the
unoccupied labels.

\subsection{The cube single-mark estimate}
\label{sec:cube-mark}

A single \emph{mark} here is one insertion of a fine activity
\(F_X\). Its high Taylor remainder gains a power from each field
direction. For the lower coefficients, the gain comes from the
spatial moments removed by the projection. The proposition combines
these estimates with the quartic damping on the other labels,
including when the mark spans more than one coarse tile.

\begin{proposition}[Contraction of the leading small-mark operator]
\label{prop:cube-mark}
Use the raw test norms and adapted quotients in
Section~\ref{sec:geometry}, and choose \(0<p<1\), \(2p>s-1\). Let \(U\)
be the fixed coarse closure of a small fine mark. Run over all
small marks \(X\) with this closure. Let \(F_X\) be an even fine activity
on the enlarged stencils of \(X\), with field jets through a
fixed order \(M\ge9\) and fine mixed norm at most \(\nu\). Use
the growing small-field regulator at \(\ell_b=b^{-s/2}\), and the
decaying and quotient regulator at \(h_b\asymp g_b^{-1/4}\ell_b\),
with large-field weight \(g_b^{9/4}\). The coarse radii are
defined using \(B=Lb\) and \(g_B\asymp g_b\).
Assume a Gaussian exponential moment at radius \(\ell_b\)
on the fixed fine stencils, with a reserve for the growing
small-field regulator there. Put
\begin{equation}
 P_X=T_8\,\mathbb E F_X(\phi+\zeta),\qquad
 R_X=(1-T_8)\mathbb E F_X(\phi+\zeta),
\end{equation}
where \(T_8\) is the Taylor jet through degree eight.

Let \(I_X^{\rm comp}\) carry the stable filtered quartic on
the unoccupied fine labels, with positive real part of its
quartic coefficient comparable to \(g_b\). Use the recovery
filter of Section~\ref{sec:geometry}. One may include vacuum-
normalized perturbative factors on those labels formed from
quadratic exponent shifts of radius-\(\ell_b\) size \(O(g_b)\)
and polynomial factors \(1+W\) of degrees at most six with
\(\|W\|_{\ell_b,0}=O(g_b^2)\).
Set
\begin{equation}
 {\cal M}_{U}F=\sum_{X:\overline X=U}
 I_X^{\rm comp}\bigl((1-\Pi_X)P_X+R_X\bigr),
\end{equation}
with the projection of Proposition~\ref{prop:cube-jet-projection}.
For sufficiently large \(L\), small filter recovery error and
\(R/b\), and then small \(g_b\), this operator satisfies
\begin{equation}
 \|{\cal M}_{U}F\|_{\mathrm{mix},B}
 \le C\bigl(L^{s-1-2p}+L^{s-3}+L^{-1}+L^{-s}+L^{-s/2}
                    +C_Lg_b^{1/4}\bigr)\nu .
\end{equation}
The coefficient of the inverse powers of \(L\) is independent
of \(L\), and the final constant \(C_L\) may depend on \(L\).
\end{proposition}

\begin{proof}
A fixed \(U\) for a small mark has at most a bounded number of
coarse labels. There are \(O(L^d)\) placements, by the bounded
number of shapes through any fine tile. For \(C\in U\), write
\(Z_C=\|\phi\|_{2,C}/(B^{d/2}h_B)\) and
\(Q_C=q_{C^*}(\phi)/h_B\) for the background field quantities,
and put \(Z^2=\sum_C Z_C^2\), \(Q^2=\sum_C Q_C^2\).
The adapted quotient yields \(F_{B,C^*}(\phi)/h_B
\le C(Z_C+Q_C)\).

We first establish a uniform reserve from the complementary
reference. In the recovery inequality of
Section~\ref{sec:geometry}, the normalized occupied squared norm
in a coarse tile is now at most \(C L^{-d}(Z_C^2+Q_C^2)\), since
the mark contains only a fixed number of fine tiles.
Consequently it can be absorbed into the recovery error.
With \(S_C\) denoting the normalized filtered norm on the
unoccupied rows,
\begin{equation}
 (1-O(R/b+\epsilon^2+L^{-d}))Z_C^2
 \le C S_C^2 + C(R/b+\epsilon^2+L^{-d})Q_C^2,
\end{equation}
where \(\epsilon\) is the recovery error. The quartic is bounded
below in real part by \(c S_C^4\), since \(g_b B^d h_B^4\asymp1\).
Quartic Young's inequality then gives any fixed negative
quadratic exponent in \(Z_C\) with an arbitrarily small
positive quadratic exponent in \(Q_C\), by the choice of
\(L\) and recovery errors. Up to fixed field-jet order, the
quartic factor has derivatives equal to its exponential
times polynomials bounded by \(C(1+Z+Q)^C\) in coarse radius
units. This uses bounded filter row sums and
\(g_b B^d h_B^4\asymp1\). These polynomials are absorbed by a
fixed fraction of the reserve.

For the optional perturbative factors, each normalized
large-field jet of a quadratic exponent shift is bounded
by \(C g_b^{1/2}(1+Z+Q)^2\). Its exponential has this same
bound in the exponent, times a polynomial in the bound for its
nonzero jets. Each polynomial factor has
jet norm at most \(1+C g_b^{1/2}(1+Z+Q)^6\), by radius
conversion of its coefficient. For fixed \(L\) the number
of factors is bounded by a fixed \(C_L\). Their product,
including the tested derivatives, can use an arbitrarily
small fraction of the reserve with a bounded prefactor.
Indeed, for fixed degree \(a\),
\begin{equation}
 \sup_{t\ge0}\frac{\log(1+\delta(1+t)^a)}{1+t^2}\longrightarrow0
 \quad(\delta\to0).
\end{equation}
Apply this with the tolerance divided by the number of factors,
using the submultiplicative finite-jet seminorm for products.
This gives the large-radius bound by the coarse regulator
with spare fractions and a constant independent of \(L\)
after choosing \(g_b\). At coarse small radius and zero
field the vacuum-normalized complement differs from \(1\)
by \(O_L(g_b)\) in its tested jets.

On each fine mark \(\|P_X\|_{\ell_b,0}\le C\nu\), by
the small-field bound at \(\zeta\) and its exponential moment.
Thus the sum of the localized low jets is bounded by
Proposition~\ref{prop:cube-jet-projection} at zero field.
The varying complementary multipliers add only \(C_Lg_b\nu\)
there, using their difference from \(1\) and the product rule.
At the large radius, monomials of degree \(a\le8\) are
bounded by \(C_L\nu g_b^{-a/4}\) times a fixed polynomial
in \(Z,Q\). The reserve and coarse weight \(g_B^{9/4}\)
bound their contribution by \(C_Lg_b^{1/4}\nu\).

For \(R_X\), test the field derivative of order \(k\le8\)
by Taylor's integral remainder at \(\zeta+t\phi\), \(0\le t\le1\).
There are \(9-k\) background directions and \(k\) tested
directions. At \(k\ge9\) no Taylor subtraction remains.
The large-radius fine cost of each restricted direction is
\(C L^{-s/2}\), with background directions carrying additionally
a factor \(C(Z+Q)\). This uses the raw restriction inequality
on the enlarged stencils and the ratio \(h_B/h_b\).
The fine decaying regulator is at most one. For the quotient
at the translated field, use its triangle inequality, then
the Gaussian moment for the \(\zeta\) term. The background
term costs \(C L^{-s}Q^2\), because there are only a bounded
number of marked tiles and the coarse admissible affine fields
restrict to fine admissible affine fields. This exponent and
all fixed powers of \(Z+Q\) fit in the coarse reserve.
The Gaussian moment and this same envelope also justify
passing the tested derivatives through the expectation.
The input mixed cost \(\nu g_b^{-9/4}\) is canceled by the
output weight. Summing the remainder estimate over the
\(O(L^d)\) possible placements gives
\begin{equation}
 C L^{d-9s/2}\nu=C L^{-5s/2}\nu\le C L^{-s/2}\nu ,
\end{equation}
as required by the stated operator bound.
This argument includes derivatives of the product with the
complement, since the derivative of \(R_X\), of any order
used by the product rule, retains at least nine directions
from the Taylor formula or the test arguments alone.

At zero field \(R_X\) has no jets through order eight.
For each higher tested jet, use the fine large-radius input
bound at \(\zeta\), with scale ratio
\(\ell_B/h_b\asymp L^{-s/2}g_b^{1/4}\).
After multiplying by \(g_b^{-9/4}\), its degree \(k\ge9\)
again gives \(O(L^{-9s/2})\), before summation.
Finally the zero-field small-radius and full large-radius
bounds imply the full mixed bound. To see this for any field,
Taylor-expand each small-radius derivative through total
order eight using the zero-field jets. Control the integral
remainder by the large-radius bound, using
\((\ell_B/h_B)^9\asymp g_B^{9/4}\).
The resulting polynomial in the small-field background norm
and the quotient exponential with coefficient \(O(g_B^{1/2})\)
are absorbed by the growing small-field regulator.
\end{proof}

The annular fluctuation estimates and boundary weights give the
hypotheses of the single-mark bound in the adapted chart. To estimate
the difference between this operator and the derivative of the full
map, Proposition~\ref{prop:cube-bounded-extraction} accounts for the
terms with additional insertions and for the change of reference.
The latter comparison must hold also under variation of the boundary
and kinetic coefficients.

\subsection{Anchorwise covariance cancellation on the cube}
\label{sec:cube-covariance-cancel}

The projection assigns each local Taylor coefficient to its anchor.
A second-order covariance has two input anchors, and its subtraction
is performed before the localized polynomial is replaced by the
outgoing one. The next lemma assigns each contribution to an anchor
so that their sum reproduces the covariance, even for overlapping
field stencils.

\begin{lemma}[Oriented covariance allocation]\label{lem:cube-oriented}
Write a polynomial potential at each anchor \(i\) as \(U_i=A_i+B_i\).
Under \(P_C=\exp(\Delta_C/2)\), use as its two outgoing parts
\(P_C A_i,P_C B_i\), recording constants separately.
For any symmetric covariance parameter \(S\), define the polynomial
bilinear covariance
\begin{equation}
 \mathcal D_S(u,v)=P_S(uv)-(P_Su)(P_Sv),\qquad
 E^S_{ik}(U)=\tfrac12\mathcal D_S(A_i,A_k)
             +\tfrac12\mathcal D_S(B_i,B_k)+\mathcal D_S(A_i,B_k).
\end{equation}
Set \(X_{S,i}(U)=\sum_k E_{ik}^S(P_{-S}U)\).
Then for each individual anchor,
\begin{equation}
 P_C X_{S,i}(U)+\sum_k E^C_{ik}(U)=X_{S+C,i}(P_CU).
\end{equation}
In a coarse cell \(D\), the part of the fresh covariance assigned to
\(D\) is the sum of the \(i\)-rows for \(i\in D\).
For each pair of different coarse cells the two oriented allocations
sum to the full pair covariance. Both assertions hold before
estimating the polynomial remainders.
\end{lemma}

\begin{proof}
For any polynomial parts \(u,v\), commutation of the heat operators
gives
\begin{equation}
 \mathcal D_{S+C}(P_{-S}u,P_{-S}v)
 =P_C\mathcal D_S(P_{-S}u,P_{-S}v)+\mathcal D_C(u,v).
\end{equation}
It applies to each of the three summands defining \(E_{ik}\).
This proves the row identity. Bilinearity and symmetry give
\(E^C_{ik}+E^C_{ki}=\mathcal D_C(U_i,U_k)\) for \(i\ne k\) and
\(E^C_{ii}=\tfrac12\mathcal D_C(U_i,U_i)\).
Summation over the two disjoint sets of labels gives the asserted
pair allocations. When finite-range fluctuation variables on the
two actual field stencils are independent the corresponding term
is zero. The algebra uses the labels for the allocations and
allows overlapping stencils.
\end{proof}

For the filtered primary interaction one can put the uniform mass,
quartic, and first filtered quadratic head in \(A_i\), and the
ordinary kinetic density and fixed quadratic head in \(B_i\).
The local density of the first filtered head is the rowwise
difference of the filtered square and the uniform square. Thus
the quartic first heat splits into the uniform mass and that
same shape, at the same anchors. The mixed quartic--kinetic
two-line term is charged to the quartic anchor by the orientation
in Lemma~\ref{lem:cube-oriented}. In particular the position of
the surviving external filtered square is fixed even at a
coarse interface. This identity requires that the outgoing densities
on those anchors be kept exactly. The comparison of the resulting
polynomial with the outgoing reference is made in
Section~\ref{sec:cube-joint-reblocking}.

\begin{lemma}[Boundary quadratic generated by the quartic]
\label{lem:new-boundary-quadratic}
Write \(b=b_j\), \(B=Lb\), and \(G=g_j^{\rm loc}\) for the raw
quartic coefficient. Use the finite filter and frozen covariance of
Sections~\ref{sec:ext-cov} and~\ref{sec:bubble}. For a coarse tile
\(D\), let \(\mathcal T_{j,D}\) be the quadratic from the Gaussian
average of the quartic after assigning its common bulk term.
The contributing anchors lie within \(O(B)\) of the boundary, and
\begin{equation}
 \|\mathcal T_{j,D}\|_{\ell_B,0}
 \le C L^s\lambda_j^{\rm tad}|G|,
 \qquad
 \lambda_j^{\rm tad}=
 \begin{cases}
 O_L(\lambda),&\text{on an ordinary step},\\
 O_L(\lambda M_R^s),&\text{on the merged step}.
 \end{cases}
 \label{eq:boundary-tadpole-source}
\end{equation}
On the shooting box its value and its derivative in the normalized quartic coordinate
are bounded in the outgoing weight \(\Lambda_R x_{j+1}\), uniformly
in \(j\) and \(n\). Its normalized mass derivative is \(o(1)\) at fixed
\(L\) as \(R\to\infty\).
\end{lemma}

\begin{proof}
The heat of an incoming quadratic adds only a
constant. The heat of the quartic produces a
quadratic proportional to \(G\gamma(x)(T\phi)_x^2\), where
\(\gamma(x)=(T\Gamma_jT^*)_{xx}\) uses the actual finite filter.
The cube and bulk rows agree beyond the propagation range of this
slice and the filter stencils, hence outside a layer of width \(O(B)\).
To estimate the two diagonal entries, covariance Cauchy--Schwarz gives
\begin{equation}
 (T\Gamma_jT^*)_{xx}
 \le \left(\sum_y |T_{xy}|\sqrt{\Gamma_j(y,y)}\right)^2 .
\end{equation}
The filter has bounded absolute row sum for both the cube and the
bulk. Thus the diagonal bounds of Sections~\ref{sec:ext-cov}
and~\ref{sec:bubble} give \(C\lambda_j^{\rm tad} b^{-s}\) in both
cases. One may take \(\lambda_j^{\rm tad}=O_L(\lambda)\) at ordinary
scales. At the merged step the sum of the diagonals from scale \(R\)
to \(b_0\) is \(O_L(\lambda R^{-s})\), so
\(\lambda_0^{\rm tad}=O_L(\lambda M_R^s)\) suffices.

In the polarization of \(\mathcal T_{j,D}\), each actual or bulk
filtered field direction costs at most \(C\ell_B\) on the enlarged
stencil, since the raw test norm controls the supremum. This bound applies to each of
the two squares in a filter restriction difference and to the bare
square subtracted on a promoted anchor. There are at most \(CB^d\)
sites in \(D\), including at edges and corners. Their total cost is
at most \(C |G|\lambda_j^{\rm tad} b^{-s} B^d\ell_B^2\).
Since \(\ell_B^2=B^{-s}\) and \(d=2s\), this is
\eqref{eq:boundary-tadpole-source}.

The outgoing quadratic has weight \(\Lambda_R x_{j+1}\).
On the shooting box, \(|G|+|\partial_{\gamma_j}G|\le Cx_j\)
for the normalized quartic coordinate \(\gamma_j=(g_j-x_j)/x_j\).
At fixed covariance and filter the source is linear in \(G\).
Its value and its \(\gamma_j\)-derivative in the outgoing weight
are therefore bounded by
\begin{equation}
 C L^s\frac{\lambda_j^{\rm tad}}{\Lambda_R}
       \frac{x_j}{x_{j+1}}\le C_L .
\end{equation}
Here \(\lambda_j^{\rm tad}\le C_L\Lambda_R\), including the merged
step, and \(x_j/x_{j+1}\) is uniformly bounded. The source is independent
of the incoming collar tensors. Its dependence on the normalized mass
\(\mu_j=\nu_j/(\Lambda_R x_j)\) comes through \(G\).
Lemma~\ref{lem:cube-inverse-mass} gives
\(|\partial_{\nu_j}G|\le C\lambda x_j\). Thus its normalized
\(\mu_j\)-derivative is at most \(C_L\lambda_j^{\rm tad}\lambda x_j=o(1)\)
for fixed \(L\) as \(R\to\infty\). The full quadratic is retained on
its anchors in the outgoing boundary potential. At the following step
its \(O(B)\) support is a collar of the incoming scale \(B\).
\end{proof}

Contractions in the Gaussian integral lower the degree of a retained
collar tensor. The next estimate measures each such term in the
weight of its output degree, including a quartic contribution to
the quadratic coefficient.

\begin{lemma}[Heat and reblocking on the collar]
\label{lem:cube-boundary-heat}
Let \(s=d/2\), \(2p>s-1\), and let \(P_{U,k}\), \(k=2,4,6\),
be homogeneous polynomials on fixed enlarged stencils of fine tiles
\(U\) in a collar of width \(C_0 b\). Let \(a_k\) bound their
fine multilinear norms at radius \(\ell_b\). Suppose the fixed
fluctuation covariance has
\begin{equation}
 \sup_U\mathbb E\left(F_{b,U^*}(\zeta)/\ell_b\right)^q
     \le C_q\epsilon_j^{q/2}\qquad(q=2,4,6).
\end{equation}
Sum the heated polynomials with fine anchors in a given coarse
tile and remove the vacuum. The coarse norm of degree \(r=2,4,6\)
is at most
\begin{equation}
 C L^{d-1-r(s/2+p)}
              \sum_{k\ge r}\epsilon_j^{(k-r)/2}a_k .
\end{equation}
The same estimate bounds the operator on variations of the
incoming polynomials at fixed covariance. The outgoing nonvacuum
polynomials remain on stencils in a physical collar of width
\(O(b)\) covered by the enlarged coarse tiles.
\end{lemma}

\begin{proof}
In the collar both depth weights are comparable on a fixed
fine stencil, so
\(F_{b,U^*}(v)\le C L^{-p}F_{B,D^*}(v)\) on a coarse direction.
This follows by taking the ratio of the two weights in each
difference seminorm and using \(\rho_b\le\rho_B\le C\rho_b\).
The ratio \(\ell_B/\ell_b=L^{-s/2}\) contributes the remaining
factor for each coarse direction. For the symmetric \(k\)-linear
form of \(P_{U,k}\), the degree-\(r\) term of the heat substitutes
\(r\) coarse directions and \(k-r\) copies of the fluctuation
field. Its multinomial factor depends only on \(k\), and the
fluctuation copies cost at most \(C\epsilon_j^{(k-r)/2}\).
The odd terms vanish by centered Gaussian symmetry. The number
of fine anchors in the collar of a coarse tile is \(O(L^{d-1})\),
by covering the collar with the strips from the physical faces.
This count includes anchors at edges and corners. The heated
tensor has its original deterministic arguments, so its support
does not increase. Applying the same calculation to the variation
proves the operator bound.
\end{proof}

\begin{corollary}[Accumulation of the boundary quadratic]
\label{cor:accumulated-boundary-source}
Use the comparison sequence of Section~\ref{sec:path}. Give the
quadratic collar tensor weight \(\Lambda_R x_j\) and the quartic
and sextic tensors weight \(x_j^2\). Under the hypotheses of
Lemmas~\ref{lem:new-boundary-quadratic} and~\ref{lem:cube-boundary-heat},
suppose \(\epsilon_j=O_L(\lambda)\) on ordinary steps and
\(\epsilon_0=o(1)\) on the merged step as \(R\to\infty\) at fixed \(L\).
For sufficiently large \(L\) and then sufficiently large \(R\),
the linear transport of the inherited collar tensors is strictly
contractive in these weights. In particular, the sum of the
quadratics in \eqref{eq:boundary-tadpole-source} transported from
all preceding steps is bounded in units of \(\Lambda_R x_j\),
uniformly in the scale and in the side length of the cube.
\end{corollary}

\begin{proof}
For the original reference, the merged-scale stencil estimate in
Section~\ref{sec:ext-cov} gives
\(\epsilon_0\le C_L\lambda M_R^{2P+2p+s}(1+\log M_R)=o(1)\).
The smaller diagonal estimate \(O_L(\lambda M_R^s)\) suffices for
the newly generated term in \eqref{eq:boundary-tadpole-source}.
For inherited tensors the extra factor in the degree-six to degree-four block is
\(O(\epsilon_j)\), and in the degree-four to degree-two block
it is \(O(\epsilon_j x_j/\Lambda_R)\). Both are uniformly
small since \(x_j\le x_0\) and \(x_0/\Lambda_R=o(1)\). The diagonal ratios from successive comparison
couplings tend to one as \(R\to\infty\). Thus the linear transport
on these weights is strictly contractive by taking \(L\) large
and then \(R\) large. In particular the inherited quadratic has
diagonal factor \(C L^{s-1-2p}x_j/x_{j+1}\). Fix \(L\) so that
this factor is bounded by \(\kappa<1\) for all sufficiently large
\(R\). A source as in \eqref{eq:boundary-tadpole-source}, generated
at each step and then transported as an inherited quadratic, contributes
at most \(C_L\sum_{m\ge0}\kappa^m=C_L/(1-\kappa)\) in the normalized
quadratic coordinate.
\end{proof}

The quartic derivative of this source is a bounded entry from
the quartic to the stable coordinate in Assumption~\ref{ass:path}.
The radii and relative weights in the path contraction may be chosen
after \(L\) is fixed. A change of outgoing reference can introduce further coefficient
increments, for which this linear estimate gives the transport
only after the increments have been bounded in the same weights.

\subsection{A local estimate for the moving reference}
\label{sec:cube-moving-reference}

Extracting a Taylor polynomial from a polymer changes the reference
on each label to which a coefficient is assigned. The linear term in
this change is used in the extraction. The estimate below controls the
remainder at both field radii.

\begin{lemma}[Remainder in a local reference change]
\label{lem:cube-reference-change}
Let \(I=e^{-A}(1+W)\) be a normalized reference on one label, with
the stable filtered quartic and the allowed quadratic term in \(A\).
Assume \(0<g\le1\), \(\deg W\le6\), and
\(\|W\|_{\ell_b,0}\le Cg^2\). Let \(q_A,q_B\) be the extracted
polynomials, each of norm at most \(C\tau\) at \(\ell_b\), where
\(\tau=\|K\|\le c g^{9/4}\). On a boundary label \(q_A\) has degree
at most two and \(q_B\) has terms of degrees four and six. On an
ordinary label \(q_A\) has degree at most four, with quartic term in
the primary coefficient, and \(q_B=0\).
Suppose \(\deg\Delta W\le6\), \(\|\Delta W\|_{\ell_b,0}\le Cg\tau\)
and put
\begin{equation}
 \widehat I=e^{-A+q_A}(1+W+q_B+\Delta W),\qquad
 H=\widehat I-I-I(q_A+q_B).
\end{equation}
Assume the primary quartic remains in its stable sector along the
segment \(A-tq_A\), \(0\le t\le1\), with the field reserve of the
single-mark estimate or of the joint product containing the label.
Then the contribution of \(H\) to the mixed norm is at most
\(Cg\tau\). If the coefficients are holomorphic and satisfy these
bounds on the enlarged activity and potential domains, then on
fixed smaller concentric domains the activity derivative is bounded
by \(Cg\), and the derivative in a unit normalized potential direction
by \(Cg\tau\).
\end{lemma}

\begin{proof}
The exact difference is
\begin{equation}
\begin{aligned}
 H=e^{-A}\bigl[ &(1+W)(e^{q_A}-1-q_A)
              +(e^{q_A}-1)q_B\\
              &{}-Wq_B+e^{q_A}\Delta W\bigr].
\end{aligned}
\end{equation}
Apply the integral Taylor formula to the two differences of
exponentials. Each resulting exponential is evaluated on the given
segment, so the field reserve absorbs the polynomial background powers
and their fixed derivatives. At the small radius the resulting bound
is \(C(g\tau+\tau^2)\).

At the large radius \(h_b\asymp g^{-1/4}\ell_b\), a polynomial of
degree \(k\) costs at most \(Cg^{-k/4}\) times its small-radius
coefficient bound. In particular \(W\) costs \(Cg^{1/2}\), and
\(Wq_B\) and \(\Delta W\) together cost at most \(C\tau g^{-1}\).
The remaining terms contain either \(q_A^2\) or \(q_Aq_B\). Both
products have degree at most eight under the stated allocation, so
they cost \(C\tau^2g^{-2}\). Multiplication by the mixed weight
\(g^{9/4}\), and the bound \(\tau\le c g^{9/4}\), give \(Cg\tau\).
Cauchy's formula on a disc of radius comparable to \(g^{9/4}\) in
the activity variable gives the first derivative bound. Applying the
same formula at fixed \(K\) on the enlarged normalized potential
domain gives the second.
\end{proof}

We also use a Cauchy estimate to pass from the derivative at zero to
a nonlinear contraction. If the exact connected output \(\Psi_j(V,K)\)
is norm holomorphic for \(\|K\|<r=cg^{9/4}\) and satisfies
\(\|\Psi_j(V,K)\|\le C(g^3+\|K\|)\) there, two Cauchy integrals
give \(\|D_K^2\Psi_j(V,K)\|\le C'(g^3+r)/r^2\le C''/r\) on the
half-size ball. Integration along the segment from zero to \(K\)
therefore gives
\begin{equation}
 \|D_K\Psi_j(V,K)-D_K\Psi_j(V,0)\|\le C''\|K\|/r .
\end{equation}
A strict derivative bound at zero persists on a sufficiently small
concentric ball. The next two subsections establish the connected
size count and the estimates on its factors.

\subsection{The connected size sum on cube labels}
\label{sec:cube-linked-size}

The next lemma sums the configurations contributing to one connected
coarse polymer. Its input is a bound for each product after the
field regulators have been estimated. This distinction matters when
two factors use overlapping field stencils.

\begin{lemma}[Size estimate for a connected output]
\label{lem:cube-connected-size}
Write \(P(X)=(|X|-N_s)_+\). Consider configurations of coarse closure
\(U\) consisting of \(m\) disjoint connected fine inputs \(X_i\),
\(q\) local extraction covers and \(r\) ordinary decorations.
The latter two types have bounded fine support, a bounded number of
possible objects per fine centre and bounded multiplicity at a centre.
Take \(N_s\) sufficiently large, in particular larger than the size of
the coarse closure of any local cover, and then take \(L\) sufficiently
large. Set \(t=\sum_i |X_i|+q+r\).
Suppose the cost of a configuration at each field radius, before the
output mixed weight, is at most
\begin{equation}
 C_L^t a^{m+q}\epsilon^r
                \prod_{i=1}^m\varrho^{-P(X_i)}.
\end{equation}
At the small radius take \((a,\epsilon)=(\tau,C_Lg)\) and at the
large radius take \((a,\epsilon)=(\tau g^{-9/4},C_Lg^{1/4})\).
For each \(A>1\), choose \(\varrho\) sufficiently large and then
\(g\) and \(\tau g^{-9/4}\) sufficiently small. With the output
size weight \((A\varrho)^{P(U)}\), the mixed norm of the sum with
\(m+q\ge1\) is at most \(C_{L,A,\varrho}\tau\). The corresponding
sum with \(m+q=0\) and \(r\ge3\) is bounded by
\(C_{L,A,\varrho}g^3\). Both bounds are uniform in \(U\) and the
volume.
\end{lemma}

\begin{proof}
For \(m+q\ge1\), the union bound for the supports and the definition
of \(P\) give
\begin{equation}
 P(U)\le\sum_iP(\overline X_i)+N_s(m+q-1)+C r ,
\end{equation}
where \(C\) bounds the closure size of a decoration. For a connected
fine input \(X\), a traversal of a spanning tree visits all its
labels in at most \(2|X|\) steps. Each segment of at most \(L\)
steps meets only a bounded number of coarse tiles, since its diameter
is \(O(Lb)\). Thus \(|\overline X|\le C_d(1+|X|/L)\).
When \(|X|>2N_s\), the choices of \(N_s\) and \(L\) give
\(P(X)-P(\overline X)\ge c|X|\).
The resulting inverse power of \(\varrho\) pays for
\(A^{P(\overline X)}\) and an arbitrarily large fixed exponential
in \(|X|\) when \(\varrho\) is sufficiently large. Each smaller input
and each cover or decoration costs at most a fixed factor in the
size weights.

To count configurations with fixed \(U\) and \(t\), there are at most
\(C_L|U|\) fine positions for occupied sites and centres, counting
the allowed multiplicities. Since \(|U|\le Ct\), choosing these
positions costs at most \(C_L'{}^t\). Choosing an edge subset of
the induced bounded-degree graph specifies a partition into connected
inputs, at an additional cost exponential in the number of occupied
labels. The choices of objects at their centres satisfy the same
bound. Part of the exponential gain on large inputs pays these
costs; a further part sums their sizes. The remaining factors are
bounded by a fixed constant per input, cover or decoration. Keep
one factor \(a\) when \(m+q\ge1\); summation of the remaining factors
is geometric for small \(a,\epsilon\). At the large radius the
output multiplier \(g^{9/4}\) converts the retained \(a\) to \(\tau\).

For the pure configurations, \(|U|\le Cr\), so the count and the
weight cost at most a fixed exponential in \(r\). Keeping three
factors of \(\epsilon\) gives \(O(g^3)\) at the small radius and
\(O(g^{3/4})\) at the large radius. The large-radius multiplier
\(g^{9/4}\) completes the estimate.
\end{proof}

The pure terms with one or two decorations still require cancellation
on their assigned labels. Section~\ref{sec:cube-joint-reblocking}
proves the product estimate and performs this subtraction.

\subsection{Joint bounds for the reblocked reference expansion}
\label{sec:cube-joint-reblocking}

There are three steps in estimating the connected output.
Lemma~\ref{lem:cube-joint-reblocking} bounds products under the
Gaussian integral. An extraction cover may enlarge the connected
support, but only the originally occupied labels lose their primary
reference. This leaves quartic damping on all other labels, even
when two factors depend on the same field values.
Proposition~\ref{prop:cube-pure-cubic} uses this bound after the
first two orders of the reference have been subtracted.
Proposition~\ref{prop:cube-bounded-extraction} adds the terms
containing an activity and estimates the resulting derivative.

\begin{lemma}[Joint regulator for reblocking]
\label{lem:cube-joint-reblocking}
Use the adapted quotient and recovery filter of
Section~\ref{sec:geometry}. Suppose \(0<p<1\), \(2p>s-1\),
and \(g_B\asymp g_b=g\). Let the Gaussian fluctuation have range
\(O(B)\) and, on a fixed enlargement of each fine tile \(P\), satisfy
\begin{equation}
 Y_P=F_{b,P^*}(\zeta)/\ell_b,\qquad
 \sup_P\E\exp(aY_P^2/\sigma)\le C_0
\end{equation}
for a fixed \(a>0\). The enlargement may contain all the stencils of
a finite menu of local insertions at \(P\). Let \(U\) be a set of
coarse labels, each the disjoint union of its fine labels. On
pairwise disjoint originally occupied fine polymers use their
large-radius activity bounds. On the other fine labels use the
stable quartic with real row coefficient at least \(c g\).
Individual factors may be evaluated at
\(\phi+t\zeta\), \(0\le t\le1\), with different parameters on
different factors. The quartic exponent may also be a rowwise
convex combination of two such quartic exponents. Allow quadratic
exponent perturbations of small-radius jet size \(C g\) per label,
and a bounded number of polynomial factors per fine label.
Assume each polynomial insertion, after factoring out its stated
coefficient cost, has fixed degree in the normalized coarse field
and in the \(Y_P\) on its stencil. The same bounds are to hold
through the fixed tested field-jet order of the mixed norm.

Choose \(L\) large, then the recovery and strip errors small,
and then the upper bounds on \(\sigma\) and \(g\) small.
The expected product in the coarse large-radius jets is bounded
by the product of the component and insertion costs times
\begin{equation}
 C_L^{|U|}
 \exp\left[-2\kappa\sum_{C\in U} Z_C^2+
                      \tfrac12 c_h\sum_{C\in U}Q_C^2\right],
 \qquad
 Z_C=\frac{\|\phi\|_{2,C}}{B^{d/2}h_B},
 \quad Q_C=\frac{q_{C^*}(\phi)}{h_B}.
\end{equation}
The analogous zero-background bound at \(\ell_B\) has prefactor
\(C_L^{|U|}\) and uses the fine small-radius activity bounds and
insertion costs. A factor \(Y_P^k\) in an insertion gives the
additional cost \(C_{L,k}\sigma^{k/2}\) when retained in the bound.
The moment and reserve hold uniformly over the segment parameters.
\end{lemma}

\begin{proof}
Let \(A_C\) denote the normalized \(L^2\) norm of the
background on the occupied sites of \(C\), and \(S_C\) the
normalized \(L^2\) norm of the unoccupied filtered rows in their
shifted backgrounds. For a rowwise convex combination apply Jensen's
inequality to the real quartics, using the lower bound on their
coefficients, so that the resulting effective shift is again in
\([0,1]\). The local recovery estimate in
Section~\ref{sec:geometry} gives, with
\(\delta=C(R/b+\|I-B_{\mathrm{rec}}T\|_{2\to2}^2)\),
\begin{equation}
 (1-\delta)Z_C^2\le A_C^2+C S_C^2+\delta Q_C^2+
        C\|T\zeta\|_{2,C}^2/(B^d h_B^2).
\end{equation}
Indeed, the exceptional strips around all fine label boundaries
have relative volume at most \(CR/b\), even when the occupied
set is disconnected. The error of the inverse filter is estimated
on a fixed enlargement of \(C\), by its operator norm and
\(F_{B,C^*}(\phi)/h_B\le C(Z_C+Q_C)\). The fluctuation term,
summed on the coarse tiles, is bounded by \(C_L\sum_PY_P^2\)
using the bounded filter row norm.

On occupied labels, \(\|\phi+t\zeta\|_{2,P}^2
\ge \|\phi\|_{2,P}^2/2-\|\zeta\|_{2,P}^2\).
Their damping supplies a negative coefficient of order
\(\kappa L^s\) on \(\sum_C A_C^2\). The triangle inequality
for each adapted quotient and the summed restriction estimate give
the background quotient cost
\(C c_h L^{s-1-2p}\sum_C Q_C^2\). The remaining occupied
fluctuation costs are bounded by \(C_L\sum_P Y_P^2\).
On each coarse tile the stable quartic is bounded below by
\(c S_C^4\), by Cauchy--Schwarz and \(g B^d h_B^4\asymp 1\).
It supplies any prescribed negative multiple of \(S_C^2\) with
a constant cost per coarse tile. Take the coefficients on \(A_C^2\)
and \(S_C^2\) large enough to get strictly more than \(2\kappa\)
on \(Z_C^2\). By the choice of \(L\) and then \(\delta\),
the positive coefficient on \(Q_C^2\) is strictly below
\(c_h/2\). A complex quadratic term of normalized size \(C g\)
uses at most \(C_L\sqrt g\sum_C (Z_C^2+Q_C^2)
+C_L\sum_PY_P^2\) in the exponent. Its background part fits in
the strict reserves for small \(g\).

The fine radius of a restricted coarse unit direction is bounded
by \(C_L\), at both radii. By the product rule and bounded degrees,
the polynomial and field-jet factors, after taking out their costs,
are bounded by a polynomial in \(Z_C,Q_C,Y_P\) per fine label,
of degree bounded uniformly per label. There are at most \(C_L\)
fine labels in a coarse tile and the enlargement graph has bounded
degree. Thus the background polynomials are absorbed by the
remaining quadratic reserves, at fixed cost per coarse tile.
Equivalently, one can first use the submultiplicative finite-jet
seminorm with factorial weights, and then absorb each local
polynomial. An activity on a large polymer uses its tested
seminorm on that entire polymer. It contributes no extra degree
per site, and the direction restriction bound still holds.

Color the fine stencil graph by grid residue classes with spacing
larger than the covariance range, increased by twice the stencil
radius. Within each of the \(\chi=O_L(1)\) colors the Gaussian
vectors are independent. For bounded \(k_P\), Hölder and the
single-tile moment give
\begin{equation}
 \E\prod_P Y_P^{k_P}e^{C_LY_P^2}
 \le C_L^{\#\{P\}}\sigma^{\sum_P k_P/2}.
\end{equation}
In fact factor out the powers of \(\sqrt\sigma\), absorb the
rescaled polynomial by \(\exp(aY_P^2/(2\chi\sigma))\), and
use \(\chi C_L\le a/(2\sigma)\). The number of fine labels in
the Gaussian bound is at most \(C_L|U|\), proving the asserted
prefactor. At zero background the stable quartic has nonnegative
real part, and all derivatives, quadratic terms and growing
activity regulators have the same Gaussian majorant. This proves
the small-radius bound. The bounds also justify Gaussian
integration of the tested derivatives.
\end{proof}

\paragraph{Centered differences and the pure reference.}
After centering, a first-order quartic difference contains at most
three background fields. At the large radius this degree bound
leaves enough powers of the coupling for a cubic remainder.
The subtraction uses the exact assignment to anchors in
Lemma~\ref{lem:cube-oriented}.

\begin{proposition}[Cubic residual for a bounded anchored split]
\label{prop:cube-pure-cubic}
Normalize the incoming vacuum. Write its reference as a product
over fine labels with factors \(I_i=e^{-U_i}(1+W_i)\), where
\(U_i\) is the stable quartic on its rows and a quadratic term.
Assume \(\|U_i\|_{\ell_b,0}\le C g\) and
\(\|W_i\|_{\ell_b,0}\le C g^2\), with \(\deg W_i\le6\).
A bounded product of such polynomial factors is also allowed
with \(W_i\) their sum to first order in their coefficients.
Put
\begin{equation}
 P_i=\E\theta U_i,\qquad
 R_i=\E\theta W_i+\sum_k E^\Gamma_{ik},\qquad
 J_i=e^{-P_i+D_i}(1+R_i-D_i),
\end{equation}
where \(\theta\) is Gaussian translation and \(E^\Gamma_{ik}\)
is the oriented row of Lemma~\ref{lem:cube-oriented}.
Suppose \(D_i\) has degree at most four, size \(C_Lg^2\) and
quartic part proportional to the primary rows on \(i\).
Constants can equivalently be removed on the labels.
Use the joint bound of Lemma~\ref{lem:cube-joint-reblocking}.
Then on zero independent activity the reblocking expansion about
\(J_i\), followed by the zero-sum covariance reapportioning, has
connected remainder \(O_{L,\varrho}(g^3)\) in the mixed norm,
with a strict coarse field reserve. The bound holds on the
enlarged complex coefficient domain whenever the displayed
size and sector assumptions hold there.
\end{proposition}

\begin{proof}
Use \(\xi_i=\theta U_i-P_i\) and put \(\gamma=g^{1/4}\) for
the large-radius estimate. In normalized fine small-radius jets
\(\xi_i\) has cost \(C_Lg\). In large-radius jets it has cost
\(C_L\gamma\) under the product bound. In fact
\begin{equation}
 \xi_i(\phi,\zeta)=
 [U_i(\phi+\zeta)-U_i(\phi)]
      -\E[U_i(\phi+\zeta)-U_i(\phi)].
\end{equation}
Each monomial of the difference has at least one fluctuation
direction measured at \(\ell_b\), leaving at most three powers
of the large background or tested directions at \(h_b\).
The coefficients of the centered expectation follow by the
same Gaussian moments. The translated \(W_i\) and all
\(R_i,D_i\) have small-radius costs \(C_Lg^2\) and
large-radius costs at most \(C_L\gamma^2\): their degree is
at most six, including for the covariance of two quartics,
since every surviving term has a cross pairing. The constants
are uniform in scale because the number of fine labels
interacting with an anchor is \(O_L(1)\).

Taylor's integral formula for the exponential, used along the
segment between \(P_i\) and \(\theta U_i\), gives
\begin{equation}
 \theta I_i=e^{-P_i}(1-\xi_i+\tfrac12\xi_i^2+\theta W_i)
                 +\mathcal E_i,
 \qquad J_i=e^{-P_i}(1+R_i)+\mathcal F_i.
\end{equation}
The errors have costs \(C_L g^3,C_L\gamma^3\) at the two
radii. For the first identity the errors contain either three
\(\xi_i\), one \(\xi_i\) and one \(W_i\), or two polynomial
factors of order \(g^2\). For the second, use the segment
between \(P_i\) and \(P_i-D_i\) and keep two factors of order
\(g^2\). The exponent remains stable on each segment, and
Lemma~\ref{lem:cube-joint-reblocking} applies after multiplying
by the other factors. In particular, a difference
\(\Delta_i=\theta I_i-J_i\) has the leading term
\(-e^{-P_i}\xi_i\) and a residual of costs \(C_Lg^2,C_L\gamma^2\).
These conclusions also follow directly from the first-order
integral formula when no cancellation is being used.

For one difference, its expected coefficient relative to the
common stable factor, through order two, is
\begin{equation}
 \tfrac12\E\xi_i^2+\E\theta W_i-R_i
    =-\sum_{k\ne i} E^\Gamma_{ik}.
\end{equation}
For two differences it is \(\E\xi_i\xi_k
=E^\Gamma_{ik}+E^\Gamma_{ki}\). The error in the latter identity
for the actual product keeps one first-order factor and one
second-order factor or two second-order factors. In either
calculation the deterministic perturbative fillers may be
replaced by their common stable parts for the displayed
second-order terms. A filler difference has second-order cost,
so the resulting error still has at least cubic order.

For a coarse cell \(C\), the pairs internal to \(C\) cancel
the internal singleton row contributions. Between two coarse
cells \(C,D\) assign to the ordered cover \((C,\{C,D\})\)
the sum \(\sum_{i\in C,k\in D}E^\Gamma_{ik}\), and to its
singleton the negative sum of those rows. The row has zero
sum for each \(C\). Using the circle identity for that row
cancels the displayed second-order coefficient on each
connected output polymer; the sign is selected to subtract
that coefficient. The replacement of the stable filler by
the perturbative reference adds a product of two second-order
factors. Multiplicative factors commute on labels, so this
subtraction uses the original assignments even on stencils
shared by two labels.

All terms with at least three fine differences have at least
three factors of \(g\) at small radius and \(\gamma\) at large
radius. The product bound and the pure connected count of
Lemma~\ref{lem:cube-connected-size} sum them with exponential
size reserve. The one- and two-difference errors are on
bounded-size coarse polymers with \(O_L(1)\) choices, hence
have the same estimates. In the exact zero-sum reapportioning
all further terms have either two inserted covariance rows
or a row and an already estimated residual. The same-scale
connected count pays the finite menu of row covers: they
have uniformly bounded size and the extra \(g^2\), respectively
\(\gamma^2\), pays their size weight. Existing residuals use
a fixed part of their stronger exponential weight for the
partition count. At common scale, products of the residual
bounds on disjoint occupied labels concatenate their
regulators, and the stable reference on uncovered labels
supplies any fixed quadratic damping, by the local recovery
with \(b=B\) and the small strip error. The spare field reserve
pays for the polynomial insertions, of bounded multiplicity
per label. Including the output weight \(g^{9/4}\) in the
large-radius bound gives the asserted \(g^3\).
\end{proof}

\paragraph{Conversion to the outgoing reference.}
A local comparison follows directly once an outgoing polynomial
representative is available on the same anchors. Suppose the
outgoing exponential is the literal sum of \(P_i-D_i\) on a
coarse label, after recording the vacuum. On a boundary label
this formula can be implemented by retaining the full quadratic
in the inner collar and using the common scalar quartic
coefficient on its actual filter rows. Assign the difference
of that coefficient from the incoming one to \(D_i\) and
leave the remaining quartic and sextic terms in \(R_i-D_i\).
This is admissible when the coefficient difference is
\(O_L(g^2)\). The sum of the remaining polynomials can be
stored as a single factor per collar, while factors on the
ordinary margin use their original anchored formulas.

In multiplying the fine \(1+(R_i-D_i)\) factors, the linear
terms sum to the retained polynomial. Every discarded product
contains at least two degree-at-most-six factors of size
\(O_L(g^2)\), hence has costs \(C_Lg^4\) and \(C_Lg\) at the
two radii under the reserve. If a bounded bilinear ordinary
reference is recomputed after an \(O_L(g^2)\) change of its
potential, the coefficient difference is \(O_L(g^3)\) of
degree at most six. Its large-radius cost is \(C_L g^{3/2}\).
Thus the literal cellwise reference error is \(O_L(g^3)\)
in the mixed norm. The estimate extends to moved inputs
when their reference sizes and the anchor equality hold on
the enlarged extraction disc. It uses neither a pointwise
inverse of a reference polynomial nor cancellation between
different coarse labels.

\paragraph{The localization--Gaussian commutator.}
On a single mark, the linearized extraction contains a difference
between applying the local projection before and after Gaussian
integration. With \(\nu\) its mixed norm, the Gaussian shift estimate
on each zero-field derivative of order \(k\le 6\) is
\begin{equation}
 \bigl|D^k(\E F_X(\zeta+\,\cdot\,) -F_X)(0)\bigr|_{\ell_b}
       \le C \sqrt{\sigma}\,\nu .
\end{equation}
Here and below the norm denotes the supremum over normalized
test directions. Integrate the \((k+1)\)-st derivative along
\(t\zeta\), use the fine growing regulator, and retain the
factor \(Y_P\) in the Gaussian bound on the bounded collection
of mark stencils. The same shift estimate holds for the
polynomial \(\Pi_X T_8 F_X\), by its bounded coefficients.
The projection uses only degrees through six, so
\begin{equation}
 \|\Pi_X T_8\E F_X(\zeta+\,\cdot\,)
              -\E(\Pi_X T_8 F_X)(\zeta+\,\cdot\,)\|_{\ell_b,0}
       \le C\sqrt{\sigma}\,\nu .
\end{equation}
The resulting polynomial has degree at most six. After insertion
in the complementary reference, passage to the large radius
costs at most \(g^{-3/2}\), which is paid by the mixed weight.
Summing over the \(O(L^d)\) small-mark placements gives
\(C_L\sqrt{\sigma}\,\nu\) in the coarse mixed norm. This
estimate includes the merged step whenever its raw
fluctuation amplitude obeys the stated moment assumption.

We now combine the commutator estimate with the connected count.
The hypotheses below specify the comparison of reference factors
on a single coefficient domain. The resulting derivative estimate
is the one used in the shooting argument.

\begin{proposition}[Extraction disc with a bounded outgoing split]
\label{prop:cube-bounded-extraction}
Suppose the fine reference and the moved reference on
\(\tau=\|K\|\le c g^{9/4}\) have the bounds of
Proposition~\ref{prop:cube-pure-cubic} on an enlarged complex
coefficient domain. Suppose their outgoing splits satisfy the
literal matching and the bounded bilinear ordinary-reference
comparison above on that domain. If extraction affects the
ordinary reference through such a bilinear factor, use its
difference bound \(C_Lg\tau\). Include among the boundary
labels a fixed halo containing the nonbulk supports. Use the
same nested grid and definitions of ordinary data in the
finite-volume and bulk maps.

For small \(c\) the connected output \(\Psi(V,K)\) of the
extraction, reblocking, covariance redistribution and outgoing
reference change is norm-holomorphic and satisfies
\begin{equation}
 \|\Psi(V,K)\|\le C_{L,\varrho}(g^3+\|K\|).
\end{equation}
At zero activity its derivative differs from the leading
localized small-mark operator by
\(C_{L,\varrho}(g^{1/4}+\sqrt\sigma)+\varepsilon_{\varrho,L}\),
where the large-set term \(\varepsilon_{\varrho,L}\) tends to
zero as \(\varrho\to\infty\) for fixed sufficiently large \(L\).
The extraction change of each full boundary tensor is
\(O_L(\tau)\). The extraction change of each ordinary margin
polynomial is \(O_L(g\tau)\). Their derivatives on smaller
domains obey the same bounds in a unit normalized potential
direction and, with \(\tau\) replaced by \(\|\delta K\|\), in
an activity direction \(\delta K\).
\end{proposition}

\begin{proof}
On each originally uncovered fine label \(i\), insert the
constant \(1-\sum_{X\ni i}J_i(X)=1\) from the zero-sum
extraction, with \(J_i(X)\) assigned to the cover \(X\).
For this formula only, denote the perturbative filler by
\(J_i^{\mathrm{pt}}\) and the moved incoming reference by
\(\widehat I_i\). After Gaussian translation the resulting
uncovered factor is exactly
\begin{equation}
 \begin{aligned}
 \theta I_i\Bigl(1-\sum_{X\ni i}\theta J_i(X)\Bigr)
 ={}&J_i^{\mathrm{pt}}
   +(\theta\widehat I_i-J_i^{\mathrm{pt}})
   +\theta(I_i-\widehat I_i)
   -\sum_{X\ni i}(\theta I_i)\theta J_i(X).
 \end{aligned}
\end{equation}
In each selection group the covers and the originally
occupied polymers into their coarse connected components.
Inside a selected cover away from its selected centre and
the occupied labels, the first three terms sum to
\(\theta I_i\); summing them first implements the original
reference fill on that cover. Outside the selected covers,
the same three choices give the ordinary multiplicative
reference change and reblocking. Gaussian expectation
factors on separated coarse components by the range
condition. Thus the expansion has the exact circle
product of the original input. Its occupied-label
damping comes from the original activities, and every
other fine label keeps one stable exponential, possibly
evaluated along a segment.

The local extraction costs are \(C_L\tau\) and
\(C_L\tau g^{-3/2}\). For the cover terms this is the
degree-six bound on the extracted jet. For the reference
differences use a coefficient segment with a stable
primary quartic and one differentiated polynomial or
exponent, of respective degrees at most six and four.
The ordinary differences have costs \(C_Lg\) and
\(C_Lg^{1/4}\), by the centered expansion. All insertions
have bounded multiplicity on a fine stencil because their
centres do and their assigned sizes are bounded. Lemma~
\ref{lem:cube-joint-reblocking} proves the joint product
estimate required by Lemma~\ref{lem:cube-connected-size}. At the large radius each fine component has coefficient
cost at most \(\tau g^{-9/4}\) times its inverse size weight, and each extraction mark costs an additional
\(g^{3/4}\) relative to this bound. The prefactor \(C_L^{|U|}\)
is allowed since the number of coarse labels is at most a
fixed multiple of the total occupied size and number of
decorations. This gives \(C_{L,\varrho}\tau\) for all
connected terms with an activity or extraction mark. The
pure terms at each fixed extracted reference are handled
by Proposition~\ref{prop:cube-pure-cubic}. At common scale
the covariance and outgoing-reference changes use only
bounded-size covers and retain the same strict field
reserve as in that proposition, with a fixed part of the
stronger size reserve. Their pure reference errors have
mixed size \(O_L(g^3)\). This gives the claimed bound.

For the derivative, hold the extracted reference fixed
while assigning each term to its independent marked
component. The derivative through the pure remainder at
the extracted reference is \(C_{L,\varrho}g^{3/4}\), by
Cauchy's formula on the disc of size \(c g^{9/4}\) and its
pure \(O(g^3)\) bound. A linear term with an ordinary
decoration retains \(g^{1/4}\) at the large radius and a
factor \(g\) in its small-radius jets. Later same-scale
decorations have second-order cost or smaller. A single
large mark without decorations has the large-set size
gap and gives the stated \(\varepsilon_{\varrho,L}\).

The remaining first variation is a single small mark.
At the fine scale the zero-sum extraction and the
singleton change have the linear combination
\(F_X-I^Xq_X\), where \(q_X=\Pi_X(F_X/I^X)\) in the
finite jet algebra. Indeed, the nonsingleton row is
\(-I^X\sum_{i\in X}q_{X,i}\). On a singleton \(i\) the
negative singleton row and the reference change combine
to \(-I_i q_{\{i\},i}\). The remainder from the literal
reference change has derivative \(O_L(g)\) in the mixed
norm, by Lemma~\ref{lem:cube-reference-change}. In the
small-radius jets, replacing \(I^Xq_X\) by \(\Pi_XT_8F_X\)
costs \(O_L(g\|F_X\|)\) since the normalized reference
is \(1+O(g)\). At the large radius both inserted
polynomials cost at most \(C_Lg^{-3/2}\|F_X\|\) before
the mixed weight, with the reserve supplied by the
small-mark complement. This gives an error bounded by
\(C_Lg^{3/4}\|F_X\|\) there. The localization--Gaussian
commutator then replaces the incoming projection by the
projection after integration at cost
\(C_L\sqrt{\sigma}\|F_X\|\). The surviving operator is
the one in Proposition~\ref{prop:cube-mark}, with the
unoccupied perturbative factors as complement. Its local
quadratic and polynomial perturbations satisfy the sizes
in that proposition. The small-mark argument needs a
quartic-only complement, so all perturbations are
absorbed after \(L\) is fixed and \(g\) is small.

For the retained differences, the jet inverse gives
\(C\tau\) per extracted fine coefficient. Gaussian heat
on a bounded collection of stencils is a bounded linear
operator on their finite jet norms. The covariance
correction is bilinear, so its change is bounded by
\(C_L(g\tau+\tau^2)\). The complete degree-two, four and
six differences in the inner collar are finite sums
of these coefficients. A margin has the ordinary
bilinear polynomial, and its outgoing ordinary potential
varies by at most \(C_L\tau\) under the bounded split.
This gives the additional \(g\) there. A boundary
variation does not change an ordinary factor or an
ordinary margin read outside its support and
propagation neighborhood. On an enlarged domain all
these differences are holomorphic, including their
vacuum normalization. A one-variable Cauchy integral
in each normalized potential direction preserves the
factor \(\tau\). On a smaller activity ball a Cauchy
integral uses numerator \(O_L(g^{9/4})\) and denominator
of order \(g^{9/4}\), preserving the extra \(g\) for
margin differences.

The finite algebraic maps are analytic on their domains.
Gaussian integration admits the uniform reserve majorant
on every smaller complex ball. Cauchy coefficients of
the integrated field jets are estimated using the same
mixed norm bounds, so their series converge in the
output norm. The connected counts give uniform bounds
in the volume. This proves norm holomorphy and justifies
the activity derivatives used above.
\end{proof}

Iteration also requires the retained coefficients to lie in the
allowed domain at the next scale. For \(d=1,2,3\),
Proposition~\ref{prop:cube-retained-chart} checks that transport.
Its first input is a uniform bound for the quadratic and quartic
contractions on an ordinary stencil.

\paragraph{The ordinary second-order polynomial for \(d=1,2,3\).}
On an ordinary stencil all coefficients are translation invariant
and the covariance and filter agree with the bulk ones. This makes
it possible to bound the cumulative reference with constants
independent of the scale. In the three-line quadratic contraction
the relevant spatial cancellation is evenness in the internal
displacement. A first-order absolute difference would be
insufficient for this purpose in dimension three.

\begin{lemma}[Bilinear ordinary reference bound below \(s=2\)]
\label{lem:cube-ordinary-W}
Let \(d=2s\), \(s<2\), and use the normalized, even bulk finite
filter \(T\) with row sum one. Let \(S\) be a cumulative bulk
covariance supported at distance \(O_L(b)\), with, for the orders
used below,
\begin{equation}
 |\nabla^k S(x)|\le C_{L,k}\lambda(R+|x|)^{-s-k}.
\end{equation}
Use the exponentially weighted absolute row bounds of \(T\)
at length \(R\). Write \(u_x=(T\phi)_x\), and take the ordinary
potential on anchor \(x\) to be a linear combination of
\begin{equation}
 u_x^4,\qquad b^{-s}\phi_x^2,\qquad
 b^{2-s} k_x(\phi),\qquad b^{2-s}R^{-2}(u_x^2-\phi_x^2),
 \qquad
 k_x(\phi)=\tfrac14\sum_{|e|_1=1}(\phi_{x+e}-\phi_x)^2 .
\end{equation}
Write \(G\) for the sum of the absolute coefficients, with the
coefficient of the last density weighted by
\(1+\log(2+b/R)\). In the oriented covariance row of
Lemma~\ref{lem:cube-oriented}, put the third density in its
second type and the other three in the first type. Define
\(\mathrm{Loc}\) to keep the vacuum, the constant-field
quartic, and the mass and kinetic moments of the bulk
quadratic, represented on the same anchors by the displayed
vacuum, filtered quartic, mass and nearest-neighbour kinetic
densities. Then
\begin{equation}
 \bigl\|(1-\mathrm{Loc})X_{S}(U)\bigr\|_{\ell_b,0;P}
       \le C_L(\lambda+\lambda^2+\lambda^3)G^2
\end{equation}
on each ordinary tile \(P\). The bound is bilinear, including
first variations with this strengthened coefficient norm.
The remaining polynomial has degree at most six.
\end{lemma}

\begin{proof}
Set \(P_1=TS\) and \(D=TST^*\). They obey the same
bounds and range as \(S\), by the absolute exponential filter
moments. The shell sums give
\begin{equation}
 \|S\|_1+\|P_1\|_1+\|D\|_1\le C_L\lambda b^s,\qquad
 \|\nabla P_1\|_2^2\le C_L\lambda^2/R^2 .
\end{equation}
When \(\phi\) is tested on this stencil, a spatial difference of
order \(k\) costs \(C_L b^{-k}\ell_b\). An even filter of mass
one thus has \((T-I)\phi=O_L((R/b)^2\ell_b)\), also after
any fixed tested derivative of order \(k\), with another
factor \(b^{-k}\). This follows by subtracting the constant and
linear terms before using the second absolute moment of the filter.
Convolutions commute with these differences in the ordinary chart.

The Gaussian Wick product identity for \(P_{-S}\) in the oriented
row leaves precisely the contractions with at least one cross
pairing between the two vertices. It follows, for instance, by
applying \(P_S\) to two separately Wick-ordered polynomial
generating exponentials and expanding the cross exponential.
For two quartics the one-cross term has norm \(C_L\lambda G^2\),
using \(b^d\ell_b^6=b^{-s}\). In the two-cross term,
subtract the constant-field representative \(u_x^4\) at the
anchor \(x\). The remaining external-leg factor is bounded
by \(C_L\ell_b^4\min(1,|r|/b)\) in multilinear tests.
Summing over the internal displacement \(r\) gives
\(C_L\lambda^2\) after the factor \(b^d\ell_b^4=1\):
\begin{equation}
 \sum_r |D(r)|^2\min(1,|r|/b)\le C_L\lambda^2 .
\end{equation}
The three-cross term at \(x\) contains \(D(r)^3u_xu_{x+r}\).
Replace \(u_{x+r}\) in the even row by the average of
\(u_{x+r}\) and \(u_{x-r}\). On subtracting \(u_x^2\),
the external-leg bound is now \(C_L\ell_b^2 |r/b|^2\).
Since
\begin{equation}
 \sum_r |D(r)|^3 |r|^2
 \le C_L\lambda^3 b^{2-s}\qquad (s<2),
\end{equation}
multiplication by \(b^d\ell_b^2/b^2=b^{s-2}\)
gives \(C_L\lambda^3\). For replacing the filtered on-site
square by its mass representative, use
\(\sum_r|D(r)|^3\le C_L\lambda^3 R^{-s}\)
and the filter difference \((R/b)^2\).
The normalized kinetic moment of the quadratic has
the same bound, since it is bounded by
\begin{equation}
 C_L b^{s-2}\sum_r |D(r)|^3(|r|^2+R^2)
 \le C_L\lambda^3 .
\end{equation}
The \(R^2\) term accounts for the external filter legs.
The symbol computation of these moments uses translation
invariance and cubic symmetry. Their nearest-neighbour
representatives have the same local norm bound.

For clarity, let \(Q\) denote the Hessian of any sum of the
displayed quadratic densities with its stated coefficients.
The mass kernel is of size \(G b^{-s}\), and the other kernels
are even zero-row kernels with second absolute moment at most
\(C G b^{2-s}\). Hence \(Q\phi\) costs
\(C_L G b^{-s}\ell_b\) on the ordinary stencils, by subtracting
the constant and linear terms. Its fixed differences of
order \(k\) satisfy the same estimate times \(C_{L,k}b^{-k}\).
For a one-cross quartic--quadratic contraction rooted at the
quartic, the contracted direction is \(TSQ\phi\), bounded
by \(C_L\lambda G\ell_b\). At a mass root use \(P_1\) on
the smooth external cubic from the quartic. At a structured
head root the expression with that smooth field \(v\)
has the form \(u_x(TSv)_x-\phi_x(Sv)_x\).
Its decomposition
\begin{equation}
 (u_x-\phi_x)(TSv)_x+\phi_x((T-I)Sv)_x
\end{equation}
recovers the factor \((R/b)^2\) on the head coefficient.
The difference-square root in a mixed quartic--kinetic
contraction does not occur, by orientation. These bounds,
with \(b^d\ell_b^4=1\), control the quartic and its
constant-field projection.

For the two-cross quartic--mass contraction the row is a
multiple of \(P_1(r)^2\) times the filtered square at the
quartic anchor. When the root is the mass anchor, moving that
square to its root and replacing it by the mass density
costs \(C_L\ell_b^2(\min(1,|r|/b)+(R/b)^2)\).
The first factor eliminates the marginal logarithm of the
absolute row, and the second times \(1+\log(2+b/R)\) is
bounded. The difference-square root is again avoided by
orientation in the kinetic contraction. At the quartic root
the row is exactly the filtered square there times a
coefficient bounded by
\(C_L G^2 b^{2-s}\sum_r|\nabla P_1(r)|^2\).
Subtracting the mass and kinetic representatives costs
\(C_L(R/b)^2\) in its normalized coefficient bound, which
cancels the factors from the kinetic input and the derivative
row. These statements also control the kinetic subtraction:
with both surviving legs at a quartic row, the relative
second moment of the external kernel is \(O(R^2)\) times its
scalar row coefficient, independently of the displacement of
that row from the selected anchor.

For the head--quartic two-cross term put
\(F(r)=D(r)^2-P_1(r)^2\). Even filtering and the second
difference estimate give
\begin{equation}
 |D(r)-P_1(r)|
     \le C_L\lambda R^2(R+|r|)^{-s-2}.
\end{equation}
One obtains this by using the second difference for shifts
of size at most \((R+|r|)/2\) and the exponentially
weighted absolute row bound on the other shifts. Thus
\(\sum |F|\le C_L\lambda^2\) and
\(\sum |r|^2|F(r)|\le C_L\lambda^2R^2[1+\log(2+b/R)]\).
A head root requires moving the square to its root. Pairing
\(r\) with \(-r\) makes this cost \(C_L\ell_b^2|r/b|^2\).
The filter and kinetic comparison cost
\(C_L\ell_b^2(R/b)^2\) against the unweighted row.
The resulting normalized bounds use precisely the
strengthened head coefficient. At a quartic root there
is no need to move the square.

For the quadratic--quadratic one-cross term the summed
partner contributes \(SQ\phi\), with test norm
\(C_L\lambda G\ell_b\). A mass root uses its coefficient
\(G b^{-s}\). At a difference-square root the external
field and \(SQ\phi\) each supply a difference at scale
\(b^{-1}\). At a head root use the preceding
\((T-I)\) decomposition with \(v=Q\phi\). All three
cases thus give \(C_L\lambda G^2\) after multiplication
by \(b^d\ell_b^2\). The mass and the relative second
moment of their spatial sum are bounded by differentiating
\(\widehat Q_1\widehat S\widehat Q_2\) at zero, using the
row and second moment estimates for \(S,Q_1,Q_2\).
There is only a vacuum from all higher cross orders
of these contractions. This exhausts the retained degrees
and proves the stated bound, with bilinearity giving the
same estimates on polarized coefficients.
\end{proof}

For this ordinary coefficient space, Lemma~\ref{lem:cube-oriented}
gives the literal heat split on each anchor. If \(X_i^+\) is
the anchored \(X_{S+\Gamma,i}(P_\Gamma U)\), set
\begin{equation}
 D_i=(\mathrm{Loc}X^+)_i
             -P_\Gamma(\mathrm{Loc}X_S(U))_i.
\end{equation}
Then \(R_i-D_i=(1-\mathrm{Loc})X_i^+\) for the incoming
polynomial \(W_i=(1-\mathrm{Loc})X_{S,i}(U)\). The heat of the
quartic contributes the fixed filtered square, which is the
sum of the chosen mass and head densities on that anchor.
The heat of a quadratic contributes only a vacuum. These
identities make the split use the same anchored mass, quartic,
kinetic and head densities on the two scales. In particular
\(D_i=O_L(g^2)\) wherever the strengthened coefficients of
\(P_\Gamma U\) are \(O_L(g)\), since \(R_i\) has that bound
and the outgoing polynomial is bounded by the lemma.
If the literal outgoing potential differs from \(P_\Gamma U\)
by \(O_L(g^2)\) in the strengthened norm, recomputation
of its ordinary polynomial costs \(O_L(g^3)\), with
the corresponding bilinear variation bound. This proves
the ordinary comparison for changes within this coefficient
space. Additional localized divergences or changes of
boundary representatives must still be accounted for
by their actual allocations.

\paragraph{Exact removal of extracted quadratic divergences.}
Ordinary localization produces second-moment densities that are
lattice divergences. To remove them without changing the linearized
activity map, choose the ordinary kinetic representative proportional
to
\begin{equation}
 k_x^{\rm c}(\phi)
   =\tfrac14\sum_i(\phi_{x+e_i}-\phi_{x-e_i})^2 .
\end{equation}
The preceding proof applies to this representative as well. Its
difference-square span is two, and its bulk quadratic symbol has
leading Taylor term \(|k|^2\), up to the fixed normalization of
the coefficient.

In a translation-invariant cubically symmetric ordinary region,
the extracted coefficients of \(J_{P,i}\) and all off-diagonal
\(Z_{P,ij},H_{P,ij}\) vanish. The diagonal \(Z\) and \(H\)
coefficients are independent of \(i\). Their contribution
can therefore be written as a multiple of the central kinetic
density and a multiple of \(\sum_i(Z_{P,ii}+H_{P,ii})\).
The latter can be carried by the boundary collar. In fact
the commutator identity in Proposition~\ref{prop:cube-jet-projection}
writes its sum over ordinary sites as a sum of local interface
tensors. Each interface tensor has normalized fine coefficient
at most \(C\tau\), by the interface-bond count in that proof.
The same reasoning applies to a first-moment divergence on
any symmetrically assigned exceptional region. Keep the
nonuniform extracted coefficients in the fine halo literally.

To implement this change at the level of all circle products,
write the literal extracted and the new reference factors
as \(I_i=e^{-\Delta_i}\widehat I_i\), with the same polynomial
factors, where \(\Delta_i\) is the quadratic density removed on
the ordinary labels and its negative is inserted on the halo
as given by the commutator identity. Then \(\sum_i\Delta_i=0\)
as a polynomial. For every occupied set with components \(X\),
\begin{equation}
 \prod_{i\ {\rm uncovered}} I_i\ \prod_X K(X)
 =
 \prod_{i\ {\rm uncovered}}\widehat I_i\
   \prod_X\left(e^{\sum_{i\in X}\Delta_i} K(X)\right).
\end{equation}
This holds on the actual fine labels with their stencils,
without dropping local divergences from any single label.
On each label the newly introduced quadratic has
small-radius size \(C\tau\). In the fused reblocking
expansion it uses the same quadratic reserve as in
Lemma~\ref{lem:cube-joint-reblocking}, including when the
label belongs to an originally occupied component.
In particular this step need not map an unrestricted
fine activity into the identical fine regulator before
Gaussian integration.

Write the correction on each occupied component by integrating
its exponential difference:
\begin{equation}
 e^{\sum_{i\in X}\Delta_i}-1
  =\int_0^1 e^{t\sum_{i\in X}\Delta_i}
            \sum_{i\in X}\Delta_i\,dt .
\end{equation}
There is one additional local insertion marked at \(i\).
Its costs are \(C\tau\) and \(C\tau g^{-1/2}\) at the two
radii. The choice of \(i\) adds only a factor \(|X|\),
paid by the stronger size count on large \(X\), or by a
fixed constant on small \(X\). The factor it multiplies
already contains an independent activity or a local
extraction mark. Thus this exact removal changes no
first variation at \(K=0\). Its connected bound is included
in the enlarged-disc estimate. The ordinary bilinear
polynomial can then be computed in the kinetic
representative \(k^{\rm c}\), because the uniform divergence
has been removed before that polynomial is formed.

\subsection{The inverse bulk mass chart on the cube}
\label{sec:cube-mass-chart}

The mass and kinetic terms in the second-order covariance polynomial
can be included in the definitions of the bulk couplings. At the
terminal scale we vary the resulting mass coordinate and need its
derivative in the original local potential.

Let \(W_j^{\rm cov}\) be the cumulative bulk covariance before step
\(j\). Define its normalized row sum and second moment by
\begin{equation}
 a_j=b_j^{-s}\sum_x W_j^{\rm cov}(x),\qquad
 m_j^{(2)}=b_j^{-s-2}\sum_x x_1^2 W_j^{\rm cov}(x).
\end{equation}
Summing the slice \(\ell^1\) bounds as a geometric series gives
\(|a_j|\le C\lambda\). Their range \(O(b_j)\) gives
\(|m_j^{(2)}|\le C\lambda\). At the terminal scale the cumulative
covariance excludes the remaining tail and satisfies the same bounds.

\begin{lemma}[Compensated bulk coordinates and their inverse]
\label{lem:cube-inverse-mass}
Suppose the bulk covariance has cubic symmetry and the quadratic
Taylor coefficient of the local kinetic symbol is \(|k|^2\). Set
\begin{equation}
 \begin{aligned}
 \nu_j &= \nu_j^{\rm loc}+a_j(\nu_j^{\rm loc})^2,\\
 g_j &=g_j^{\rm loc}(1+4a_j\nu_j^{\rm loc}),\\
 z_j &=(1+2a_j\nu_j^{\rm loc})z_j^{\rm loc}
               -\tfrac12m_j^{(2)}(\nu_j^{\rm loc})^2 .
 \end{aligned}
\end{equation}
These changes compensate the second-order cumulative mass--mass,
mass--kinetic and mass--quartic projections in the uniform bulk
coordinates. On the complex box \(|\nu_j^{\rm loc}|\lesssim |g_j|\)
with \(\lambda |g_j|\) sufficiently small, the branch of the inverse
through zero local mass is holomorphic. Holding \(g_j,z_j\), the
covariance and the separate tensor coefficients fixed gives
\begin{equation}
 \partial_{\nu_j}\nu_j^{\rm loc}=1+O(\lambda |g_j|),\qquad
 |\partial_{\nu_j}g_j^{\rm loc}|\le C\lambda |g_j|,\qquad
 |\partial_{\nu_j}z_j^{\rm loc}|
       \le C\lambda (|\nu_j^{\rm loc}|+|z_j^{\rm loc}|).
\end{equation}
\end{lemma}

\begin{proof}
For a quadratic Hessian \(Q\) the covariance polynomial has quadratic
part \(\tfrac12\langle Q\phi,W_j^{\rm cov}Q\phi\rangle\).
Its constant and quadratic Fourier coefficients are those of
\begin{equation}
 \bigl(b_j^{-s}\nu_j^{\rm loc}
       +b_j^{2-s} z_j^{\rm loc}|k|^2+O(|k|^4)\bigr)^2
                \widehat W_j^{\rm cov}(k).
\end{equation}
Cubic symmetry gives the coefficient of \(|k|^2\) in
\(\widehat W_j^{\rm cov}\) as
\(-\tfrac12\sum_x x_1^2 W_j^{\rm cov}(x)\).
In mass and kinetic units the two coefficients are therefore
\(a_j(\nu_j^{\rm loc})^2\) and
\(2a_j\nu_j^{\rm loc}z_j^{\rm loc}
-\tfrac12m_j^{(2)}(\nu_j^{\rm loc})^2\).
The one-line mass--quartic contraction gives
\(4a_j\nu_j^{\rm loc}g_j^{\rm loc}\), as in
Section~\ref{sec:bubble}. Heating the extracted uniform quadratic
changes only its constant term. Thus the changes absorb the
differences of the cumulative projections to second order.
The separate zero-row terms retain their own polynomial increments.

Write \(u=\nu_j^{\rm loc}\). The inverse through zero for
\(\nu_j=u+a_j u^2\) has derivative
\(\partial_{\nu_j}u=(1+2a_j u)^{-1}\).
The remaining two equations are linear in \(g_j^{\rm loc}\)
and \(z_j^{\rm loc}\), respectively. Differentiating them gives
\begin{equation}
 \partial_{\nu_j}g_j^{\rm loc}
   =-\frac{4a_jg_j^{\rm loc}}{1+4a_j u}\,\partial_{\nu_j}u,\qquad
 \partial_{\nu_j}z_j^{\rm loc}
   =\frac{m_j^{(2)}u-2a_jz_j^{\rm loc}}{1+2a_j u}\,
      \partial_{\nu_j}u .
\end{equation}
The denominators are uniformly bounded away from zero by the
hypothesis on the mass box. The moment bounds then give the three
derivative estimates.
\end{proof}

Thus the inverse mass variation reproduces the uniform mass to
leading order, also in \(d=4\) with a small marginal kinetic
coefficient. The finite-volume boundary tensors are assigned
separately in the outgoing reference. Their bounds and those of
the activity enter the terminal large-field comparison.

\subsection{An invariant retained-coefficient chart in dimensions one through three}
\label{sec:cube-retained-chart}

To apply the preceding activity estimate in \(d=1,2,3\), we now specify
the outgoing coefficients on the cube. An ordinary site retains its
bulk density even when it lies in the outer part of a cell meeting
the boundary. In the inner collar the full nonuniform quadratic
is retained. When a bulk difference-square density enters this
collar, its contribution is thus measured in the boundary tensor
norm. The same assignment enters the coefficient recursion and
the change of reference in the activity map.

On an interval, a layer at a fixed distance from the boundary has at
most two sites. Each coarse interval meets \(O(1)\) fine intervals
in the old collar and contains \(O(L)\) fine intervals altogether.
Take, for example, \(p=1/8\) at \(s=1/2\). The old quadratic collar
tensor then has transport factor \(C L^{-3/4}\) by
Lemma~\ref{lem:cube-boundary-heat}, while the bulk kinetic
coefficient transports with \(L^{-3/2}\). A newly generated
tadpole may occupy \(O(Lb)\) sites. Its norm is bounded by
\(C|G|\lambda_j^{\rm tad}b^{-1/2}(Lb)\ell_{Lb}^2
=C L^{1/2}|G|\lambda_j^{\rm tad}\) by
\eqref{eq:boundary-tadpole-source}. On the next step it belongs to
the old collar. The initial boundary head costs
\(C\lambda G_0 M_R^{-3/4}\) by
Proposition~\ref{prop:initial-cube-heads}. For the ordinary quadratic
projection there is just one coordinate \(z_1\), so the four moments
are \(A(1,1),A(1,z_1),A(z_1,z_1),A(1,z_1^2)\).
The centered difference and its square give the same invertible
moment matrix as in Proposition~\ref{prop:cube-jet-projection}.
Each interface contributes \(O(1)\) crossing bonds to its commutator.
These counts also apply to the terminal interval when a leftover
piece is merged with the last full cell.

\begin{proposition}[Retained coefficient estimates below the local endpoint]
\label{prop:cube-retained-chart}
Let \(d\in\{1,2,3\}\) and \(\max\{0,(s-1)/2\}<p<s/2\),
with \(p\ne1/2\). Use the frozen finite-range
covariance, the initially truncated filter, the joint initialization,
and the comparison sequence of Sections~\ref{sec:ext-cov},
\ref{sec:cube-joint-initialization} and~\ref{sec:path}. Suppose
the activity agrees with its bulk comparison off the grouped boundary
labels, and norm the finite-volume and bulk activities by the maximum
of their mixed norms. There is a single outgoing chart with ordinary mass, quartic,
local kinetic, and fixed-shape quadratic coefficients, full collar
tensors of degrees two, four and six, and ordinary margin factors.
In this chart the ordinary and collar reference comparisons required
in Proposition~\ref{prop:cube-bounded-extraction} hold on its enlarged
extraction disc. On the shooting box, the retained coefficient map
obeys the quartic, mass and stable rows of
Assumption~\ref{ass:path}, including their normalized first
variations. Boundary comparisons preserve equality with the bulk off
the new grouped boundary labels.
\end{proposition}

\begin{proof}
Use \(b=b_j\), \(B=Lb\). Away from the physical boundary use the site
densities of Lemma~\ref{lem:cube-ordinary-W}, with the central kinetic
representative in the subsequent divergence identity. Include also
the fixed initial zero-row head \(H_0\) of
Proposition~\ref{prop:initial-cube-heads}, in difference-square form.
Use the notation \(Z_T=T^*T-I\) for the structured head, with \(T\)
normalized by its bulk row sum. Its anchored density is
\((T\phi)_x^2-\phi_x^2\). The actual filtered rows, and the uniform
mass density, are shared by the ordinary and the collar factors.
The other quadratic terms inside the collar are represented by a full
quadratic \(Q_{\partial,2}\). The additional terms of degrees four
and six there are stored in \(1+B_{\partial,4}+B_{\partial,6}\).

The inner collar at scale \(B\) consists of the fine-layer
anchors within a fixed multiple of \(B\) of the physical boundary.
This accommodates the covariance range \(O(B)\): in the scalar
polynomial construction the largest time in an ordinary slice
is at most a constant times the coarse length in units of \(R\).
Choose the collar width with a fixed margin over this range and
the fine-scale localization and support enlargements. The
incoming collar is correspondingly within a fixed multiple of
\(b\), so only \(O(L^{d-1})\) old collar labels are reblocked
in a coarse collar tile. New promotions may use all \(O(L^d)\)
fine labels of a coarse tile. They are sources from the incoming
ordinary coefficients, with constants allowed to depend on \(L\).
The stencils used by every new inner factor lie within a fixed
enlargement of its coarse cell and within depth \(O(B)\).
All remaining anchors in that cell use the ordinary factors,
including their polynomial corrections. This set of ordinary
anchors is the margin of the cell. Their polynomial corrections
can be added on the cell, with the cost of the omitted products
as in Section~\ref{sec:cube-joint-reblocking}. At initialization
the collar contains the finite boundary head and the bulk difference
squares with anchors inside it. The ordinary second-order
polynomial at that scale is zero. The cutoff defining which fine
anchors enter the outgoing collar can be rounded to fine layers
with no change in these bounds. For sides not divisible by the
block length, start the grid at the same corner at every scale.
On each coordinate axis merge the leftover interval with its last
full interval and use the Cartesian products. The merged interval
at the fine scale lies within the merged interval at the coarse
scale, since the earlier coordinate of the coarse merged interval
is a fine grid coordinate still to its left. This gives nested
labels whose volumes and stencil enlargements are comparable to
those of ordinary cubes.

We first check the ordinary comparison in this chart. The bulk
Fourier estimate in Section~\ref{sec:bubble} with an extra multiplier
bounded by \(|k|^r\) gives
\(\|\nabla^r\Gamma_i^{\rm bulk}\|_\infty
 \le C_r\lambda b_i^{-s-r}\) for each fixed order needed in the
comparison. On the low-frequency ball the extra multiplier is
handled by scaling \(Rk/\sqrt{r_i}\), where
\(r_i=(R/b_i)^2\) is the spectral scale in that section. On the
complement, the fixed gap gives the factor \(r_i^{2m-1}\) of the
higher-resolvent estimate. The observation weight has integral
with the extra factor \(|k|^r\) bounded by \(C_r R^{-d-r}\).
Indeed the product formula for the spread-out box gives
\(\lvert\widehat D_R(k)\rvert\le C(1+R|k|_\infty)^{-1}\);
the excluded zero site contributes only \(O(R^{-d})\).
Choose the fixed smoothing power larger than \(d+r\).
The white channel contributes \(O(e^{-R})\), which obeys the
same bound. Increasing the fixed higher-resolvent order makes
the large-frequency part no greater than the claimed slice bound.
Finite propagation then gives the cumulative bounds in
Lemma~\ref{lem:cube-ordinary-W}, by summing scales of size at least
a constant multiple of the observed displacement. This includes the
unmerged pieces of the first covariance.

The ordinary bilinear estimate allows \(H_0\) in the second part
of the orientation in Lemma~\ref{lem:cube-oriented}. For this
extension use the size
\begin{equation}
 z_{H_0}(b)=b^{s-2}\sum_y |y|^2|H_0(y)|
\end{equation}
in the coefficient norm. At each difference-square root, the
external and contracted legs contribute a difference over the same
span \(y\), so the product costs \(C|y|^2/b^2\) in the ordinary
field norm. Summing with \(|H_0(y)|\) gives exactly \(z_{H_0}(b)\).
For the mixed quartic--head term with two cross pairings, the
orientation charges it to the quartic root. Both exterior legs
there form the same on-site filtered square. The coefficient is
bounded by
\begin{equation}
 C \sum_y |y|^2 |H_0(y)|\,\|\nabla(TS)\|_2^2 ,
\end{equation}
where \(S\) is the cumulative covariance. This follows by applying
the two endpoint differences to the two contracted legs and using
\(\|f(\cdot+y)-f\|_2\le C |y|\|\nabla f\|_2\).
The subtraction of the on-site mass and kinetic moments uses
\((R/b)^2\), as in Lemma~\ref{lem:cube-ordinary-W}.
These are all changes involving the fixed head. They also apply to
its variation by polarization.

Write \(G\) for the raw quartic, \(u=b^{-s}\nu^{\rm loc}\) for the
raw mass row, \(S_j^0=\sum_y S_j(y)\) for the row of the cumulative
covariance, \(S_j\) for its kernel, and
\(\gamma_j=(T\Gamma_jT^*)(0)\) in the bulk. In particular
\(3B^s|\gamma_j|\le C_L\lambda_j^{\rm tad}\),
where \(\lambda_j^{\rm tad}\le C_L\Lambda_R\).
The first heat gives the raw mass row
\(u_1=u+3\gamma_jG\) and the structured term with Hessian
\(3\gamma_jG Z_T\), on the same anchors. Use
\begin{equation}
 D_i=(\mathrm{Loc}\,X_{S_j+\Gamma_j}(P_{\Gamma_j}U))_i
                    -P_{\Gamma_j}(\mathrm{Loc}\,X_{S_j}(U))_i
\end{equation}
on the outgoing ordinary anchors. This is the polynomial subtracted
from the first-heat potential. The constant terms are recorded in
the vacuum. The ordinary projection keeps scalar quartic, mass and
kinetic coefficients. There is one further structured density
in \(D_i\) because the localized quartic in
\(\mathrm{Loc}\,X_{S_j}(U)\) heats to a filtered square.
If \(a_{4,j}\) is the coefficient of that localized quartic in
the normalization \(G(T\phi)^4/4\), the extra contribution to
the outgoing quadratic has Hessian
\(3\gamma_j a_{4,j}(Z_T-\kappa_T K_{\rm loc})\).
Here \(K_{\rm loc}\) is the kinetic representative with
second Taylor coefficient \(|k|^2\) and \(\kappa_T\) is that
coefficient for \(Z_T\) in the same units.
Its mass and second Taylor coefficients vanish. This formula follows
by subtracting the scalar mass and kinetic projections from the heat
of \(a_{4,j}(T\phi)^4/4\). In particular the actual outgoing
structured coefficient is updated by first heat and this explicit
correction alone.

Put \(\theta=L^{s-2}<1\) and choose \(a_1\) with
\(\theta<a_1<1\), also \(L^{-2}<a_1\). Give the normalized
structured head the weight \(h_j=\Lambda_R x_0 a_1^j\); the fixed
initial head can use the same weight. Give the local kinetic and
full collar quadratic the weight \(\Lambda_Rx_j\), and
\(B_{\partial,4},B_{\partial,6}\) the weights \(x_j^2\).
On fixed boxes in these weights,
\begin{equation}
 h_j[1+\log(2+b_j/R)]\le C_Lx_j .
\end{equation}
To check this, \(b_j/R=M_RL^j\), the reciprocal sequence estimate
implies \(x_0/x_j\le C(1+\lambda^2x_0j)\), and
\(\Lambda_R(1+\log(2+M_R))=o(1)\). The factor \(a_1^j\)
absorbs any fixed power of \(1+j\). The ordinary bilinear estimate
is therefore \(O_L(x_j^2)\) on the common enlarged box,
including the variation in these coordinates.

For the structured source itself the constant-field quartic
projection gives, by the two-line quartic and the one-line
quartic--mass contractions,
\begin{equation}
 a_{4,j}=9G^2\sum_y (TS_jT^*)(y)^2+4Gu S_j^0.
\end{equation}
Every other quadratic vertex has zero row and contributes zero to
this quartic projection. The cumulative bubble bound gives
\(\sum_y(TS_jT^*)(y)^2
 \le C_L\lambda^2(j+\log(2+M_R))\), by summing its increments.
With \(|S_j^0|\le C_L\lambda b^s\), this gives
\(\lvert a_{4,j}\rvert\le C_L x_j\) on the enlarged box, and
the same bound on first variations in a unit normalized quartic
direction. In a unit normalized mass direction the bound is
\(C_L\lambda\Lambda_R x_j^2\), by differentiation of this formula
and the inverse mass chart. The old structured density transports
with \(\theta\). Both its first-heat and corrected source are
bounded in outgoing units by
\begin{equation}
 C_L\lambda_j^{\rm tad}(R/b_j)^2
      (|G|+|a_{4,j}|).
\end{equation}
Dividing by \(h_{j+1}\) works uniformly since
\begin{equation}
 \frac{\lambda_j^{\rm tad}(R/b_j)^2}{\Lambda_R a_1^{j+1}}
 \le C_L M_R^{-2}(L^{-2}/a_1)^j .
\end{equation}
The normalized stable diagonal from the fixed-head transports is
\(\theta/a_1<1\). The outgoing structured change from first heat
is included in \(P_{\Gamma_j}U\) when comparing its polynomial
with the outgoing ordinary polynomial. The additional \(D_i\)
correction has size \(O_L(x_j^2)\) even in the log-strengthened
norm: before the logarithm is applied it is bounded by
\begin{equation}
 C_L\lambda_j^{\rm tad}(R/b_j)^2
   [\lambda^2(j+\log(2+M_R))|G|^2+\lambda |G\nu^{\rm loc}|] .
\end{equation}
The geometric decay of \((R/b_j)^2\) absorbs the additional
logarithm. All other changes of ordinary representatives from
\(P_{\Gamma_j}U\) are \(O_L(x_j^2)\) in the strengthened norm by
the bounded projection and covariance estimate. Recomputing the
ordinary second-order factor therefore costs \(O_L(x_j^3)\).

For completeness the scalar quartic projection gives the raw update
\begin{equation}
 G_+-G=-\mathfrak b_jG^2-4G(u_1S_{j+1}^0-uS_j^0).
\end{equation}
Indeed the quartic part of the localized row of
\(X_{S_j+\Gamma_j}(P_{\Gamma_j}U)\) is the formula for \(a_{4,j}\)
at the new covariance and with \(u_1\). Heat preserves the quartic
part of the old localized row, and the difference of the two bubbles
is \(\mathfrak b_j/9\). In compensated coordinates
\(g=G(1+4uS_j^0)\), \(g_+=G_+(1+4u_+S_{j+1}^0)\) this implies
\begin{equation}
 g_+-g+\mathfrak b_jG^2
 =4S_{j+1}^0\bigl[(G_+-G)u_1+G_+(u_+-u_1)\bigr].
\end{equation}
There is no estimate here on the unprojected size of the cumulative
bubble; only its increment is used. The right side is \(O_L(x_j^3)\)
since \(B^s|u_+-u_1|+|G_+-G|\le C_Lx_j^2\).
Replacing \(G^2\) by \(g^2\) costs \(O_L(x_j^3)\).
The same bounds give
\begin{equation}
 |\nu_+-L^s\nu-3B^s\gamma_jg|
       +|z_+-L^{s-2}z|\le C_L x_j^2
\end{equation}
in the compensated mass and local kinetic chart. Its quadratic
compensation has uniformly bounded coefficients by
Section~\ref{sec:cube-mass-chart}. Since the first filtered tadpole
was assigned to \(Z_T\), it does not contribute to \(z_+\).
All these error bounds hold on the enlarged boxes, so Cauchy's
inequality gives the stated orders in a unit normalized potential
direction.

We next describe the outgoing split on the inner collar. Use the
same oriented covariance row and the sum of the actual incoming
polynomials on each of its site or tensor anchors. A fine
correction may be assigned to the collar whenever one of its
stencils meets the region of disagreement. By finite propagation
all such anchors are in the chosen new inner collar.
In \(D_i\) keep the constant and full quadratic part of \(R_i\)
there, and the unique primary quartic change giving the common
outgoing scalar quartic on those filter rows. In forming the full
outgoing collar quadratic, subtract the common outgoing mass on
the same anchors. The nonprimary parts of degrees four and six
of \(R_i-D_i\) are grouped into \(B_{\partial,4}^+,B_{\partial,6}^+\);
the quartic discrepancy from changing the coefficient on an
anchor is included in its fourth-degree polynomial.
This gives the literal exponent with an outgoing ordinary margin
and a full inner quadratic. The uniformly bounded \(O_L(x_j^2)\)
polynomials on the newly promoted fine anchors are moved into
this same collar polynomial. Products discarded in merging them
are exactly the products of at least two perturbative factors
estimated in Section~\ref{sec:cube-joint-reblocking}.

On the retained old collar tensors, linear heat and reblocking
are those of Lemma~\ref{lem:cube-boundary-heat}. A newly promoted
fine anchor carries an ordinary kinetic or head term in the
linear quadratic part. Its old stencil is at depth at least
\(c b_j\), where the fine tested norm of these terms is at most
\(C_L\) times their ordinary coefficient norms. Restrict the
coarse tests to the fine stencils and sum the \(O(L^d)\) fine
labels per coarse tile; this bounds the contribution by
\(C_L\Lambda_R x_j\). It is a triangular linear arrow from
the stable ordinary coefficients to the collar. The common raw
mass has no such arrow: its linear uncontracted square is the same
before and after engineering rescaling at each anchor. The
primary first heat on the \(O(B)\) collar satisfies
\eqref{eq:boundary-tadpole-source}. At a newly promoted anchor,
subtracting the common first-heat mass leaves the zero-row
filtered square and the covariance discrepancy in that estimate.
Both are retained in the outgoing quadratic, so the connected
remainder has no first-order term from this promotion.

The nonlinear collar terms have norm \(C_L x_j^2\) in the
small-radius units, by the finite-stencil covariance bound.
In detail, subtract the independent-copy Gaussian expectation
from the covariance. Each factor then has at least one fluctuation
direction, and the Gaussian moment and polarization bound that
bilinear map by \(C_L\sigma_j\), where the raw jet amplitude
\(\sigma_j=o(1)\) also for the merged step. Each normalized local
coefficient direction is bounded by its own weight in raw units.
An ordinary margin polynomial is a bounded bilinear function
of the outgoing ordinary coordinates and has no collar derivative.
Its derivative in a raw mass or stable quadratic direction is
at most \(C_Lx_j\) times that direction in the strengthened norm.
These bounds also treat the old ordinary polynomials promoted to
the new collar. Thus mass-to-collar quadratic derivatives are
\(O_L(x_j)\) after normalization, mass-to-higher tensors are
\(O_L(\Lambda_R)\), and quadratic-to-higher nonlinear stable
derivatives are \(O_L(\Lambda_R)\).
For the linear heat from degree six to degree four the extra
factor is \(O_L(\sigma_j)\). From degrees four and six to degree
two it is \(O_L((\sigma_j+\sigma_j^2)x_j/\Lambda_R)\).
The diagonal collar factors are at most
\(C L^{d-1-r(s/2+p)}\) for \(r=2,4,6\), times ratios of
successive weights tending to one. The normalized quartic-to-collar
derivatives are bounded, by the \(C_L\lambda_j^{\rm tad}|G|\)
bound for the quadratic source and the \(O_L(x_j^2)\) higher
sources. The mass row receives only the bulk derivatives.
In particular its mass derivative after removal of \(L^s\nu\)
is \(o(1)\) in normalized units.

Finally, the extraction of a small mark changes the ordinary and
collar coefficients by \(O_L(\tau)\), where \(\tau=\|K\|\).
Use the exact divergence reassignment preceding
Section~\ref{sec:cube-mass-chart}. On ordinary anchors this is a
common scalar displacement; the nonuniform part and the actual
interface terms are assigned to the widened collar. Its quadratic
terms there use the full tensors. The corresponding difference of
ordinary margin factors is \(O_L(x_j\tau)\) by the bilinear
comparison on the same stencils. The structured head has no
independent extracted coefficient. Its extraction dependence is
through the displayed two source formulas. A change of size
\(O_L(\tau)\) in the extracted coefficients changes \(G+a_{4,j}\)
by at most \(C_L\tau\) on this box: the factor
\(\lambda^2(j+\log(2+M_R))|G|\) is uniformly bounded.
The normalized head variation is therefore at most
\(C_L M_R^{-2}\tau/x_0\). The connected expansion only needs the
bounded coefficient changes on the enlarged disc
\(\tau\le c g^{9/4}\). In the shooting norm its potential
derivatives and its variation in a unit activity direction
are bounded by the same estimates with \(\tau\) replaced by
\(C f_j\), using Cauchy's inequality on the enlarged boxes.
The relevant uniform small factors are
\begin{equation}
 x_j/\Lambda_R,\qquad f_j/(\Lambda_R x_j),\qquad f_j/x_j^2 .
\end{equation}
For the transient part use
\(x_j^{-1}\le C(x_0^{-1}+\lambda^2j)\) and the geometric factor
\(d_0^j\). The initial ratios \(U/G_0^2\) and
\(U/(\Lambda_R G_0)\) tend to zero by
Section~\ref{sec:cube-joint-initialization}; the remaining factors
follow from the microscopic scalings.

The only linear stable coupling with a bound not tending to zero
besides the diagonals is from the ordinary coordinates to the
collar. Use an equivalent fixed product norm giving the collar
a sufficiently small relative weight. Choose \(L\) large for the
diagonals first, and then \(R\) large for the small reverse and
off-diagonal terms, with the radii of the boxes fixed in this order.
All scalar and tensor formulas use identical ordinary stencils
away from the halo. Finite propagation of the reblocking and the
same choice of localization on the common grid preserve bulk
agreement there, by grouping every disagreeing label into the
new boundary family. This proves the coefficient bounds and
compatibility with the outgoing split.
\end{proof}

The preceding proposition controls the retained coefficients. The
following corollary combines their transport with the contraction
of the activity and identifies the data passed to the terminal
integration.

\begin{corollary}[Coefficient and activity bounds along the fractional cube path]
\label{cor:fractional-cube-path}
In dimensions \(d=1,2,3\), the chart in
Proposition~\ref{prop:cube-retained-chart} satisfies the one-step
estimates in Assumption~\ref{ass:path} on the shooting boxes.
At its terminal scale, its coefficient bounds and inverse-mass
variation are those required in
Propositions~\ref{prop:terminal-packing} and~\ref{prop:terminal-graph}.
\end{corollary}

\begin{proof}
Use the paired bulk and finite-volume activity charts throughout
Proposition~\ref{prop:cube-bounded-extraction}. The bulk version has
the same range and product estimates, with no collar.
Proposition~\ref{prop:cube-mark} bounds the leading linear map by a
quantity tending to zero as \(L\to\infty\) and then \(R\to\infty\).
Choose \(L\), the size weight, and then \(R\) so that this bound
and the errors in Proposition~\ref{prop:cube-bounded-extraction}
are less than \(d_0/4\) in total. The same choice applies in the
bulk. The uniform disc bound there and Cauchy's inequality give
\begin{equation}
 \|D_K\Psi_j(V,K)-D_K\Psi_j(V,0)\|
       \le C\|K\|/g_j^{9/4}
\end{equation}
on a smaller concentric disc. The shooting box is contained in
this smaller disc and its relative radius tends to zero, since
\(U/G_0^{9/4}=o(1)\) and \(f_j=x_j^3+Ud_0^j\).
Since \(f_{j+1}\ge d_0 f_j\), the normalized activity diagonal
has a strict contraction. The bound on the source term at zero
is \(C x_j^3\), and the potential derivatives on the enlarged
shooting boxes are \(C(x_j^3+\|K\|)\) by Cauchy's inequality.
This gives the full activity row with values and derivatives in
units \(f_{j+1}\). The other rows are supplied by
Proposition~\ref{prop:cube-retained-chart}.

At the terminal scale the ordinary mass, local kinetic and
collar quadratic sizes are \(O(\Lambda_R x_N)\).
The fixed heads satisfy
\(h_N\le C\Lambda_R x_N\). The margin and the collar polynomial
factors have sizes \(O(x_N^2)\) and fixed field degrees.
The full coefficient box is holomorphic, including its terminal
inverse mass chart. The full collar tensors, fixed heads,
and the activity are independent raw coordinates in that inverse.
The raw quartic and local kinetic coefficients have derivative in
\(\nu_N\) of size \(O(\lambda x_N)\) by
Lemma~\ref{lem:cube-inverse-mass}. In the total quadratic, the
raw mass contributes the specified uniform square with multiplicative
derivative \(1+O(\lambda x_N)\). The difference from the target
uniform mass therefore has derivative \(o(1)\) at the small radius.
The remaining hypotheses of the two terminal propositions are the
raw Gaussian moments in Section~\ref{sec:ext-cov}, the fixed-filter
bounds, and the already normalized activity estimate. This proves
the asserted terminal bounds.
\end{proof}

\section{The marginal kinetic coefficient in dimension four}
\label{sec:four-dimensional-flow}

At the local endpoint a quadratic term with two spatial derivatives has
zero scaling dimension. We first isolate its coefficient in the bulk
quadratic interaction. We then estimate its increments and its transport
together with the terms supported near the boundary.

\subsection{Local representatives for quadratic terms}

\paragraph{An anchorwise cumulative bound in dimension four.}
The sunset kernel has a logarithmic cumulative second moment
at \(d=4\). We now treat this endpoint. Removing the second
moment at the original anchor gives a uniform ordinary
second-order remainder on a coefficient box where the retained
kinetic coefficient is \(O(g)\).

\begin{lemma}[Anchored ordinary reference in \(d=4\)]
\label{lem:cube-ordinary-W-d4}
Use the bulk covariance and filter assumptions of
Lemma~\ref{lem:cube-ordinary-W} with \(s=2\), and use the same
degrees and oriented rows. Represent the ordinary kinetic term
on anchor \(x\) by
\begin{equation}
 k_x^{\mathrm a}(\phi)=-\phi_x\Delta_{\mathrm c}\phi_x,\qquad
 \Delta_{\mathrm c}\phi_x
   =\tfrac14\sum_{i=1}^4
     (\phi_{x+2e_i}-2\phi_x+\phi_{x-2e_i}).
\end{equation}
The quadratic Fourier coefficient of its spatial sum is
\(|k|^2\). Keep the mass and kinetic moments in the local
representatives \(\phi_x^2,k_x^{\mathrm a}\), with the
prescribed quartic and vacuum also kept. Then the ordinary
remainder \((1-\mathrm{Loc})X_S(U)\) has the bilinear bound of
Lemma~\ref{lem:cube-ordinary-W}, with the same strengthened
filtered-head coefficient. The bound includes first variations
in those coordinates. It also holds after the exact removal
of ordinary extraction divergences, with the \(H\)-representative
kept instead of \(Z\).
\end{lemma}

\begin{proof}
The sums of \(k^{\mathrm a}\) and the central difference squares
\(k^{\mathrm c}\) coincide. Indeed, for a fixed coordinate
their sums both equal
\(\tfrac14\sum_x(\phi_{x+2e_i}-\phi_x)^2\).
The ordinary \(Z+H\) divergence has the same commutator
identity as before. Thus the exact exponential change uses
the ordinary \(H\)-representative with its sign and the
interface tensor from that identity.

In the sunset row put \(D=TS T^*\) and \(u=T\phi\).
Expand \(u_{x+r}\) to spatial order two at \(x\). The
discrete Taylor polynomial can use central first and
second differences, since these are exact on polynomials of
degree two. Its tested error is
\(C_L\ell_b |r/b|^3\) for \(r\ne0\), using three bounded
differences on the ordinary stencil. The odd linear term has
zero row because \(D(r)^3\) is even. Cubic symmetry makes the
quadratic term a scalar multiple of \(\Delta_{\mathrm c}u_x\).
The normalized test bound after subtracting these terms is
\begin{equation}
 C_L b^{s-3}\sum_r |D(r)|^3|r|^3
       \le C_L\lambda^3,\qquad s=2 .
\end{equation}
To replace \(u_x\Delta_{\mathrm c}u_x\) by
\(\phi_x\Delta_{\mathrm c}\phi_x\), each filtered leg adds
\((R/b)^2\) in the corresponding tested difference norm.
The logarithm in the removed coefficient therefore costs
only \(C_L\lambda^3(R/b)^2[1+\log(2+b/R)]\).
The constant term of this spatial expansion is \(u_x^2\)
times an absolute row sum \(O(\lambda^3R^{-2})\).
The difference from its mass representative has the same
\((R/b)^2\) filter gain and the kinetic moment is
\(O(R^2)\) times this row sum. Thus both give bounded
normalizations. These replacements agree exactly with the
quadratic Taylor moment of the summed anchor polynomial:
the relative \(k^2\) moment of the external filter legs is
the one removed from the constant row, and their row sum
is one on the spatial quadratic term.

The mixed rows admit the bounds of
Lemma~\ref{lem:cube-ordinary-W} with this kinetic density.
In the two-cross mixed row with a kinetic vertex, the chosen
anchor is the quartic and the internal sum is a multiple of
\(\sum_r (TS)(r)(-\Delta_{\mathrm c})(TS)(r)\).
By summation by parts over the bulk partner it is the
same difference-square sum as before. In a quadratic--quadratic
one-cross row rooted at \(k^{\mathrm a}\), the two contracted
possibilities leave \(\Delta_{\mathrm c}\) either on the
external field or on the smooth summed partner \(SQ\phi\).
Both give \(b^{-2}\). The Fourier moments of the sum and
all the remaining shell estimates are those in the
ordinary proof. The finite-jet and outgoing-reference
comparison therefore holds on the \(O(g)\) kinetic box.
\end{proof}

\paragraph{A fixed short head of vanishing second moment.}
The local power admits a microscopic quadratic remainder
whose coefficients need not decay in the original core
units. Its normalized effect decays in the ordinary
polynomial after removing its constant and second moments
on the actual anchors. The following bound also applies to
the bulk head produced by a short inverse in
Proposition~\ref{prop:kinetic-reference-family}.

\begin{lemma}[Ordinary estimate for a higher-order short head]
\label{lem:cube-fixed-short-head}
In the setting of Lemma~\ref{lem:cube-ordinary-W-d4}, add to
the second type of the incoming quadratic potential the
fixed ordinary density
\begin{equation}
 H_x(\phi)=\tfrac12\sum_v a_v(\phi_{x+v}-\phi_x)^2,\qquad
 a_{-v}=a_v,\qquad \sum_v a_v v_i v_k=0\quad(1\le i,k\le4).
\end{equation}
Assume it has finite range \(O(b_0)\) and
\begin{equation}
 \sum_v |a_v|\,|v|^2\le M,\qquad
 \sum_v |a_v|\,|v|^4\le C M R^2 .
\end{equation}
Then in the bilinear ordinary bound at scale \(b\ge b_0\)
this head contributes as a coefficient of size
\(C M(R/b)^2\). In the ordinary kinetic increment, the
quartic--head two-cross term from the change of the
covariance at an ordinary scale has magnitude at most
\(C_L\lambda^2 g M(R/b)^2\).
For a merged covariance with cumulative gradient bound
\(\|\nabla TS\|_2^2\le C\lambda^2 R^{-2}\), it has size
\(C_L\lambda^2 g M\).
These statements are bilinear, with the corresponding
first variations.
\end{lemma}

\begin{proof}
In the homogeneous bilinear form of \(H_x\), pair \(v\) and
\(-v\) and Taylor-expand the two differences to degree two.
Their cubic contributions cancel. The two linear differences
give the tensor \(\sum_v a_v v_i v_k\), which vanishes.
The remaining tests cost
\(C_L \ell_b^2 |v/b|^4\), by the same discrete Taylor
interpolation with bounded differences through order three.
Consequently \(H_x\) has the bound of a
quadratic coefficient \(C M(R/b)^2\) on the tile,
since \(b^d\ell_b^2=b^2\).
The Hessian of the sum has bounded absolute fourth moment
\(C M R^2\), zero row and zero second moment. As a convolution
on ordinary fields it costs at most \(C_L M R^2 b^{-4}\ell_b\)
and the corresponding bound with any fixed extra
spatial differences. This follows by pairing \(v,-v\)
and subtracting the Taylor polynomial through degree three.

For a one-cross quartic--head term the chosen anchor is
the quartic, and the contracted direction is the convolution
of the cumulative \(TS\) with this Hessian applied to the
field. The bound gains \(M(R/b)^2\).
In a one-cross quadratic--quadratic term the same
convolution argument applies if the head is the summed
partner. If the head is the root, the contraction has
one difference of the external field and one difference of
the smooth summed field. The paired local estimate just
proved applies to these two fields, giving the same gain.
The summed mass and second coefficients of this term
vanish by the zero row and zero second moment of the
head Hessian.

In the two-cross quartic--head term the quartic is again
the root. Write \(P_1=TS\). The scalar multiplying its
external square \(u_x^2\) is bounded by
\begin{equation}
 C g\sum_v|a_v|\sum_r|P_1(r+v)-P_1(r)|^2
      \le C_L g M\lambda^2 R^{-2},
\end{equation}
using telescoping and the cumulative
\(\|\nabla P_1\|_2^2\) bound.
The full ordinary projection of the external square is
its mass plus its kinetic moment on that anchor. In fact
evenness of the fixed filter, cubic symmetry and its
fourth absolute moment imply
\begin{equation}
 u_x^2=\phi_x^2+c_T\phi_x\Delta_{\mathrm c}\phi_x+
             O_L(\ell_b^2(R/b)^4)
\end{equation}
as a tested polynomial, with \(c_T=O(R^2)\).
This follows by Taylor-expanding the filter to degree
three, with cancellation of both odd degrees; the square
of the second-degree correction also has order
\((R/b)^4\). The second-degree term has exactly the
relative moment of the external square. Thus its
unprojected remainder improves on the mere
\((R/b)^2\) filter estimate. Multiplication by
\(b^2(R/b)^4\) in the preceding scalar row bound
gives \(C_Lg\lambda^2 M(R/b)^2\).
There is no other nonconstant degree from two quadratics,
so the claimed strengthened estimate follows.

Finally write \(P_C=TC\) for the filtered ordinary
slice. In the kinetic increment the quartic and this
fixed short input coefficient do not change under first
heat. The scalar difference of the two covariance rows
uses two products of differences where at least one
difference is on \(P_C\). By the ordinary slice and
cumulative bounds,
\begin{equation}
 \|\nabla P_C\|_\infty\le C_L\lambda b^{-3},\qquad
 \|\nabla P_1\|_1+\|\nabla P_C\|_1\le C_L\lambda b ,
 \qquad s=2 .
\end{equation}
Every shifted \(v\)-difference in a product is therefore
bounded in the sum over \(r\) by \(C_L|v|^2\lambda^2 b^{-2}\),
using telescoping, one supremum and one absolute row sum.
Multiply by \(g M\) and by the \(O(R^2)\) kinetic
moment of the filtered external square. The merged case
uses the earlier cumulative gradient estimate. This is the
asserted increment and derivative bound.
\end{proof}

\paragraph{A kinetic representative bounded up to the physical faces.}
A common choice of kinetic density up to the physical faces is
also possible. In the following estimate take
\(J_R=R^{-2}P_\Lambda L_RP_\Lambda^*\) and normalize
\(J_R\) by a bounded positive factor so that its bulk symbol
has leading term \(|k|^2\), and use the anchored density
\(\phi_x(J_R\phi)_x\) with zero extension. Its sum over an
enlarged fine tile has bounded bilinear \(F_b\)-norm at
radius \(\ell_b\), uniformly up to the faces, edges and
corners. For each tested direction \(v\), at distance more
than \(C R\) from the exterior, the symmetry of the walk
gives
\begin{equation}
 |(J_Rv)_x|\le C\ell_b(\rho_b(x)/b)^p\rho_b(x)^{-2},
\end{equation}
by summing centered second differences over displacements
of size at most \(R\). At smaller distances the same
estimate comes directly from the absolute jump-row bound
and \(\rho_b(x)\asymp R\). Multiply by the field bound
on the second direction and use
\begin{equation}
 \sum_{x\in P^*}(\rho_b(x)/b)^{2p}\rho_b(x)^{-2}
 \le C b^2 ,
\end{equation}
which follows by summing the layers of at most \(Cb^3\)
sites and using \(2p>1\). With \(\ell_b^2=b^{-2}\)
this proves the uniform norm. Freezing one normalized
kinetic coordinate therefore gives the same local bound
for its first variation.

In the ordinary region the alternative representative differs
from \(k_x^{\mathrm a}\) in Lemma~\ref{lem:cube-ordinary-W-d4}
by a tested coefficient of size \(O((R/b)^2)\) after summing
over a tile. The difference kernel is even and has zero
constant and second moments, and fourth absolute moment
\(O(R^2)\), so the fourth-order discrete Taylor bound
gives this estimate. The same argument gives the
corresponding bound on the difference of smooth contracted
legs in the ordinary bilinear estimate. The spread-out
representative thus keeps the ordinary kinetic moment and
admits uniform tests on an anchor newly assigned to the
collar.

\subsection{Summable perturbative sources for the local kinetic coordinate}
\label{sec:kinetic-increments}

The cumulative sunset has a second moment of logarithmic size in
dimension four. The shooting argument of
Corollary~\ref{cor:marginal} uses its successive increments.
At least one covariance line in each increment belongs to the
current slice, and this gives a scale-uniform bound.

\begin{lemma}[Kinetic moment of the sunset increment in \(d=4\)]
\label{lem:kinetic-increment}
Write \(W\) for a cumulative filtered bulk covariance and \(C\) for
the next ordinary filtered slice, at side length \(b\ge R\).
Assume they are supported in \(|x|\le C_1 b\) and obey the
bounds from Section~\ref{sec:bubble},
\begin{equation}
 |W(x)|\le C_2\lambda (R+|x|)^{-2},\qquad
 |C(x)|\le C_2\lambda b^{-2}.
\end{equation}
The kernels are taken even. The displayed cumulative envelope
also holds after filtering: the finite Neumann filter has bounded
absolute exponential row and column norms at scale \(R\), and
\((R+|x-y|)^{-2}\le (1+|y|/R)^2(R+|x|)^{-2}\).
The zero-row kinetic second moment
of the three-line quadratic increment, before multiplication by
its quartic factors, is bounded by
\begin{equation}
 \sum_x |x|^2\big|(W+C)(x)^3-W(x)^3\big|\le C\lambda^3 .
\end{equation}
For a merged ultraviolet slice with \(W=0\), the bound is
\(C\lambda^3\log(2+b_0/R)\).
\end{lemma}

\begin{proof}
Expand the difference as \(3W^2 C+3 W C^2+C^3\).
By counting lattice shells the first term has moment at most
\begin{equation}
 C\lambda^3 b^{-2}
 \sum_{1\le r\le Cb} r^{5}(R+r)^{-4}\le C\lambda^3 .
\end{equation}
For the second term the factor \(b^{-4}\) multiplies the shell
sum with exponent \(-2\). For the third the factor \(b^{-6}\)
multiplies the sum with exponent zero. These sums are
\(O(b^4)\) and \(O(b^6)\), respectively. Making the kernel
zero-row subtracts a multiple of \(\delta_0\), which does not
affect its second moment. If the two external legs carry the
normalized filter, the kinetic coefficient of an even zero-row
kernel is unchanged: multiply its symbol by
\(\lvert\widehat{\overline T}\rvert^2=1+O(R^2|k|^2)\) and
use that the kernel symbol vanishes to second order at zero.
The merged envelope \(C\lambda(R+|x|)^{-2}\) gives the logarithm
on taking its cube and the second moment.
\end{proof}

The quartic factors in the polynomial increment are bounded by
\(C x_j^2\), and have the same first-variation bound in a unit
normalized quartic direction. Thus the sunset contribution has
summable size and variation, at most
\begin{equation}
 C\lambda^3\left(x_0^2\log(2+b_0/R)+\sum_{j<N}x_j^2\right)
 \le C\lambda^3 x_0^2\log(2+b_0/R)+C\lambda x_0 .
\end{equation}
Here we used the reciprocal-increment comparison of
Section~\ref{sec:path}. The first-heat structured head has an
additional geometric gain: its raw coefficient is bounded by
\(C\lambda_j^{\rm tad}g_j b_j^{-2}\) since the filtered slice
diagonal gives the tadpole. The normalized zero-row second
moment of its fixed shape is \(O(R^2)\) at \(s=2\).
Consequently the incremental bound in that marginal coordinate
is \(C\lambda_j^{\rm tad}g_j(R/b_j)^2\), with the same bound
for its normalized quartic derivative. Its sum is finite
and tends to zero on the microscopic initial scales. These are
the polynomial sources in the marginal flow. The connected
remainder and its variation require an activity estimate on the
chosen kinetic domain. Proposition~\ref{prop:residual-coupled-chart}
supplies it on the shrinking domain used for shooting.

Appendix~\ref{app:fixed-kinetic} gives the corresponding identities
when the kinetic coefficient is held in a fixed disc. For the
four-dimensional trajectory we now choose a family of finite-range
covariances and bound the residual kinetic coefficient scale by scale.

\subsection{A finite-range family of kinetic Gaussian references}
\label{sec:kinetic-reference-family}

For the four-dimensional shooting argument we vary the Gaussian
reference before decomposing its covariance. The local split
allows a kinetic parameter in the physical block of the auxiliary
precision, so that each ordinary slice still has finite propagation.
The compensating quadratic density is explicit and the physical
spin measure and its pressure are unchanged. The kinetic moment of
this density will supply the second adjustable bare coordinate
alongside the temperature.

\begin{proposition}[A tunable finite-propagation Gaussian family]
\label{prop:kinetic-reference-family}
Work in \(d=4\) with the local split of
Lemma~\ref{lem:local-split}. Let \(\zeta=(1+t_R)^{-1}\),
\(J_\Lambda=P_\Lambda L_RP_\Lambda^*\), and fix a sufficiently
small real \(\delta\) in an interval independent of \(R,n\).
Set
\begin{equation}
 B^\delta=B+\delta L_R,\qquad
 \mathcal Q_{c,\delta}=
 \begin{pmatrix}
 I&-\mathcal L_\Lambda\\
 -\mathcal L_\Lambda^*& B^\delta_\Lambda/c
 \end{pmatrix},
 \qquad
 V_{c,\delta}=\bigl(A_\Lambda^{-1}
                         -c(B^\delta_\Lambda)^{-1}\bigr)^{-1}.
\end{equation}
Use the original observation \(\mathcal O=(\mathcal L_\Lambda^*,0)\).
For \(c\) below the principal pole \(c_{*,\Lambda}^\delta\),
\begin{equation}
 V_{c,\delta}=\mathcal O\mathcal Q_{c,\delta}^{-1}\mathcal O^*.
\end{equation}
Define
\begin{equation}
 \xi_i=\text{the \(i\)-th increasing eigenvalue of }
 A_\Lambda^{-1/2}J_\Lambda A_\Lambda^{-1/2}.
\end{equation}
Then \(c_{*,\Lambda}^\delta=1+(\zeta+\delta)\xi_1\).
The pole observation direction is independent of \(\delta\)
up to scalar normalization. Set
\begin{equation}
 c^\delta=c_{*,\Lambda}^\delta
             -\theta_0\lambda^2\bar g_N (R/n)^2.
\end{equation}
The polynomial finite-range decomposition of
\(C^\delta=a_R V_{c^\delta,\delta}\) has the annular Gaussian
moments and bulk filtered bubble bounds of
Sections~\ref{sec:ext-cov} and~\ref{sec:bubble}, uniformly in
this real interval. The physical Gaussian density change from
\(C^0\) to \(C^\delta\) inserts the quadratic Hessian
\begin{equation}
 D^\delta=(C^0)^{-1}-(C^\delta)^{-1}
   =a_R^{-1}\bigl[c^\delta(B^\delta_\Lambda)^{-1}
                         -c^0 B_\Lambda^{-1}\bigr]
\end{equation}
in the residual potential. Its bulk kernel and its
boundary defect, in the exponentially weighted absolute
row and column norms at length \(R\), are bounded by
\(C a_R^{-1}|\delta|\). The boundary defect has an additional
factor \(e^{-a\,\dist(x,\Lambda^c)/R}\) in row \(x\).
These estimates also hold for one derivative in \(\delta\),
with \(|\delta|\) replaced by one. All coefficients and
finite-horizon Gaussian data are continuous in the real
parameter.
\end{proposition}

\begin{proof}
The physical Schur complement at shift \(r\ge0\) is
\begin{equation}
 H_{r,\delta}=B_\Lambda^\delta/c+rI
                         -A_\Lambda/(1+r).
\end{equation}
For \(r=0\) the block inverse tested through the original
observation is \(A_\Lambda+A_\Lambda H_{0,\delta}^{-1}
A_\Lambda\). The matrix inverse identity gives \(V_{c,\delta}\).
The inverse exists when \(H_{0,\delta}\) is positive, which
is equivalent to \(c<1+(\zeta+\delta)\xi_1\), since
\begin{equation}
 B^\delta_\Lambda=A_\Lambda+(\zeta+\delta)J_\Lambda .
\end{equation}
The generalized eigenvalues of \(A_\Lambda\) against
\(B^\delta_\Lambda\) are
\([1+(\zeta+\delta)\xi_i]^{-1}\). If \(v\) satisfies
\(J_\Lambda v=\xi_1 A_\Lambda v\), then
\(\big[A_\Lambda^{-1}
        -c_{*,\Lambda}^\delta(B^\delta_\Lambda)^{-1}\big]A_\Lambda v=0\).
So the pole direction in the observation space is \(A_\Lambda v\).
The established two lowest generalized levels at \(\delta=0\)
give \(\xi_1\asymp (R/n)^2\) and
\(\xi_2-\xi_1\asymp (R/n)^2\), with the upper bound on
\(\xi_2\) from the same min--max comparison. The formula for
the generalized eigenvalues therefore preserves the gap
uniformly. The exact displacement in \(c^\delta\) gives the
same principal variance in the adapted source coordinate,
namely \(n^2/(\theta_0\lambda\bar g_N)\).

To verify finite propagation and the local estimates, put
\(Q^\delta=Q+\delta D_R\), so that
\(B^\delta=(1+\delta)I-Q^\delta\).
The choice of the smoothing amplitude in the proof of
Lemma~\ref{lem:local-split} gives \(Q(x)\ge aD_R(x)\)
off the diagonal for a fixed \(a>0\). Thus \(Q^\delta\)
is nonnegative with zero diagonal, range \(O(R)\), and
supremum \(O(R^{-4})\) for sufficiently small signed
\(\delta\). Its row sum is \(q+\delta\), and
\((q+\delta)/(1+\delta)\le 1-c_1\) uniformly. The block
operator has the original finite propagation and uniform
spectral upper bound. In the polynomial decomposition the
degree bounds are unchanged, so each covariance slice has
finite range in the same scale units.

Here are the specific uniform inputs for its boundary estimates.
In the physical Schur complement the bulk convolution can
be written as
\begin{equation}
 H_{r,\delta}
  =\bigl((1+\delta)/c+r\bigr)I
       -Q^\delta/c-A/(1+r)
  =\frac{\zeta+\delta}{c} L_R+rI+
                   \left(c^{-1}-(1+r)^{-1}\right)A .
\end{equation}
The off-diagonal jump kernel in the first expression has
the same nonnegative projected spread-out component used
in the barrier proof. Thus the normal-coordinate barrier
and the loss on jumps killed outside the cube give the same
\(r^{-1}w(x)\) resolvent row bound, for
\(w(x)=(\min(b,R+\dist(x,\Lambda^c))/b)^{p_0}\),
\(p_0<1\), with \(b=Rr^{-1/2}\).
The bounded correction multiplying \(A\) is absorbed by
choosing the spectral shift as in the unperturbed proof.
In the auxiliary Duhamel comparison the diagonal entries
are \(1+r\) and \((1+\delta)/c+r\); all the remaining
couplings are nonnegative. Its no-coupling subtraction and
the bulk domination are unchanged. The bulk Schur symbol in
the second expression is comparable to \(r+\omega(k)\)
after the same enlarged regularization. The high-frequency
gap, the \(O(R^{-4})\) squared integral of the physical
off-diagonal multiplier, and the observation smoothing
are uniform, since \(Q^\delta=Q+\delta D_R\) and \(A\)
is fixed. The interior block equations and cutoff
derivative estimates in Lemma~\ref{lem:local-annulus}
use a finite-range even physical convolution with bounded
row norms, second moment \(O(R^2)\), and these low- and
high-frequency bounds. All hold with \(Q^\delta\), including
the Fourier derivative bounds on the scaled torus because
\(D_R\) has fixed scaled support. This verifies the same
annular variance and exponential moment proof, including
the microscopic collar and merged scales. The tail
comparison, principal source bounds and bubble bounds
then follow by the same spectral polynomial estimates:
the lower bubble band uses the Schur comparison at two
fixed regularizations and the observation lower bound at
low frequency, and its upper bound uses the just established
bulk diagonal bound. These arguments use \(\delta\) only
in fixed compact bounds.

For the local density change use the original and modified
short inverse series
\begin{equation}
 (B^\delta_\Lambda)^{-1}
 =\frac1{1+\delta}\sum_{m\ge0}
                  \left(\frac{Q^\delta_\Lambda}{1+\delta}\right)^m .
\end{equation}
Their weighted absolute row sums converge geometrically
with a fixed exponential distance weight at length \(R\).
A killed-path defect at row \(x\) requires at least
\(c\,\dist(x,\Lambda^c)/R\) steps. Reserving part of the
geometric decay for those steps proves the stated boundary
bound. A derivative in \(\delta\) marks one step or
differentiates the scalar normalization and changes the
per-length bound by at most a fixed polynomial in \(m\).
The path bound with a slightly weakened exponential
is uniform for the derivative chain. Since \(c^\delta-c^0
=\delta\xi_1\), the mean value formula using these estimates
gives the additional \(|\delta|\) for \(D^\delta\).
Symmetry gives the column estimates. Finally,
substituting the two precision matrices into the density
identity gives the displayed sign of \(D^\delta\);
the scalar Gaussian normalization is a vacuum.
\end{proof}

\paragraph{The initial kinetic and mass coordinates.}
Let \(p=\widehat B(0)\),
\(B_2=\tfrac12\sum_xx_1^2Q(x)\),
and \(j_2=\tfrac12\sum_xx_1^2D_R(x)\asymp R^2\).
The bulk two-jet of this density change is explicit:
\begin{equation}
 \widehat D^\delta(k)
 = \frac{\delta\xi_1}{a_R p}
 - \frac{\delta\xi_1 B_2+\delta c^\delta j_2}
           {a_R p^2}\,|k|^2 +O_{R,\delta}(|k|^4).
\end{equation}
Indeed both short inverse denominators have row sum \(p\)
and \(\widehat B^\delta(k)=p+(B_2+\delta j_2)|k|^2
+O_R(|k|^4)\). At fixed \(c\), the mass term is exactly
zero. With \(c^\delta\) chosen above its magnitude is
\(|\delta|\xi_1/(a_Rp)\). In the scale-\(b_0\)
mass coordinate this is
\(O(|\delta|\lambda^{-1}(b_0/n)^2)\).
The derivative of the kinetic coefficient is
\begin{equation}
 -\frac{\xi_1 B_2+(c^\delta+\delta\xi_1)j_2}{a_R p^2},
\end{equation}
negative and of order \(\lambda^{-1}\), uniformly at
large \(n/R\). In particular a change
\(|\delta|=O(\lambda^2G_0)\) changes the initial kinetic
coefficient by \(O(\lambda G_0)\). Before canonical
field conversion the added Hessian has size
\(O(\lambda^2G_0)=O(g_{\rm mic})\) in the exponentially
weighted row norm. Its bulk mass may be assigned to
the scalar mass chart, its zero-row remainder to the
bulk quadratic heads, and its clipping defect to the
boundary quadratic heads. These assignments retain the
same short-tail path estimates as physical initialization.
Thus the initial kinetic coefficient can be varied while keeping
the ordinary fluctuations within the finite-propagation
construction.

\paragraph{Exact transfer from the unperturbed scale decomposition.}
For comparison of the terminal distributions across this family,
it is convenient to express the perturbed covariance directly in
terms of the unperturbed one. This also gives control of the
ordinary slices and their first variations. Write
\(d_\delta=(\zeta+\delta)/\zeta\) and
\(c=1+d_\delta(c'-1)\), where \(c'<c_{*,\Lambda}^0\).
For the zero-shift Schur complements one has
\begin{equation}
 H_{0,\delta}(c)=\frac{d_\delta c'}{c}H_{0,0}(c'),\qquad
 V_{c,\delta}=t_\delta V_{c',0}+(1-t_\delta)A_\Lambda,
 \qquad t_\delta=\frac{c}{d_\delta c'}.
\end{equation}
Indeed both the coefficient of \(L_R\) and the coefficient
of \(A\) in the first identity have the displayed common
factor. Substitution in
\(V=A_\Lambda+A_\Lambda H_0^{-1}A_\Lambda\) gives the second
identity. In the new reference the distance to its pole is
\(d_\delta\) times the original pole displacement. In the bulk
two-jet formulas for this choice of \(c\), replace \(\xi_1\) by
\((c'-1)/\zeta=O((R/n)^2)\). The microscopic tail bounds,
strict kinetic slope and bound on the scalar perturbation
are unchanged in size.

To obtain a positive finite-range decomposition for both signs
of \(\delta\), fix \(\eta>0\) with
\(\eta\|\mathcal Q_{c',0}\|\le 1/2\). Multiply each original
auxiliary slice polynomial by \(1-\eta t\) in its spectral
variable \(t\), including the functional-calculus terminal
remainder. Functional calculus preserves positivity and bounds
each new slice between one-half and one times the original
slice. The ordinary propagation range increases by at most
the microscopic range of the auxiliary precision. The sums
of the resulting slices and remainder give
\(\mathcal Q_{c',0}^{-1}-\eta I\).
On applying the observation, multiplying by \(t_\delta\), and
adding \([1-t_\delta+t_\delta\eta]A_\Lambda\) to the merged
slice, one obtains \(V_{c,\delta}\) exactly. The coefficient
of \(A_\Lambda\) is positive on a fixed neighbourhood of
\(\delta=0\). The positive remainder of the merged slice is
comparable to its unperturbed counterpart, and the added
\(A_\Lambda\) has fixed microscopic range.

The ordinary annular variance bounds transfer by positive
operator comparison, since a stencil variance is a quadratic
form in the covariance. The ordinary first \(\delta\) variation
is the same scale kernel times \(\partial_\delta t_\delta\),
which is uniformly bounded when \(c'\) is fixed. The bubble
lower bound transfers by the positive slice comparison on the
bulk Fourier fibers. The upper bound transfers by the absolute
row and diagonal estimates. In the merged slice the additional
\(A_\Lambda\) may be bounded directly. Its canonical diagonal
is \(O(\lambda R^{-2})\) and its supported test differences of
order \(k\) have standard deviation
\(O(\sqrt\lambda R^{-1-k})\), using the smoothing power in
Lemma~\ref{lem:local-split} chosen larger than the tested order
and the exponentially small white channel. For the supremum of
the weighted stencil tests one may even use a union bound on
a merged tile: the extra fluctuation parameter is at most
\(C\lambda M_R^{2P+2p+2}\log(2+b_0)\), which still tends
to zero. This gives the growing fluctuation reserve at the
merged scale. Its canonical diagonal alone yields the
tadpole bound with \(\lambda M_R^{2}\). These statements
also hold for one variation of the scalar multiplier of this
added covariance. The resulting decomposition varies the
terminal reference and is available uniformly on the fixed
external kinetic domain \(|\delta|/\lambda\le c\).

\paragraph{The kinetic source on the shrinking ordinary box.}
Backward shooting for the residual kinetic coefficient requires
a summable bound on its increment and its first variation.
It is important here to retain the mass projection. Its
cancellation removes the tadpole multiplied by the cumulative
bubble, which would otherwise exceed the cubic error after
quadratic compensation.

\begin{proposition}[Bulk kinetic increment in the shrinking chart]
\label{prop:bulk-kinetic-source}
Assume the exact ordinary split of
Lemma~\ref{lem:cube-oriented} and
Lemma~\ref{lem:cube-ordinary-W-d4}, with the mass and kinetic being
the complete quadratic Fourier two-jet. Use the kinetic
representative \(k^{\mathrm a}\). Let the bulk quadratic Hessian
before the step have symbol
\begin{equation}
 Q(k)=u+z^{\mathrm{loc}}|k|^2+O(|k|^4).
\end{equation}
Assume \(|u|\le C b^{-2}x_j\), \(|z^{\rm loc}|\le Cx_j\) and
\(G\asymp x_j\) on the enlarged potential box. Include the
zero-two-jet short heads as in Lemma~\ref{lem:cube-fixed-short-head}.
Let \(M_j\) bound their total second absolute moment
and \(C R^2 M_j\) their total fourth absolute moment.
Require \(M_j R/b_j\le C x_j\) and \(M_0\le C x_0\).
This includes the fixed microscopic remainder and the
zero-two-jet part of the filtered square. The comparison sequence
and the activity satisfy the bounds of Section~\ref{sec:path}.
For the enlarged extraction chart, assume the retained
scalar coefficient and its first variations change by
\(O_L(\tau)\) and the corresponding bounds in
Proposition~\ref{prop:cube-bounded-extraction}, with
\(\tau\le C f_j\). The Gaussian covariance is held fixed in
these variations.

Then for \(z\) in the compensated bulk quadratic chart of
Section~\ref{sec:cube-mass-chart},
\begin{equation}
 |z_+-z|+\|D(z_+-z)\|
  \le C_L\bigl[\lambda^3x_j^2+f_j
                     +\lambda x_j(R/b_j)^2\bigr].
\end{equation}
Here \(D\) measures directions in the enlarged normalized
potential and activity boxes, including \(z/x_j\). At the
merged step the right side may be replaced by
\begin{equation}
 C_L\bigl[\lambda^3x_0^2\log(2+M_R)+f_0+\lambda x_0\bigr].
\end{equation}
This proposition estimates the ordinary bulk source. It
assumes availability of the stated local outgoing split and
does not construct the coupled boundary chart.
\end{proposition}

\begin{proof}
Use \(S=S_j\) for the cumulative bulk covariance and \(C\)
for the incoming slice. Put \(D=TST^*\), \(E=TCT^*\),
\(\gamma=E(0)\), \(P=ST^*\delta_0\), \(P_C=CT^*\delta_0\).
A superscript \(+\) on a kernel means that \(S\) is replaced
by \(S+C\). First heat replaces the quadratic Hessian by
\(Q_1=Q+3\gamma G\,T^*T\). The kinetic coefficient of
this heat is \(O(R^2|\gamma G|)\).

The quadratic part of the Wick covariance projection has
three kinds of terms. The three-cross quartic pair has
Hessian obtained from \(6G^2D(r)^3\) between the two
external filtered legs. In the two-cross quartic--quadratic
pair both external legs are at the quartic anchor. Its
Hessian on that anchor is
\begin{equation}
 3G\,\langle P,Q P\rangle\,
       (T^*\delta_x)(T^*\delta_x)^* .
\end{equation}
Indeed the second derivative of \(G(T\phi)_x^4/4\) along
two filtered directions has scalar \(3G(T\phi)_x^2\),
and the Hessian of the quadratic is \(Q\). The Wick
factor \(1/2!\) cancels on passing to the Hessian. This is
independent of the anchor assignment when summed in the
bulk. Lastly, the one-cross quadratic pair has summed
Hessian \(Q^2 S\). All other nonconstant degrees are
quartic or sextic. The quartic projection before the
step has coefficient
\begin{equation}
 a_4=9G^2\sum_r D(r)^2+4Gu\sum_r S(r).
\end{equation}
The old localized quartic is heated in the anchor
subtraction, supplying the same filtered square with
Hessian \(3\gamma a_4T^*T\). These formulas use \(G/4\)
as the quartic normalization throughout.

For the mixed pair, the difference of scalar rows is
exactly
\begin{equation}
 \langle P^+,Q_1P^+\rangle-\langle P,QP\rangle
 =\langle P^+,(Q_1-Q)P^+\rangle
       +2\langle P_C,QP\rangle+\langle P_C,QP_C\rangle .
\end{equation}
The mass part \(uI\) of the last two summands is
bounded by \(C_L|u|\lambda^2\), since
\(\|P\|_1\le C_L\lambda b^2\) and
\(\|P_C\|_\infty\le C_L\lambda b^{-2}\), and the same
bounds hold for the corresponding factors with \(P_C\).
For an even zero-row kernel \(Z\) of second absolute
moment \(M_Z\), use
\(\langle f,Zh\rangle=-\tfrac12\sum_v Z(v)
 \langle f(\cdot+v)-f,h(\cdot+v)-h\rangle\).
The bounds \(\|\nabla P_C\|_\infty\le C_L\lambda b^{-3}\)
and \(\|\nabla P\|_1+\|\nabla P_C\|_1\le C_L\lambda b\)
give \(C_L M_Z\lambda^2 b^{-2}\) for the scalar row. This
applies to the local kinetic and to any zero-row head.
For a fixed zero-second-moment head one can improve it:
put \(Z\) on the fresh smooth factor and subtract its
Taylor polynomial through order three. With fourth
absolute moment \(M_{Z,4}\) the bound becomes
\(C_L M_{Z,4}\lambda^2 b^{-4}\), since
\(\|\nabla^4P_C\|_\infty\le C_L\lambda b^{-6}\).
These estimates apply to both mixed terms in the display.
The first summand in that display is
\(3\gamma G\sum_r (D^+(r))^2\), whose absolute value is
at most \(C_L|\gamma G|\lambda^2[1+\log(2+b/R)]\).

Every mixed scalar row is converted to its kinetic
coefficient by multiplying by the relative quadratic
moment of the two external filter legs, of size \(O(R^2)\).
Including \(3G\), its increment is bounded by
\begin{equation}
 C_L\lambda^2 x_j (R/b)^2
    \bigl(x_j+M_j(R/b)^2
             +\lambda x_j[1+\log(2+b/R)]\bigr)
\end{equation}
at an ordinary step. This is \(O_L(\lambda x_j(R/b)^2)\).
Indeed the comparison sequence and \(b/R=M_RL^j\)
give \(\lambda^2x_j[1+\log(2+b/R)]\le C_L\).
At the merged step \(S=0\). The bounds
\(\|P^+\|_2^2\le C_L\lambda^2[1+\log(2+M_R)]\) and
\(\|\nabla P^+\|_2^2\le C_L\lambda^2R^{-2}\) give the
asserted merged envelope for the mixed increment.
The kinetic coefficient of the quartic pair is
\(O_L(\lambda^3 x_j^2)\), by
Lemma~\ref{lem:kinetic-increment}. The external filter
moment contributes only
\(C_LR^2 G^2\sum_r |(D+E)(r)^3-D(r)^3|\);
that sum is at most \(C_L\lambda^3 b^{-2}
[1+\log(2+b/R)]\) at an ordinary step, which is
also sufficient. The merged bound is the stated
logarithmic second moment. Finally,
\(R^2|\gamma| (|G|+|a_4|)\) is at most
\(C_L\lambda x_j(R/b)^2\), using
\(|a_4|\le C_L x_j\), and at most \(C_L\lambda x_0\)
at the merged step.

We next give the raw mass bound required for the
quadratic compensation. In its projection the sunset
increment has potentially singular term
\(18G^2\gamma\sum_r D(r)^2\). The difference of its
\(D^2E\) row from that singular part is \(O_L(\lambda^3
b^{-2}G^2)\): use
\(|E(r)-\gamma|\le C_L\lambda b^{-2}|r|/b\) and
\(\sum_r |D(r)|^2|r|/b\le C_L\lambda^2\) on the
covariance support. The terms with at least two \(E\)
have the same bound. The mixed row from the first-heat
change \(Q_1-Q\) has singular mass
\(9G^2\gamma\sum_r D(r)^2\). Its difference from the
full mixed contribution with \(Q_1-Q\) is bounded by
\(C_L G^2|\gamma|\lambda^2\), since
\(\sum_r|(D^+(r))^2-D(r)^2|\le C_L\lambda^2\).
The old localized quartic heat subtracts precisely
\(27G^2\gamma\sum_r D(r)^2\) in the potential correction.
Thus all three singular coefficients cancel. The
remaining part of \(3\gamma a_4\), involving \(Gu\sum S\),
has size \(C_Lb^{-2}x_j^2\). The previously estimated
non-heat mixed rows also have this bound in raw mass
units: for the fixed zero-second-moment head use
\(b^{-4}M_{Z,4}\) and \(M_j(R/b)^2\le C_Lx_j\).
For the merged step the uncorrected cubic row bound
\(C\lambda^3R^{-2}\) replaces this cancellation.
After multiplication by \(b_0^2\) its additional factor
\(\lambda^3 M_R^2\) is bounded. The merged non-heat and
heat mixed rows similarly use
\(\lambda^2 M_R^2\) and \(\lambda^3 M_R^2[1+\log(2+M_R)]\)
times bounded normalized inputs, so give the same
conclusion.

For the pure quadratic part and its compensation, work
in the algebra of Fourier two-jets. The reference update
subtracts \(Q_1^2(S+C)-Q^2S\) as well as the non-quadratic
projected correction just bounded. Each quadratic
quantity here is in dimensionful mass and kinetic units.
The pure projected kinetic correction is \(O_L(x_j^2)\),
using \(S(0), (S+C)(0)=O_L(\lambda b^2)\) and their
second moments \(O_L(\lambda b^4)\).
Its mass correction is \(O_L(b^{-2}x_j^2)\) by the
same calculation. Together with the preceding mixed,
sunset and tadpole estimates this proves
\begin{equation}
 |u_+-u_1|\le C_L b^{-2}x_j^2,\qquad
 |z^{\rm loc}_+-\operatorname{kin}Q_1|\le C_L x_j^2 .
\end{equation}
In the kinetic estimate the old localized quartic heat is
\(O_L(x_j^2)\) even though its rough bound above was
linear in \(x_j\): retain
\(|a_4|\le C_L[\lambda^2(1+j+\log(2+M_R))+\lambda]x_j^2\)
and use the geometric factor \(R^2|\gamma|\).
Here \(\operatorname{kin}\) denotes the coefficient of
\(|k|^2\). The merged bound holds since \(a_4=0\) there.

The compensated coordinate is the kinetic two-jet of
\(Q+Q^2S\). Substituting the reference update into this
expression cancels the pure quadratic projection
exactly to second order. Its only remaining term in
this calculation is the corresponding two-jet of
\begin{equation}
 (Q_+^2-Q_1^2)(S+C)
     = (Q_+-Q_1)(Q_++Q_1)(S+C).
\end{equation}
Its kinetic coefficient is \(O_L(\lambda x_j^3)\)
by the last displayed bounds. It fits inside \(f_j\).
Using the exact rational quadratic two-jet instead
changes these estimates only by a cubic expression on
this shrinking domain. The extraction changes cost
\(O_L(\tau)\) in the retained coefficient and have the
stated variation bound. All the polynomial estimates
hold also on the complex potential box. Applying
Cauchy's inequality on a fixed smaller normalized box
gives their derivative bounds with the same envelopes.
This proves the assertion for the bulk step.
\end{proof}

\paragraph{Orientation when a filtered quadratic is split into two jets.}
The first heat of the quartic produces a kinetic two-jet and a
short quadratic remainder. Both pieces can participate in a
two-cross contraction with a quartic. Assigning that contraction
to the quartic anchor requires a redistribution of the individual
covariance rows. The next lemma bounds this redistribution and
uses its antisymmetry to preserve the sum for each pair.

\begin{lemma}[Reorientation of the ordinary two-cross rows]
\label{lem:kinetic-reorient}
Use the covariance, test kernels and ordinary enlarged shrinking box
of Proposition~\ref{prop:bulk-kinetic-source}. In addition to the
fixed difference-square head in Lemma~\ref{lem:cube-fixed-short-head},
allow a short filtered quadratic of the form
\begin{equation}
 q_T\bigl[(T\phi)_x^2-\phi_x^2-c_T\phi_x\Delta_{\mathrm c}\phi_x\bigr],
 \qquad |q_T|R^2\le M_j,\qquad
 M_jR/b_j\le Cx_j,\qquad M_0\le Cx_0 ,
\end{equation}
where \(c_T\) is the quadratic coefficient of the summed
external filtered square in this representation. Use the symmetric
covariance row for the terms other than the mixed two-cross term
of a quartic and a zero-row quadratic. Assign the latter term to
the quartic anchor.

This ordinary representative has the bilinear bounds of
Lemma~\ref{lem:cube-ordinary-W-d4}, with an effective strength
\(C M_j(R/b)^2\) for the zero-two-jet heads, including
the fixed microscopic head if its moments are included in
\(M_j\). In the
anchorwise heat identity it suffices to replace the symmetric
fresh covariance row by a row with the same pair totals and an
\(O_L(x_j^2)\) bound per fine label after summing its neighbors.
The row is supported on pairs at distance \(O_L(b)\), and the
bounds hold uniformly in the enlarged complex potential domains.
\end{lemma}

\begin{proof}
Let \(P=TS\), \(D=TST^*\), and \(A_T=T-I\) as bulk
convolution kernels. If \(T=I+(c_T/2)\Delta_{\mathrm c}\)
to quadratic order, \(A_T-(c_T/2)\Delta_{\mathrm c}\)
has zero row and second moment. Its absolute moments of orders
two and four are \(O(R^2)\) and \(O(R^4)\).
For the mixed two-cross term of this filtered head and a
quartic the scalar radial row is, up to fixed Wick factors,
\begin{equation}
 F_S=D^2-P^2-c_T P\Delta_{\mathrm c}P
      =2P\big(A_T-(c_T/2)\Delta_{\mathrm c}\big)P+(A_TP)^2 ,
\end{equation}
where the square and the first multiplication are pointwise.
By fourth-order filter cancellation and the cumulative kernel
difference bounds,
\begin{equation}
 |F_S(r)|\le C_L\lambda^2 R^4 (R+|r|)^{-8}.
\end{equation}
For this estimate use paired Taylor subtraction for shifts at
most \((R+|r|)/2\), and the exponentially weighted row bound for
the complementary shifts. The anchored external filtered square
has a fourth-order remainder after mass and kinetic subtraction,
as in Lemma~\ref{lem:cube-fixed-short-head}. The two-cross term
assigned to the quartic anchor therefore costs
\(C_Lg\lambda^2 M_j(R/b)^2\). A one-cross term rooted at the head
is the polarized difference of the filter product and its
degree-two spatial Taylor polynomial, evaluated on the external
field and a smooth summed partner. Its fourth-order bound
gives the same geometric gain. If the head is the summed
partner the bulk convolution Hessian has zero row and second
moment, so the gain follows by putting its four differences
on that smooth partner. These also handle quadratic--quadratic
rows. A local kinetic one-cross term may be kept symmetric:
at its root its two alternatives put the Laplacian on the
external field or on the summed smooth partner. Both cost
\(b^{-2}\). This proves the ordinary reference bounds with
the chosen convention.

Denote the correction of the ordinary row on the pair
\((i,k)\) by \(J_{ik}^{S}(U)\). For each mixed term it is half
the corresponding scalar times the difference of the two
external filtered squares; thus \(J_{ki}^{S}=-J_{ik}^{S}\).
The exact heat row is the symmetric fresh row plus
\begin{equation}
 J^{S+C}_{ik}(P_CU)-P_C J^S_{ik}(U).
\end{equation}
This follows by applying the heat identity of
Lemma~\ref{lem:cube-oriented} to the original symmetric row
and adding the correction on both sides. Its pair totals are
unchanged. In the bulk \(P_C\) of the difference of the two
external filtered squares adds zero because the diagonal heat
is uniform on the anchors.

We give estimates for this difference. Write \(P_C=TC\) for
the fresh filtered kernel. For the filtered head above the
change of \(F_S\) is a sum of the cross products in its last
formula with one or both kernels replaced by \(P_C\).
With \(v(r)=\min(1,|r|/b)\), its weighted absolute row sum is
\begin{equation}
 \sum_r v(r)\bigl|F_{S+C}(r)-F_S(r)\bigr|
 \le C_L\lambda^2 R^3 b^{-3}.
\end{equation}
Indeed the term with \(P_C\) multiplying
\((A_T-(c_T/2)\Delta_{\rm c})P\) uses
\(\|P_C\|_\infty\le C_L\lambda b^{-2}\)
and \(\sum |r||((A_T-(c_T/2)\Delta_{\rm c})P)(r)|
 \le C_L\lambda R^3\). The latter follows from the
pointwise bound \(C_L\lambda R^4(R+|r|)^{-6}\).
For the oppositely placed fourth difference use
\(C_L\lambda R^4 b^{-6}\) on \(P_C\) and
\(\|P\|_1\le C_L\lambda b^2\).
In the cross product with \(A_T\) on each factor use
\(\|A_TP\|_1\le C_L\lambda R^2[1+\log(2+b/R)]\)
and \(\|A_TP_C\|_\infty\le C_L\lambda R^2 b^{-4}\).
The terms with two fresh factors satisfy the same bound.
Use \((R/b)[1+\log(2+b/R)]\le C\) to conclude.
The factor \(v(r)\) is supplied in multilinear tests
by the difference of the external filtered squares.
Multiplication by \(g|q_T|b^2\) gives
\(C_Lg\lambda^2 M_j R/b=O_L(g^2)\).

For a fixed difference-square head of zero second moment,
the scalar cross row is
\begin{equation}
 \sum_v a_v(\Delta_vP)(\Delta_vP_C),
 \qquad \Delta_v f(r)=f(r+v)-f(r).
\end{equation}
Pair opposite \(v\) and
Taylor-expand both differences to order two. The degree-two
product cancels by the moment condition and the cubic product
by symmetry. The absolute summed remainder is bounded by
\begin{equation}
 C_L M_j R^2\bigl(
  \|\nabla^3P\|_1\|\nabla P_C\|_\infty+
  \|\nabla^2P\|_1\|\nabla^2P_C\|_\infty+
  \|\nabla P\|_1\|\nabla^3P_C\|_\infty\bigr).
\end{equation}
To see this without a pointwise restriction on the span,
subtract the linear and quadratic products in an order that
pairs a third-order remainder with the full first difference,
a second-order remainder with the second-order term, and a
third-order remainder with the first-order term. The
telescoping remainder bounds in \(\ell^1\) and \(\ell^\infty\)
give \(|v|^4\) in each case. The cumulative derivative
row norms of orders one, two and three are respectively
\(C_L\lambda b\), \(C_L\lambda[1+\log(2+b/R)]\)
and \(C_L\lambda/R\). The fresh derivative bounds then
give \(C_L\lambda^2 M_j R b^{-3}\).
The same calculation applies when both factors are fresh,
and multiplication by \(g b^2\) again yields
\(O_L(g^2)\).

For the local kinetic with coefficient \(O(g)\), there is no
need for a zero second moment. One supremum of a fresh
first difference and one cumulative \(\ell^1\) first
difference give \(C_L\lambda^2 |z| b^{-2}\), as in the
kinetic source estimate. Its row correction is therefore
\(O_L(g^2)\) in the fine jets. The remaining change of the
quadratic coefficients under the first heat has zero-row
part \(3\gamma G(T^*T-I)\). Here the scalar row is
\((D^+)^2-(P^+)^2\), whose absolute row norm is
\(O_L(\lambda^2)\) by even-filter second differences.
The coefficient is \(O_L(\lambda g b^{-2})\), so its
changed orientation also costs \(O_L(g^2)\).
At the merged step use \(S=0\), the cumulative kernel
bounds from length \(R\), and the factor \(b_0^2/R^2\)
in raw mass normalization. The resulting coefficients are
bounded by \(C_L\lambda^2 M_R^2 g^2\), including the
filtered head and first heat after their factors of
\(\lambda\) are included. This preserves the bound.

All these transfer rows have the finite support of the
cumulative and fresh kernels and their stencils. The finite
cube subtraction may use their bulk versions on the pairs
read by any new ordinary anchor, and zero or their
antisymmetric completion on the other pairs inside the
boundary chart. A fixed coarse collar contains all
nonordinary neighbors. Antisymmetry leaves the pair
covariance sums unchanged, so the same connected row
cancellation is valid; no equality of bulk and clipped
coefficients at an individual collar anchor is needed.
\end{proof}

\subsection{The coupled cube chart for a residual kinetic coefficient}
\label{sec:residual-coupled-chart}

Fix a member of the Gaussian reference family during the RG
iteration. The bulk source estimate controls the small residual
kinetic coefficient in this reference. It remains to use the
same quadratic coefficients in the ordinary and collar parts of
the activity map. The ordinary reference retains the full kinetic
two-jet. When a fine anchor enters the coarse collar, its actual
fine quadratic is assigned to the collar tensor. In particular,
the central difference stencil is used only on the ordinary side
of the physical boundary.

\begin{proposition}[Coupled anchored split on the shrinking kinetic box]
\label{prop:residual-coupled-chart}
Let \(d=4\), \(1/2<p<1\), and use the canonical finite-range
covariance, including any fixed member of
Proposition~\ref{prop:kinetic-reference-family}.
Use the comparison sequence of Section~\ref{sec:path}, and take
\(\Lambda_R\asymp\lambda M_R^2\). Assume the initialized
free and ordinary activities agree off the grouped boundary family
and have the bounds of Proposition~\ref{prop:cube-joint-initialization}
with the chosen quadratic inputs. Keep the Gaussian reference fixed
when defining derivatives in the shooting variables.

There is one outgoing anchored split as required in
Proposition~\ref{prop:cube-bounded-extraction}, on the enlarged
complex domains
\begin{equation}
 |\nu_j|+|z_j|+\|Q_{\partial,2}\|_{\ell_j,0}
                    \le C\Lambda_R x_j,\qquad
 \|B_{\partial,4}\|_{\ell_j,0}+\|B_{\partial,6}\|_{\ell_j,0}
                    \le Cx_j^2 ,
\end{equation}
where the displayed output tensor bounds describe their
coordinate units. The incoming structure also allows the fixed
short zero-two-jet head from physical initialization and a
filtered zero-two-jet head in the geometric units specified
below. The stable-coordinate derivative is strictly contracting
in an equivalent product norm, with the marginal kinetic
coordinate held fixed. The normalized mass and quartic rows
obey Assumption~\ref{ass:path}. The additional kinetic row
obeys Proposition~\ref{prop:bulk-kinetic-source}, with bounded
normalized arrows from \(z_j/(\Lambda_R x_j)\) to the other
rows. The output retains agreement with the bulk off the new
boundary family.
\end{proposition}

\begin{proof}
\emph{Representatives and weights.}
Use ordinary anchored mass, the common positive filtered primary
quartic, and the kinetic density \(k_x^{\mathrm a}\) of
Lemma~\ref{lem:cube-ordinary-W-d4}. This density is used only
on stencils a fixed positive fine depth from the physical boundary.
The ordinary quadratic Hessian is then its full quadratic
Fourier two-jet plus the two short heads of zero two-jet.
The fixed microscopic head is put in difference-square form,
with its mass and kinetic moments removed before it is assigned
to ordinary anchors. The finite truncation can be performed first
and its second moment then removed exactly. Subtract a central
difference square with the same cubic second moment, and transfer
that square to \(k^{\mathrm a}\) by the exact divergence identity
in Section~\ref{sec:cube-joint-reblocking}. Its interface piece
belongs to the initialized boundary tensor. The exponentially
small discarded quadratic tail is assigned as in physical
initialization. The retained fixed head has bounded span
\(O(b_0)\), second absolute moment \(O(\lambda x_0)\), and fourth
absolute moment \(O(\lambda x_0R^2)\).

For the initial off-shell kinetic variable, multiply the reference
on each ordinary initial anchor by the exponential of its assigned
kinetic change, and multiply an occupied component by the product
of these exponentials on its actual labels. Keep the full boundary
quadratic fixed in this operation. On the real physical parameter
the kinetic change is zero relative to the bare coefficient, so
this is the original initialized interaction. For all off-shell
values the additive distribution makes the circle product identity
exact, just as for initial mass reweighting. In the initial
ordinary region all central stencils have depth comparable to
\(b_0\), so a change of size \(O(x_0)\) has small-radius quadratic
cost \(O(x_0)\) per label and large-radius cost \(O(\sqrt{x_0})\).
The uniform initial field reserve absorbs this exponent on
the enlarged disc, with an additional fixed fraction if the
mass is varied simultaneously. Thus the new initial activity
has size \(O(U)\), and a unit normalized kinetic variation
has the same bound by Cauchy's inequality. Ordinary anchors
outside the boundary family are assigned in the same way in
the two initializations, preserving their agreement.

Use the filtered quadratic density
\begin{equation}
 H_x^T(\phi)=(T\phi)_x^2-\phi_x^2-c_T\phi_x\Delta_{\mathrm c}\phi_x
\end{equation}
from Lemma~\ref{lem:kinetic-reorient}. Write its raw Hessian
coefficient as \(q_{T,j}\), so its contribution to the potential
has coefficient \(q_{T,j}/2\). Its second absolute-moment
parameter is at most \(C R^2 |q_{T,j}|\). At scale \(b=b_j\)
define the effective coefficient
\(m_{T,j}=(R/b_j)^2R^2q_{T,j}\).
Its unit is
\begin{equation}
 h_j=\Lambda_Rx_0 a_1^j,\qquad L^{-2}<a_1<L^{-1}.
\end{equation}
On the enlarged boxes this gives
\begin{equation}
 \frac{R}{b_j} R^2|q_{T,j}|
  \le C\Lambda_R M_R x_0 (La_1)^j\le C_Lx_j.
\end{equation}
The bound uses the reciprocal comparison for \(x_j\) and
\(\Lambda_R M_R=o(1)\). At the merged step
\(R^2|q_{T,0}|\le C\Lambda_R M_R^2x_0\le Cx_0\).
These are the strengthened estimates for the effective orientation
row and the mixed kinetic source. In particular the
enlarged domains allow \(R^2|q_{T,j}|\) to be larger than its
value on the physical trajectory. The effective quartic
coefficient norm of Lemma~\ref{lem:kinetic-reorient} is \(O(x_j)\)
in these domains.

\emph{Ordinary reassembly.}
Write \(G\) for the raw primary coupling, \(u=b^{-2}\nu^{\rm loc}\),
\(S_j\) for the cumulative bulk covariance and
\(\gamma_j=(T\Gamma_jT^*)(0)\) in the bulk. First heat on the
ordinary quartic provides the full filtered square
\(3\gamma_jG(T\phi)_x^2\) in Hessian units. The old localized quartic
provides another such square with coefficient \(3\gamma_j a_{4,j}\),
where
\begin{equation}
 a_{4,j}=9G^2\sum_r (TS_jT^*)(r)^2+4Gu\sum_r S_j(r).
\end{equation}
The first term follows from the two-cross quartic pair, the second
from the one-cross mass pair. A zero-row quadratic has zero
constant-field quartic projection. The first-heat square splits
literally into its mass, \(k^{\mathrm a}\) part, and \(H_x^T\);
the same split is used for the square from the old localized
quartic. The ordinary outgoing Hessian two-jet is the projection
of the full \(P_{\Gamma_j}U-D\), with the total kinetic coefficient
including those two filtered squares. Thus the remaining head
coefficient has the exact scalar equation
\begin{equation}
 q_{T,j+1}=q_{T,j}+3\gamma_j(G+a_{4,j}).
\end{equation}
This is the equation before extraction when the input is \(U\),
and the same formula with the extracted ordinary coefficients
is used on the enlarged activity disc. The effective transport of
the old term is \(L^{-2}\). The absolute effective forcing is at most
\begin{equation}
 C_L\lambda^{\rm tad}_j (R/b_j)^4x_j,\qquad
 \lambda^{\rm tad}_j=
 \begin{cases}C\lambda M_R^2,&j=0,\\C\lambda,&j>0.\end{cases}
\end{equation}
Here \(3|\gamma_j|\le C_L\lambda^{\rm tad}_j b_j^{-2}\) and
\(|a_{4,j}|\le C_Lx_j\) by the cumulative bubble sum and the
reciprocal comparison. Dividing by \(h_{j+1}\) gives a uniformly
small bound, since \(L^{-4}<a_1\). The same estimate holds for
a unit normalized potential derivative by the explicit formula
and Cauchy's inequality. Thus this head has a strict stable
diagonal \(L^{-2}/a_1\), with small extra normalized entries.

Let \(\widetilde X_{S,i}\) be the anchored ordinary covariance
polynomial using the antisymmetric reorientation of
Lemma~\ref{lem:kinetic-reorient}, and let
\(W_i=(1-\mathrm{Loc})\widetilde X_{S_j,i}(U)\).
In the ordinary region let
\begin{equation}
 D_i=\bigl(\mathrm{Loc}\,\widetilde X_{S_j+\Gamma_j}(P_{\Gamma_j}U)\bigr)_i
       -P_{\Gamma_j}\bigl(\mathrm{Loc}\,\widetilde X_{S_j}(U)\bigr)_i .
\end{equation}
Use the effective fresh row of Lemma~\ref{lem:kinetic-reorient}
in the covariance reapportioning. The anchorwise identity gives
\(R_i-D_i=(1-\mathrm{Loc})\widetilde X_{S_j+\Gamma_j,i}(P_{\Gamma_j}U)\).
Here all ordinary coefficient formulas are evaluated on the bulk
comparison of the extracted input where appropriate. These
identities use the actual anchored representatives, including
the split of the first-heat square. If a fine label contains both
site densities and a previously grouped ordinary polynomial, use
the latter's sum of anchored densities in this linear calculation.
The product errors from grouping are counted only after the two
orders have been matched.

Both \(R_i\) and the displayed remainder have \(O_L(x_j^2)\)
fine jet bounds. For \(R_i\), use the fixed-stencil Gaussian
moment for \(P_{\Gamma_j}W_i\), and the bounded effective
fresh row. For the remainder, use
Lemmas~\ref{lem:cube-ordinary-W-d4}, \ref{lem:cube-fixed-short-head}
and~\ref{lem:kinetic-reorient}. Hence \(D_i=O_L(x_j^2)\).
Its quartic is along the primary density. As a consequence
the ordinary positive sector persists throughout the real
coefficient segments used for the product estimates. The change
of \(q_T\) in the ordinary first-heat input and in the literal
outgoing potential differs only by \(3\gamma_j a_{4,j}\).
In the strengthened units for the row reorientation its size
is at most \(C_L\lambda^{\rm tad}_j(R/b_j)^3|a_{4,j}|\).
Using
\begin{equation}
 |a_{4,j}|\le
 C_L[\lambda^2(1+j+\log(2+M_R))+\lambda]x_j^2
\end{equation}
makes this \(O_L(x_j^2)\), including the merged input.
The other ordinary coefficients in the literal outgoing
potential differ from \(P_{\Gamma_j}U\) by \(O_L(x_j^2)\),
by the raw mass and kinetic estimates in
Proposition~\ref{prop:bulk-kinetic-source} and the quartic
projection. Thus replacing the ordinary bilinear factor
evaluated at \(P_{\Gamma_j}U\) by the one evaluated at the
literal outgoing coefficients has size \(O_L(x_j^3)\)
on the same anchors. The local comparison of
Section~\ref{sec:cube-joint-reblocking} applies.

\emph{The common collar and ordinary margin.}
At each scale mark a fixed number of cell layers along the
physical boundary. For a side length not divisible by the grid
size, start from the same corner and merge the last incomplete
interval into the previous full interval in each coordinate.
The fine merged interval starts within the last coarse interval,
and all other fine boundaries respect the integer ratio of the
two grids. The Cartesian products thus give nested rectangular
labels. The clipped labels belong to the boundary family. Keep the disjoint nested cell labels and, if
grouping at a face or an edge is used, group only a bounded
number of adjacent cells per label. Thus each label has volume
comparable to \(b_j^4\). Inside that boundary family keep a
collar of thickness \(C_*b_j\) at scale \(b_j\); the remaining
anchors there are its ordinary margin. Choose \(C_*\) larger than the
range and support constants in coarse units, and leave a fixed
additional buffer of coarse labels for ordinary suffixes of
the enlarged stencils. Thus at the step to \(B=Lb\) the new
collar can have thickness \(C_*B\). It contains the old
boundary stencils, their localization halo, and every anchor
whose incoming covariance row or polynomial support reads
nonbulk data. This uses finite propagation at range \(O(B)\);
it does not require the range constant in fine units to be
independent of \(L\). The incoming boundary tensors still
occupy only a fixed width in fine units. The ordinary expression
applies to every margin anchor outside the buffer. One can
implement the buffer by using the slightly enlarged collar in
the boundary labels and reducing the margin accordingly;
the width and stencil enlargement constants remain fixed in
coarse units.

Use the actual finite-cube covariance row on these supports.
Add the ordinary antisymmetric correction on pairs for which
it is required by a new ordinary anchor; complete it
antisymmetrically on the adjoining collar anchor. Elsewhere in
the collar one may use the symmetric row. Both pair totals
and support bounds are unchanged. All added rows have fine
jet norm \(O_L(x_j^2)\) by
Lemma~\ref{lem:kinetic-reorient}; they need agree with the actual
finite covariance only in their symmetric pair totals.
The definition of \(R_i\) then uses the whole \(P_{\Gamma_j}W_i\)
on each old support. On each new inner anchor take the complete
constant and quadratic part of \(R_i\) into \(D_i\).
The quartic part put in \(D_i\) is just that needed to change
the local primary coefficient on its actual filter row to
the outgoing common value. The difference in these scalar
quartic coefficients is \(O_L(x_j^2+\tau)\).
Put the remaining homogeneous degree-four and degree-six
coefficients in the corresponding inner polynomial factor,
and record constants separately. Define the new inner
quadratic by summing the literal \(P_i-D_i\) there and
subtracting only the common mass on those same anchors.
All kinetic and short-shape quadratics inside the collar are
therefore in its full tensor. Exponents match exactly, and
the linear polynomial factors match through second order.
Grouping their higher products has the mixed cubic cost
proved in Section~\ref{sec:cube-joint-reblocking}.
No summation-by-parts equality of distinct finite-volume
cell factors is used.

For the linear coefficient bounds on this chart, inherited
boundary tensors are supported within \(O(b)\) of the
physical boundary on disjoint fine labels. They therefore
have the gains of the full collar heat estimate, namely
\(C L^{1-2p}\), \(C L^{-1-4p}\), and \(C L^{-3-6p}\)
for the matched degrees two, four, and six. The constants
are independent of \(L\). They follow by counting at most
\(C L^3\) fine boundary labels in one coarse enlargement,
and restricting each external direction with factor
\(C L^{-1-p}\). The Gaussian heat off-diagonal terms with
\(r-m\) contracted directions add the factor
\(C_L\sigma_j^{(r-m)/2}\), where \(\sigma_j=o(1)\) also at
the merged step. These estimates apply equally to face,
edge, and corner stencils since the union of the face
layers overcounts their intersections.

There are also genuine linear terms on newly promoted
ordinary anchors. For those stencils the smallest physical
depth is a fixed positive multiple of \(b\); shallower
stencils were already stored in the old tensor or its
fixed halo. A coarse test direction on a stencil at
depth \(t\) has value bounded by \(C\ell_B(t/B)^p\)
and a difference of order \(a\) bounded by the
same value times \(t^{-a}\), with \(t\) capped at \(B\).
For a kinetic density the two legs thus cost
\(C\ell_B^2(t/B)^{2p}t^{-2}\). There are at most
\(C B^3\) sites in a unit depth layer. The kinetic
contribution in normalized coarse units is at most
\begin{equation}
 C |z^{\mathrm{loc}}| B^{1-2p}
       \left(\sum_{c b\le t\le B}t^{2p-2}
                  + B^{2p-1}\right)
       \le C|z^{\mathrm{loc}}|.
\end{equation}
The term at depth comparable to or larger than \(B\)
uses the capped bound. The analogous polarized bound for
\(H^T\) or the fixed short head has \(t^{-4}\) in
place of \(t^{-2}\) and coefficient \(C M_jR^2\),
where \(M_j=R^2|q_{T,j}|\) for \(H^T\) and \(M_j\)
bounds the second absolute moment for the fixed head.
The polarized fourth-order angular cancellation is the
same as in Lemma~\ref{lem:cube-fixed-short-head}.
It gives
\begin{equation}
 C M_j R^2 B^{1-2p}
       \left(\sum_{c b\le t\le B}t^{2p-4}
                       +B^{2p-3}\right)
       \le C M_j(R/b)^2 L^{1-2p}.
\end{equation}
At the lowest scale choose the finite-head cutoff inside
the fixed initial ordinary stencil buffer, as permitted by
the exponentially small tail bound. Thus the line segments
used in the Taylor estimate lie at comparable depth even
at that scale. If a stencil in the fixed halo reads a
nonordinary anchor, the whole halo is kept in a full tensor
instead. These linear arrows are bounded in the quoted
units, because \(|z^{\rm loc}|=O(\Lambda_R x_j)\) and the
effective head weights are \(h_j=O(\Lambda_R x_j)\).
They are arrows into the collar. There is no reverse
linear arrow into an ordinary coefficient.

The uncontracted bare mass density and its engineering
rescaling are identical in the two charts on every anchor.
Its linear contribution cancels exactly upon subtracting
the common outgoing mass. The first-heat
quadratic compares \(3G(T\Gamma_jT^*)_{xx}\) with its bulk
coefficient on the corresponding filter rows. Applying
\eqref{eq:boundary-tadpole-source} with \(s=2\) on the entire
\(O(B)\) collar gives \(C L^2\lambda_j^{\rm tad}|G|\) per coarse
tile, including the difference of the filtered and bare squares
at promoted anchors. Its normalized quartic derivative is at most
\(C_L\lambda_j^{\rm tad}/\Lambda_R\). In the inverse mass chart \(G\) depends only on \(g_j,\nu_j\).
The derivative of this source with respect to the residual kinetic
coordinate at fixed covariance is therefore zero.

All remaining nonlinear sources and the old ordinary
polynomial factors promoted into the collar have fine
coefficient size \(C_Lx_j^2\). This follows from
polarization and the ordinary comparison on the same
anchors. In the genuinely clipped covariance rows use the
two centered Gaussian increments, which together cost
\(C_L\sigma_j\) times the product of the fixed-stencil
potential norms. The effective antisymmetric rows are
bounded by their ordinary analytic estimates. The resulting
coarse tensor bound is \(C_Lx_j^2\) since the number of
fine labels inside the new fixed coarse collar is \(O_L(1)\)
per coarse label. The same statement holds for the residual
mass after its common linear density has been removed.
Cauchy's inequality in the normalized boxes gives
\(O_L(x_j/\Lambda_R)\) for its contribution to the
quadratic row after division by \(\Lambda_R x_{j+1}\).
For derivatives from a quadratic coefficient into a
quartic or sextic factor, bilinearity gives \(O_L(\Lambda_R)\)
after division by \(x_{j+1}^2\). The bound uses the effective
norm of a head on its ordinary polynomial stencils. The
ordinary margin factors have no boundary derivative and
are subject to the same ordinary bilinear variation bound.
The normalized heat \(6\) to \(4\) is \(O_L(\sigma_j)\);
the heat \(4,6\) to \(2\) is
\(O_L((\sigma_j+\sigma_j^2)x_j/\Lambda_R)\).
Thus all reverse and nonlinear variations in the stable
block are small, apart from the stated triangular
linear arrows. Choose \(L\) large for the diagonal
gains, and then give the collar a smaller relative weight
than the ordinary stable coordinates. This absorbs the
bounded one-way stable arrows. The marginal input is kept
independent in this stable-block estimate.

\emph{Scalar rows and the activity domain.}
The raw ordinary scalar quartic obeys
\begin{equation}
 G_+-G=-\mathfrak b_jG^2-4G(u_1\widehat S_{j+1}(0)
                                          -u\widehat S_j(0)),
 \qquad u_1=u+3\gamma_jG .
\end{equation}
The proof is the constant-field projection used for \(a_{4,j}\),
since the kinetic and both short heads have zero row.
Multiplying by the compensated factor
\(1+4u\widehat S_j(0)\) cancels the second term to cubic
order, by the raw coefficient bounds in
Proposition~\ref{prop:bulk-kinetic-source}. The mass
equation after subtracting its engineering term has forcing
\(3B^2\gamma_jG+O_L(x_j^2)\) in outgoing units.
The only potentially unbounded cumulative bubble term
in that error cancels exactly as proved in the same
proposition. The compensated quadratic projection removes
the leading pure-quadratic variation. Its inverse changes
a unit raw mass variation by \(1+O(\lambda x_j)\) and the
remaining raw coefficients by \(O(\lambda x_j)\), as in
Section~\ref{sec:cube-mass-chart} on this shrinking box.
The normalized mass entries from the residual mass error
are small since \(x_j/\Lambda_R=o(1)\). The kinetic
source estimate itself is Proposition~\ref{prop:bulk-kinetic-source}.
All these scalar maps are holomorphic on the larger
coefficient box, with the asserted bounds for their
first variations by Cauchy's inequality.

On the extraction disc the uniform ordinary displacement
of mass, quartic and kinetic is \(O_L(\tau)\) in normalized
raw fine-scale units. Implement the extra spatial-moment
densities by the exact exponential divergence reassignment
of Section~\ref{sec:cube-joint-reblocking}, using
\(k^{\mathrm a}\) instead of the central difference
square. The scalar kinetic span is unchanged. The
nonuniform displacement stays local in the fine halo
and is assigned to the new collar. No independent
filtered head displacement is extracted. Its formula
above changes by \(C_L\lambda^{\rm tad}_j(R/b_j)^4\tau\)
in its effective units, because
\(\lambda^2[1+j+\log(2+M_R)]|G|\le C_L\).
The new collar extraction difference is \(O_L(\tau)\)
and the ordinary margin difference is \(O_L(x_j\tau)\),
by the local matching estimates and their bilinear
variation. These statements hold on the entire enlarged
disc, and their first variations follow as in
Proposition~\ref{prop:cube-bounded-extraction}. On the
shooting boxes the additional normalized factors are
\(f_j/(\Lambda_R x_j)\) and \(f_j/x_j^2\), both \(o(1)\)
by the microscopic comparison. The same bounded matching
proves the connected disc bound and the strictly
contracting activity derivative by that proposition.
Thus extraction and transport use the same outgoing reference on
each anchor.
\end{proof}

The number of cells on a face grows with its area. To obtain a
uniform bound when an ordinary label enters the boundary family,
we sum over connected covers containing one fixed cell.

\begin{lemma}[Anchored boundary promotion on a cube]
\label{lem:anchored-boundary-promotion}
Use the bounded-degree graph of cube cells of scale \(b_0\).
Let \(K^0(X)\) be its connected
activity at zero boundary defects, with
\(\sup_D\sum_{X\ni D} e^{A_0|X|}\|K^0(X)\|_{r,*}\le C\tau_r\).
Assume the single-label reference factors have uniformly bounded
norm in the same reserve. Suppose the first marked
terms obtained by Cauchy extraction on the boundary-link
discs have bounds
\begin{equation}
 \begin{aligned}
 \|D_{X,e}\|_{r,*}\le C R^{-4}\tau_r
   e^{-A_0|X|}
   e^{-a d_\partial(B,C)/R},\qquad e=\{B,C\}\subset X,\\
 \text{where}\qquad
 d_\partial(B,C)=\operatorname{dist}(B,C)
                +\min(\delta(B),\delta(C)).
 \end{aligned}
\end{equation}
The norm includes the spare field reserve, \(X\) is graph connected,
and \(\delta(B)\) is its distance to the physical boundary.
For every fixed output exponent \(A_1\), take \(A_0\) and
\(b_0/R\) sufficiently large and \(\tau_r\) sufficiently small.
There are connected covers \(X\cup J_e\) reaching the grouped
boundary layer with
\begin{equation}
 \sup_D\sum_{X,e:\ D\in X\cup J_e}
 e^{A_1|X\cup J_e|}\|D_{X,e}\|_{r,*}\le C_{A_1}R^{-4}\tau_r
\end{equation}
uniformly in the size of the cube. Regrouping the circle product on
these covers preserves the field density. The activity of components
with at least one mark has anchored norm at most \(C_{A_1}R^{-4}\tau_r\).
On components avoiding the boundary labels the activity equals \(K^0\).
\end{lemma}

\begin{proof}
If \(K_X(\mathbf z)\) is the interpolated activity with physical
defects \(\mathbf z\), integration along \(t\mathbf z\) gives
\(K(X)=K^0(X)+\sum_e D_{X,e}\), where
\(D_{X,e}=\int_0^1 z_e\,\partial_e K_X(t\mathbf z)\,dt\).
Attach to \(e\) a path from one of its cells to a nearest face and,
if needed, a path between its endpoints. These paths form \(J_e\).
They use at most \(C(1+d_\partial(B,C)/b_0)\) cells, and the
distance in the cell graph from either endpoint to any cell of
\(J_e\) obeys the same bound.

For a root \(D\in X\), choose \(X\) rooted at \(D\) and sum over at most
\(|X|^2\) marks, dropping the remaining defect weight. The connected
set count through \(D\) is at most a fixed exponential in \(|X|\).
For \(D\in J_e\), both endpoints lie within \(C(1+k)\) grid
units of \(D\) if \(k=\lceil d_\partial(B,C)/b_0\rceil\).
There are at most \(C(1+k)^{2d}\) choices. For each pair sum over
connected \(X\) rooted at \(B\). The cover size weight uses at
most \(C A_1(1+k)\) from the exponential decay in the distance,
which leaves a convergent sum in \(k\).

Expand the original circle product into packings of its
nonadjacent connected actual supports. Each occupied support is
now assigned either to \(K^0(X)\) with cover \(X\) or to its
selected marked term with cover \(X\cup J_e\). Group the covers
that overlap or touch and insert the old reference factors on
their unused labels. This is an exact regrouping because
the actual occupied supports are disjoint, independently of
their covers, and the filler labels are exactly the unused
labels consumed by grouping. The field reserve multiplies on
the disjoint actual labels, even if their tested neighborhoods
intersect. A grouped component having no mark is just one
unmarked activity, since the original actual supports are
mutually nonadjacent.

For completeness, the local norm bound for the regrouping
follows from a rooted tree sum. Absorb filler constants and the
desired output size weight into the spare exponents of the
cover weights. A parent cover \(Y\) can attach each child at
one of \(O(|Y|)\) possible labels, the labels in its closed
graph neighborhood. A connected family containing a given
root has a spanning tree of touching covers, and may be
overcounted by specifying an ordered list of children at
each such label. If \(u\) bounds the resulting total weight
of descendant trees rooted at any label, the children's lists
cost at most \((1-u)^{-C|Y|}\). A fixed spare size exponent pays
this factor for sufficiently small \(u\). Induction on the
tree height bounds \(u\) by \(C\tau_r\) from the cover-anchored
input. To bound the subseries with at least one mark, multiply
all marked weights by \(R^4\) in this positive majorant.
Its anchored input is still \(O(\tau_r)\), so the same
induction applies. Returning to the original weights gives
the gain \(O(R^{-4}\tau_r)\) for the marked subseries.
All its output covers reach the boundary, so components avoiding
the boundary labels retain their unmarked activity.
\end{proof}

\section{Physical initial data and terminal covariance}
\label{sec:physical-entry}

To apply the one-step estimates, the effective interaction at the initial
scale must represent the spin measure at the chosen temperature. We
construct the initial polymer activity with the mass left free, then
identify the mass at which the representation is physical. In dimension
four the same construction leaves the kinetic coefficient free as well.

\subsection{Joint initialization in the cube chart}
\label{sec:cube-joint-initialization}

To initialize the bulk and boundary charts in the same activity
norm, we must account for the two truncations in the microscopic
density. The pressure terms use the full inverse of the short
filter, whereas the reference on unused labels uses its finite
Neumann truncation. We keep this reference fixed and expand the
filter tail together with the quadratic tail. Recording both cell
labels of each filter entry, including the column label, permits
an exact comparison with the bulk off the grouped boundary labels.

\begin{proposition}[Joint initial activity and its mass derivative]
\label{prop:cube-joint-initialization}
Use initial cells of side comparable to \(b_0=RM_R\), with
\(M_R\) a sufficiently large fixed multiple of \(\log R\). Use
fixed enlarged cell stencils and the raw field tests \(F_{b_0}\)
of Section~\ref{sec:geometry}. On the boundary use the same tests
on the clipped stencils. Clipped cells can be joined to adjacent
cells so that the resulting labels have volume comparable to
\(b_0^d\). Let \(T\) be a finite Neumann truncation of the frozen
primary filter \(H_F=B_\Lambda^{-1}\). Decompose the primary-field
Hessian into its bulk scalar row, a finite head, and the remainder,
as specified below. There is an exact circle-product initialization
with the quartic reference at \(T\) and the finite quadratic head.
Its activity \(K_{\rm init}(\beta,\nu)\) is equal to the corresponding
bulk activity on every polymer avoiding the grouped boundary labels.

Set \(g=g_{\mathrm{mic}}\), \(G_0\asymp A_{\rm sc}^4g\), and
\begin{equation}
 U=A_{\rm sc}^6g^2+G_0^{9/4}g^{1/2}+G_0^3 .
\end{equation}
At the two initial radii, use the growing regulator and the
decaying and quotient regulator, respectively. With any fixed
exponential polymer weight and a fixed reserve in both regulators,
the mixed activity norm obeys
\begin{equation}
 \sup_{|\nu|\le cG_0}\|K_{\rm init}(\beta,\nu)\|_{\mathrm{mix},0}
 \le C U .
 \label{eq:cube-joint-init}
\end{equation}
The activity is holomorphic in this disc, and on a fixed smaller
disc its mass derivative is at most \(CU/G_0\) in the same norm.
The estimates are uniform in the volume and in a fixed physical
window \(|\beta-1|\le A_\beta g\) for \(n\ge n_0(R)\).
\end{proposition}

\begin{proof}
\emph{Filter and Hessian tails.}
All calculations of filter kernels below are in the physical field
\(h=A_{\rm sc}R^{-d/4}\phi\). Write \(\delta(x)=\dist(x,\Lambda^c)\).
The Neumann expansion, the short-range row gap and the weighted
Hessian estimate of Lemma~\ref{lem:extension-cavity} give
\begin{align}
 H_F&=T+E^0+E^\partial, \\
 \sup_x\sum_y e^{a|x-y|/R}
 \bigl(|E^0(x,y)|+e^{a\delta(x)/R}|E^\partial(x,y)|\bigr)
 &\le \epsilon_R . \label{eq:cube-init-filter-tail}
\end{align}
The same estimate holds with rows and columns exchanged. Here
\(E^0\) is the restriction of the bulk Neumann tail. In the
difference \(E^\partial\) every walk has length exceeding the
truncation order and visits the exterior. The geometric row bound
therefore pays both the truncation and the extra boundary weight
after reducing \(a>0\). The truncation order and \(M_R\) can be
increased by fixed multiples of \(\log R\) so that
\(\epsilon_R\le R^{-M}\) for any fixed \(M\), with all finite
stencils still in a fixed number of cell collars.

Remove the bulk scalar row from the primary-field Hessian.
Write its bulk zero-row kernel \(Z^0\) in difference-square form
\(-\frac14\sum_{x,y} Z^0(x,y)(h_x-h_y)^2\), summing over the
actual cube sites. The correction to the diagonal is then the
missing row \(-\sum_{y\notin\Lambda} Z^0(x,y)\). Include it with
the finite-volume Hessian discrepancy. The bulk kernel has
exponentially weighted absolute rows and columns of size \(C g\).
The boundary discrepancy, including the missing row, has the
same bound with the extra factor \(e^{a\delta(x)/R}\) in
\eqref{eq:cube-init-filter-tail}, with \(\epsilon_R\) replaced by
\(C g\). For the Hessian pullback this follows also by expanding
the two conjugating factors in short walks. In a discrepancy one
factor is a Hessian discrepancy or a walk reaches an exterior
summation site. In either case the weighted bound and
\(\delta(x)\le\delta(y)+|x-y|\) retain the boundary weight at
the initial site. Take the bulk head for \(|x-y|\le c_1 b_0\)
and the boundary head for
\(|x-y|+\min(\delta(x),\delta(y))\le c_1 b_0\), for a fixed
\(c_1>0\). Assign these heads by fixed symmetric local weights
to their incident labels, using grouped boundary labels for the
boundary head. The remainder has the same estimates as the
original kernels with bound \(C g e^{-c_2 M_R}\) after reducing
\(a\). The diagonal correction belongs entirely to the boundary
part. These conventions leave the bulk pair density unchanged
when both endpoints are present.

\emph{Interpolation of cell pairs.}
For a cell pair \(e=(B,D)\), let \(d_e=\dist(B,D)\). For
boundary terms use \(d_e^\partial=d_e+
\min(\delta(B),\delta(D))\), where \(\delta(B)=\min_{x\in B}\delta(x)\).
Choose a tail variable for each ordered block of \(E^0\) or
\(E^\partial\), and independently for each unordered cell pair
of the bulk or boundary quadratic remainder, including the
diagonal pairs. Each variable multiplies its block or its exact
signed pair contribution. Write \(D_e=d_e\) or
\(D_e=d_e^\partial\) according to the type of variable. For a
small fixed \(c_3>0\) use two polydiscs, with radii
\begin{equation}
 \rho_e^+=\epsilon^{-1/4}e^{c_3D_e/R},
 \qquad
 \rho_e^-=\epsilon^{-1/8}e^{c_3D_e/R},
\end{equation}
where \(\epsilon=R^{-M}\). By increasing \(M_R\) and the
Neumann cutoff before fixing the cell collars, we can take both
weighted tail bounds at most \(\epsilon\). Hence on the larger
polydisc the common absolute majorant of the amplified tails has
weighted row and column bounds \(C\epsilon^{3/4}\). This applies
also when several types of variables have the same pair of cells.
For a quadratic difference-square term the estimates control the
absolute pair coefficient, and hence the diagonal contribution
as well as the off-diagonal contribution.

Apply the cell BKAR formula to the Hessian-subtracted short
pressure. On each nonsingleton tree component \(C\), integrate
the tree derivative on its induced short graph, evaluated on
the filter rows of \(C\). On a singleton \(B\), subtract the
fixed quartic filler \(I_{4,T}(B)=\exp[-V_4((Th)_B)]\)
from the exact singleton pressure weight. Denote these two
types of pressure atom by \(p_C\). Since the Hessian-subtracted
pressure weight is a quadratic exponential times a finite spin
partition function divided by its zero-field value, it is entire
in the external field. Thus the pressure formula holds for
complex tail variables, with \(p_C\) holomorphic. In the
pressure packing the marked atoms are disjoint, and all its other
labels carry the quartic fillers. The finite quadratic head
factor on a label is present whether or not the label is marked.

We first estimate the product for a fixed packing, restricted to
a set of labels \(W\) that contains all its marked atoms and
both endpoints of every selected tail variable. Other tail
variables are zero. Insert the filler, with its finite head, at
every unused label of \(W\). Let \(F_W\) be this product times
the exponential of the selected quadratic tails. Denote the
two physical radii by
\begin{equation}
 e=A_{\rm sc}R^{-d/4}b_0^{-d/4},\qquad
 a=b_0^{-d/4}.
\end{equation}
In the remainder of this proof \(r\) denotes either physical
radius. Define
\(\tau_r=b_0^d(r^6+gr^4)\). We use normalized squares
\begin{equation}
 Z_B^2=\frac{\|h\|_{2,B}^2}{b_0^d r^2},\qquad
 Q_B^2=\frac{q_{B^*}(h)^2}{r^2},
\end{equation}
where \(q_{B^*}\) is the adapted affine quotient on the fixed
stencil \(B^*\). Its defining affine decomposition and the
finite-dimensional affine norm comparison on \(B\) give
\begin{equation}
 \frac{\sup_{B^*}|h|^2}{r^2}\le C(Z_B^2+Q_B^2).
 \label{eq:cube-init-quotient}
\end{equation}
Indeed the residual is controlled in supremum by \(q_{B^*}\),
and the \(L^2\) norm of the affine part on \(B\) controls its
supremum on \(B^*\). These comparisons are uniform for a clipped
label joined to a neighbour because it contains a comparable
cube and is contained in a fixed enlargement of that cube.

For real tail variables in the larger polydisc all marked rows
use \(T\) plus an error with absolute row and column bounds
\(C\epsilon^{3/4}\), and the fillers use \(T\). The same is true
of any convex mixture of those rows, with the same absolute
majorant. For \(x\) in \(W\) at distance more than the range
of \(B_\Lambda\) from the artificial interfaces of \(W\), the
identity \(B_\Lambda T=I-Q_\Lambda^{m+1}\) recovers \(h_x\) using
only those rows that are available in \(W\). The operator norm
error is \(C\epsilon^{3/4}\) by the row and column bounds. Its
input consists of the fixed stencils and the selected column
labels. The omitted strip in any cell has volume
\(O((R/b_0)b_0^d)\). Squaring, summing and using
\eqref{eq:cube-init-quotient} in these strips and for the
stencil input gives, for the available filtered rows \(u\),
\begin{equation}
 \sum_{B\in W} Z_B^2
 \le C\frac{\|u\|_{2,W}^2}{b_0^d r^2}
    +C(R/b_0+\epsilon^{3/2})\sum_{B\in W}(Z_B^2+Q_B^2).
 \label{eq:cube-init-stitching}
\end{equation}
Here \(m\) is the order of the truncation. Physical-face rows
are handled by the restricted short inverse and need no
extension of \(h\) through the physical face.

Lemma~\ref{lem:pressure-singleton} and Proposition~\ref{prop:pressure-trees}
give, at a real filter, damping by \(e^{-c\sum\Psi(u_x)}\)
on all the marked rows with a factor \(\tau_r\) per marked atom.
One may take an arbitrarily small additional loss in the
normalized cell supremum squares, and any fixed exponential size
margin on the nonsingleton atoms. To use the fixed singleton
filler, split the singleton difference at \(e^{-V_4(u)}\).
The difference of the remaining quartic exponentials, by
integration along their segment, costs at most
\(C\epsilon^{3/4}b_0^d r^4\) times a fixed polynomial in the
normalized cell suprema. This is smaller than \(\tau_r\).
The exponential in the segment formula damps a convex
combination of the two row fields, by convexity of the quartic.
The estimates through a fixed number of field jets hold in the
same way because filtered test directions have supremum at most
\(Cr\). When the derivative falls on the outer finite head or on
the selected quadratic exponential, its polynomial factors
have quadratic coefficient bounded by \(C(g+\epsilon^{3/4})\).
At either radius
\begin{equation}
 b_0^d r^2(g+\epsilon^{3/4})=o(1).
 \label{eq:cube-initial-quad-small}
\end{equation}
Both the finite heads and the tail exponents therefore consume
only an arbitrarily small fixed multiple of the sum in
\eqref{eq:cube-init-quotient}. The cell supremum penalties of
the pressure estimates have the same property, by the
\(\ell^2\) bound on the matrix of cell suprema of the common
filter majorant. That matrix has bounded row and column sums:
the finite head has fixed cell range, and the tails have an
exponentially weighted row and column bound.

At the large radius \(r=a\), \(b_0^d r^4=1\).
For every fixed \(K>0\) and all sufficiently large \(R\),
\(c\Psi(t)\ge K t^2/(b_0^d r^2)-C_K/b_0^d\). For \(|t|\le1\)
this is quartic Young's inequality, and for \(|t|>1\) it
follows from \(\Psi(t)\ge t^2/2\). Take \(K\) to leave a
strict reserve in the required negative \(Z_B^2\) exponents.
In \eqref{eq:cube-init-stitching} the coefficient of \(Q_B^2\)
then tends to zero and fits in any prescribed positive quotient
allowance. At the small radius one keeps a fraction of the growing
regulator directly, using \eqref{eq:cube-init-quotient} and
discarding the real damping. The finite jet Leibniz coefficients,
the singleton polynomial factors and the fixed losses per cell
are bounded by a fixed exponential in \(|W|\). Thus, in either
fixed weighted jet norm with a reserve, denoted by \(\|\cdot\|_{r,*}\),
the real bound takes the form
\begin{equation}
 \|F_W\|_{r,*}\le C^{|W|}
        \prod_{C\ {\rm marked}}(\tau_r e^{-A_1|C|}).
 \label{eq:cube-init-real-packing}
\end{equation}
Here \(A_1\) is arbitrary. On singleton atoms its fixed cost
can be included in \(C^{|W|}\), or retained as a separate
fixed factor per singleton.

For use in Cauchy's formula we also need a bound on the
large complex polydisc. For a pressure graph \(U_C\), with
\(J\) the interpolated short coupling, set
\(A_J=L_J^{\mathsf T}S_JL_J\). The proof of the weighted
Hessian estimate applies to small fixed complex perturbations
of the coupling variables around the tree interpolants:
the absolute row gap and the bound on the largest coefficient
are preserved. Indeed the proof of that estimate only uses the
absolute weighted gap of the complex hyperbolic-tangent matrix
and the small maximum entry. It gives \(A_J=I-J+O(g)\) in absolute rows
and columns. Integrating the zero-field coupling derivative
from zero along the segment \(tJ\) gives
\(|\log(Z_J(0)/2^{|U_C|})|\le C|U_C|g\), since the derivative
uses only off-diagonal spin covariances and each is \(O(g)\).
For \(u=v+iw\), positivity of the real spin weight gives
\(|Z_J(L_Ju)|\le Z_{\operatorname{Re}J}(\operatorname{Re}L_Ju)\).
Compare its field to \(L_{\operatorname{Re}J}v\), using the
bound one on every coordinate of the real pressure gradient.
The difference costs \(C\|w\|_1\).
The quadratic subtractions differ in absolute value by
\(C(g\|v\|_2^2+\|w\|_2^2+\|v\|_2\|w\|_2)\) using
\(\operatorname{Re}(I-J)=I-\operatorname{Re}J\).
The complex pressure weight is consequently bounded by
the real \(\Psi\) damping at \(v\) times the exponential of
these costs and \(C|U_C|g\). For the quartic filler use \(\operatorname{Re}(x+iy)^4\ge
\tfrac12 x^4-Cy^4\) and \(|y|\le Cr\). For a real
displacement \(v_1\), the inequality
\(\Psi(v_0+v_1)\ge c\Psi(v_0)-C\Psi(v_1)\) follows from the
triangle inequality and the bounded ratio \(\Psi(2t)/\Psi(t)\).
Thus displacing the real field by at most \(Cr\) consumes
only a fixed multiple of \(b_0^d r^4\) per label.

For the finitely many tested field jets, evaluate every
factor on \(h+\sum_i z_i v_i\) with \(|z_i|\le1\) and
\(\|v_i\|_\infty\le Cr\) on the stencils. The imaginary part
has \(\ell^2\) norm at most \(C\epsilon^{3/4}\) times the norm
of \(h\) on the stencils and selected columns, plus
\(Cr\sqrt{b_0^d|W|}\); the same bound controls the displacement
between \(v\) and the real row substitution at \(h\), with the
tail term omitted for that displacement. Sum over the atoms
before using Young's inequality, as their row labels are
disjoint. The part quadratic in the background uses a coefficient
\(C(g+\epsilon^{3/4})\) or an arbitrarily small fixed multiple
of the normalized squares in \eqref{eq:cube-init-quotient}.
Young's inequality for the remaining terms costs at most
\(C b_0^{2d}|W|\). Thus the real recovery bound
\eqref{eq:cube-init-stitching} at the real row substitution
still leaves the prescribed reserves. Cauchy's formula on fixed
coupling circles costs at most a constant per tree edge, and
there are at most a fixed exponential number of trees. Cauchy's
formula on the field-jet circles costs a fixed factor. We obtain
the uniform bound \(e^{C b_0^{2d}|W|}\) in \(\|\cdot\|_{r,*}\)
on the large polydisc, with constants independent of the number
of tail variables.

For a point \(\boldsymbol z\) of the small polydisc consider
the common scalar line
\(\operatorname{Re}\boldsymbol z+w\operatorname{Im}\boldsymbol z\).
It stays in the larger polydisc for \(|w|\le c\epsilon^{-1/8}\)
and has real entries on the diameter. At \(w=i\) the semicircle
has harmonic measure \(O(\epsilon^{1/8})\). Apply the two-constants
theorem to each fixed normalized field-jet evaluation and take
the supremum. For \(M\) sufficiently large,
\(\epsilon^{1/8} b_0^{2d}=o(1)\) and
\(\epsilon^{1/8}|\log\tau_r|=o(1)\). The bound
\eqref{eq:cube-init-real-packing} follows on the small polydisc
with a different \(C\) and half the size margin. Thus taking
any fixed desired size margin is possible also there.

We give the exact support expansion and its count. For a fixed
packing and for a subset \(S\) of the tail variables, let
\(\Delta_S F=\sum_{D\subset S}(-1)^{|S|-|D|} F(\boldsymbol1_D)\).
The evaluation at all variables one is \(\sum_S\Delta_S F\),
including \(S=\varnothing\). In a finite volume this is an
identity of finite sums. The term uses precisely the Taylor
monomials with positive exponent in each variable of \(S\)
and zero exponents outside \(S\). By Cauchy's formula and the
small polydisc bound each selected variable contributes a
factor at most
\begin{equation}
 \sum_{k\ge1}(\rho_e^-)^{-k}
 \le 2\epsilon^{1/8}e^{-c_3D_e/R}.
\end{equation}
Associate a connected grid path to the pair of each selected
variable, and in a boundary variable also include a path to
its selected physical face. The supports used to form connected
components are the actual marked atom sets and these paths.
Group touching as well as overlapping supports, and include a
filler for every otherwise unused label in the resulting
component \(W\). An atom reads only the selected filter columns
connected to it by these paths and its fixed finite head stencil.
The latter need not be added to the polymer labels, since field
dependence on these stencils is allowed for the reference and
for every activity. As functions of the selected variables,
different components factor exactly. This is also true of
\(\Delta_S\), since the variable sets in different components
are disjoint. An unused label outside the components keeps its
original reference. Conversely all pairwise nonadjacent choices
of output components and their internal data give precisely the
corresponding terms of the product. This constructs the circle
product on the enlarged support graph without any division
by a reference factor. In particular, at zero adjustable mass,
\begin{equation}
 e^{-Q^{\mathrm{head}}-Q^{\mathrm{tail}}}
 e^{-\mathcal U_\Lambda(H_Fh)}=(I_T\circ K_{\rm init}(\beta,0))(\Lambda).
\end{equation}

\emph{Sum over connected components.}
Regard each marked atom or selected edge with its path as an
object with a

connected cover \(Y\). Include the fixed \(C^{|W|}\) and any
desired output polymer size weight in the per-object exponent
\(e^{c|Y|}\), using \(|W|\le\sum |Y|\). The sum of atom
weights through one root label is \(O(\tau_r)\), by choosing
its size margin larger than the exponential number of
connected sets through that label. For each edge type the
anchored sum through one root is \(O(\epsilon^{1/8})\) after
any fixed size margin, by the connected-cover count of
Section~\ref{sec:geometry}. For completeness, convergence of
the connected-object sum follows by choosing a rooted spanning
tree of touching covers. Root at an object covering the given
label, and assign each child to the first label in the closed
neighbourhood of its parent cover that belongs to the child.
There are at most \(C|Y|\) possible attachment labels at a
parent of cover \(Y\). Allow an ordered list of children at
each attachment label. If the anchored sum of all trees of
height at most \(h-1\) is \(T_{h-1}\), the extra factor at
this parent is bounded by \((1-T_{h-1})^{-C|Y|}\).
Reserve an additional fixed \(e^{|Y|}\) in the input
weights and take the anchored sum \(v\) of the inflated object
weights small. Induction gives \(T_h\le 2v\) for every \(h\),
since \((1-2v)^{-C|Y|}\le e^{|Y|}\). Every connected family
is counted at least once. The resulting rooted activity bound
is \(C(\tau_r+\epsilon^{1/8})\), with the same field reserve.
The bound is uniform in the cube size. Norm tests on a
component restrict to all the stencils used above, and the
estimate retained its reserves on the disjoint actual labels.
In particular the counting step incurs no extra field cost.

A component of the exact activity that avoids boundary labels
contains no boundary variable, because such a selected
variable was assigned a face-reaching path. On interior rows
its finite Neumann head is the bulk head. Its pressure atoms
use the same induced short graph as the corresponding bulk
atoms, and its bulk difference-square tails are exactly the
same as in the bulk calculation. Use the same path assignments
for bulk pairs and average over tied paths or faces when
enforcing lattice symmetry. The ordinary paths can be taken
coordinate-monotone, so a path between two cells of the cube
stays in the cube. No bulk coefficient on a boundary-free
polymer changes due to clipping, because its local stencils
and all retained labels lie outside the clipping collar.
Thus the boundary-free coefficients agree term by term.
The bulk activity on a fixed finite polymer is defined by
these same finite component formulas.

Finally give each label the factor \(e^{-\nu m(B)}\), where
\(m(B)=\frac12 b_0^{-s}\sum_{x\in B}\phi_x^2\) is the
normalized scalar mass density. The product over an occupied
component is exactly \(e^{-\nu m(W)}\), by additivity, including
the fillers in that component. The reweighting preserves the
circle product and the bulk agreement. Its quadratic cost
in the two radii, for \(|\nu|\le cG_0\), is \(O(G_0)\) and
\(O(G_0^{1/2})\), respectively, and fits inside the retained
reserve. It therefore defines a holomorphic family in the
fixed norm. Shrinking the mass disc gives the claimed derivative
by Cauchy's formula. Since
\begin{equation}
 \tau_e\le C A_{\rm sc}^6 g^2,\qquad
 \tau_a\le Cg^{1/2},
\end{equation}
and \(\epsilon^{1/8}\) can be made smaller than \(G_0^3\),
the weighted sum of the two estimates proves
\eqref{eq:cube-joint-init}. All bounds used here persist on
the real thermal window. For fixed finite volume the same family
is continuous in \(\beta\) in the activity norm. The finite coefficient
formulas are continuous on bounded field sets, uniformly over the
tested directions. Outside a bounded field set the strict regulator
reserve makes the norm contribution uniformly small, by
\eqref{eq:cube-init-quotient}. As in the pressure estimates, for
each fixed window width the quartic defect proportional to
that width is absorbed by the sextic seed when \(R\) is
sufficiently large.
\end{proof}

\begin{corollary}[Physical entry into the mass chart]
\label{cor:physical-mass-entry}
For \(d=1,2,3,4\), take the finite heads of
Proposition~\ref{prop:initial-cube-heads} and the joint activity of
Proposition~\ref{prop:cube-joint-initialization}. Fix \(A_0\) and
\(\lvert\beta-1\rvert\le A_0g_{\rm mic}\). For sufficiently large
\(R\) and \(n\ge n_0(R)\), their circle product equals the physical
density at the scalar mass
\(\nu_0^{\rm bare}=\lambda^{-1}M_R^s r_h(\beta)\). This mass is
\(O_{A_0}(\Lambda_R G_0)\) and lies on the smaller mass disc of the
proposition. In the units of Assumption~\ref{ass:path},
\(U/(\Lambda_R G_0)=o(1)\) and \(U/G_0^2=o(1)\) as \(R\to\infty\).
\end{corollary}

\begin{proof}
At the stated mass, the reference factor and the activity reweighting
together restore the scalar term subtracted from the physical Hessian.
At other \(\nu\) the same circle product gives the off-shell family.
The crossing estimate in Section~\ref{sec:micro-extension} gives
the bound on \(\nu_0^{\rm bare}\), and \(\Lambda_R=o(1)\) places it
inside the smaller disc. Substituting the microscopic scales into
\(U\) gives
\begin{equation}
 \frac{U}{\Lambda_R G_0}
 \le \frac{C}{M_R^s}
       \left(g_{\mathrm{mic}}^{1-4a_*}
             +g_{\mathrm{mic}}^{7/4-7a_*}
             +g_{\mathrm{mic}}^{2-10a_*}\right)=o(1),
 \qquad
 \frac{U}{G_0^2}
 \le C\left(g_{\mathrm{mic}}^{2a_*}
             +g_{\mathrm{mic}}^{3/4-a_*}+G_0\right)=o(1).
\end{equation}
Together with \eqref{eq:cube-joint-init} and the derivative bound
of Proposition~\ref{prop:cube-joint-initialization}, these ratios give
the initial activity and mass-derivative estimates in the shooting chart.
\end{proof}

\subsection{Initial and terminal data for the kinetic reference family}
\label{sec:kinetic-reference-endpoints}

The shooting construction uses two endpoint properties of the
kinetic reference family. At the initial end, matching both bare
coordinates must recover the microscopic density. At the terminal
end, changing the reference must preserve the Gaussian estimates
for curvature and magnetic response. The two lemmas verify these
properties before the nonlinear matching.

\paragraph{Exact physical entry with a free initial kinetic coordinate.}
We first extract the kinetic moment from the bulk quadratic and
keep the full boundary quadratic fixed. This determines the
variation of the bare kinetic coordinate while confining the
central difference stencil to the ordinary region.

\begin{lemma}[Initial kinetic extraction and off-shell reweighting]
\label{lem:kinetic-initial-extraction}
Let \(d=4\), \(1/2<p<1\), and use
Proposition~\ref{prop:kinetic-reference-family} with
\(|\delta|\le C\lambda^2G_0\).
After the bulk row is removed, let \(z_0^{\rm bare}\)
be the \(k^2\) coefficient of the retained bulk
quadratic head in Hessian units. Then
\(|z_0^{\rm bare}|\le C\lambda G_0\).
The remainder of its ordinary bulk head has zero row
and zero second moments and absolute moments of degrees
two and four bounded respectively by \(C\lambda G_0\)
and \(C\lambda G_0R^2\). The full physical collar
quadratic, including the interface terms from this
extraction, has norm \(C\lambda G_0\) at radius
\(\ell_{b_0}\).

There is an exact initial activity in the outgoing chart
with independent mass and kinetic coordinates \(\nu,z\),
holomorphic on \(|\nu|\le cG_0\),
\(|z-z_0^{\rm bare}|\le cG_0\), with norm \(CU\).
On a fixed smaller disc both partial derivatives are
at most \(CU/G_0\). It equals the physical interaction
at the bare mass and kinetic coordinates and
agrees with the bulk activity off the grouped
boundary labels. The initial kinetic coefficient has
uniformly negative derivative in the scaled external
parameter \(a=\delta/\lambda\), bounded away from zero
for sufficiently large \(R\).
\end{lemma}

\begin{proof}
The canonical weighted row of the zero-row bulk
head is bounded by \(C\lambda G_0R^{-2}\):
combine Proposition~\ref{prop:initial-cube-heads}
with the density-change bound
\(C a_R^{-1}|\delta|\) in
Proposition~\ref{prop:kinetic-reference-family}.
Its second and fourth absolute moments are consequently
bounded by \(C\lambda G_0\) and \(C\lambda G_0R^2\).
Keep the finite head in anchored difference-square form.
Subtract its scalar second moment times a unit-curvature
difference-square representative with spans \(2e_i\).
The cut head is cubically symmetric, so this removes
all its second moments. The subtraction preserves zero
row and satisfies the same absolute moment bounds.
The starting coefficient has size \(C\lambda G_0\).

For exact placement, separate the ordinary sites \(S\)
from the initial collar so that their stencils are in
the bulk-comparison region and the interface is at
depth at least a fixed positive multiple of \(b_0\).
At those sites only, replace the kinetic difference-square
density by \(k_x^{\rm a}\) of
Lemma~\ref{lem:cube-ordinary-W-d4}, using the same
normalization for the potential and its Hessian. The
difference of the two unit-curvature representatives
is a fixed multiple of
\begin{equation}
 \sum_{i=1}^4
 \bigl(\phi_{x+2e_i}^2-2\phi_x^2+\phi_{x-2e_i}^2\bigr).
\end{equation}
For example the representative
\(\tfrac18\sum_i[(\phi_{x+2e_i}-\phi_x)^2+
(\phi_{x-2e_i}-\phi_x)^2]\) differs from
\(-\phi_x\Delta_{\rm c}\phi_x\) by \(1/8\) times this
display. Summing over \(S\) cancels every internal bond
and leaves the crossing-bond differences of \(\phi^2\).
On the interface the raw first difference bound gives
\(|\phi_y^2-\phi_x^2|\le C\ell_{b_0}^2/b_0\)
for tested directions at radius \(\ell_{b_0}\); the
bilinear bound follows by polarization. At most
\(Cb_0^3\) bonds contribute per interface label.
Since \(\ell_{b_0}^2=b_0^{-2}\), the complete interface
tensor has normalized size \(C|z_0^{\rm bare}|\).
Assign its actual crossing stencils to the collar.
The entire original head in the inner collar is
retained there, including its clipping defect. Its
norm has the claimed bound by
Proposition~\ref{prop:initial-cube-heads} for
\(p>1/2\). No ordinary nearest-neighbour density is
tested at the physical face.

Denote by \(\Delta_i\) the changes in local quadratic
exponents from this rearrangement, including the
interface density. They have zero total as a polynomial.
On each affected label their jet norm at \(\ell_{b_0}\)
is at most \(C\lambda G_0\). On an originally occupied
component \(X\), multiply its activity by
\(e^{\sum_{i\in X}\Delta_i}\) with the convention that
the old uncovered exponent is the new one minus
\(\Delta_i\). This is precisely the finite circle-product
identity used for the extraction divergences. The same
rearrangement on the bulk comparison uses the
translation-invariant divergence with no interface
density. Hence the transformed activities still agree
on all polymers avoiding the boundary labels.

Now allow \(z\) to differ from \(z_0^{\rm bare}\). Multiply
on each ordinary label by the exponential of the
ordinary kinetic density with coefficient
\(-(z-z_0^{\rm bare})\), in the potential normalization.
Put this same factor on every site label in an occupied
component, as well as on each uncovered label. This
defines the off-shell family exactly, and keeps the
raw collar tensor fixed. The off-shell mass is introduced
by the uniform scalar multiplication in
Proposition~\ref{prop:cube-joint-initialization}.
Any of these multipliers of quadratic coefficient
size \(A\) at the small radius has large-radius
exponent bounded by
\begin{equation}
 C A G_0^{-1/2}\sum_{i}
       \left(Z_i^2+Q_i^2\right),
\end{equation}
where \(Z_i,Q_i\) are the local normalized \(L^2\)
and quotient quantities. This follows by testing the
quadratic on its fixed enlarged stencil and using
\(F_{b_0}(\phi)/h_{b_0}\le C(Z_i+Q_i)\).
The same bound at the small radius uses \(A\) instead
of \(A G_0^{-1/2}\). For the rearrangement take
\(A=C\lambda G_0\), and for the off-shell multipliers
take \(A=C G_0\). These costs and the fixed number of
tested jets use only a fixed fraction of the spare
regulator and size margins of the joint initialization.
Thus the norm remains \(CU\) on the two complex discs.
Cauchy's inequality gives the two derivative bounds.
The equality at the bare coordinates is exact.

Finally the derivative in \(a=\delta/\lambda\) of the
bulk kinetic coefficient of the full density difference
is negative and bounded away from zero by
Proposition~\ref{prop:kinetic-reference-family}.
Its discarded tail has derivative in the same kinetic
moment at most \(C e^{-c b_0/R}\), by the weighted
short inverse estimate for one \(\delta\) derivative
and multiplication by \(\lambda\). The old microscopic
quadratic is independent of \(a\). Hence the same sign
and lower bound hold for the bare coordinate of the
cut head. All exact rearrangements and initial
norms are continuous on the physical temperature
and counterterm windows.
\end{proof}

For the terminal comparison we must control both the remaining
covariance in the principal coordinate and the magnetic response of
the slices already integrated. The transfer formula gives both
covariances in terms of the same unperturbed decomposition.
\begin{lemma}[Terminal Gaussian bounds under covariance transfer]
\label{lem:kinetic-reference-terminal}
Use the transfer decomposition of
Section~\ref{sec:kinetic-reference-family}, with the same
scale horizon as its unperturbed decomposition. Denote the
unperturbed terminal and integrated covariances by \(C_T^0,C_<^0\).
There is a common choice of the dual principal direction
\(\ell\), up to scalar normalization, for all the references.
Let \(T_\delta,P_\delta\) denote their terminal and integrated
covariances. Uniformly in the physical counterterm window,
the terminal variance, conditioned Gaussian exponential
moment, magnetic profile, and integrated magnetic direction
satisfy the bounds used in Propositions~\ref{prop:terminal-graph}
and~\ref{prop:terminal}.
\end{lemma}

\begin{proof}
Write \(q_T\) for the unperturbed auxiliary terminal kernel
in spectral variables and
\begin{equation}
 M=a_R\mathcal O\mathcal Q_{c',0} q_T(\mathcal Q_{c',0})
                    \mathcal O^* .
\end{equation}
Since \(0\le q_T(t)\le t^{-1}\) on the positive spectrum, one has
\(0\le M\le a_R A_\Lambda\) and \(M\le \|\mathcal Q_{c',0}\| C_T^0\).
The transferred covariances are exactly
\begin{equation}
 T_\delta=t_\delta(C_T^0-\eta M),\qquad
 P_\delta=t_\delta C_<^0+(1-t_\delta)a_RA_\Lambda
                         +t_\delta\eta M .
 \label{eq:kinetic-transferred-split}
\end{equation}
In particular at \(\delta=0\) the transfer moves \(\eta M\)
from the prescribed terminal covariance into the integrated covariance.
The comparison \(\tfrac12 t_\delta C_T^0\le T_\delta\le t_\delta C_T^0\) follows from the multiplier bound.
The generalized principal vector solves
\(J_\Lambda\ell=\xi_1 A_\Lambda\ell\), independently of
\(\delta\). Its finite-volume normalizations are
uniformly equivalent since \(B_\Lambda^\delta\asymp I\).
Use one and the same vector in the comparisons below.

For \(\sigma_0^2=\ell^{\mathsf T}C_T^0\ell\) and
\(u_0=C_T^0\ell/\sigma_0^2\) the old bounds give
\(\sigma_0^2\asymp n^2/(\lambda x_N)\) and
\(\|u_0\|_{F_{n,p}}\le Cn^{-2}\).
If \(D\) is a weighted stencil test in \(F_{n,p}\), the
unperturbed hard covariance bound and the decomposition
along \(\ell\) imply
\begin{equation}
 D C_T^0D^*\le C\bigl(\lambda n^{-2}+\sigma_0^2n^{-4}\bigr).
\end{equation}
By covariance Cauchy--Schwarz and \(M\le a_R A_\Lambda\),
\begin{equation}
 |DM\ell|\le C\sqrt{\lambda}R
               \bigl(\sqrt\lambda n^{-1}+\sigma_0n^{-2}\bigr).
\end{equation}
It follows that the variance of \(T_\delta\) along \(\ell\) has
the same order as \(\sigma_0^2\). Dividing the last display by
\(\sigma_0^2\) and accounting for the change in the denominator
shows \(\|u_\delta-u_0\|_{F_{n,p}}=o(n^{-2})\) as
\(n\to\infty\) at fixed \(R\) and terminal aspect. In
particular the magnetic overlap has the same lower bound.
The conditioned covariance is monotone under positive
covariance comparison, by the infimum formula in the proof
of Lemma~\ref{lem:kinetic-terminal-energy}. Its Gaussian norm
moment follows by adding an independent Gaussian to attain a
fixed multiple of the old conditioned covariance and applying
conditional Jensen's inequality.

For the integrated magnetic direction,
\(\|A_\Lambda\one\|_{F_{n,p}}\le C_R n^p\). Outside a strip of
microscopic width the field is constant, and inside the strip
the local tests have only the boundary weight loss and
microscopic differences. Also
\begin{equation}
 \|M\one\|_{F_{n,p}}
 \le C\sqrt{\lambda}R n^2
               \bigl(\sqrt\lambda n^{-1}+\sigma_0 n^{-2}\bigr)
 \le C_R n^2 .
\end{equation}
For the last inequality use \(x_N^{-1}=O_R(1+\log n)\).
Together with the old bound for \(C_<^0\one\), the exact formula
for \(P_\delta\) proves the integrated-source estimate.
All these covariances and normalizations are continuous
in the real external parameter.
\end{proof}

\section{A dimension-independent shooting theorem}\label{sec:path}

The linearized bulk mass grows by a factor \(L^s\) at each step, whereas the
quartic coupling decreases only logarithmically. We therefore prescribe
the quartic coupling and the stable coordinates at the microscopic
scale and determine the mass from the zero-curvature condition at
scale \(N\). Theorem~\ref{thm:path} implements these two boundary
conditions by a contraction on paths. Its proof allows the spaces
of boundary tensors and polymer activities to vary with the scale.

Let \(b_j=b_0 L^j\) be the step scales, \(N\asymp \log n\),
\(x_0=G_0\asymp A_{\rm sc}^4 g_{\mathrm{mic}}\), and set
\begin{equation}
 x_{j+1}=x_j-\mathfrak b_jx_j^2,\qquad f_j=x_j^3+U d_0^j,
\qquad 0<d_0\le\tfrac14 .
\end{equation}
Assume \(0\le\mathfrak b_j\le C\lambda^2\) and
\(\mathfrak b_j\ge c\lambda^2\) except on a bounded number of
initial or terminal steps. The merged first step may instead have
\(\mathfrak b_0\le C\lambda^2\log(2+M_R)\), where
\(M_R=b_0/R\asymp\log R\). Assume
\(\sup_j \mathfrak b_jx_j=o(1)\) and
\begin{equation}
 \Lambda_R=C_0\lambda M_R^s,\qquad
 x_0/\lambda^2+U/x_0+\Lambda_R=o(1).
\end{equation}
For the physical initialization of
Proposition~\ref{prop:cube-joint-initialization},
\(x_0/\lambda^3=O(g_{\mathrm{mic}}^{1-10a_*})\to0\) and
\begin{equation}
 U=A_{\mathrm{sc}}^6 g_{\mathrm{mic}}^2+
        G_0^{9/4}g_{\mathrm{mic}}^{1/2}+G_0^3
       =o(G_0^{9/4}) .
\end{equation}
The latter follows since the first term divided by \(G_0^{9/4}\)
is \(O(g_{\mathrm{mic}}^{3a_*-1/4})\).
Thus the microscopic estimate has the smallness required at entry into
the map.

Write \(\gamma_j=(g_j-x_j)/x_j\),
\(\mu_j=\nu_j/(\Lambda_R x_j)\), and \(k_j=K_j/f_j\), and let
\(w_j\) denote \(W_j\) expressed in its scale-dependent stable weights.
The one-step map must contract the stable coordinates in these weights.
We state its estimates on fixed enlarged convex boxes
\(|\gamma|\le\varepsilon\) and with fixed sufficiently large radii
in \(w,\mu,k\). All derivatives below are in these normalized
input coordinates, holding the frozen covariance fixed.

\begin{assumption}[Weighted RG and endpoint data]\label{ass:path}
The step equations have the form
\begin{align}
 g^+ &= g-\mathfrak b_jg^2+e_j,\qquad
   |e_j|+\|D e_j\|\le C f_j,\\
 \nu^+ &= L^s\nu+\Lambda_R x_j M_j(\gamma,w,\mu,k),\\
 w^+ &= \mathcal W_j(\gamma,w,\mu,k),\qquad
 k^+ = \mathcal K_j(\gamma,w,\mu,k).
\end{align}
The maps \(M_j,\mathcal W_j,\mathcal K_j\) are bounded at zero in
the normalized input and output coordinates. For \(\kappa<1\) and \(\delta_R=o(1)\)
they have the derivative bounds
\begin{equation}
\begin{array}{c|cccc}
 &\gamma&w&\mu&k\\ \hline
 \mathcal W_j&C&\kappa&\delta_R&\delta_R\\
 M_j&C&C&\delta_R&\delta_R\\
 \mathcal K_j&C&C&C&\kappa .
\end{array}
\end{equation}
All constants are uniform in \(j,N\). The fixed initial \(w_0\) is
bounded and \(g_0=x_0\). The initial family has
\(\|K_{\rm init}\|\le CU\) and
\(\|D_{\nu_0} K_{\rm init}\|\le CU/x_0\).
The exact terminal curvature equation has a graph
\(\mu_N=B_N(\gamma_N,w_N,k_N)\) on these boxes,
bounded at zero, with derivative bounds \((C,C,\delta_R)\).
The maps and data are real on real arguments and continuous in the
external data in the same charts. Uniform constants and the estimates
hold on the boxes with the successive radii used in the proof below.
\end{assumption}

The term \(M_j\) includes the mass generated by Gaussian integration. With a
fixed bulk filter of row sum one its actual first-order coefficient is
\(3 b_{j+1}^s(\overline T\Gamma_j^{\rm bulk}\overline T^*)(0)\).
Centering at the actual slice keeps the Gaussian reference
independent of the mass variable. At the terminal scale the natural
profile in the curvature calculation is \(\rho_N=n^{s/2}u_T\);
the derivative of the uniform quadratic variation is
\(b_N^{-s}\|\rho_N\|_2^2\asymp 1\) if
\(\|u_T\|_2\asymp1\). Proposition~\ref{prop:terminal-graph} derives the graph and
its first variations from the terminal coefficient and activity
bounds and a small inverse-chart derivative. At the other end of the path, the reference and the polymer packing
must be evaluated at the same trial value of \(\nu_0\).

\begin{theorem}[A trajectory with prescribed terminal curvature]\label{thm:path}
For \(d\in\{1,2,3,4\}\), under Assumption~\ref{ass:path} and the
stated comparison-sequence bounds, there is a unique trajectory in
the controlled box, for all sufficiently large \(R\) and every
horizon with these bounds. It has \(g_0=x_0\), the prescribed
\(W_0\), \(K_0=K_{\rm init}(\beta,\nu_0)\), zero terminal curvature,
and
\begin{equation}
 |g_j-x_j|\le\varepsilon x_j,\qquad
 |\nu_j|\le C\Lambda_R x_j,\qquad
 \|w_j\|\le C,\qquad \|K_j\|\le Cf_j .
\end{equation}
The constants and contraction factor are uniform in the horizon, and
\(\nu_0^*\) is continuous in the external data. For the canonical
scales at large \(n\),
\begin{equation}
 g_N\asymp (G_0^{-1}+\lambda^2N)^{-1}\asymp_{d,R}(\log V)^{-1},
 \qquad
 \|K_N\|\le C(g_N^3+Ud_0^N).
\end{equation}
\end{theorem}

\begin{proof}
Put \(a_j=1-2\mathfrak b_j x_j\). Since
\(0\le a_j\le (x_{j+1}/x_j)^2\), the forward Green function of the
linearized quartic equation is bounded by
\((x_j/x_{i+1})^2\) for a source at step \(i<j\).
Reciprocal increments satisfy
\begin{equation}
 x_{i+1}^{-1}-x_i^{-1}=\frac{\mathfrak b_i}{1-\mathfrak b_i x_i}.
\end{equation}
Thus \(x_j^{-1}\asymp x_0^{-1}+\lambda^2 j\), with constants
uniform under the small first-step correction. If \(a_0'\in(0,1)\)
is fixed sufficiently close to one, then \(x_{i+1}\ge a_0'x_i\)
and \(d_0<(a_0')^3\). It follows that
\begin{equation}
 \sup_j x_j\sum_{i<j}\frac{f_i}{x_{i+1}^2}
       \le C\bigl(x_0/\lambda^2+U/x_0\bigr)=o(1).
\end{equation}
For the cubic part the sum is bounded by
\(C x_j\sum_{i<j}x_i\); the reciprocal comparison and
\(\log(1+t)/(1+t)\le1\) give the bound.
For the transient use \(x_{i+1}\ge (a_0')^{i+1}x_0\) and
\(x_j\le x_0\). The residual
\(-\mathfrak b_i(g_i-x_i)^2\) has normalized size \(O(\varepsilon^2)\)
and Lipschitz norm \(O(\varepsilon)\) after the same Green function.
Indeed
\begin{equation}
 x_j\sum_{i<j}\frac{\mathfrak b_i x_i^2}{x_{i+1}^2}
 \le C x_j\sum_{i<j}(x_{i+1}^{-1}-x_i^{-1})\le C.
\end{equation}
The forward quartic equation thus maps an error of size \(f_i\)
at scale \(i\) to a uniformly small relative error at every later
scale.

Define the mass component by solving the linear recurrence backward
with the terminal graph evaluated at the input path. The normalized
forcing at step \(i\ge j\) is multiplied by
\(L^{-s(i-j+1)}x_i/x_j\), and the terminal graph by
\(L^{-s(N-j)}x_N/x_j\). Their sums are uniformly bounded.
Use one-step forward maps for \(w,k\), with the prescribed \(w_0\)
and \(k_0=K_{\rm init}(\beta,\nu_0)/f_0\). The latter has
normalized mass derivative at most
\(CU\Lambda_R/f_0\le C\Lambda_R\).

Choose \(\varepsilon\) small for the quartic self-coupling.
Fix successively large radii for \(w,\mu,k\), in this order.
Each radius absorbs its bounded input from the preceding
coordinates and its strict diagonal transport. The backward mass
component has only a small derivative in \(\mu\). Terms from the
later coordinates into earlier ones are absorbed by increasing
\(R\) after the radii are fixed. For a contraction, use the
supremum norm on paths with coordinate weights
\(1,t,t^2,t^3\) in order \((\gamma,w,\mu,k)\).
For \(t>0\) sufficiently small, the strictly lower triangular
bounded derivatives are small in this norm. The diagonal bounds
are strictly below one and the upper triangular bounds tend to
zero by the Volterra estimate and \(\delta_R\to0\).
This gives a uniform strict contraction of the closed box.
Its fixed point satisfies the exact recurrences. The contraction
estimate in common charts yields continuity, and the reciprocal
increment comparison gives the bounds at \(N\).
\end{proof}

\begin{corollary}[Matching the bare mass by varying the temperature]
\label{cor:physical-mass-crossing}
Suppose the hypotheses of Theorem~\ref{thm:path} hold for the
physical initial family of Proposition~\ref{prop:cube-joint-initialization}
on the real window \(\lvert\beta-1\rvert\le A_0g_{\rm mic}\), with
the constants in the bound for the target mass fixed before \(A_0\).
If \(A_0\) is sufficiently large, then for all sufficiently large
\(R\) and \(n\ge n_0(R)\) there is a \(\beta\) in the window such that
\(\nu_0^{\rm bare}(\beta)=\nu_0^*(\beta)\). At this temperature the
physical initial family follows the trajectory with the prescribed
terminal curvature.
\end{corollary}

\begin{proof}
By Corollary~\ref{cor:physical-mass-entry}, the bare scalar
is \(\nu_0^{\rm bare}=\lambda^{-1}(b_0/R)^s r_h(\beta)\), where
\(r_h\) is the bulk row of the interaction Hessian in the primary
field. By Corollary~\ref{cor:bare-thermal-crossing},
\(r_h=h_0^2 r\), with \(h_0=p^{-1}\) bounded above and below,
\(r_h(1)=O(g_{\rm mic})\), and \(r_h'(\beta)=-h_0^2(1+O(g_{\rm mic}))\).
The fixed box and target constants must be chosen before \(A_0\).
The nonscalar variations are controlled by
Proposition~\ref{prop:initial-cube-heads} and the estimates following it,
with the threshold on \(R\) allowed to depend on \(A_0\).
The target from Theorem~\ref{thm:path} is
\begin{equation}
 r_h^*(\beta)=\lambda(R/b_0)^s\nu_0^*(\beta),\qquad
 |r_h^*(\beta)|\le C\lambda^2 x_0\le C' g_{\mathrm{mic}}.
\end{equation}
Choosing \(A_0\) large gives opposite signs of
\(r_h-r_h^*\) at the endpoints of the window. At a zero of this
continuous function, the full initialization, including its
reference and polymer activity, is physical and has the prescribed
terminal curvature.
\end{proof}

A marginal coefficient can also be prescribed at the microscopic
scale if the sum of its increments and first variations is small.
The next corollary treats that case, including the dependence of
the bubble coefficient on the additional coordinate.

\begin{corollary}[A summable marginal coordinate]\label{cor:marginal}
Theorem~\ref{thm:path} permits an additional finite-dimensional
coordinate with initial value \(z_0\) and step equation
\(z^+=z+\zeta_j\) in fixed normalized units. Suppose the weighted
step and terminal bounds are uniform on an enlarged box for
\(z\), with bounded derivatives in \(z\) of the normalized
forcing in the stable and mass rows and the terminal graph,
and with \(\|D_z e_j\|\le C f_j\).
The quartic equation may use a coefficient \(\mathfrak b_j(z)\).
In that case use \(\mathfrak b_j=\mathfrak b_j(z_0)\) for the
comparison sequence and assume
\(\|D_z\mathfrak b_j(z)\|\le C\mathfrak b_j\) on the enlarged box.
If \(\sum_{j<N}(\sup|\zeta_j|+\sup\|D\zeta_j\|)=o(1)\), with
the derivative in all normalized coordinates on the box, the
same conclusion holds with
\(\sup_j|z_j-z_0|=o(1)\).
\end{corollary}

\begin{proof}
Use the cumulative equation \(z_j=z_0+\sum_{i<j}\zeta_i\)
in the path map. Its size and all its derivatives tend to zero
uniformly in the horizon, so place this coordinate first in
the weighted order in the proof of Theorem~\ref{thm:path}.
For the variable bubble, the additional quartic source at step
\(i\) is \(-[\mathfrak b_i(z_i)-\mathfrak b_i(z_0)]g_i^2\).
Its normalized forward Green bound is at most
\begin{equation}
 C\|z-z_0\|_\infty\sup_j x_j
                  \sum_{i<j}\frac{\mathfrak b_i x_i^2}{x_{i+1}^2}
 \le C'\|z-z_0\|_\infty .
\end{equation}
The derivative in \(z\) has a bounded forward norm, and the
derivative in the normalized quartic path is
\(O(\|z-z_0\|_\infty)\) by the same sum. Use a ball
\(\|z-z_0\|_\infty\le r_R=o(1)\) containing the image of the
cumulative equation. The quartic source preserves the quartic
box there. The other rows have bounded derivatives in this
earlier coordinate, so the same weighted norm gives the
contraction. For example, unnormalized step bounds of order
\(x_j^2+Ud_0^j\) have summable total
\(O(x_0/\lambda^2+U)\) under the reciprocal comparison.
\end{proof}

For a kinetic coefficient in dimension four, the corollary requires
uniform control of the remainder and its first variation on an
enlarged domain of fixed radius. The coupled chart of
Section~\ref{sec:residual-coupled-chart} controls the residual
kinetic coefficient on a domain that shrinks with \(g_j\).

\subsection{Backward shooting for a small residual kinetic coefficient}
\label{sec:backward-kinetic}

To use this shrinking domain, prescribe zero residual kinetic
coefficient at the terminal scale and sum its increments backward.
The sum from scale \(j\) to the endpoint must then be small in units
of \(x_j\). An external parameter in the Gaussian reference will
allow the resulting initial value to be matched along with the mass.

\begin{proposition}[Backward kinetic shooting on the shrinking box]
\label{prop:kinetic-backward}
Use the comparison sequence and the normalized coordinates of
Theorem~\ref{thm:path}. In dimension four suppose the maps of
Assumption~\ref{ass:path} are available on
\(\sup_j |z_j|/(\omega_R x_j)\le c_z\), with
\(0<\omega_R\le 1\) and an additional residual kinetic
equation \(z_{j+1}=z_j+\xi_j\) and terminal condition \(z_N=0\).
Assume the unaltered bounds of Assumption~\ref{ass:path} are
uniform for these kinetic inputs. Derivatives of the normalized
stable, mass and activity maps, and of the terminal graph, with
respect to \(z_j/(\omega_R x_j)\) are bounded. In the quartic row suppose
the added derivative is bounded by
\(C(f_j+\mathfrak b_j x_j^2)\) in the same direction. The initial
activity may depend on \(z_0/(\omega_R x_0)\) with bounded derivative after
division by \(f_0\), and the prescribed other initial coefficients
retain their bounds. If
\begin{equation}
 \sup_{j\le N}\frac1{x_j}\sum_{i=j}^{N-1}
       \bigl(\sup |\xi_i|+\sup\|D\xi_i\|\bigr)
 \le \vartheta_R=o(\omega_R),
\end{equation}
where \(D\) uses all the normalized coordinates, there is a
unique trajectory in the controlled box satisfying both terminal
conditions. It satisfies
\(\sup_j |z_j|/x_j\le\vartheta_R\).
The bare target values \(z_0^*,\nu_0^*\) depend continuously
on external data when the stated maps and data do.
\end{proposition}

\begin{proof}
In the path map of Theorem~\ref{thm:path} insert the component
\begin{equation}
 (z_j/(\omega_R x_j))^{\mathrm{out}}
       = -(\omega_R x_j)^{-1}\sum_{i=j}^{N-1}\xi_i .
\end{equation}
On the full controlled box it has supremum bound and Lipschitz
constant at most \(\vartheta_R/\omega_R\). Place this component first in
the weighted path norm. In the quartic component the new
derivative is bounded uniformly, since its forward Green
bound in normalized coordinates is
\begin{equation}
 C\sup_j x_j\sum_{i<j}
        \frac{f_i+\mathfrak b_i x_i^2}{x_{i+1}^2}\le C'.
\end{equation}
With \(z=0\) the original self-map bounds apply. The value
of the other arrows from this component on the smaller ball
\(\sup |z_j|/x_j\le 2\vartheta_R\) is
\(O(\vartheta_R/\omega_R)\).
Fix the other box sizes with the strict margins in the proof
of Theorem~\ref{thm:path}, and then take \(R\) large. This
gives a self-map. All arrows from the kinetic component
are bounded, so weighted contraction in the other rows
is proved in the same order as before. All arrows into the
kinetic component are small after that choice of weights,
by the bound \(\vartheta_R/\omega_R\). The fixed point is the desired
trajectory. Every controlled solution belongs to the smaller
ball by the cumulative kinetic equation, which gives uniqueness
in the controlled box. Continuity follows from the same
uniform contraction in common charts.
\end{proof}

The one-step bounds in the coupled chart contain a sunset term, a
cubic remainder and a contribution from the short-range quadratic
kernel. The following estimate sums these three sources in the units
required by Proposition~\ref{prop:kinetic-backward}.

\begin{lemma}[A summable source in the backward kinetic norm]
\label{lem:kinetic-backward-sum}
Under the comparison-sequence bounds of Theorem~\ref{thm:path},
suppose the ordinary step bounds for
\(\sup|\xi_j|+\sup\|D\xi_j\|\) are
\begin{equation}
 C\bigl[\lambda^3x_j^2+f_j
              +\lambda x_j(R/b_j)^2\bigr].
\end{equation}
At the merged step one may replace the first coefficient by
\(\lambda^3\log(2+M_R)\) and the last by
\(\lambda x_0 M_R^2(R/b_0)^2\).
These bounds imply the cumulative hypothesis of
Proposition~\ref{prop:kinetic-backward}, uniformly in the horizon,
with \(\omega_R=\Lambda_R\) in the microscopic scalings of this
paper.
\end{lemma}
\begin{proof}
On all ordinary scales with the bubble lower bound,
\(\lambda^2 x_i^2\le C(x_i-x_{i+1})\).
The bounded number of other scales contribute at most
\(Cx_j^2\) to any sum starting at \(j\). Thus
\(\sum_{i\ge j}x_i^2\le Cx_j/\lambda^2\), and
\(\sum_{i\ge j}x_i^3\le Cx_j^2/\lambda^2\).
The transient obeys
\(\sum_{i\ge j}Ud_0^i/x_j\le C U/x_0\), because
\(x_{i+1}/x_i\ge a_0'>d_0\) as in the proof of
Theorem~\ref{thm:path}. The bulk head gives a geometric
sum of size at most \(C\lambda x_j(R/b_j)^2\).
The merged terms divided by \(x_0\) are bounded by
\(C\lambda+C\lambda^3x_0\log(2+M_R)\).
Consequently the cumulative bound divided by \(x_j\)
is at most
\begin{equation}
 C\bigl(\lambda+x_0/\lambda^2+U/x_0
                +\lambda^3 x_0\log(2+M_R)\bigr).
\end{equation}
Here \(x_0/\lambda^2=o(\lambda)\) follows from \(a_*<1/10\),
and the initial ratios in Section~\ref{sec:cube-joint-initialization}
give \(U/x_0=o(\lambda)\). The merged logarithmic term is
\(o(\lambda)\). Thus the bound is \(O(\lambda)=o(\Lambda_R)\)
in dimension four.
\end{proof}

The backward solution specifies two microscopic target values. To
obtain a physical trajectory, both must agree with their bare values
at the same temperature and counterterm. Signs on the four faces
of a rectangle suffice for this two-parameter crossing.

\begin{corollary}[Simultaneous thermal and kinetic matching]
\label{cor:kinetic-two-matching}
In Proposition~\ref{prop:kinetic-backward} let the maps depend
continuously on a real temperature \(\beta\) and a real
external counterterm parameter \(a\) in a closed rectangle.
Suppose the two bare coefficients are continuous there.
Let the thermal bare coefficient minus \(\nu_0^*\) have
strict opposite signs on the two temperature faces, uniformly
in \(a\). Let the kinetic bare coefficient minus \(z_0^*\)
have strict opposite signs on the two counterterm faces,
uniformly in \(\beta\). Then there are parameters in the
rectangle at which the trajectory is physical and both
terminal conditions hold. The kinetic sign condition holds,
for example, if the kinetic bare coefficient is
\(z_{\rm bare}(\beta)+k(\beta,a)a\), with
\(|z_{\rm bare}(\beta)|\le C\lambda x_0\) and \(k\) of fixed
sign with \(k_*=\inf |k|>0\), provided the counterterm interval
\([a_-,a_+]\) has \(a_-<0<a_+\) and
\begin{equation}
 k_*\min\{-a_-,a_+\}>
             (C\lambda+\vartheta_R)\sup_{(\beta,a)}x_0 .
\end{equation}
The supremum is over the chosen rectangle.
\end{corollary}
\begin{proof}
Multiply the two continuous matching functions by signs so
each is negative on its lower coordinate face and positive
on its upper coordinate face. Subtract this vector of
functions from the identity on the rectangle and project
coordinatewise to the rectangle. Brouwer's theorem gives a
fixed point. The strict face signs force each coordinate
of the fixed point to be interior, where it must be a zero
of its matching function. At this point the full physical
initialization agrees with the shot trajectory. For the sufficient
criterion, \(|z_{\rm bare}(\beta)-z_0^*|\le(C\lambda+\vartheta_R)x_0\)
by Proposition~\ref{prop:kinetic-backward}. At each endpoint the
term \(k(\beta,a_\pm)a_\pm\) strictly dominates this error. Its
signs on the two faces are opposite and uniform in \(\beta\).
\end{proof}

For the finite-propagation family, the sign on the counterterm
faces comes from the bulk kinetic coefficient of the microscopic
density. The curvature matched by this construction uses the
terminal slice of the chosen reference. Section~\ref{sec:d4-prescribed}
will impose the condition for the original split.

\begin{corollary}[Physical matching for the tunable reference]
\label{cor:kinetic-reference-match}
Use the reference family in Proposition~\ref{prop:kinetic-reference-family} with \(c'=c^0\)
fixed and the transfer decomposition of
Section~\ref{sec:kinetic-reference-family} at the original cut.
Suppose a uniform one-step chart for this family satisfies
Proposition~\ref{prop:kinetic-backward} and the terminal
curvature hypotheses of Proposition~\ref{prop:terminal-graph},
and use the source bound of Lemma~\ref{lem:kinetic-backward-sum}.
Then the physical parameters can be matched with
\(|\delta|\le C\lambda^2 x_0\) for every sufficiently large
\(n\) at fixed \(R\). The prescribed terminal curvature is
the one using the Gaussian slices of this family.
\end{corollary}

\begin{proof}
Set \(a=\delta/\lambda\). The initial moment estimates of
Proposition~\ref{prop:initial-cube-heads} give a bare bulk
kinetic coefficient of size \(C\lambda x_0\) at \(a=0\).
For the transfer choice of \(c\), the bulk two-jet calculation in
Section~\ref{sec:kinetic-reference-family} uses
\((c'-1)/\zeta=O((R/n)^2)\) in place of \(\xi_1\). Its kinetic
coefficient therefore has slope in \(a\) negative and bounded away
from zero, uniformly on the rectangle. The initial
cumulative covariance is zero, so the compensated and
local initial kinetic coefficients agree. The exact
initialization uses the same heads and tails after addition
of this density difference, by its weighted short-kernel
bound. The change in the normalized bare mass is bounded by
\(C|a|(b_0/n)^2\). The thermal crossing of
Section~\ref{sec:path} therefore holds uniformly in the
counterterm window by enlarging the minimum size \(n_0(R)\)
at fixed \(R\). Lemma~\ref{lem:kinetic-backward-sum} gives
\(\vartheta_R=O(\lambda)\), so the kinetic slope and the
bound on \(z_0^*\) give opposite signs on the faces
\(a=\pm A\lambda G_{\rm ref}\) for fixed sufficiently large \(A\),
since \(x_0\asymp G_{\rm ref}=A_{\rm sc}^4g_{\mathrm{mic}}\)
uniformly on the thermal window.
Corollary~\ref{cor:kinetic-two-matching} matches the two
coordinates. Multiplying the counterterm window by
\(\lambda\) gives the asserted bound on \(\delta\).
\end{proof}

\section{Terminal integration and return to spins}\label{sec:terminal}

At the last scale, the quartic potential confines the principal
coordinate to a range of order \((V/g)^{1/4}\). The principal profile
has overlap of order \(\sqrt V\) with the uniform source. Together they give a
magnetic second moment of order \(V^{3/2}g^{-1/2}\).
Proposition~\ref{prop:terminal} makes this argument precise and
restores the source in the integrated slices. Its dependence on the
dimension enters through \(n^s=\sqrt V\). In this section \(a_*\)
denotes the terminal amplitude \((V/g)^{1/4}\).

\begin{assumption}[Terminal estimates]\label{ass:terminal}
At the matched point let \(g=g_N\to0\),
\(h=g^{-1/4}V^{-1/4}\), and \(a_*=(V/g)^{1/4}\).
There is a test norm \(E_n\) controlling the supremum norm and a
quotient seminorm \(q_n\le C E_n\). For the conditional
decomposition of Section~\ref{sec:model}, assume
\begin{equation}
 \sigma_T^2\asymp_{d,R}\sqrt V/g,\qquad
 \|\ell\|_2\le C,\qquad
 \|u_T\|_{E_n}\le C V^{-1/2},\qquad
 |\one^{\mathsf T}u_T|\ge c\sqrt V,\qquad
 \E e^{c\sqrt V\|Y\|_{E_n}^2}\le C.
\end{equation}
Constants in this assumption may depend on \(d,R\). Assume also that
\(v=C_<\one\) satisfies \(\|v\|_{E_n}\le C\sqrt V\), where
\(C_<\) is the integrated covariance. The exact real terminal
interaction is nonnegative and can be written
\begin{equation}
 F_T=e^{-V_4}+\mathcal E,\qquad
 V_4=\frac g4\sum_{x\in\Lambda}(T\phi)_x^4 .
\end{equation}
The filter \(T\) is uniformly bounded in \(\ell^\infty\), with
uniformly bounded inverse in \(\ell^4\).
For \(0\le k\le2\), assume the pointwise bound
\begin{equation}
 |D_{v_1}\cdots D_{v_k}\mathcal E(\phi)|
 \le \epsilon_n H(\phi)G(\phi)\prod_{i=1}^k
      \frac{\|v_i\|_{E_n}}{h},\qquad \epsilon_n=o(1).
\end{equation}
Here \(H=e^{-\kappa\|\phi\|_2^2/(Vh^2)}\),
\(G=e^{c_h q_n(\phi)^2/h^2}\), and \(c_h/\kappa\) is sufficiently small in terms of the
principal-profile constants. The physical initial interaction
and \(F_T\) are related by exact Gaussian integration, up to a
field-independent positive factor.
\end{assumption}

Proposition~\ref{prop:terminal-packing} below gives
\(\epsilon_n\le C_R(\sqrt g+g^{-9/4}\|K_N\|)\) from the terminal
cell factors. For the source, Section~\ref{sec:ext-cov} gives
\(\|C_j\one\|_{F_{n,p}}\le C_R n^p b_j^{s-p}\).
Since \(p<s\), the geometric sum up to the terminal scale is
\(O_R(n^s)=O_R(\sqrt V)\).

\subsection{Terminal comparison from a finite-cell RG output}
\label{sec:terminal-output}

At fixed terminal aspect the cube contains a bounded number of
cells. Expanding their circle product gives the pointwise error
in Assumption~\ref{ass:terminal} from the retained coefficients
and the activity norm. We use the adapted chart,
where the terminal test norm is \(E_n=F_{n,p}\) and the terminal
quotient can be taken to be \(q_n=F_{n,p}\). Here \(p<1\) and
\(2p>s-1\).

The regulator ratio \(c_h/\kappa\) must be fixed before the terminal
aspect is enlarged to control the sum over cells. This order requires
a principal-profile bound uniform in the choice of aspect. Whenever
the annular estimates hold in a weight \(p'>p\) and to a few more
orders than needed by \(F_{n,p}\), the bound
\(\|u_T\|_{F_{n,p}}\le C V^{-1/2}\) can use a constant independent
of the terminal aspect, for all sufficiently large \(n\) at each
fixed aspect. To see this, first use one fixed comparison aspect
and let \(c\uparrow c_{*,\Lambda}\) with \(n,R\) fixed. The
conditional covariance and low-tail variance bounds apply
uniformly in this limit. Since the integrated variance stays
bounded while the full soft variance diverges, the conditional
profile tends to \(u_\star=B_\Lambda^{1/2}e\). The profile bound
in the higher norm \(F_{n,p'}\) therefore holds for \(u_\star\)
with the constant from this fixed comparison aspect. For the
actual aspect the higher norm of \(V^{1/2}(u_T-u_\star)\) is
bounded, possibly with a larger constant. Its normalized
\(\ell^2\) norm tends to zero by the exact profile identity in
Section~\ref{sec:ext-cov}. At depth at most \(\delta n\), the
difference of the weights gives a factor \(C\delta^{p'-p}\)
in the \(F_{n,p}\) norm. In the interior, the higher norm
gives bounded rescaled first differences. A value bounded
away from zero would persist on a cube of side \(c_\delta n\)
and contradict the normalized \(\ell^2\) bound. To include
higher derivatives of order \(k\), apply the discrete Taylor
interpolation on an interior cube with mesh \(\lfloor\eta n\rfloor\)
and a fixed order \(P'>k\). In rescaled units the resulting
bound is \(C_{\delta}(\eta^{-k}\|\cdot\|_\infty+
\eta^{P'-k})\), by the same fixed-order interpolation matrix
as in the annular estimate. Send \(n\to\infty\), then
\(\eta\downarrow0\), then \(\delta\downarrow0\). Thus
\(\|V^{1/2}(u_T-u_\star)\|_{F_{n,p}}\to0\).
This permits the choice of the output coefficient \(c_h/\kappa\)
before enlarging the terminal aspect.

\begin{proposition}[Pointwise terminal error from the terminal coordinates]
\label{prop:terminal-packing}
Let \(b\asymp n/A\) with fixed large \(A\), and use the adapted
cell quotient and actual disjoint cell labels of
Section~\ref{sec:geometry}. Suppose the terminal circle product,
with vacuum normalized to one, has the following output.
\begin{enumerate}
\item The primary quartic is \(V_4=\tfrac g4\sum_x(T\phi)_x^4\)
with cellwise anchor factors. The filter \(T\) has a uniformly
bounded absolute row norm, bounded \(\ell^4\) inverse, and a
bounded finite-range forward filter \(B_{\rm rec}\) of range
\(O(R)\) with \(\|I-B_{\rm rec}T\|_{2\to2}=o_R(1)\).
\item The retained quadratic coefficients have size \(O_R(g)\)
in the cell field norm at radius \(\ell_b=b^{-s/2}\).
The nonprimary polynomial factors are \(1+P\), of degree
at most six, with coefficients of size \(O_R(g^2)\) in that
norm. There are a bounded number per cell.
\item With \(h_b=g^{-1/4}b^{-s/2}\) the large-field bound on each
connected activity, through two derivatives, is
\begin{equation}
 \begin{aligned}
 |D_{v_1}\cdots D_{v_k} K(X,\phi)|
 &\le C\delta_K\prod_{U\in X}
   \exp\!\left\{-\kappa_0\frac{\|\phi_U\|_2^2}{b^d h_b^2}
                  +c_0\frac{q_{U^*}(\phi)^2}{h_b^2}\right\}\\
 &\qquad{}\times\prod_{i=1}^k\frac{\sup_{U\in X}F_{b,U^*}(v_i)}{h_b},
 \qquad 0\le k\le2,
 \end{aligned}
\end{equation}
where \(\delta_K=o(1)\) and \(\kappa_0,c_0>0\) are fixed.
Here \(U^*\) is a fixed field-support enlargement of \(U\).
\end{enumerate}
Then the pointwise terminal bound of Assumption~\ref{ass:terminal}
holds with \(\epsilon_n\le C_{A,R}(\sqrt g+\delta_K)\). The
output constants \(\kappa,c_h>0\) can be fixed before taking
\(A\) sufficiently large. The conclusion applies to the
terminal value of the controlled path whenever
\(\delta_K=O(g^{-9/4}\|K_N\|)\) and its retained coefficients
have the stated sizes.
\end{proposition}

\begin{proof}
For the truncated Neumann filter used in the microscopic and
bulk charts the filter assumptions follow directly from
\(T=\sum_{j=0}^{M}Q_\Lambda^j\). Up to the fixed scalar
normalization, \(B_{\rm rec}=I-Q_\Lambda\) and
\(T^{-1}=(I-Q_\Lambda)(I-Q_\Lambda^{M+1})^{-1}\).
The row and column bounds give
\(\|Q_\Lambda\|_{4\to4}\le q<1\), and
\(\|I-B_{\rm rec}T\|_{2\to2}\le q^{M+1}\).

Set \(h=g^{-1/4}n^{-s/2}\),
\(Z=\|\phi\|_2/(\sqrt Vh)\), and \(Q_n=F_{n,p}(\phi)/h\).
Restriction to cell neighbourhoods does not increase the
raw norm by more than a fixed factor. The radius conversion
from \(\ell_b\) to \(h_b\) is \(g^{-1/4}\).
Thus a terminal quadratic and its first two direction
derivatives cost \(C_A\sqrt g(1+Q_n)^2\), with each tested
direction contributing \(F_{n,p}(v)/h\). A degree-\(k\)
polynomial correction costs \(C_A g^{2-k/4}(1+Q_n)^k\)
in these same jets, hence a power at least \(g^{1/2}\).
There are at most \(C_A\) cells. In the term with no occupied
cells, \(\ell^4\) invertibility gives \(V_4\ge c Z^4\).
Writing \(e^{-U}-1=-U\int_0^1e^{-tU}\,dt\) for the
quadratic corrections and expanding the finite product
of \(1+P\)'s about one yields the error
\(C_{A,R}\sqrt g\,e^{-\kappa Z^2+c_h Q_n^2}\) in the two
tested jets. For this absorption use the quartic decay,
the bound \(e^{C_A\sqrt g(1+Q_n)^2}\), and a fixed
fraction of the output quotient reserve.

Let \(X\) be the set of occupied cells in a nonempty term,
and write \(Z_X^2=\|\phi_{\cup X}\|_2^2/(Vh^2)\) and
\(S^2=\|(T\phi)_{(\cup X)^c}\|_2^2/(Vh^2)\).
Outside \(\cup X\) and its range-\(O(R)\) strips, recovery
of the original field uses only the unoccupied filtered
rows. Those strips have volume \(O((R/b)V)\). On their
interior sites use \(\|\phi\|_\infty\le h Q_n\);
on ordinary unoccupied sites use the operator estimate
for \(I-B_{\rm rec}T\). These estimates give
\begin{equation}
 (1-2\|I-B_{\rm rec}T\|_{2\to2}^2) Z^2
    \le Z_X^2+C S^2+C(R/b)Q_n^2 .
\end{equation}
The coefficient of \(Z_X^2\) comes from retaining the
occupied part of the \(\ell^2\) norm exactly. The activity
damping costs \(-\kappa_0 (n/b)^s Z_X^2\). The unoccupied
quartic supplies any fixed negative multiple of \(S^2\)
at a fixed constant cost, since \(Vgh^4=1\).
For the quotient factors, the summed restriction estimate
of Section~\ref{sec:geometry} gives
\begin{equation}
 \sum_U \frac{q_{U^*}(\phi)^2}{h_b^2}
       \le C A^{s-1-2p} Q_n^2 .
\end{equation}
Choose \(A\) so the decay on \(Z_X^2\) exceeds the
required multiple and \(c_0 C A^{s-1-2p}<c_h/4\).
Then take \(R\) large and \(R/b\) small. The recovery
inequality leaves a strict surplus in both exponents,
which absorbs \(C_A\sqrt g(1+Q_n)^2\) and the finite-degree
polynomials from the two field derivatives of the reference.
The term is bounded in these jets by
\(C_{A,R}\delta_K^m e^{-\kappa Z^2+c_h Q_n^2}\), where \(m\)
is the number of connected occupied components.
Summing over the bounded number of subsets proves the claim.
\end{proof}

\subsection{The terminal mass graph from small-field coefficients}
\label{sec:mass-graph}

The shooting theorem requires the mass as a function of the other
terminal coordinates. Varying the mass at fixed activity and fixed
Gaussian covariance gives a leading curvature derivative equal to
the squared norm of the rescaled principal profile in mass units.
The next proposition bounds the error in this derivative, which
allows us to solve the scalar curvature equation uniformly.

\begin{proposition}[Terminal curvature estimate in the cube chart]
\label{prop:terminal-graph}
Use a fixed terminal aspect and the canonically normalized tail.
Let \(x=x_N\), \(f=f_N\), and \(\Lambda_R\) be as in
Section~\ref{sec:path}, and assume
\((x^2+f)/(\Lambda_R x)=o(1)\). Suppose the exact terminal
circle-product chart is holomorphic and even in the field on
an enlarged complex coefficient box and satisfies
the following small-field estimates, with the vacuum removed.
\begin{enumerate}
\item The primary quartic uses a fixed real filter \(T\) with
uniformly bounded absolute rows, and its coefficient is of order
\(x\) in the stable sector. The total quadratic has the form
\(\tfrac12 \nu b^{-s}\|\phi\|_2^2+U_{\rm st}\), where
\(U_{\rm st}\) has size \(O(\Lambda_R x)\) at
\(\ell_b=b^{-s/2}\). The remaining polynomial factors are
\(1+P\), of bounded degree and coefficient size \(O(x^2)\).
\item The activity and its first two field jets satisfy the
small-field bound \(Cf\) times the product of growing cell
regulators. In a unit normalized activity direction the same
estimate, with \(Cf\), bounds the variation. The growing
regulator on this fixed number of cells is bounded by
\(\exp(C\|\phi\|_{E_n}^2/n^{-s})\).
\item In the terminal chart, the \(\nu\) derivative of
\(V_4+U_{\rm st}\) has size \(o(1)\) at the small-field
radius. This holds if these raw coordinates are independent
of \(\nu\), or if the inverse normal-form chart has a small
first variation there. The test directions satisfy
\(\|\rho_N\|_{E_n}=O(n^{-s/2})\) for
\(\rho_N=n^{s/2}u_T\).
\end{enumerate}
For each fixed stable box, take the mass radius to be
\(C_M\Lambda_R x\), with \(C_M\) larger than a fixed multiple of
the bound for \(|q_N(0)|/(\Lambda_R x)\) divided by the lower
bound on \(c_N=b^{-s}\|\rho_N\|_2^2\). The coefficient bounds must
hold with a further fixed-fraction collar for Cauchy's estimates.
Suppose the hard field has the exponential moment
\(\E e^{c n^s\|Y\|_{E_n}^2/\lambda}\le C\), and
\(n^s/\sigma_T^2=O(\lambda x)\).
Then the exact terminal curvature has the graph in
Assumption~\ref{ass:path}, with normalized activity derivative
\(O(f/(\Lambda_R x))\).
\end{proposition}

\begin{proof}
Write \(F_T\) for the terminal interaction, and
\(Z(a)=\E F_T(a\rho_N+Y)\), using the rescaled coordinate \(a\).
At \(a=0\) we use two field jets in the direction \(\rho_N\).
The terminal number of cells is bounded, and
\(\|\rho_N\|_{E_n}/\ell_b=O(1)\). If
\(y=\|Y\|_{E_n}/\ell_b\), the integrated jets are bounded
uniformly whenever their pointwise majorant is a fixed polynomial
in \(y\) times \(e^{C y^2}\), by the hard Gaussian moment
and small \(\lambda\).

The empty-activity term of \(F_T\) is
\(e^{-V_4-U}\prod(1+P)\), with
\(U=\tfrac12\nu b^{-s}\|\phi\|_2^2+U_{\rm st}\).
In the complex box, the primary quartic exponential has modulus
at most one on real fields. Taylor's integral formula for the
exponential bounds the integrated two-jet difference from
\(1-V_4-U\) by \(C x^2\). Indeed each term in that remainder
carries two factors of \(V_4+U\) or their field derivatives,
and the adverse exponential from \(U\) is at most
\(e^{C\Lambda_R x(1+y)^2}\). The factors \(1+P\) produce
errors of order \(x^2\) by the same integrated jet bound.
Every nonempty activity term is \(O(f)\) in these jets,
by the product rule and the small-field regulator.
The entire error is holomorphic on the enlarged box, and
changing only the normalized activity by a vector of unit norm
changes its first variation by at most \(C f\).
This follows by differentiating the finite circle product,
using the stated activity variation and integrating the same
majorant.

In particular \(Z(0)=1+O(x+f)\), so its scalar inverse is
uniform. Let
\(c_N=b^{-s}\|\rho_N\|_2^2\asymp1\).
The second \(a\) derivative of the uniform mass is \(c_N\nu\).
The leading quartic contribution is bounded by \(C\lambda x\)
since
\begin{equation}
 \left|\E D_{\rho_N}^2 V_4(Y)\right|
 \le Cx\sum_{z\in\Lambda}
      |(T\rho_N)_z|^2 \E |(TY)_z|^2
 \le C\lambda x\, V n^{-2s}
 =C\lambda x .
\end{equation}
By evenness \(Z'(0)=0\). Thus the curvature
\(q_N=n^s/\sigma_T^2-Z''(0)/Z(0)\) has the form
\(c_N\nu+r_N\), where \(r_N=O(\Lambda_R x)\).
At \(\nu=0\) the bound depends on the fixed stable box and
has activity contribution \(O(f)=o(\Lambda_R x)\).

For the \(\nu\) derivative of \(Z''(0)\), Cauchy's formula on the
larger mass box differentiates the \(O(x^2+f)\) error with
cost \(C/(\Lambda_R x)\). The assumed small first variation of
\(V_4+U_{\rm st}\) shows
\(\partial_\nu Z''(0)=-c_N+o(1)\) and
\(\partial_\nu Z(0)=O(1)\). Consequently
\(\partial_\nu q_N=c_N+o(1)\).
The normalized \(\gamma,w\) derivatives of \(r_N/(\Lambda_R x)\)
are bounded by Cauchy's formula on their fixed enlarged boxes.
The normalized activity derivative is \(O(f/(\Lambda_R x))\)
by the integrated jet variation. These bounds give a contraction
of \(\nu\mapsto-c_N^{-1}r_N\) on a disc of radius a fixed
multiple of \(\Lambda_R x\), after choosing that multiple
from the zero-mass bound and then taking \(R\) large.
Its fixed point is holomorphic and real for real data. Implicit
differentiation gives the required normalized derivatives.
\end{proof}

The ratio in the proposition is implied by
\(x_0/\Lambda_R\to0\) and \(U/(\Lambda_R x_0)\to0\), both of which
follow from the initial scales in Section~\ref{sec:path}. Indeed
\(f_j/(\Lambda_R x_j)\le x_0^2/\Lambda_R+
[U/(\Lambda_R x_0)](d_0/a_0')^j\), where
\(x_{j+1}\ge a_0'x_j\). The frozen resonance and the
terminal low-tail variance estimate supply
\(n^s/\sigma_T^2=O(\lambda x_N)\) before the mass is matched.
The inputs still to be checked in each application are the retained
coefficient estimates, the first variation of the inverse mass
coordinate, and the small-field activity bound.

We can now use the quartic scale to estimate the magnetic second
moment. The first bound below concerns the terminal field. The
second estimates the change when a uniform source is coupled to
the full auxiliary field, including the slices already integrated.

\begin{proposition}[The terminal moment and Gaussian source restoration]
\label{prop:terminal}
Under Assumption~\ref{ass:terminal},
\begin{equation}
 \E_T(\one^{\mathsf T}\phi)^2\asymp V^{3/2}g^{-1/2},\qquad
 |\E_0(\one^{\mathsf T}\phi)^2-\E_T(\one^{\mathsf T}\phi)^2|
 \le C V^{3/2}.
\end{equation}
Here \(\E_T\) and \(\E_0\) use the normalized terminal and
initial auxiliary laws.
\end{proposition}

\begin{proof}
Put \(a=a_*t\). Since \(\ell^{\mathsf T}Y=0\) and
\(\ell^{\mathsf T}u_T=1\), one has
\(\|a u_T+Y\|_2\ge |a|/\|\ell\|_2\).
Also \(a_*=\sqrt V h\); the quotient bound on \(u_T\)
and the triangle inequality give
\begin{equation}
 H(a_*t u_T+Y)G(a_*t u_T+Y)
 \le e^{-c t^2}\exp(C c_h h^{-2}\|Y\|_{E_n}^2).
\end{equation}
Since \(h^{-2}=g^{1/2}\sqrt V\), the latter exponential has a
uniform expectation with room for Cauchy--Schwarz. The Gaussian
variance of \(\one^{\mathsf T}Y\) is \(O(V^{3/2})\), from the
supremum norm bound and the exponential moment. Consequently the
same expectation with insertion \(|\one^{\mathsf T}Y|^k\),
\(k\le2\), is bounded by \(C V^{3k/4} e^{-ct^2}\).

The same field-derivative estimate assumed for \(\mathcal E\) holds
with constant \(C\) for \(e^{-V_4}\). Indeed
\(\|T\phi\|_4^4\ge c\|\phi\|_4^4\ge
cV^{-1}\|\phi\|_2^4\). Differentiation produces fixed-degree
polynomials in \(g^{1/4}\|T\phi\|_4\) and in the factors
\(g^{1/4}\|Tv_i\|_4\le C\|v_i\|_{E_n}/h\).
Absorbing the former by a fraction of \(e^{-V_4}\) leaves \(H\).
Thus the preceding Gaussian envelope applies to the integrand
and its two field derivatives.

For compact \(t\)-sets,
\(\E V_4(a_*t u_T+Y)\le C(t^4+g)\) by the supremum norm
bound on \(T\) and the fourth moment of \(\|Y\|_\infty\).
Jensen's inequality yields a uniform positive lower bound for the
quartic weight integrated over \(Y\), and the error is \(o(1)\) there.
The Gaussian factor \(e^{-a_*^2t^2/(2\sigma_T^2)}\) is bounded
below on compact \(t\)-sets since
\(a_*^2/\sigma_T^2=O(\sqrt g)\). Integrating the upper envelope
and using, for the lower second moment, \(1\le |t|\le2\),
gives \(Z_T=\E_{C_T}F_T\asymp a_*/\sigma_T\) and
\(\E_T a^2\asymp a_*^2\).
The weighted envelope also gives
\(\E_T(\one^{\mathsf T}Y)^2=O(V^{3/2})\).
The \(L^2\) triangle inequalities and the overlap of \(u_T\)
give the first conclusion.

For the second conclusion put \(S=\one^{\mathsf T}\phi\).
Gaussian translation in the earlier covariance turns the source
\(J\one\) into a factor \(e^{J^2\one^{\mathsf T}C_<\one/2}\)
and the shift \(F_T(\phi+Jv)\), with the terminal source
\(e^{JS}\) retained. The same envelope permits two derivatives
at \(J=0\). For fixed \(n\), the shift incurs only a constant
and an arbitrarily small loss in the two quadratic exponents
as \(J\) varies near zero, and the linear source is absorbed
by the remaining Gaussian reserve. Thus the difference of moments is
\begin{equation}
 \one^{\mathsf T}C_<\one+
   2 Z_T^{-1}\E_{C_T}[S D_v F_T]+
     Z_T^{-1}\E_{C_T} D_v^2 F_T .
\end{equation}
The first term is at most \(CV^{3/2}\). With
\(B_v=\|v\|_{E_n}/h\le C V^{3/4}g^{1/4}\),
the two other terms are bounded by
\(C V^{3/4}g^{-1/4}B_v\) and \(C B_v^2\), respectively, by
the derivative envelope and the normalization bound. Both are
\(O(V^{3/2})\).
\end{proof}

To return to spins we must account for the variation of
\(A_\Lambda\one\) near the boundary. The ferromagnetic two-point
function is nonnegative, so uniform upper and lower bounds on
the entries of this vector suffice.

\begin{lemma}[Positive-row transfer to the Ising susceptibility]
\label{lem:spin}
Let \(A_\Lambda\) be a fixed positive Gaussian split of a
ferromagnetic Ising interaction such that \(A_\Lambda\) has
nonnegative entries, \(\inf_x(A_\Lambda)_{xx}\ge a_R'>0\)
and \(\sup_x\sum_y(A_\Lambda)_{xy}\le p_R<\infty\), uniformly in
\(n\) at fixed \(R\). Let the unscaled auxiliary field be
\(h_{\rm phys}=\sqrt\beta A_\Lambda\sigma+\eta\) with independent
\(\eta\sim N(0,A_\Lambda)\). At zero field,
\begin{equation}
 \beta (a_R')^2 \E_\beta(\one^{\mathsf T}\sigma)^2
 \le \E_\beta(\one^{\mathsf T}h_{\rm phys})^2
 \le \beta p_R^2 \E_\beta(\one^{\mathsf T}\sigma)^2+p_R V.
\end{equation}
\end{lemma}

\begin{proof}
The vector \(A_\Lambda\one\) has entries in \([a_R',p_R]\).
The two-point functions \(\E_\beta\sigma_x\sigma_y\) are
nonnegative, as follows by expanding the zero-field Ising
partition functions in \(\tanh(\beta J_{xy})\) and retaining
the parity constraint. Multiplying these expectations by
the two row weights and summing gives the bounds on the
weighted spin term. The independent noise adds the
nonnegative variance \(\one^{\mathsf T}A_\Lambda\one\le p_R V\).
\end{proof}

\subsection{The four-dimensional matched point for the tunable reference}
\label{sec:d4-alternative-theorem}

We first prove the susceptibility bound for the terminal split
obtained by tuning the Gaussian reference. The temperature and
counterterm solve the two matching equations on a common rectangle.
Fix the original cut \(J\), the pole parameter \(c'=c^0\) and the
transfer parameter \(\eta>0\) of Section~\ref{sec:kinetic-reference-family}.
For each \(\delta\) use \(d_\delta=(\zeta+\delta)/\zeta\),
\(c=1+d_\delta(c'-1)\) and \(t_\delta=c/(d_\delta c')\). The
terminal covariance is \(T_\delta\) from \eqref{eq:kinetic-transferred-split}.
At the matched parameters this decomposition is held fixed when taking
the curvature.

\begin{theorem}[Susceptibility at the finite-propagation matched point]
\label{thm:d4-finite-propagation}
In dimension four, for all sufficiently large fixed \(R\)
and \(n\ge n_0(R)\), there are an external parameter
\(|\delta_*|\le C\lambda^2x_0\) and a physical temperature
\(\tau_{\Lambda}^{\mathrm{pc},\delta_*}\) for which the
terminal effective potential of the transferred split
\eqref{eq:kinetic-transferred-split} has zero curvature in its
principal coordinate \(a=\ell^{\mathsf T}\phi\), with terminal
Gaussian covariance \(T_{\delta_*}\). The physical Ising susceptibility
at this temperature satisfies
\begin{equation}
 c_R n^2\sqrt{\log n}
 \le \chi_\Lambda^F(\tau_{\Lambda}^{\mathrm{pc},\delta_*})
 \le C_R n^2\sqrt{\log n}.
\end{equation}
The constants are independent of \(n\).
\end{theorem}

\begin{proof}
Use the covariance transfer decomposition in
Section~\ref{sec:kinetic-reference-family}. Its ordinary
annular covariance bounds, filtered bubble bounds and
merged amplitude hold uniformly across the real matching
window. The ordinary and transferred terminal data can
use the same terminal aspect, chosen in the order of
Section~\ref{sec:terminal-output}. The off-shell activity
for the two scalar coordinates is given by
Lemma~\ref{lem:kinetic-initial-extraction}.
Its normalized initial mass and kinetic derivatives
after division by \(f_0\) are at most
\begin{equation}
 C(U/x_0)\Lambda_R x_0/f_0\le C\Lambda_R .
\end{equation}
The remaining prescribed stable coefficients have the
entry bounds of the coupled chart. In particular after
removing the bulk kinetic two-jet the fixed short head
has effective size \(O(\lambda x_0/M_R^2)\) at entry.
The cut and the boundary tensors are fixed at each choice
of the external parameters and continuously depend on
those parameters in the common label charts.

Proposition~\ref{prop:residual-coupled-chart} and
Proposition~\ref{prop:cube-bounded-extraction} give the
ordinary, collar and activity bounds on the same enlarged
extraction disc. At radius \(f_j\) the activity contraction
is the one in Corollary~\ref{cor:fractional-cube-path}:
choose \(L\) and the size margin so that the leading operator
and the large-set error sum to less than \(d_0/4\), then
take \(R\) large for the mixed-norm decoration and
Cauchy errors. The ratio \(f_j/g_j^{9/4}\) tends to zero
uniformly, and \(f_{j+1}\ge d_0 f_j\). The normalized
kinetic coordinate is \(z_j/(\Lambda_R x_j)\). Its arrows
into the stable, mass and activity rows are bounded by
the coupled chart, and the derivative of the quartic
error in that coordinate is \(O(f_j)\).
Proposition~\ref{prop:bulk-kinetic-source} and
Lemma~\ref{lem:kinetic-backward-sum} give the backward
kinetic norm \(O(\lambda)=o(\Lambda_R)\), including
first variations. Thus the estimates for the dynamics in
Proposition~\ref{prop:kinetic-backward} have the stated
strict margins.

At the terminal scale the stable head and collar
quadratic sizes are \(O(\Lambda_R x_N)\).
The raw kinetic coefficient satisfies the same bound by
the inverse compensated chart and the terminal condition
\(z_N=0\). Both collar polynomial degrees and the ordinary
margin factors have size \(O(x_N^2)\).
At fixed stable and activity coordinates the inverse mass
chart has the derivative of Lemma~\ref{lem:cube-inverse-mass}, so
the scalar mass derivative is \(1+O(\lambda x_N)\) and
the induced variation of the quartic and local kinetic
terms is \(O(\lambda x_N)\) in the small-radius units.
The fixed heads and full collar tensors are independent
raw coordinates in this inverse. The terminal activity
and its first variations are in the stated mixed norm,
on a fixed number of labels. The Gaussian estimates are
given by Lemma~\ref{lem:kinetic-reference-terminal}.
Hence Proposition~\ref{prop:terminal-graph} applies,
also for kinetic inputs on the shrinking box by the same
finite-jet argument. It gives the required continuous
terminal mass graph and its bounded normalized kinetic
derivative.

We may now apply the backward shooting proposition and
Corollary~\ref{cor:kinetic-reference-match}. The latter
uses the exact finite-head kinetic coefficient, which
differs from the bulk two-jet used for its sign comparison
by the exponentially small initial tail. This difference
fits in the strict crossing margin. At the resulting
temperature and counterterm both scalar bare coordinates
match the original density, as in
Lemma~\ref{lem:kinetic-initial-extraction}. The
spin model at that temperature is the physical one,
since the covariance change was compensated in its
auxiliary density.

Finally \(g_N\asymp_R(\log n)^{-1}\).
The terminal bounds in the chart and
Proposition~\ref{prop:terminal-packing} give the pointwise
terminal error with
\(g_N^{-9/4}\|K_N\|=o(1)\).
Lemma~\ref{lem:kinetic-reference-terminal} supplies all
the Gaussian bounds for Proposition~\ref{prop:terminal},
including the integrated-source estimate. It follows
that the initial auxiliary magnetic second moment has
order \(V^{3/2}g_N^{-1/2}\). The density compensation
preserves this initial auxiliary law. The spin comparison
of Lemma~\ref{lem:spin} therefore applies with the same
\(O_R(V)\) error. Divide by \(V=n^4\) to get the asserted
bounds.
\end{proof}

\subsection{Comparison of the two terminal curvature conditions}
\label{sec:kinetic-curvature-comparison}

The physical law of the full auxiliary field is the same in the
two references. The terminal coordinate depends on how its
covariance is split. To recover its original marginal, we keep
its conditional distribution given the full field. Gaussian
regression then identifies the covariance that must remain after
the preliminary integrations.

\begin{lemma}[Original terminal coordinate under a changed reference]
\label{lem:original-curvature-posterior}
Let \(C^0=C^0_<+C^0_T>0\) be the original Gaussian covariance,
and let \(a_0=\ell^{\mathsf T}\phi^0_T\) be its terminal
coordinate, with \(\sigma_0^2=\ell^{\mathsf T}C^0_T\ell>0\).
Put
\begin{equation}
 w_0=(C^0)^{-1}C^0_T\ell,\qquad
 \eta_0^2=\sigma_0^2-w_0^{\mathsf T}C^0w_0 .
\end{equation}
In any Gaussian reference \(\Sigma>0\) with physical interaction
factor \(F^\Sigma\) representing the same law of the full
auxiliary field, set
\begin{equation}
 V_{0,\Sigma}=\eta_0^2+w_0^{\mathsf T}\Sigma w_0,\qquad
 v_{0,\Sigma}=\Sigma w_0,\qquad
 u_{0,\Sigma}=v_{0,\Sigma}/V_{0,\Sigma}.
\end{equation}
Suppose that \(C^\Sigma_{<,k}\) has been integrated by scale
\(k\), with exact interaction \(F^\Sigma_k\), and write
\(C^\Sigma_{T,k}=\Sigma-C^\Sigma_{<,k}\). Define
\begin{equation}
 B_{0,\Sigma,k}=C^\Sigma_{T,k}
               -\frac{v_{0,\Sigma}v_{0,\Sigma}^{\mathsf T}}{V_{0,\Sigma}}.
 \label{eq:original-curvature-hard}
\end{equation}
If \(B_{0,\Sigma,k}\ge0\), the density of the original
coordinate \(a_0\), up to a constant independent of \(a\), is
\begin{equation}
  e^{-a^2/(2V_{0,\Sigma})}
       \E_{B_{0,\Sigma,k}}F^\Sigma_k(a u_{0,\Sigma}+Y).
 \label{eq:original-curvature-density}
\end{equation}
In particular its zero-curvature condition can be imposed on
the reference-\(\Sigma\) RG chart using this profile and hard
covariance. The condition \eqref{eq:original-curvature-hard}
is equivalent to
\begin{equation}
 v_{0,\Sigma}\in\operatorname{ran} C^\Sigma_{T,k},\qquad
 v_{0,\Sigma}^{\mathsf T}
 (C^\Sigma_{T,k})^\dagger v_{0,\Sigma}\le V_{0,\Sigma},
\end{equation}
where the dagger denotes the inverse on the range.
\end{lemma}

\begin{proof}
In the original Gaussian split, the joint covariance of the
full field \(\phi\) and \(a_0\) has cross column
\(C^0_T\ell\). Gaussian regression therefore represents
\(a_0=w_0^{\mathsf T}\phi+\eta_0 Z\), with \(Z\) standard
normal independent of \(\phi\). Weighting by the physical
interaction, a function only of \(\phi\), preserves this
conditional law. As a result, under the reference \(\Sigma\)
one can first form the same linear expression with
\(\phi\) centered Gaussian of covariance \(\Sigma\) and then
weight by \(F^\Sigma(\phi)\). The proxy variance of the
linear expression is \(V_{0,\Sigma}\) and its cross column
is \(v_{0,\Sigma}\). Conditioning on its value \(a\)
gives mean \(a u_{0,\Sigma}\) and covariance
\(\Sigma-v_{0,\Sigma}v_{0,\Sigma}^{\mathsf T}/V_{0,\Sigma}\).
If \eqref{eq:original-curvature-hard} is positive semidefinite,
this covariance is the sum of \(C^\Sigma_{<,k}\) and
\(B_{0,\Sigma,k}\). Gaussian addition gives
\eqref{eq:original-curvature-density}. The range test follows
by conjugating the rank-one inequality by the inverse square
root of \(C^\Sigma_{T,k}\) on its range and then testing on its
kernel. The argument includes degenerate hard covariances.
\end{proof}

This gives a sufficient terminal estimate for matching the
original coordinate via a different reference. Suppose the
outgoing chart is available at a cut \(k\) a bounded number
of scales before the original terminal scale, and that
\(B_{0,\Sigma,k}\) in \eqref{eq:original-curvature-hard} is
positive and has the uniform hard exponential norm moment of
Proposition~\ref{prop:terminal-graph}. Suppose also that
\begin{equation}
 \|n^{s/2}u_{0,\Sigma}\|_{E_n}=O(n^{-s/2}),\qquad
 b_k^{-s}\|n^{s/2}u_{0,\Sigma}\|_2^2\asymp1,\qquad
 n^s/V_{0,\Sigma}=O(\lambda x_k).
\end{equation}
If the profile, covariance and proxy variance are continuous
in the external parameters and independent of the variable
terminal mass, the proof of Proposition~\ref{prop:terminal-graph}
applies to \eqref{eq:original-curvature-density}. Its strict
mass derivative is the displayed normalized profile square
plus the same inverse-chart error. One can therefore use this
scalar graph in the backward shooting instead of the graph
for the modified terminal coordinate. The number of steps
released into the terminal comparison is bounded, so the
comparison sequence, the terminal aspect constants and the
activity weights are unchanged up to fixed factors. The usual
terminal moment proof still uses the natural decomposition of
the modified reference. Its estimate needs the terminal
coefficient box of order \(\Lambda_R x_k\), which the new
curvature graph supplies.

A scalar example illustrates both issues in this comparison. A
small change of reference can move the curvature zero and can
make the candidate conditional covariance at the original cut
negative.

\begin{proposition}[The curvature and the hard covariance need not transfer]
\label{prop:curvature-reference-obstruction}
Even a scalar covariance change of the form
\(\Sigma_t=tC^0+(1-t)A\) can alter the natural terminal
curvature for a fixed physical law. It can also make
\eqref{eq:original-curvature-hard} indefinite at the original
terminal cut for arbitrarily small \(|t-1|\).
\end{proposition}

\begin{proof}
In dimension one for the auxiliary field, take
\begin{equation}
 C^0=2,\qquad C_T^0=1,\qquad \ell=1.
\end{equation}
Then \(w_0=1/2\) and \(\eta_0^2=1/2\). For \(t\) near one
take
\begin{equation}
 \Sigma_t=2t+\tfrac12(1-t),\qquad C_{T,t}=t .
\end{equation}
Both parts of this Gaussian split are positive. This is of
the form of a scaled covariance plus a microscopic term.
For the original conditional coordinate the proposed hard
variance at the same cut is
\begin{equation}
 H_t=t-\frac{\Sigma_t^2}{2+\Sigma_t}.
\end{equation}
It vanishes at \(t=1\) and its derivative there is
\(-1/8\). Hence it is negative for all sufficiently small
positive \(t-1\).

The same example shows that the natural curvatures can differ
even for a fixed full physical law. Take the full law to be
the even mixture of normals of variance one and means
\(\pm\sqrt3\). In a Gaussian reference of variance
\(\Sigma_t\), its positive density relative to that Gaussian
defines the interaction factor. The natural terminal
coordinate is conditionally Gaussian given the full field,
with mean \(w_t\phi\) and variance \(\eta_t^2\), where
\begin{equation}
 w_t=t/\Sigma_t,\qquad
 \eta_t^2=t-t^2/\Sigma_t.
\end{equation}
Thus its marginal is the even mixture of normals of
variance \(v_t=w_t^2+\eta_t^2\) and means
\(\pm w_t\sqrt3\). Its log density has second derivative
\(3w_t^2/v_t^2-1/v_t\) at zero. This vanishes exactly when
\(\eta_t^2/w_t^2=2\). The latter ratio is
\(\Sigma_t(\Sigma_t-t)/t\), equal to two at \(t=1\) and
with derivative \(1/2\) there. The curvature consequently
changes under an arbitrarily small reference modification.
\end{proof}

For the transfer family on the cube we can estimate the variance
and the profile throughout the counterterm window. The covariance
moment then follows whenever the subtraction in
\eqref{eq:original-curvature-hard} is positive.

\begin{proposition}[Bounds for the original-coordinate comparison]
\label{prop:original-curvature-comparison-bounds}
Use the covariance transfer of Section~\ref{sec:kinetic-reference-family}
with the unperturbed parameter fixed there, in dimension four.
Let \(k\) be a fixed number of scales before the original terminal
cut. For every \(1/2<p<1\), on the physical counterterm window and
for all sufficiently large \(n\) at fixed \(R\), the proxy variance
and the profile in Lemma~\ref{lem:original-curvature-posterior}
satisfy the bounds required for its terminal mass graph. If
\(B_{0,\Sigma,k}\ge0\), its hard Gaussian exponential moment in
\(F_{n,p}\) also has the required uniform bound. Thus positivity of
this particular rank-one subtraction is sufficient for all the
Gaussian inputs of that transfer at the chosen cut.
\end{proposition}

\begin{proof}
Write \(T=C_T^0\), \(P=C_<^0\), \(G=T+P\),
\(A_c=a_R A_\Lambda\), and \(\sigma^2=\ell^{\mathsf T}T\ell\).
Then \(\Sigma=tG+(1-t)A_c\), where \(t=t_\delta\) is bounded
above and below by positive constants. Put \(u=T\ell/\sigma^2\)
and \(r=A_c w_0\), where \(w_0=G^{-1}T\ell\).
The regression quantities satisfy
\begin{align}
 v_{0,\Sigma}&=t\sigma^2u+(1-t)r,\\
 V_{0,\Sigma}&=t\sigma^2+(1-t)
                 (\eta_0^2+w_0^{\mathsf T}A_c w_0).
\end{align}
The old conditional variance satisfies
\(\eta_0^2=\ell^{\mathsf T}(P-PG^{-1}P)\ell\le C\lambda n^2\)
by Section~\ref{sec:ext-cov}. Also \(A_c\le G\) and
\(w_0=\ell-G^{-1}P\ell\), so
\begin{equation}
 w_0^{\mathsf T}A_cw_0
 \le 2\ell^{\mathsf T}A_c\ell+2\ell^{\mathsf T}PG^{-1}P\ell
 \le C\lambda n^2 .
\end{equation}
Consequently \(V_{0,\Sigma}=t\sigma^2+O(|1-t|\lambda n^2)\)
and \(V_{0,\Sigma}\asymp n^2/(\lambda x_k)\).

For the profile estimate we use two elementary operator
bounds. First an exponentially short kernel, with the clipping
defect bounds of Proposition~\ref{prop:kinetic-reference-family},
maps \(F_{n,p'}\) to \(F_{n,p}\) with bounded norm for
\(p<p'<1\) and \(n\ge n_0(R)\). At a site of truncated boundary depth
\(\rho=\min(n,R+\dist(x,\Lambda^c))\), split its bulk row at a
small multiple of \(\rho\). Short translations in that row
commute with each tested difference and preserve the local
weights. On the long translations and the clipping defect,
use the exponential weight before taking the fixed number
of point differences. Their extra cost is bounded by
\begin{equation}
 C (\rho/n)^{p'-p}(1+\rho)^P e^{-c\rho/R}.
\end{equation}
This has supremum tending to zero as \(n\to\infty\) at
fixed \(R\). The same estimate applies to products of the
short kernels after taking an intermediate exponent.
Second, for the integrated covariance,
\begin{equation}
 \|Pf\|_{F_{n,p'}}\le C\lambda n^2\|f\|_{F_{n,p'}}
 \qquad (p'<1).
\end{equation}
To see this, apply the estimate to each individual slice at
side \(b\) before merging. On a stencil at depth at most
\(C b\), every read field value costs
\(C(b/n)^{p'}\|f\|_{F_{n,p'}}\). For a difference of order
\(a\) on the covariance row, covariance Cauchy--Schwarz
and the annular bounds give
\(C\lambda b^{-2}(\min(b,\rho)/b)^{p''}
\min(b,\rho)^{-a}\), with \(p'<p''<1\).
Summing the \(O(b^4)\) read sites and multiplying by
the raw test weight gives \(C\lambda b^2\|f\|_{F_{n,p'}}\).
At larger depth the slice equals the bulk there by finite
propagation. Move the differences to \(f\) and use the
absolute row bound \(C\lambda b^2\). The geometric sum over
the integrated slices proves the displayed estimate.

The limiting pole profile of Section~\ref{sec:terminal-output}
is \(B_\Lambda\ell\) in the common normalization.
The short inverse estimate and the pole profile bound,
using exponents strictly between \(p\) and one, therefore
give \(\|\ell\|_{F_{n,p'}}\le Cn^{-2}\) for a fixed
\(p'>p\). The exact inverse relation
\(A_cG^{-1}=I-c' A_\Lambda B_\Lambda^{-1}\) gives
\begin{equation}
 r=A_c\ell-(I-c' A_\Lambda B_\Lambda^{-1})P\ell .
\end{equation}
The two operator bounds just proved show \(\|r\|_{F_{n,p}}\le C\lambda\).
The formula for the cross column now gives
\(\|u_{0,\Sigma}-u\|_{F_{n,p}}=o(n^{-2})\), uniformly on
the counterterm window. In particular the normalized
profile square has the same upper and lower bounds as the
unperturbed one.

It remains to bound the hard moment, conditional on positivity.
Let \(D\) be the sum of the unperturbed slices released by
stopping at \(k\). In the transferred decomposition
\(C^\Sigma_{T,k}\le t(T+D)\). Subtract the regression rank-one
matrix, and write \(T_h=T-\sigma^2 uu^{\mathsf T}\).
The displayed variance and cross-column identities give
the form bound
\begin{equation}
 B_{0,\Sigma,k}
 \le t(T_h+D)+C\lambda n^2 uu^{\mathsf T}
                          +\frac{C}{\lambda n^2}rr^{\mathsf T}.
\end{equation}
Indeed the difference of the two soft rank-one coefficients
is \(O(\lambda n^2)\), and the mixed form of \(u,r\) is bounded
by the last two terms by Cauchy--Schwarz. The remaining
\(rr^{\mathsf T}\) term from the subtraction is nonpositive.
The old hard covariance has the required Gaussian moment.
The bounded number of released slices have the annular moment
at scales comparable to \(n\). The square roots of the last
two rank-one terms have profile norms at most
\(C\sqrt\lambda n^{-1}\). Combining these finitely many
moments and applying Gaussian domination by conditional
Jensen proves the required bound for \(B_{0,\Sigma,k}\)
when it is positive. All vectors and matrices in the
comparison are continuous in the external parameter.
\end{proof}

\subsection{The prescribed terminal curvature in dimension four}
\label{sec:d4-prescribed}

To impose the curvature of the prescribed split, we make two
changes in the calculation. First we integrate in a microscopic
field whose covariance with the original coordinate is exactly
\(T\ell\). This removes the term \((1-t)A_cw_0\), which in
Lemma~\ref{lem:original-curvature-posterior} requires an inverse
estimate at high frequencies. Then we stop one scale earlier. The
slice retained at this cut supplies a variance margin of order
\(\lambda n^2\), enough to make the prescribed regression covariance
positive throughout the matching rectangle.

Throughout this subsection \(T=C_T^0\), \(G=C^0\), \(A_c=a_RA_\Lambda\),
\(\sigma^2=\ell^{\mathsf T}T\ell\), and \(w_0,\eta_0^2\) are the
original regression data of Lemma~\ref{lem:original-curvature-posterior}.
The original pole parameter \(c'=c^0\), cut \(J\), and coordinate
\(a_0\) are held fixed throughout the matching rectangle.
Write \(t=1-\varepsilon\) and
\(\Sigma=tG+\varepsilon A_c\). This is the transfer family of
Section~\ref{sec:kinetic-reference-family} with
\(\delta=\zeta\varepsilon c'/(1-\varepsilon c')\).
The preterminal integrations will use its finite-range slices.

The field transformation must preserve the short-range estimates
used at initialization. Its dependence on the external parameter
must also allow the bare kinetic coefficient to cross the value
obtained by backward shooting.

\begin{lemma}[Microscopic change of field]
\label{lem:aligned-initialization}
Let \(S=G\Sigma^{-1}\). Use \(\psi\) as the field in the RG map
and represent the original full field by \(\phi=S\psi\).
For \(|\varepsilon|\le H\lambda^2x_0(0)\) with fixed \(H\), the
transformed physical interaction has the initial activity bounds
and the off-shell mass and kinetic first-variation bounds of
Lemma~\ref{lem:kinetic-initial-extraction}. Its bare kinetic
coefficient is strictly increasing in \(\varepsilon/\lambda\),
with slope bounded above and away from zero. Its scalar mass
satisfies the thermal crossing estimate of Section~\ref{sec:path}.
\end{lemma}

\begin{proof}
In the following matrix identities \(A,B\) are the actual cube
matrices. Set \(D=B-\varepsilon c'A\). From the old precision we get
\begin{equation}
 \Sigma G^{-1}=I-\varepsilon c'AB^{-1}=DB^{-1},\qquad
 S=BD^{-1},\qquad B^{-1}S=D^{-1}.
\end{equation}
If \(F^0=e^{-U^0}\) is the original full interaction relative to
the Gaussian with covariance \(G\), the interaction in the
\(\psi\) field relative to the Gaussian with covariance \(\Sigma\)
is, up to a constant,
\begin{equation}
 F_{\rm in}(\psi)=
 \exp\{-U^0(S\psi)-\tfrac12\psi^{\mathsf T}H_\varepsilon\psi\},
 \qquad
 H_\varepsilon=S^{\mathsf T}G^{-1}S-\Sigma^{-1}
 =\frac{\varepsilon c'}{a_R}D^{-1}(B-c'A)D^{-1}.
\end{equation}
Indeed \(S^{\mathsf T}G^{-1}=\Sigma^{-1}\),
\(S-I=\varepsilon c'A D^{-1}\), and
\(a_R\Sigma^{-1}A=D^{-1}(B-c'A)\).
Thus the old Hessian is conjugated by \(S\), and the old
Hessian-subtracted spin pressure is evaluated at \(D^{-1}\)
applied to the new primary field in physical units.

For both signs of \(\varepsilon\) write \(D=d_{\varepsilon}(I-K)\), where
\begin{equation}
 d_{\varepsilon}=1-\varepsilon c'A(0),\qquad
 K=d_{\varepsilon}^{-1}\big[Q_\Lambda+\varepsilon c'(A-A(0)I)\big].
\end{equation}
The bulk off-diagonal kernel of \(Q\) dominates a fixed multiple of
\(D_R\), and the off-diagonal entries of \(A\) are bounded by
a fixed multiple of \(D_R\). Hence \(K\) is nonnegative for
large \(R\). Its absolute row and column sums, including a
small fixed exponential weight at length \(R\), are bounded
by a constant strictly below one. This uses the original row
gap and \(\varepsilon=o_R(1)\).
The inverse \(D^{-1}\), its first \(\varepsilon\)-derivative
and the products in the Hessian identity consequently have
uniform weighted bounds. Each term in a clipping discrepancy
has a short walk reaching the exterior, which supplies the
additional factor \(e^{-c\dist(x,\Lambda^c)/R}\) in row \(x\).
A differentiated inverse marks one step and costs only a
polynomial in the walk length. In particular the added Hessian
has weighted size \(C|\varepsilon|/a_R\), and its derivative
has size \(C/a_R\). The derivatives of the two conjugating
factors of the old residual Hessian preserve its weighted
bound \(O_{A_0}(g_{\rm mic}/a_R)\) on a fixed thermal window
of width \(A_0g_{\rm mic}\).

The finite filter used on a label is
\(T_{\rm f}=d_{\varepsilon}^{-1}\sum_{i=0}^{m}K^i\), where \(m\) is the
microscopic Neumann cutoff of Section~\ref{sec:cube-joint-initialization}.
It has nonnegative entries, and \(DT_{\rm f}=I-K^{m+1}\). This supplies the recovery
identity on occupied and unoccupied label sets and the bounded
\(\ell^4\) inverse at the terminal scale. The bulk tail and its
clipping defect have the bounds of
\eqref{eq:cube-init-filter-tail}. The pressure atoms on each
induced spin graph depend on their filtered arguments as in Proposition~\ref{prop:cube-joint-initialization}.
The new filter bounds supply the two ingredients in its
real-field estimate, namely bounded test directions on the
enlarged stencils and recovery on the available rows. The
same exponentially small tails on the two polydiscs give its
complex-field estimate and its connected path expansion.
The bulk and boundary quadratic heads have the weighted
bounds of Proposition~\ref{prop:initial-cube-heads}, since
\(|\varepsilon|=O(g_{\rm mic})\). Cutting these heads and
then removing their bulk second moments gives the interface
terms on the same actual collar labels as in
Lemma~\ref{lem:kinetic-initial-extraction}. Multiplication on
both occupied and unoccupied labels by the off-shell mass and
kinetic factors therefore gives activity bounded by \(CU\)
and partial derivatives bounded by \(CU/x_0\) on the smaller
discs of that lemma.

Normalize the bulk truncation by its row sum, absorbing
the fourth power of this sum into the initial quartic.
The sum is bounded above and away from zero and all fixed
absolute moments of order \(j\) of the normalized filter are
\(O_j(R^j)\), by geometric summation of the walks. It has
the cubic symmetries used in the ordinary transport and
has normalized Fourier multiplier bounded below on the
low-frequency bubble band. In particular the initial quartic
\(x_0(\varepsilon)\asymp x_0(0)\) and its comparison sequence
obey the bounds used in the shrinking chart. The coefficients
are continuous in the external parameters. The pole displacement
still uses the original reference sequence.

Finally \(B-c'A=\zeta J_\Lambda+(1-c')A\). With
\(p=\widehat A(0)=\widehat B(0)\) in bulk and
\(j_2=\tfrac12\sum_z z_1^2D_R(z)\asymp R^2\), the
\(\varepsilon\)-derivative of the coefficient of \(|k|^2\)
in the added bulk Hessian is
\begin{equation}
 \frac{c'\zeta j_2}{a_Rp^2}\frac{1+\varepsilon c'}
 {(1-\varepsilon c')^3}
       +O\!\left(\frac{R^2}{a_R}(R/n)^2\right).
\end{equation}
This follows by expanding the scalar bulk symbol
\(\varepsilon c'(\widehat B-c'\widehat A)/
[a_R(\widehat B-\varepsilon c'\widehat A)^2]\).
The old conjugated Hessian contributes
\(O_{A_0}(g_{\rm mic}R^2/a_R)\) to this derivative. Cutting the
weighted tails before extracting the kinetic coefficient costs
\(O((R^2/a_R)e^{-c b_0/R})\) in the derivative. Since
\(a_R=\lambda R^2\), the slope in \(\varepsilon/\lambda\)
has the asserted strict lower bound. The old scalar row is
multiplied by \((1-\varepsilon c')^{-2}\). The new normalized
bare mass from \(H_\varepsilon\) is temperature independent
and is \(O(|\varepsilon|\lambda^{-1}(b_0/n)^2)\), since
\(|1-c'|=O((R/n)^2)\). These estimates preserve the strict
thermal crossing for \(n\ge n_0(R)\).
\end{proof}

At the last integrated scale the auxiliary precision has a spectral
interval on which both the old terminal covariance and the released
slice are large. The generalized principal eigenvector gives a test
with uniformly positive spectral mass on this interval.

\begin{lemma}[A variance margin after releasing one scale]
\label{lem:aligned-curvature-margin}
Let \(J\) be the original cut, let \(k\le J-1\) be a bounded
number of scales earlier, and write \(T=C^0_T\), \(\sigma^2=
\ell^{\mathsf T}T\ell\). Set \(W_k=t^{-1} C^\Sigma_{T,k}\)
for the transfer decomposition. Then \(T\ell\in\operatorname{ran} W_k\)
and, uniformly in the physical window,
\begin{equation}
 (T\ell)^{\mathsf T}W_k^\dagger(T\ell)
       \le \sigma^2-c\lambda n^2 .
\end{equation}
The constant can depend on the fixed terminal aspect.
\end{lemma}

\begin{proof}
Write \(\mathcal Q=\mathcal Q_{c',0}\), \(y=\mathcal O^*\ell=
(\mathcal L_\Lambda\ell,0)\), and let \(\mu\) be the spectral
measure of \(\mathcal Q\) at \(y\). Denote the old terminal multiplier by
\(f\), the sum of the released multipliers by \(d\), and set
\(h(s)=(1-\eta s)(f(s)+d(s))\). Then
\(T=a_R\mathcal O f(\mathcal Q)\mathcal O^*\) and
\(W_k=a_R\mathcal O h(\mathcal Q)\mathcal O^*\).
For every physical test \(z\), spectral Cauchy--Schwarz gives
\begin{equation}
 |\langle z,T\ell\rangle|^2
 \le \langle z,W_k z\rangle\,
        a_R\int \frac{f(s)^2}{h(s)}\,d\mu(s),
\end{equation}
with zero quotients interpreted as zero. This proves the range
assertion and bounds the inverse quadratic form by the integral.
Since \(\eta\|\mathcal Q\|\le 1/2\) and \(s f(s)\le1\),
\begin{equation}
 \sigma^2-(T\ell)^{\mathsf T}W_k^\dagger(T\ell)
 \ge a_R\int\frac{f(s)d(s)}{f(s)+d(s)}\,d\mu(s)
                -2\eta a_R\|y\|^2.
\end{equation}

Put \(r_J=L^{-2J}\asymp M_*(R/n)^2\). On
\(0<s\le\epsilon r_J\), the bound
\(\sum_{i<J}q_i(s)\le C/r_J\) gives \(f(s)\ge 1/(2s)\)
if \(\epsilon\) is small. The low-frequency lower bound for
\(q_{J-1}\) gives \(d(s)\ge c/r_J\) there.
To bound the spectral mass, take \(U=(\mathcal L_\Lambda\ell,\ell)\).
The identity \(A_\Lambda\ell=c_{*,\Lambda}^{-1}B_\Lambda\ell\)
and the normalization \(\ell^{\mathsf T}B_\Lambda\ell=1\) give
\begin{equation}
 \mathcal Q U=\left(0,\frac{c_{*,\Lambda}-c'}{c'c_{*,\Lambda}}
                   B_\Lambda\ell\right),\qquad
 \langle y,U\rangle=c_{*,\Lambda}^{-1}.
\end{equation}
Both vectors have bounded norm, and the last scalar is bounded
below. The component of \(U\) above \(\epsilon r_J\) has norm
at most \(\|\mathcal Q U\|/(\epsilon r_J)=O(\lambda^2\bar g_N)\).
Hence \(\mu([0,\epsilon r_J])\ge c>0\) for all sufficiently
large \(n\) at fixed \(R\). The integral is at least
\(c/r_J\), while \(\|y\|^2=O(1)\). Using
\(a_R/r_J\asymp\lambda n^2\) proves the claim.
\end{proof}

The preceding lemma estimates the inverse quadratic form at the
earlier cut. We now compare its margin with the change in the
variance of the original scalar. This gives a positive Gaussian
covariance on which the curvature calculation can be performed.

\begin{proposition}[Positive hard covariance for the original coordinate]
\label{prop:prescribed-positive}
Perform the finite-range integrations in \(\psi\) and stop one scale
before the prescribed cut. Write \(C_k=C^\Sigma_{T,k}\). In this
representation the original coordinate has Gaussian variance
\(\widetilde V=\eta_0^2+w_0^{\mathsf T}S\Sigma S^{\mathsf T}w_0\)
and cross covariance \(v=T\ell\) with \(\psi\). Throughout the
matching rectangle, for sufficiently large fixed \(R\) and
\(n\ge n_0(R)\),
\begin{equation}
 v\in\operatorname{ran} C_k,\qquad
 \widetilde V-v^{\mathsf T}C_k^\dagger v\ge c\lambda n^2,\qquad
 \widetilde B_k=C_k-\frac{vv^{\mathsf T}}{\widetilde V}
       \succeq \frac{c\lambda n^2}{\widetilde V}C_k .
 \label{eq:prescribed-positive}
\end{equation}
The profile \(v/\widetilde V\), variance \(\widetilde V\), and
hard covariance \(\widetilde B_k\) satisfy the Gaussian
hypotheses of Proposition~\ref{prop:terminal-graph} for the
curvature of the original terminal marginal.
\end{proposition}

\begin{proof}
Under the physical law one may sample the original coordinate
given \(\phi\) as \(w_0^{\mathsf T}\phi+\eta_0 Z\), where \(Z\)
is an independent standard normal. After the change of variable
the conditional mean is \(w_0^{\mathsf T}S\psi\), also when
the physical interaction is included. Under the unweighted
Gaussian in \(\psi\) the cross column is
\(\Sigma S^{\mathsf T}w_0=Gw_0=T\ell\). The identity
\(tG\Sigma^{-1}G=G-\varepsilon S A_c\) gives
\begin{equation}
 t\widetilde V=\sigma^2-\varepsilon
             (\eta_0^2+w_0^{\mathsf T}S A_cw_0).
\end{equation}
For the matrix in the error, put \(M=G^{-1/2}A_cG^{-1/2}\),
so \(0\preceq M\preceq I\). Then
\(S A_c=G^{1/2}(tI+\varepsilon M)^{-1}M G^{1/2}\)
is positive and bounded above by \(C A_c\). The variance
bounds in Proposition~\ref{prop:original-curvature-comparison-bounds}
give \(\eta_0^2+w_0^{\mathsf T}A_cw_0\le C\lambda n^2\).
Hence
\begin{equation}
 |t\widetilde V-\sigma^2|\le C|\varepsilon|\lambda n^2 .
\end{equation}
Lemma~\ref{lem:aligned-curvature-margin} applied to
\(W_k=C_k/t\) proves the first two estimates in
\eqref{eq:prescribed-positive}, after the terminal aspect is
fixed and \(R\) is large enough that \(C|\varepsilon|\) fits
in its margin. The range estimate and Cauchy--Schwarz on the
support of \(C_k\) give the third. The lower bound is uniform
in variance units. The relative factor \(\lambda n^2/\widetilde V\)
is of order \(\lambda^2 x_k\).

Set \(u=T\ell/\sigma^2\). The new mean profile
\(v/\widetilde V=(\sigma^2/\widetilde V)u\) is a bounded
positive multiple of the original one. Since
\(\ell^{\mathsf T}u=1\), \(\|\ell\|_2\le C\) and
\(\|u\|_{F_{n,p}}\le Cn^{-2}\), the profile at the rescaled
coordinate has \(F_{n,p}\)-norm \(O(n^{-1})\) and
\(b_k^{-2}\|nv/\widetilde V\|_2^2\asymp1\).
Moreover \(n^2/\widetilde V=O(\lambda x_k)\), since shortening
the horizon by one step preserves comparability with
\(\bar g_N\). For the hard field let \(D_{\rm rel}\) be the
released old slice and let \(T_h=T-\sigma^2 uu^{\mathsf T}\).
The transferred covariance is bounded above by
\(t(T+D_{\rm rel})\), by its spectral multiplier. Thus
\begin{equation}
 0\preceq\widetilde B_k\preceq t(T_h+D_{\rm rel})
                 +C|\varepsilon|\lambda n^2 uu^{\mathsf T}.
\end{equation}
The old hard covariance and the released slice at a fixed
fraction of scale \(n\) have the raw Gaussian norm moment of
Section~\ref{sec:ext-cov}. A Gaussian with the last rank-one
covariance has norm bounded by \(C\sqrt{\lambda}\,n^{-1}|Z|\).
Gaussian addition and conditional Jensen's inequality give
\(\E\exp(c n^2\|Y\|_{F_{n,p}}^2/\lambda)\le C\) for
\(Y\sim N(0,\widetilde B_k)\).

Gaussian regression followed by addition of the integrated
covariance now gives the original scalar density, up to a
positive constant,
\begin{equation}
 \exp[-a^2/(2\widetilde V)]\,
 \E_{\widetilde B_k}F_k(a v/\widetilde V+Y),
 \label{eq:aligned-original-density}
\end{equation}
where \(F_k\) is the exact RG output for \(F_{\rm in}\)
in the \(\psi\) field. Indeed the conditional covariance of the
full \(\psi\) equals \((\Sigma-C_k)+\widetilde B_k\).
The regression data here are continuous in the external
parameter and independent of the running mass. In the physical
\(\phi\) coordinates the same hard covariance is
\(S\widetilde B_k S^{\mathsf T}\), for the representation with
total covariance \(S\Sigma S^{\mathsf T}\). This identifies
precisely the subtraction estimated in
\eqref{eq:prescribed-positive}. In the direct comparison using
\(\phi\) itself with covariance \(\Sigma\), the column would
still be \(\Sigma w_0\).
\end{proof}

The original scalar density is now expressed in the variables of the
coupled RG chart. We use that density for mass matching and the
natural terminal field of the same chart for the magnetic moment.

\begin{corollary}[The prescribed matched point in dimension four]
\label{cor:d4-prescribed}
For all sufficiently large fixed \(R\) and \(n\ge n_0(R)\),
there is a physical temperature at which the original prescribed
terminal marginal has zero curvature and
\(\chi_\Lambda^F\asymp_R n^2\sqrt{\log n}\).
\end{corollary}

\begin{proof}
Use scales up to \(k\) for the backward shooting with terminal
kinetic coordinate \(z_k=0\). For each fixed external pair
\((\beta,\varepsilon)\), Lemma~\ref{lem:aligned-initialization}
puts the physical microscopic family and its two off-shell
coordinates in the joint initial chart. The RG map in \(\psi\)
uses the transfer slices themselves. In particular it has the
finite propagation, filtered bubble and annular estimates used in
Proposition~\ref{prop:residual-coupled-chart}. The normalized
microscopic filter of Lemma~\ref{lem:aligned-initialization}
has the same moment and recovery bounds needed for that
proposition. Its coupled collar transport, the connected
extraction of Proposition~\ref{prop:cube-bounded-extraction}
and the summable kinetic source of
Lemma~\ref{lem:kinetic-backward-sum} therefore apply as in
Theorem~\ref{thm:d4-finite-propagation}. They give the same
terminal coefficient box and inverse mass derivative, with
\(x_k\asymp_R(\log n)^{-1}\). The comparison coordinate
used at the terminal step is given by
\eqref{eq:aligned-original-density}. By
Proposition~\ref{prop:prescribed-positive} all its Gaussian
inputs satisfy Proposition~\ref{prop:terminal-graph}.
In particular the leading derivative of its rescaled curvature
with respect to the terminal mass is
\(b_k^{-2}\|nv/\widetilde V\|_2^2\asymp1\), with an \(o(1)\)
error. The terminal scalar graph has the required stable and
activity derivatives. In the residual chart the inverse compensated kinetic and mass
coordinates have the same first-variation bounds used in
Theorem~\ref{thm:d4-finite-propagation}. Cauchy estimation
on the enlarged kinetic disc gives the bounded normalized
kinetic derivative of this scalar graph.

Proposition~\ref{prop:kinetic-backward} gives continuous bare
targets \(\nu_0^*,z_0^*\). Fix the thermal width as in
Section~\ref{sec:path}, then take \(H\) large enough that the
two kinetic faces enclose \(z_0^*\), using the positive slope
in Lemma~\ref{lem:aligned-initialization} and the estimates
\(|z_0^*|+|z_0^{\rm bare}(\beta,0)|\le C_{A_0}\lambda x_0(0)\).
The constants in the backward target bound can be taken
before choosing \(H\). In the bulk source of
Proposition~\ref{prop:bulk-kinetic-source} the fixed microscopic
head enters through \(M_0/x_0\) and \(M_j R/(b_j x_j)\).
Its raw second moment is \(O_{A_0,H}(\lambda x_0)\).
The initial ratio is therefore small, and at later steps
\((x_0/x_j)R/b_j\) stays bounded by the reciprocal comparison
and geometric growth of \(b_j\). The separately generated
filtered head has the uniform bounds of
Proposition~\ref{prop:residual-coupled-chart}. In the initial
stable box a collar contribution \(O_{A_0,H}(\lambda x_0)\)
costs \(O_{A_0,H}(M_R^{-2})\). Thus each fixed choice of
widths fits the same box and source estimates by taking
\(R\) large afterwards. The added scalar perturbation tends
to zero in the thermal face margin as \(n\) grows at fixed \(R\).
The uniform thermal slope and
the bounded normalized target consequently give the opposite
thermal face signs. Corollary~\ref{cor:kinetic-two-matching}
sets both bare discrepancies to zero. At this point the
off-shell reweightings restore \(F_{\rm in}\) exactly. Its
density in \eqref{eq:aligned-original-density} has zero
curvature for the original \(a_0\).

For the magnetic moment return to the natural terminal field in
\(\psi\) at cut \(k\). Lemma~\ref{lem:kinetic-reference-terminal}
applies here because the aspect changes by a fixed factor and the
pole displacement fixed at \(J\) still has the required order at \(k\). More explicitly, conditioning
\(T+D_{\rm rel}\) on \(\ell\) gives a hard covariance bounded
above by
\begin{equation}
 T_h+(I-u\ell^{\mathsf T})D_{\rm rel}(I-\ell u^{\mathsf T}).
\end{equation}
This follows by using \(u\) as a candidate linear predictor.
The added projected slice has the terminal moment, since
\(\ell^{\mathsf T}D_{\rm rel}\ell=O(\lambda n^2)\). Adding the released slice
changes the old mean profile by \(o(n^{-2})\) in the raw norm by
covariance Cauchy--Schwarz and \(\sigma^2\asymp n^2/(\lambda x_k)\).
The transferred natural hard moment, profile and integrated
unit source bounds then follow by the multiplier argument of
that lemma. Propositions~\ref{prop:terminal-packing}
and~\ref{prop:terminal} use the coefficient bounds at the
matched mass and give
\(\E(\one^{\mathsf T}\psi)^2\asymp_R
V^{3/2}g_k^{-1/2}\) for the full, interacting field.
These moment bounds only use the size of the matched mass.
They apply with the mass chosen by the original curvature graph.

In the physical field this linear functional is paired with
\(m=S^{-\mathsf T}\one=(I-\varepsilon c'B^{-1}A)\one\).
The absolute short row bounds give
\(1-C|\varepsilon|\le m_x\le1+C|\varepsilon|\) uniformly in
the cube. Thus \(Am\) has entries bounded above and away from
zero at fixed \(R\). In physical units the conditional
Hubbard--Stratonovich field has mean \(\sqrt\beta A\sigma\)
and Gaussian covariance \(A\). Positivity of the spin
two-point functions as in Lemma~\ref{lem:spin} compares the
weighted spin variance to \(\E(\one^{\mathsf T}\sigma)^2\),
and the added noise variance is \(m^{\mathsf T}Am=O_R(V)\).
Converting the fixed field normalization and dividing by
\(V=n^4\) proves the assertion.
\end{proof}

\subsection{Proof of the susceptibility theorem}
\label{sec:proof-susceptibility}

\begin{proposition}[Physical terminal data in the fractional dimensions]
\label{prop:fractional-matched-data}
Fix \(d\in\{1,2,3\}\). For all sufficiently large fixed \(R\) and
\(n\ge n_0(R)\), there is a physical temperature at which the
prescribed terminal curvature vanishes. At this temperature the
terminal interaction satisfies Assumption~\ref{ass:terminal}
with \(g=g_N\asymp_{d,R}(\log V)^{-1}\).
\end{proposition}

\begin{proof}
In dimension one take \(p=1/8\), so \(p<s/2=1/4\) and
\(2p>s-1=-1/2\). The raw field norm with this weight admits the
slice, principal-profile and integrated-source bounds in
Section~\ref{sec:ext-cov}. The initial quadratic estimates in
Proposition~\ref{prop:initial-cube-heads} and the interval
specialization of Proposition~\ref{prop:cube-retained-chart} therefore
use the same weight as the terminal integration.

Proposition~\ref{prop:cube-retained-chart}
and Corollary~\ref{cor:fractional-cube-path} give the one-step and
terminal coefficient estimates for the exact initial family of
Proposition~\ref{prop:cube-joint-initialization}.
The terminal graph follows from Proposition~\ref{prop:terminal-graph}
and the conditioned Gaussian bounds of Section~\ref{sec:ext-cov}.
Apply Theorem~\ref{thm:path} and Corollary~\ref{cor:physical-mass-crossing}
to the joint physical initialization.
The resulting real temperature has zero terminal curvature.
The reciprocal increments give \(g\asymp_{d,R}(\log V)^{-1}\).
The activity estimate gives
\(g^{-9/4}\|K_N\|\le C g^{3/4}+C U d_0^N g^{-9/4}=o(1)\).
Proposition~\ref{prop:terminal-packing} gives the pointwise
terminal error and its first two field derivatives. The finite
Neumann filter has the row and inverse bounds proved there.
The prescribed principal variance is of order \(\sqrt V/g\) at
fixed \(R\). By Lemma~\ref{lem:terminal-gaussian-comparison}
the integrated slices contribute only \(O_R(\sqrt V)\) to it,
so \(\sigma_T^2\asymp_R\sqrt V/g\). The short-range row gap gives
\(\|\ell\|_2\le\|B_\Lambda^{-1/2}\|\le C\).
Section~\ref{sec:ext-cov} gives the exponential moment of the
conditioned Gaussian field, and Corollaries~\ref{cor:terminal-principal-profile}
and~\ref{cor:integrated-uniform-source} give the profile and the
integrated uniform-source bounds at the same weights. At the bare
mass the initial density is physical, and exact Gaussian integration
preserves positivity. This verifies Assumption~\ref{ass:terminal}.
\end{proof}

\begin{proof}[Proof of Theorem~\ref{thm:susceptibility}]
In dimensions one through three, take the physical temperature of
Proposition~\ref{prop:fractional-matched-data}. Proposition~\ref{prop:terminal}
gives the initial auxiliary moment of order \(V^{3/2}g^{-1/2}\). The fixed canonical
field rescaling changes only the constants for fixed \(R\).
The microscopic split has a positive diagonal and a bounded row,
so Lemma~\ref{lem:spin} applies with \(\beta\) bounded above and
below by positive constants. Its \(O_R(V)\) error is of lower
order. Dividing the spin moment by \(V\) proves the result.

In dimension four, Corollary~\ref{cor:d4-prescribed} supplies a
physical temperature for the prescribed coordinate and gives
\(\chi_\Lambda^F\asymp_R n^2\sqrt{\log n}\). Since \(V=n^4\), this
has the asserted order.
\end{proof}



\appendix

\section{Polynomial covariance slices and terminal conditioning}
\label{app:gaussian-decomposition}

\subsection{Scalar slices}
\label{app:gaussian-scalar}
Fix a common spectral upper bound \(S\) for the auxiliary precisions and
write \(4t/S=2-2\cos\theta\), \(0\le\theta\le\pi\).
We use the scalar periodization of \cite[Lemma~2.1]{IsingCorrelation}.
Let \(\chi\) be a nonzero, nonnegative, even smooth function supported
in \([-1/2,1/2]\), \(F=|\widehat\chi|^2\), and
\(I_F=\int_0^\infty T F(T)\,dT\). Define
\begin{equation}
 P_T(\theta)=\sum_{h\in\mathbb Z}F(T(\theta+2\pi h)),\qquad
 q_i(t)=\frac{4}{S I_F}\int_{I_i}T P_T(\theta)\,dT,
 \quad I_0=(0,L],\quad I_i=[L^i,L^{i+1}]\ (i\ge1).
\end{equation}
Put \(r_i=L^{-2i}\). For \(J\ge1\), write \(r=L^{-2J}\),
\(g_<(t)=\sum_{i<J}q_i(t)\), \(f(t)=t^{-1}-g_<(t)\), and
\(s_J(t)=tf(t)=1-tg_<(t)\).

\begin{lemma}[Polynomial covariance slices and their tail]
\label{lem:scalar-polynomial-slices}
The polynomials \(q_i\) are strictly positive on \([0,S]\),
satisfy \(q_i(t)\ge q_i(0)\) for \(t<0\), and sum to \(1/t\)
on \(0<t\le S\). For each fixed integer \(m\ge1\),
\begin{equation}
 \deg q_i\le C_L L^i,\qquad
 q_i(t)\le C_m r_i^{2m-1}(t+r_i)^{-2m},\qquad
 q_i(t)\ge c/r_i\quad(0\le t\le c_*r_i).
 \label{eq:gaussian-scalar-bounds}
\end{equation}
The upper bound holds on \([0,S]\). The polynomial \(s_J\) has
degree at most \(C_Lr^{-1/2}\), and on \(0<t\le S\),
\begin{align}
 g_<(t)&\asymp (t+r)^{-1}, &
 0\le s_J(t)&\le C_m\bigl(r/(t+r)\bigr)^m,\label{eq:gaussian-partial-scalar}\\
 0\le f(t)(1-s_J(t))&\le C_m r^{2m-1}(t+r)^{-2m}.
 \label{eq:gaussian-tail-scalar}
\end{align}
\end{lemma}

\begin{proof}
Poisson summation expresses \(P_T\) as an even trigonometric polynomial
of degree at most \(T\). Its Fourier coefficients are proportional to
the autocorrelation of \(\chi\) at integer multiples of \(1/T\), so
they are nonnegative and vanish outside this degree. Integrating each
period and using
\(\sum_h(\theta+2\pi h)^{-2}=(2-2\cos\theta)^{-1}\) gives
\(\sum_{i\ge0}q_i(t)=1/t\) for \(0<t\le S\).
The definition gives \(q_i\ge0\) on \([0,S]\). For \(t<0\) each cosine
in its expansion becomes a hyperbolic cosine, so \(q_i(t)\ge q_i(0)\).
For \(T\ge1\), rapid decay gives
\(P_T(\theta)\le C_a(1+T^2t)^{-a}\) for every \(a>0\). For
\(T<1\) it is a constant of order \(T^{-1}\).
Integration over the scale interval gives the degree and upper
bounds in \eqref{eq:gaussian-scalar-bounds}. The lower bound follows
from \(F(0)>0\). In fact the integral of the zeroth period is strictly
positive for every \(0\le\theta\le\pi\), since the entire Fourier
transform of \(\chi\) cannot vanish on a nonempty interval. Continuity
then gives \(\min_{[0,S]}q_i>0\) at each fixed \(i\).

The degree bound for \(s_J\) follows from that for the slices.
The partial sum satisfies \(g_<\le\min(t^{-1},C/r)\). For \(t\ge D r\), geometric summation
of \eqref{eq:gaussian-scalar-bounds} gives
\(f(t)\le C_m r^{2m-1}t^{-2m}\), hence \(g_<\ge(2t)^{-1}\) for fixed
large \(D\). For \(t\le Dr\) the zeroth period in the last integrated
slice gives \(q_{J-1}(t)\ge c_{D,L}/r\). After setting \(T=L^J z\),
this follows from the positive minimum of
\(\int_{1/L}^1 z F(yz)\,dz\) for \(0\le y=L^J\theta\le C_D\).
Finally \(0\le tf\le1\) and \(1-s_J=tg_<\) control the residual for
\(t\le r\), and the bound on \(f\) controls it for \(t\ge r\).
\end{proof}

\subsection{Comparison of the joint Gaussian fields}
\label{app:gaussian-joint}
In this appendix the joint Gaussian fields have unnormalized
covariances, so observed covariances are multiplied by \(a_R\) at the
end. Write \(Q_c=\mathcal Q_c\), \(D=\mathcal D\),
\(U=\mathcal L_\Lambda\), \(A=U^*D^{-1}U\), and \(B=B_\Lambda\).
Let \(Z=X_<+X_T\) be the sum of independent Gaussian fields with
covariances \(g_<(Q_c)\) and \(f(Q_c)\). Write
\(a=\ell^{\mathsf T}\mathcal O Z\), \(Y=\mathcal O X_T\), and
\(C_T=\mathcal O f(Q_c)\mathcal O^*\).

\begin{lemma}[Conditioning the joint Gaussian field]
\label{lem:joint-conditioned-covariance}
For \(1\le c<c_{*,\Lambda}\), on tests of finite inverse energy,
\(\operatorname{Cov}(Z\mid a=0)\le C Q_1^{-1}\) uniformly in
\(c\). The observed terminal field satisfies
\begin{equation}
 \mathcal S_\ell(C_T)\le\operatorname{Cov}(Y\mid a=0)
 =\mathcal O\big[
 s_J(Q_c)\operatorname{Cov}(Z\mid a=0)s_J(Q_c)
 + f(Q_c)(I-s_J(Q_c))\big]\mathcal O^*.
\end{equation}
\end{lemma}

\begin{proof}
The quadratic-form identities follow first on regularized or
spectrally truncated tests and then by Gaussian \(L^2\) limits.
Regress \(X_T\) on \(Z\) to obtain \(X_T=s_J(Q_c)Z+\eta\), where
\(\eta\) is independent of \(Z\) and has covariance
\(f(Q_c)(I-s_J(Q_c))\). This gives the displayed identity after
conditioning on \(a=0\). The lower bound follows by rank-one regression,
since \(\operatorname{Cov}(Y,a)=C_T\ell\) and
\(\operatorname{Var}a\ge\ell^{\mathsf T}C_T\ell\).

To compare the conditioned joint covariance to \(Q_1^{-1}\), condition
first on the full observed field \(h=U^*Z_a\), where \(Z_a\) denotes
the auxiliary component. Given \(h\), the two blocks are independent.
Their conditional means are \(D^{-1}UA^{-1}h\) and \(cB^{-1}h\);
their covariance blocks are
\(R_A=D^{-1}-D^{-1}UA^{-1}U^*D^{-1}\) and \(cB^{-1}\).
Indeed, completing the physical square gives a tilt of the auxiliary
Gaussian depending only on \(h\), and ordinary Gaussian regression
gives the first mean and covariance. Denote these combined mean and
covariance operators by \(T_c\) and \(J_c\), respectively. Then
\begin{equation}
 \operatorname{Cov}(Z\mid a=0)=J_c+T_c\mathcal S_\ell(V_c)T_c^*,
 \qquad Q_1^{-1}=J_1+T_1 V_1T_1^*.
\end{equation}
Section~\ref{sec:ext-cov} gives
\(\mathcal S_\ell(V_c)\le C V_1\), \(\|V_1\|\le C\varrho^{-s}\) and
\(|c-1|\le C\varrho^s\). We have \(J_c\le C J_1\), while
\((T_c-T_1) V_1(T_c-T_1)^*\) is confined to the physical block and
bounded there by \(C\varrho^s B^{-2}\le C B^{-1}\). Expanding
\(T_c=T_1+(T_c-T_1)\) and using Cauchy--Schwarz thus gives
\(\operatorname{Cov}(Z\mid a=0)\le C Q_1^{-1}\).
\end{proof}

\subsection{Inverse energy on polynomial tests}
\label{app:gaussian-krylov}
In the fractional channel let \(N=N_{\rm sm}\), and let \(W\) have
continuous auxiliary component
\(\sqrt{w_\alpha(u)}D_R^{N/2}P(u,\cdot)\), where \(P\) has degree less
than \(K\) in \(u\). Suppose the coefficient vectors of \(P\), the white
component, and the physical component are supported in a common cube
of side \(M\ge R\). With \(0<\alpha=d/4<1\), we claim
\begin{equation}
 \langle W,Q_1^{-1}W\rangle
 \le C\{K^2+K^{2-2\alpha}(M/R)^{2\alpha}\}\|W\|^2 .
 \label{eq:gaussian-krylov}
\end{equation}

For the polynomial part, write \(t=1-u\) and
\(\nu(t)=t^\alpha(1-t)^{N-\alpha-1}\). For any scalar polynomial
\(p\) of degree below \(K\),
\begin{equation}
 \int_0^1 \frac{|p|^2}{t}\nu\,dt\le CK^2\|p\|_{L^2(\nu)}^2,
 \qquad
 \int_0^1 \frac{|p|}{t}\nu\,dt\le CK^{1-\alpha}\|p\|_{L^2(\nu)} .
\end{equation}
Indeed, for \(a=N-\alpha-1\), the shifted Jacobi polynomial
\(P_j^{(\alpha,a)}(1-2t)\) has value
\(\binom{j+\alpha}{j}=O((j+1)^\alpha)\) at zero and squared norm
comparable to \((j+1)^{-1}\). Its ratio to the endpoint value is
\(\sum_{h=0}^j(-j)_h(j+\alpha+a+1)_ht^h/((\alpha+1)_h h!)\).
For \(j<K\) and \(t\le K^{-2}\), the absolute sum is bounded by
\(\sum_{h\ge0} C^h/(h!)^2\). Cauchy--Schwarz in the orthonormal
expansion gives \(\sup_{t\le K^{-2}} |p(t)|\le
CK^{1+\alpha}\|p\|_{L^2(\nu)}\). Integrate this estimate up to
\(K^{-2}\). Above \(K^{-2}\) use \(t^{-1}\le K^2\) for the squared
estimate, and Cauchy--Schwarz with
\(\int_{K^{-2}}^1 t^{\alpha-2}(1-t)^a\,dt\le CK^{2-2\alpha}\)
for the linear estimate.

For the spatial part, a vector \(v\) supported in a cube of side
\(M\ge R\) satisfies, with normalized Fourier measure,
\begin{equation}
 \int \omega(k)^{-\alpha}|\widehat v(k)|^2\,dk
 \le C(M/R)^{2\alpha}\|v\|_2^2 .
\end{equation}
On \(|k|\le M^{-1}\) use \(\|\widehat v\|_\infty^2\le CM^d\|v\|_2^2\)
and integrate \((R|k|)^{-2\alpha}\), which is integrable since
\(d>2\alpha\). On the complement use the supremum of the multiplier.
By the variational formula for inverse energy the finite-volume form
is bounded by its bulk counterpart on the zero extension of the
physical component. Completing the bulk auxiliary square gives
\begin{equation}
 \langle W_a,D^{-1}W_a\rangle+
 \big\|(B_\infty-A_\infty)^{-1/2}
        (W_p+\mathcal L^*D^{-1}W_a)\big\|_2^2 .
\end{equation}
The first polynomial inequality bounds the first term by
\(CK^2\|W_a\|^2\), since \(1-u\widehat D_R\ge1-u\).
Put \(d_k=\widehat D_R(k)\). The second polynomial inequality bounds
the squared modulus of the continuous contribution to
\(\widehat{\mathcal L^*D^{-1}W_a}(k)\) by
\begin{equation}
 CK^{2-2\alpha}|d_k|^{2N}
     \int_0^1 w_\alpha(u)|\widehat P(u,k)|^2\,du .
\end{equation}
Use \(|d_k|^{2N}\le |d_k|^N\) and apply the spatial bound to
\(D_R^{N/2}P(u,\cdot)\), supported in a cube of side \(M+C_NR\).
Since \(B_\infty-A_\infty=\zeta(I-D_R)^\alpha\), integration gives
the second term of \eqref{eq:gaussian-krylov}. The same spatial bound
applies directly to the physical and white sources.

\begin{corollary}[Resolvent bound for the conditioned terminal field]
\label{cor:terminal-resolvent-majorant}
Let \(1\le d\le4\), \(1\le c<c_{*,\Lambda}\), and take the terminal
cut \(r=L^{-2J}\asymp M_*(R/n)^2\) of
Section~\ref{sec:FinitePropagationTerminalConditioning}.
Let \(Y\) and \(C_T\) have the unnormalized convention of
Lemma~\ref{lem:joint-conditioned-covariance}. For each fixed integer
\(m\ge1\), uniformly in \(c\),
\begin{equation}
 a_R\mathcal S_\ell(C_T)\le a_R\operatorname{Cov}(Y\mid a=0)
 \le C_m a_R r^{2m-1}\mathcal O (Q_c+rI)^{-2m}\mathcal O^* .
\end{equation}
\end{corollary}

\begin{proof}
For the terminal cut choose \(K\asymp_L r^{-1/2}\) larger than the
degree of \(s_J\) plus one. For every physical test \(v\) the vector
\(W=s_J(Q_c)\mathcal O^*v\) has the stated form. Each application of
the continuous diagonal multiplies the polynomial by \(I-uD_R\), and
the off-diagonal input adds a term independent of \(u\). Each block
increases the spatial support by at most \(C_NR\), so
\(M\le n+C_N RK\le C_{M_*,L,N}RK\) at the terminal aspect.
Thus \eqref{eq:gaussian-krylov} gives
\(\langle W,Q_1^{-1}W\rangle\le Cr^{-1}\|W\|^2\).
For the local split \(D=I\), the Schur complement
\(B-A=\zeta(I-D_R)_\Lambda\) has lowest eigenvalue at least
\(c(R/n)^2\) by Section~\ref{sec:ext-cov}. The couplings are bounded,
so completing the same square yields \(\|Q_1^{-1}\|\le C(n/R)^2\),
which gives the same inequality at this aspect.
Insert this bound in the regression formula of
Lemma~\ref{lem:joint-conditioned-covariance}. Use
\(s_J(t)^2\le C_m(r/(t+r))^{2m}\) in its first term and
\eqref{eq:gaussian-tail-scalar} in its second term. Multiplication
by \(a_R\) gives the stated bound.
\end{proof}

\section{Resolvent derivatives at the local endpoint}\label{app:local-resolvent}

Lemma~\ref{lem:local-annulus} bounds discrete derivatives on a stencil
at depth \(\rho_b(x)\) from a face. The calculation below uses the notation of Section~\ref{subsec:local-annulus}.

\begin{proof}[Proof of Lemma~\ref{lem:local-annulus}]
Put \(M_r=\mathcal Q_c+rI\), \(D=B/c+rI\), and
\(H=D-A/(1+r)\) in the bulk. Write
\begin{equation}
 H=d_pI-K_r,\quad d_p=c^{-1}+r,\quad K_r=Q/c+A/(1+r).
\end{equation}
Here \(Q\) has zero diagonal. The bulk Fourier multiplier of
\(H\) is comparable to \(r+\min\{|Rk|^2,1\}\) by the local
Schur comparison, and the same lower bound holds for the
smallest eigenvalue of the bulk matrix \(M_r(k)\). In fact the
block triangular factorization giving the Schur complement
has uniformly bounded factors and inverses.

Let \(U_l=M_r^{-l}f\) with \(\|f\|_2\le1\), and write \(u_l\)
for its physical component and \(w_l=\mathcal L_\Lambda^*
(U_l)_{\rm aux}\) for its observation. Extend physical vectors
by zero. Set \(\widetilde u_l=u_l-d_p^{-l}f_{\rm phys}\).
For \(l\ge2\),
\begin{equation}
 \|\widetilde u_l\|_\infty+\|w_l\|_\infty
 \le C_l R^{-2}r^{1-l}.
 \label{eq:local-rows}
\end{equation}
At \(l=1\) the upper bound \(CR^{-2}(1+\log(1/r))^{1/2}\)
suffices. To prove these bounds, put
\(D_0=\operatorname{diag}((1+r)I,d_pI)\) on all channels.
The Duhamel expansion of \(e^{-t(D_0-K)}\) in the
entrywise nonnegative coupling \(K=D_0-M_r\) shows that
\(M_r^{-l}-D_0^{-l}\) is entrywise nonnegative and bounded by
its bulk matrix elements. Indeed, the removed term is the
path with no coupling; the other terms are positive and finite-volume
paths are a subset of bulk paths. Integrate the expansion with
weight \(t^{l-1}/(l-1)!\). At \(|Rk|\le c_1\), the squared
bulk row norm is controlled by integrating \((r+|Rk|^2)^{-2l}\),
which gives \(C_l R^{-4} r^{2-2l}\) for \(l>1\) and
\(CR^{-4}(1+\log(1/r))\) at \(l=1\). At the remaining
frequencies the uniform matrix gap and the resolvent identity
bound \(\|M_r(k)^{-l}-D_0^{-l}\|\) by
\(C_l (|\widehat Q(k)|+\|\widehat{\mathcal L}(k)\|)\).
Its squared integral is \(O_l(R^{-4})\) by the size and
smoothing of the split. This proves the row estimates for
\(U_l-D_0^{-l}f\). The absolute row norm of \(\mathcal L^*\)
is bounded and the \(\ell^2\) norm of each row is \(O(R^{-2})\);
applying \(\mathcal L_\Lambda^*\) gives the asserted bound on
\(w_l\).

Let \(d_x=\min\{b,\operatorname{dist}(x,\Lambda^c)\}\).
The absolute-leg bound proved above already suffices if
\(d_x\le C_{m,\gamma}R\). Its proof also applies to every site
when \(b=O(R)\), since only the local support and weights of
the observation leg are used. At interior sites the block
equations give
\begin{equation}
 \begin{aligned}
 H\widetilde u_l &=
    \widetilde u_{l-1}+(1+r)^{-1}w_{l-1}
                  +d_p^{-l}K_r f_{\rm phys},\\
 H w_l &=(1+r)^{-1}
     \bigl(A\widetilde u_{l-1}+D w_{l-1}
                        +d_p^{-(l-1)}A f_{\rm phys}\bigr).
 \end{aligned}
\end{equation}
For the second equality first use
\((1+r)w_l=A u_l+w_{l-1}\) and
\(D u_l=u_{l-1}+w_l\), and commute the bulk scalar convolutions
where the stencils lie in the cube. Let \(k_0>|\gamma|/2\) be
a fixed integer. Iterating the displayed equations for
\(j\le k_0\), in a ball of radius \(d_x/2\) with a reserved
\(O(k_0R)\) collar, gives
\begin{equation}
 \|H^j\widetilde u_l\|_{\infty,\,B(x,d_x/2)}
 \le C_{j,l} R^{-2}r^{1-l+j},\qquad l-j\ge2.
\end{equation}
Indeed, the finite-range convolutions multiplying the two
previous components commute with \(H\) in this interior
calculation and have bounded absolute row norms. The source
terms, including any subsequent \(H\)-factors on them, are
bounded by \(C_{j,l}R^{-2}\) using the row \(\ell^2\) norms of
\(K_r,A\) and the bounded absolute row norms of \(H\).
Thus the exponent follows from \eqref{eq:local-rows}.

We convert the interior bounds on powers of \(H\) to derivative
bounds using a cutoff Green function. Set \(\delta=(R/d_x)^2\ge r\)
and \(\widetilde H=H+\delta I\), and use a smooth lattice
cutoff \(\chi\) equal to one on \(B(x,d_x/4)\) with support
in \(B(x,d_x/2)\) and differences of order \(j\) bounded by
\(C_jd_x^{-j}\). Let
\(h_x=\widetilde H^{-k_0} A\nabla^{\gamma,*}\delta_x\)
in the bulk. Then
\begin{equation}
 \|h_x\|_1\le C\delta^{-k_0}d_x^{-|\gamma|},
 \qquad
 \|\nabla^j[(1-\chi)h_x]\|_1
 \le C_j\delta^{-k_0}d_x^{-|\gamma|-j},\quad 0\le j\le 2k_0.
 \label{eq:local-parametrix}
\end{equation}
For these kernel bounds, decompose near zero in the Fourier
variable \(\theta=Rk\) into smooth dyadic bands of size
\(T\ge\sqrt\delta\) and a final ball of size \(\sqrt\delta\).
The symbol of \(\widetilde H\) is comparable to
\(\delta+|\theta|^2\) there; its derivatives at scale \(T\)
are bounded by \(C_a T^{2-a}\), with the zero-th order bound
also applying on the final ball. The resulting row
\(\nabla^j h_x\) on one band is bounded by
\begin{equation}
 C_JR^{-4-|\gamma|-j} T^{4+|\gamma|+j-2k_0}
           (1+T|y-x|/R)^{-J}.
\end{equation}
This follows by integration by parts after rescaling the band,
including the discrete difference multipliers. At high
frequencies \(\widetilde H\) has a uniform gap and bounded
derivatives in \(\theta\); the factor \(\widehat A\) gives
arbitrarily many fixed powers of \((1+|\theta|)^{-1}\)
by choosing the split smoothing order large enough.
The exponentially small white term is harmless on the
scaled torus. Thus the high-frequency row is bounded by
\(C_JR^{-4-|\gamma|-j}(1+|y-x|/R)^{-J}\).
Summing the low-frequency \(\ell^1\) bounds at \(j=0\)
uses \(2k_0>|\gamma|\). For \(j>0\), in the tail
\(|y-x|\gtrsim d_x\) reserve a large fixed power of
\((1+Td_x/R)^{-1}\) before summing. The discrete product
rule with the cutoff gives \eqref{eq:local-parametrix}.

Since \(K_r\) is even and has range \(O(R)\),
\(\widetilde H\) is a constant of size \(O(\delta)\) plus
a bounded weighted sum of centered second differences over
jumps of size \(O(R)\). A product of \(j\) such differences
costs at most \(C_jR^{2j}\) times the largest \(\ell^1\)
norm of a lattice difference of order \(2j\), by telescoping
each jump into coordinate steps. Thus the second bound in
\eqref{eq:local-parametrix} gives
\begin{equation}
 \|\widetilde H^{k_0}[(1-\chi)h_x]\|_1
 \le C \sum_{j=0}^{k_0}
   \delta^{k_0-j}R^{2j}\delta^{-k_0}d_x^{-|\gamma|-2j}
 \le C d_x^{-|\gamma|}.
\end{equation}
The cutoff identity
\begin{equation}
 \nabla^\gamma A\widetilde u_l(x)
 =\langle \chi h_x,\widetilde H^{k_0}\widetilde u_l\rangle
  +\langle \widetilde H^{k_0}[(1-\chi)h_x],
                  \widetilde u_l\rangle
\end{equation}
now bounds its left side by \(C_lR^{-2}r^{1-l}d_x^{-|\gamma|}\)
when \(l\ge k_0+2\). Indeed the first term uses
\(\|\widetilde H^{k_0}\widetilde u_l\|_{\infty,\,\supp\chi}
 \le C_l\delta^{k_0}R^{-2}r^{1-l}\) by the interior powers
and \(r\le\delta\), and the second term uses
\eqref{eq:local-rows}. All pairings are with zero extensions.

Finally, exactly on the cube,
\begin{equation}
 w_m=(1+r)^{-m}w_0+
             \sum_{l=1}^m(1+r)^{-(m-l+1)}A_\Lambda u_l .
\end{equation}
For \(l\ge k_0+2\), use the cutoff bound on the
\(\widetilde u_l\) part. For smaller \(l\) the bound
\(\|\nabla^\gamma A\|_1\le C_\gamma R^{-|\gamma|}\)
and the first-resolvent or higher-resolvent row bounds
suffice: the loss \((d_x/R)^{|\gamma|}\le r^{-|\gamma|/2}\)
and the first-resolvent logarithm are absorbed by taking
\(m>k_0+2+|\gamma|/2\). The direct-source rows
\(\nabla^\gamma w_0\) and
\(d_p^{-l}\nabla^\gamma A f_{\rm phys}\) are bounded by
\(C R^{-2-|\gamma|}\), by the smoothing of the split.
Altogether \(|\nabla^\gamma w_m(x)|\le
C_m R^{-2}r^{1-m}d_x^{-|\gamma|}\).
Squaring this dual norm and multiplying by \(a_Rr^{2m-1}\)
gives the claim since \(a_R=\lambda R^2\) and
\(\lambda R^{-2}r=\lambda b^{-2}\).
\end{proof}

\section{Gaussian integration at a fixed kinetic coefficient}
\label{app:fixed-kinetic}

We collect the identities and bounds for integrating a quadratic
interaction when its kinetic coefficient ranges over a fixed disc.
The square-completion formula also gives the terminal comparison in
Section~\ref{sec:kinetic-terminal-absorption}. The trajectory of
Section~\ref{sec:four-dimensional-flow} uses the finite-propagation
reference on its shrinking kinetic domain.

\subsection{Exact bulk compensation of the pure quadratic terms in dimension four}
\label{sec:kinetic-exact-compensation}

For a purely quadratic input, completing the Gaussian square gives
an exact version of the second-order compensation in
Section~\ref{sec:cube-mass-chart}. The first two Taylor coefficients
of the bulk Fourier symbol determine the mass and kinetic
coordinates. This calculation allows a fixed kinetic coefficient as
the quartic coupling tends to zero and identifies the corresponding
change in the covariance.

\begin{lemma}[Exact compensation of the bulk quadratic two-jet]
\label{lem:kinetic-exact-compensation}
Let \(C\) be a finite-range bulk slice and \(W=W_j^{\rm cov}\)
the integrated bulk covariance at step \(j\). Write
\begin{equation}
 \widehat H(k)=b_j^{-2}u+z^{\rm loc}|k|^2+O(|k|^4)
\end{equation}
for the symbol of the incoming quadratic Hessian. Put
\(a=b_j^{-2}\widehat W(0)\) and
\(m=b_j^{-4}\sum_x x_1^2 W(x)\). The compensated mass and
kinetic coefficients are
\begin{equation}
 \nu=\frac{u}{1-au},\qquad
 z=\frac{z^{\rm loc}-\tfrac12 m u^2}{(1-au)^2}.
 \label{eq:kinetic-rational-compensation}
\end{equation}
Under pure quadratic Gaussian integration and with
\(W^+=W+C\), the compensated dimensionful mass
\(b_j^{-2}\nu\) and the compensated kinetic coefficient \(z\)
are unchanged. In particular there is only the engineering
rescaling of \(\nu\) when \(b_j\) is replaced by \(b_{j+1}\).
The identities are uniform and analytic on the low-jet domain
\(|a u|\le 1/4\), \(|z^{\rm loc}|\le z_0\), wherever the
Gaussian quadratic integration exists. The corresponding
compensated coefficient of a single quartic insertion,
to first order in that insertion, is
\(g=g^{\rm loc}(1-au)^{-4}\).
\end{lemma}

\begin{proof}
For a finite Gaussian vector with covariance \(C\) and
symmetric Hessian \(H\), write \(X=C^{1/2}Z\) with \(Z\)
standard Gaussian. Completing the square in \(Z\) gives
\begin{equation}
\begin{aligned}
 \E_C\!\Big[e^{-\langle\phi+X,H(\phi+X)\rangle/2}
                    F(\phi+X)\Big]
 &=\det(I+C^{1/2}HC^{1/2})^{-1/2} \\
 &\quad\times e^{-\langle\phi,H^+\phi\rangle/2}
   \E_{\widetilde C}F\big((I+CH)^{-1}\phi+X\big),\\
 H^+&=H(I+CH)^{-1},\\
 \widetilde C&=C^{1/2}(I+C^{1/2}HC^{1/2})^{-1}C^{1/2}.
\end{aligned}
\label{eq:kinetic-heat}
\end{equation}
This holds when \(I+C^{1/2}HC^{1/2}\) is positive definite
and requires no inverse of \(C\). It gives the bulk symbol
identities by working with periodic convolution matrices.
For scalar symbols, the identity
\begin{equation}
 \frac{\widehat H^+}{1-(\widehat W+\widehat C)\widehat H^+}
 =\frac{\widehat H}{1-\widehat W\widehat H}
\end{equation}
holds near \(k=0\) whenever the denominators at zero are
nonzero. Taking the Taylor coefficients and using
\(\widehat W(k)=b_j^2a-\tfrac12 b_j^4m|k|^2+O(|k|^4)\)
gives \eqref{eq:kinetic-rational-compensation}. For the
quartic, the degree of the inserted polynomial is unchanged
only by the term of the expectation without Gaussian
contractions. Each zero-momentum external leg is then
multiplied by
\((1+\widehat C(0)\widehat H(0))^{-1}\).
The same denominator identity to the fourth power yields
the stated quartic compensation. Inversion of
\eqref{eq:kinetic-rational-compensation} gives
\begin{equation}
 u=\frac{\nu}{1+a\nu},\qquad
 z^{\rm loc}=\frac{z}{(1+a\nu)^2}+\tfrac12 m u^2.
\end{equation}
Since \(|a|+|m|\le C\lambda\), differentiation gives
\(\partial_\nu u=1+O(\lambda|u|)\),
\(|\partial_\nu z^{\rm loc}|\le C\lambda(|z|+|u|)\), and
\(|\partial_\nu g^{\rm loc}|\le C\lambda|g|\) uniformly on
this domain.
\end{proof}

For the canonical bulk slice, the rational chart has a fixed
kinetic disc when the Hessian is represented by the local
Laplacian. Indeed let \(\omega(k)=1-\widehat D_R(k)\) and
\(r=(R/b_j)^2\). The shifted-precision bound in
Section~\ref{sec:bubble} bounds the bulk auxiliary inverse
by \(C(r+\omega(k))^{-1}\) in operator norm in dimension
four. Applying the scalar cutoff bound before integrating
over \(k\), and using boundedness of the observation
vector, gives
\begin{equation}
 0\le\widehat C(k)\le
 C_m\lambda R^2 r^{2m-1}(r+\omega(k))^{-2m}
 \le C_m\lambda R^2(r+\omega(k))^{-1}.
\end{equation}
The last bound also holds for the merged slice at its largest
scale. To see this, sum the preceding bound with the
corresponding geometric \(r_i\ge r\). For \(r_i\ge\omega(k)\)
sum \(r_i^{-1}\), and for \(r_i<\omega(k)\) sum
\(r_i^{2m-1}\omega(k)^{-2m}\).
Consequently, for
\(\widehat H(k)=b_j^{-2}u+z^{\rm loc}R^{-2}\omega(k)\)
with a bounded adjustment of the normalization of
\(z^{\rm loc}\),
\begin{equation}
 \sup_k |\widehat C(k)\widehat H(k)|
 \le C\lambda(|u|+|z^{\rm loc}|).
\end{equation}
For small \(\lambda\), this bounds the denominator of the
Gaussian identity uniformly away from zero in a kinetic disc
of radius independent of \(g_j\), including for the
holomorphic finite-dimensional Gaussian integrals on that disc.

In particular the pure mass--kinetic contraction need not be
summed as a source proportional to \(\lambda|u||z|\). It is
absorbed exactly by the chart. However, completing the square in
a bulk slice with a quadratic input generally changes its
propagation. For a nonzero even finite-range convolution \(C\)
and a nonzero even finite-range Hessian \(H\) of zero row sum,
\(\widehat C/(1+\widehat C\widehat H)\) cannot be a Laurent
polynomial. If it were, \(1+\widehat C\widehat H\) would divide
\(\widehat C\) in the Laurent polynomial ring, and hence divide
\(1=(1+\widehat C\widehat H)-\widehat H\widehat C\).
A Laurent polynomial unit is a monomial. Symmetry under
inversion of all variables makes this monomial constant, and
its value at \(k=0\) is one, which contradicts
\(\widehat C\widehat H\ne0\). Thus this exact compensation of
the projected quadratic contribution does not justify using the
dressed covariance in the connected finite-range step.

\begin{lemma}[Scale-local expansion of the dressed ordinary slice]
\label{lem:kinetic-dressed-slice}
Let \(C\) be an ordinary covariance slice on a four-dimensional
cube at scale \(b\ge R\). Fix \(1/2<p<1\), put
\(\rho(x)=\min\{b,R+\dist(x,\Lambda^c)\}\) and
\(w(x)=(\rho(x)/b)^p\), and use the raw test norm \(F_b\).
Suppose the slice has range \(K b\) and the annular stencil
variance bounds of Section~\ref{sec:ext-cov} with this
weight and up to at least \(\max\{2,P\}\) orders. With
\(J_R=R^{-2}P_\Lambda L_RP_\Lambda^*\), set
\begin{equation}
 H=zJ_R+b^{-2}u I,\qquad t=|z|+|u|.
\end{equation}
For \(C_0\lambda t<1\), the dressed covariance in
\eqref{eq:kinetic-heat} has the absolutely convergent expansion
\begin{equation}
 \widetilde C=\sum_{m\ge0} C^{(m)},\qquad
 C^{(m)}=(-CH)^mC.
\end{equation}
Each term has range at most \(K_0(m+1)b\) and satisfies
\begin{equation}
 |\nabla_x^\gamma\nabla_y^\eta C^{(m)}(x,y)|
 \le C_0\lambda b^{-2}(C_0\lambda t)^m
 \frac{w(x)w(y)}{\rho(x)^{|\gamma|}\rho(y)^{|\eta|}},
 \qquad |\gamma|,|\eta|\le P.
 \label{eq:kinetic-dressed-slice}
\end{equation}
The same assertion for one unit coefficient derivative in
\((u,z)\) holds with \((C_0\lambda t)^m\) replaced by
\(m(C_0\lambda)^m t^{m-1}\), for \(m\ge1\).
In particular both series converge with an arbitrarily
prescribed fixed exponential weight in the range index \(m\)
when \(\lambda t\) is sufficiently small. For real coefficients
in this disc the Gaussian with covariance \(\widetilde C\)
has the following joint fluctuation bound on any finite
family \(\mathcal P\) of different scale-\(b\) tiles. For
\(Y_P=F_{b,P^*}(X)/\ell_b\) on fixed tile enlargements
and bounded nonnegative integers \(k_P\), one has
\begin{equation}
 \E_{\widetilde C}\!
    \left[\prod_{P\in\mathcal P}Y_P^{k_P}
            \exp\left(\frac{c}{\lambda}
                         \sum_{P\in\mathcal P}Y_P^2\right)\right]
 \le C^{|\mathcal P|}\lambda^{\sum_P k_P/2}.
\end{equation}
\end{lemma}

\begin{proof}
Let \(\|f\|_w=\sup_x|f(x)|/w(x)\). Covariance
Cauchy--Schwarz gives the product of the two annular stencil
standard deviations in every mixed entry. In particular
\begin{equation}
 |\nabla_x^\gamma C(x,y)|
 \le C\lambda b^{-2}\frac{w(x)}{\rho(x)^{|\gamma|}}w(y).
\end{equation}
For \(\nabla_x^\gamma(CJ_R)(x,y)\), extend the second
variable by zero. At depth larger than \(2\sqrt d R\),
write the symmetric walk as an average of centered second
differences of steps of length at most \(\sqrt d R\).
Each such difference is bounded, by twice telescoping in
coordinate directions, by \(C R^2\) times the largest
second stencil bound within distance \(\sqrt d R+2\)
of \(y\), where \(w\) and \(\rho\) are comparable to their
values at \(y\). Closer to the boundary, use the zeroth
bound on the extended row and
\(\rho(y)\asymp R\). The normalization \(R^{-2}\) gives
in both cases
\begin{equation}
 |\nabla_x^\gamma(CJ_R)(x,y)|
 \le C\lambda b^{-2}\frac{w(x)}{\rho(x)^{|\gamma|}}
                    \frac{w(y)}{\rho(y)^2}.
\end{equation}
The scalar mass term has the same bound since
\(b^{-2}\le\rho(y)^{-2}\).
All these rows have support within a fixed multiple of \(b\)
from the stencil anchor. A ball of this radius meets
\(O(b^3)\) sites at each given boundary depth below \(b\),
and has at most \(O(b^4)\) sites of larger depth. Therefore
\begin{equation}
 \sum_{|y-x|\le K_1 b}\frac{w(y)^2}{\rho(y)^2}
 \le C b^3 b^{-2p}\sum_{r=1}^{\lceil b\rceil}
            (r+R)^{2p-2}+C b^2\le C b^2.
\end{equation}
Multiplying the row estimate by \(|f(y)|\le w(y)\|f\|_w\)
gives \(\|CH f\|_{F_b}\le C\lambda t\|f\|_w\),
and hence the same bound with \(F_b\) on the left
replaced by \(\|\cdot\|_w\).

At the right end of \((CH)^m C\), the \(\|\cdot\|_w\)
norm of the \(\eta\)-stencil of column \(y\) of \(C\) is
at most \(C\lambda b^{-2}w(y)\rho(y)^{-|\eta|}\).
Apply the weighted bound to the intervening factors, and
apply the \(F_b\) bound for the leftmost factor. This proves
\eqref{eq:kinetic-dressed-slice}, including the original
mixed bound for \(m=0\). Each factor has range bounded
by a fixed multiple of \(b\). Marking one of the \(m\)
factors and using the same estimate on it with \(t=1\)
proves the derivative assertion.

For real coefficients, the nonzero eigenvalues of \(CH\)
are those of the symmetric \(C^{1/2}HC^{1/2}\) by the
determinant identity for products. They have absolute value
at most the weighted operator norm of \(CH\). It follows
that the hard resolvent is positive and
\(\widetilde C\le (1-C_0\lambda t)^{-1}C\).
Absorb each power of \(Y_P/\sqrt\lambda\) into
\(C\exp(cY_P^2/\lambda)\). The remaining exponential of the
sum of squared norms is convex. Add an independent centered
Gaussian to attain the dominating covariance and apply
conditional Jensen to that exponential. For the dominating
multiple of \(C\), color the tiles in a bounded number of
residue classes so that the fixed enlargements within one
color are separated by more than the finite range. Hölder
over the colors and independence within colors give a product
of single-tile bounds for small enough \(c\), proving the
joint bound.
\end{proof}

The expansion controls a single dressed covariance link and its
coefficient variation at ordinary scales. For real coefficients it
also gives a joint Gaussian majorant. To estimate products with
activities, we express the quadratic tilt as a product of local
densities under the original Gaussian and keep track of the shifted
background separately.

\subsection{Local kinetic insertions in the Gaussian integral}
\label{sec:kinetic-density}

A local insertion of the quadratic density contains two fluctuation
fields. Its cost is small at either field radius, uniformly as the
quartic coupling tends to zero. This permits us to use the original
finite-range Gaussian integral and the regulator estimates of
Lemma~\ref{lem:cube-joint-reblocking} for products of insertions.

\begin{lemma}[A kinetic insertion costs \(O(\lambda)\) at both radii]
\label{lem:cube-kinetic-density}
Let \(d=4\), \(1/2<p<1\), and use the covariance, raw norm and
tile geometry of Lemma~\ref{lem:kinetic-dressed-slice}. Write
\(H=zJ_R+b^{-2}uI\), with \(J_R=R^{-2}P_\Lambda L_RP_\Lambda^*\),
and let \(|z|+|u|\le t_0\) on a fixed complex domain. There are
quadratic polynomials \(E_P\), supported on the fixed fine stencils
\(P^*\), such that
\begin{equation}
 \tfrac12\langle\zeta,H\zeta\rangle=\sum_P E_P(\zeta).
\end{equation}
Let \(v_P=e^{-E_P}-1\). In any product satisfying the hypotheses
of Lemma~\ref{lem:cube-joint-reblocking}, insert \(v_P\) on each
tile of a set \(A\) of distinct fine labels, without changing the
field factor on that label. At both radii the product bound of
that lemma holds, with its strict reserve and an additional factor
\((C_L\lambda (|z|+|u|))^{|A|}\). The derivative of the inserted
product in a unit \((z,u)\) direction has the bound obtained
by replacing this factor, for \(|A|\ge1\), by
\begin{equation}
 |A|(C_L\lambda)^{|A|}(|z|+|u|)^{|A|-1}.
\end{equation}
If the product of the original field factors depends on the
coefficients, its derivative is estimated separately by the
same product bound.
\end{lemma}

\begin{proof}
Extend \(\zeta\) by zero off \(\Lambda\). Assign to \(P\) the
contribution
\begin{equation}
 \frac{z}{4R^2}\sum_{x\in P,y\in\Lambda}
       D_R(x-y)(\zeta_x-\zeta_y)^2+
 \frac{z}{2R^2}\sum_{x\in P,y\notin\Lambda}
       D_R(x-y)\zeta_x^2+
 \frac{u}{2b^2}\sum_{x\in P}\zeta_x^2 .
\end{equation}
The ordered interior pairs and the two orientations of
the exterior pairs give the claimed exact sum. Write
\(M=F_{b,P^*}(\zeta)\), \(Y_P=M/\ell_b=bM\).
For pairs within \(O(R)\) of the physical boundary, both
fields are bounded by \(C M(R/b)^p\), including the zero
extension. A tile meets at most \(C b^3R\) sites in this
region. Away from this collar, with depth parameter
\(\rho=\min\{b,R+\dist(x,\Lambda^c)\}\), summing first
differences bounds the pair difference by
\(C M R b^{-p}\rho^{p-1}\).
The row sum of \(D_R\) is one. As in the energy estimate
of Lemma~\ref{lem:kinetic-terminal-energy}, layer summation
bounds the kinetic contribution by
\begin{equation}
 C|z|M^2\bigl[b^2(R/b)^{2p-1}
                 + b^{3-2p}\textstyle\sum_{r=1}^b(R+r)^{2p-2}
                 + b^2\bigr]\le C|z|Y_P^2 .
\end{equation}
The mass term is at most \(C|u|Y_P^2\), by the tile volume
bound. This proves both \(|E_P|\le C (|z|+|u|)Y_P^2\) and
\(|D_{z,u}E_P|\le C Y_P^2\). The integral formula for an
exponential difference gives the extra polynomial cost
\(C (|z|+|u|)Y_P^2\) times \(e^{C t_0 Y_P^2}\) at each mark.
Its first variation replaces the coefficient by \(C\).
These insertions have no background field derivatives.
In the proof of Lemma~\ref{lem:cube-joint-reblocking}, the
additional exponential in the fluctuation seminorm is absorbed
by its \(O(\lambda)\) Gaussian moment. The same coloring and
Hölder argument there retains a factor \(C_L\lambda\) for
each marked \(Y_P^2\). No background regulator is used by
the new insertions. Summing the mark that carries the coefficient
derivative proves the last estimate.
\end{proof}

Here the complex variation is in the density at a fixed field
argument. For real coefficients set \(S=(I+CH)^{-1}\).
The square-completion identity can be written as
\begin{equation}
 \E_C\big[e^{-\langle\phi+\zeta,H(\phi+\zeta)\rangle/2}
                     F(\phi+\zeta)\big]
 = e^{-\langle\phi,H^+\phi\rangle/2}
       \E_C\big[e^{-\sum_P E_P(\zeta)}F(S\phi+\zeta)\big].
\end{equation}
The expectation of the density at zero background gives the
determinant in the square-completion formula. Holding
\(\psi=S\phi\) fixed as an independent real field permits complex
variation of the coefficients in the density. The expansion then
uses finite-range links. Returning to the original field requires
an additional estimate on \(S\phi\), beginning with its restriction
to a single tile.

\begin{lemma}[Local tests of the transformed background]
\label{lem:cube-kinetic-background}
Under the assumptions of Lemma~\ref{lem:kinetic-dressed-slice},
fix a target fine tile \(P\). Let \(P^{(m)}\) contain the
\(O(mb)\)-neighbourhood of its tested stencils, clipped at the
physical boundary. With the same fine depth and scale weights
on these sets,
\begin{equation}
 F_{b,P^*}\big((S-I)\phi\big)
 \le \sum_{m\ge1}(C\lambda (|z|+|u|))^m
                      F_{b,P^{(m)}}(\phi).
\end{equation}
One coefficient derivative has the corresponding bound with
the summand prefactor replaced by
\(m(C\lambda)^m(|z|+|u|)^{m-1}\).
\end{lemma}

\begin{proof}
The proof of Lemma~\ref{lem:kinetic-dressed-slice} bounds
\(CH\) from the weighted supremum norm to \(F_b\), with
operator size \(C\lambda (|z|+|u|)\) and range \(O(b)\).
Apply the weighted supremum bound to the intermediate factors
of \((CH)^m\), and the derivative bound to the outermost operator
on the output stencil. All intervening field sites are in
\(P^{(m)}\). Summation of the Neumann series gives the
estimate, and differentiating the series marks one of its
\(m\) factors.
\end{proof}

This bound controls the field on a tile by the field in successive
neighbourhoods. A neighbourhood of radius \(mb\) can contain order
\(m^4\) labels, while the decay in the estimate is exponential in
\(m\). To pull the background substitution back into the polymer
norm one would have to assign its external field arguments to
connected sets of comparable total length, with the same fixed
number of tested derivatives.

\subsection{Connected Gaussian integration on a fixed kinetic disc}
\label{sec:kinetic-vacuum-graphs}

For independent background \(\psi\), the local densities of
Lemma~\ref{lem:cube-kinetic-density} give a connected expansion
under the original finite-range Gaussian. Tilting the Gaussian
also changes its vacuum normalization. Taking a joint cumulant
of the active factors cancels the vacuum contribution on every
nonempty coefficient and preserves the fixed order of field
derivatives.

\begin{lemma}[Normalized connected expansion for a quadratic tilt]
\label{lem:dressed-connected-expansion}
Write \(C\) for an ordinary finite-range covariance and
\(H=\sum_i H_i\) for localized symmetric quadratic matrices.
The label \(i\) of each matrix has a fixed-size stencil, and
the stencil graph of the fine labels has bounded degree and
includes the covariance range. Put
\begin{equation}
 Q_i(\zeta)=\tfrac12\langle\zeta,H_i\zeta\rangle,
 \qquad v_i(\zeta)=e^{-Q_i(\zeta)}-1.
\end{equation}
Assume that the quadratic bounds are uniform on a fixed complex
parameter disc and that on fixed stencils
\begin{equation}
 |Q_i(\zeta)|\le C t_0\sum_{P\sim i}Y_P^2,\qquad
 \sup_P\E_C e^{aY_P^2/\sigma}\le C_0,
 \qquad Y_P=F_{b,P^*}(\zeta)/\ell_b .
\end{equation}
Here \(t_0\) is fixed, \(a>0\), and \(\sigma\) is sufficiently
small. Let the active factors be a finite family of translated
activities or local insertions, each with a connected
cover containing its field stencil. Their actual occupied
labels are pairwise disjoint. Any label included only for the
cover keeps its assigned reference. Use active factors and fillers of the class of
Lemma~\ref{lem:cube-joint-reblocking}, with a strict regulator
reserve. A product of Gaussian expectations uses independent
copies of the original Gaussian, one in each expectation. The
multiplicity of polynomial insertions per label in each copy
remains bounded.

Then normalized expectation with covariance
\(\widetilde C=C^{1/2}(1+C^{1/2}HC^{1/2})^{-1}C^{1/2}\)
has an exact expansion into connected covers in which a
quadratic-tilt insertion costs \(C\sigma\) at either radius.
A component may use such insertions at unoccupied labels.
The sum of all terms involving a tilt insertion gains
\(C\sigma\), uniformly in volume, after any fixed additional
exponential cover weight, provided the active factors also
have a sufficiently small weighted anchored sum. Only terms
with at least one active factor occur in the normalized
connected expansion. The estimate is holomorphic in the
coefficients on the indicated disc and has the same first
variation bound on smaller fixed discs. All bounds hold
through the original fixed order of background field jets.
\end{lemma}

\begin{proof}
Each \(v_i\) satisfies
\begin{equation}
 |v_i(\zeta)|\le |Q_i(\zeta)|e^{|Q_i(\zeta)|}.
\end{equation}
For a set of distinct tilt insertions in one expectation the
local quadratic powers have bounded multiplicity at a given
stencil. Finite-range coloring and H\"older's inequality
therefore retain one factor \(C\sigma\) per insertion, as
in Lemma~\ref{lem:cube-joint-reblocking}. More explicitly,
color the stencil graph so vectors in any one color are
independent. Absorb each bounded power of \(Y_P/\sqrt\sigma\)
in a fixed fraction of its single-stencil exponential
moment, and absorb the added quadratic exponential in
another fixed fraction for small \(\sigma\). The constants
are exponential in the total number of read labels. Each
tilt factor is independent of the background field, so
differentiation in a tested background direction does not
increase its Gaussian polynomial loss.

For clarity, the reserve estimate works with the multiple
copies just described. On the actual occupied fine labels
use the occupied decay and quotient estimate of
Lemma~\ref{lem:cube-joint-reblocking}; the disjointness
of those labels permits each background square to be counted
only once. On each unoccupied row use the stable quartic on
that row, either at the background or with its assigned
fluctuation. A convex combination of real quartic exponents
gives the same lower bound with a convex combination of the
row displacements. In the recovery estimate, bound the sum
of squares of these displacements by the sum of
\(C_L(Y_P^{(\alpha)})^2\) over pairs consisting of a read
tile \(P\) and a Gaussian copy \(\alpha\). This follows
from the absolute filter row bound and bounded local
stencil size, without requiring one displacement on all
rows. Recovery therefore gives the same background
inequality and strict reserve. The occupied-field
triangle inequalities are unchanged, and the additional
fluctuation exponents and polynomial powers are paid by
coloring within each copy. Independence across copies
multiplies the moment constants. Their product is bounded
by a constant exponential in the total number of read
labels, counted with multiplicity over the copies.

Fix a nonempty family \(M\) of distinct active factors
\(f_A\), and introduce scalar variables \(t_A\). Their joint
cumulant in the normalized tilted expectation is
\begin{equation}
 \kappa_{\widetilde C}(f_A:A\in M)
   =\left[\prod_{A\in M}t_A\right]
       \log\E_C\left[
         \prod_i(1+v_i(\zeta))
         \prod_{A\in M}(1+t_A f_A)
       \right].
 \label{eq:kinetic-normalized-cumulant}
\end{equation}
For real parameters the finite-dimensional Gaussian
identity follows from completing the square, also if
\(C\) is singular. The identity extends holomorphically
within the tilted disc, using domination by the indicated
Gaussian moments. Taking a nonempty coefficient in
\eqref{eq:kinetic-normalized-cumulant} subtracts the
logarithm of the vacuum denominator automatically.
No logarithm of an activity or of a reference factor is
used. At fixed real background the logarithm at \(t=0\)
can be evaluated as a power series in \(t\); it suffices
to bound its squarefree coefficients in the field norm.

Here is the connected support expansion for that coefficient.
Call a tilt insertion or an active factor an elementary
object, with its chosen connected cover. Expand the two
products inside the expectation. Components of the resulting
objects are determined by touching covers. Distinct
components use independent Gaussian stencils. Consequently
the expectation is the hard-core gas of polymers \(G\),
where \(G\) is a nonempty connected set of elementary objects
and its weight is their joint expectation. The hard-core
condition is that two covers do not touch. One instance
of any elementary object is allowed within a polymer.
In the Mayer expansion of its logarithm, a term is a
connected family of incompatible polymers. An elementary
object can occur in distinct members of this list. Taking
the coefficient in \eqref{eq:kinetic-normalized-cumulant}
makes every active factor occur exactly once. Tilt
insertions can repeat in different members. These members
are products of expectations, so the repeated insertions
use different Gaussian copies. Within each expectation
the tilt multiplicity is bounded as before.

For completeness one can obtain absolute summability from
an anchored estimate alone. Give each elementary object
\(O\) a nonnegative majorant \(r(O)\) containing its
two-radius cost and the fixed exponential in read labels.
Write \(Y(O)\) for its cover. For any required exponential
cover weight, leave a fixed extra margin and suppose that
\begin{equation}
 \sup_i\sum_{O:i\in Y(O)}
      r(O)e^{a_1|Y(O)|}\le\delta_1
\end{equation}
with \(a_1\) sufficiently large and \(\delta_1\) small. The
sum of connected elementary configurations through a root
label, weighted by the product of the object majorants,
is at most \(2\delta_1\) after consuming one fixed part
of the size margin. Indeed choose a rooted spanning tree
of touching objects and charge each child to a site in
the closed neighbourhood of its parent cover. A parent
of cover size \(k\) has at most \(Dk\) possible attachment
sites. Allow ordered lists of children at each site. If
the rooted sum at the previous tree height is bounded by
\(2\delta_1\), their cost is at most
\((1-2\delta_1)^{-Dk}\), which is absorbed by the
reserved factor \(e^k\) for small \(\delta_1\).

For the Mayer coefficient of a list of \(q\) polymers,
the absolute value of the alternating sum over connected
subgraphs of its incompatibility graph is at most the
number of its spanning trees. To see this, fix an order
on the edges and group connected graphs by the greedy
spanning tree. For a given tree the graphs are obtained
by including an arbitrary subset of extra edges of larger
order than every edge on their replacement tree path.
The alternating sum on such an interval is zero or has
absolute value one. Root the tree at a polymer containing
the prescribed label. The factor \(1/q!\) in the Mayer
formula then becomes \(1/(q-1)!\) after paying for the
choice of root. Ordering the children of every rooted
tree allows this last factor to pay for the labels of
nonroot vertices. The same rooted-tree bound, with the
remaining size margin, now applies to the polymers as
objects. This proves absolute convergence uniformly in the
volume. The formal Mayer coefficients agree with the
Taylor coefficients at the vacuum by the usual expansion
of the logarithm into connected incompatibility graphs.
Equivalently that identity follows by partitioning every
graph into its connected components and exponentiating
the connected series. Absolute convergence after
slightly inflating the object majorants justifies
evaluation of all tilt variables at one in
\eqref{eq:kinetic-normalized-cumulant}.

To turn these connected coefficients into reblocked weights,
use the moment-cumulant identity
\begin{equation}
 \E_{\widetilde C}\prod_{A\in M} f_A
     =\sum_{\mathcal P\text{ a partition of }M}
              \prod_{E\in\mathcal P}
                \kappa_{\widetilde C}(f_A:A\in E).
\end{equation}
This follows by exponentiating the polynomial logarithm in
the same squarefree coefficient algebra. Give each cumulant
term its connected union cover and group touching covers.
Their products are bounded by one more application of
the rooted-tree bound with a further fixed size margin.
For each such group include the reference factors at the
unused labels in its cover, exactly once. The reference
on labels outside the groups stays unchanged. With this
assignment, the original multiplicative expansion on
the fine labels gives the outgoing circle product. Each cover
includes the labels that determine compatibility of its
activity or extraction choice. Thus separated groups
can be assembled independently in the reverse expansion.
The construction uses actual occupied labels for regulator
damping and uses covers for connected support. In particular
a cover joining two activities across vacant labels never
removes their fillers.

All terms at zero tilt are the original finite-range
expectations. In the difference there is at least one
tilt insertion. Inflate its activity from \(C\sigma\)
to a fixed small constant when making the anchored
estimate and keep one factor \(C\sigma\) outside the
sum. This proves the difference bound. The same operation
can keep any prescribed fixed number of costs belonging
to other kinds of objects when those costs are factored out
before making the anchored estimate. Uniform majorants on
the complex disc justify coefficientwise Cauchy
differentiation. Since its radius is independent of
\(\sigma\), the first coefficient variation of the tilt
difference still has factor \(C\sigma\). Gaussian
expectations and products of expectations require only
the tested background jets of their factors. This proves
the lemma.
\end{proof}

\paragraph{Effect on the mixed activity estimates.}
One can use this expansion in the chart in which the background
argument of a fluctuation factor is the independent variable
\(\psi\) of the Gaussian square completion. The connected estimate
for terms containing an activity or an extraction cover can then
be obtained by the size comparison of Lemma~\ref{lem:cube-connected-size}.
In the large-radius bound the activity cost is
\(\tau g^{-9/4}\), and the insertion cost is \(O(g^{1/4})\)
for each first-order centered ordinary difference. In the
small-radius bound these costs are \(\tau\) and \(O(g)\).
The output size weight \((A\varrho)^{P(U)}\) is paid by
multiplying the object weights by
\((A\varrho)^{|\overline{Y(O)}|+c}\), with a fixed \(c\)
per connecting edge. Indeed a contact across the covariance
range needs only a bounded coarse path. A root on such a
path has a participating cover within bounded coarse distance.
For a long input polymer \(X\), take \(L\) large so that
\(|\overline{Y(X)}|\le C+|X|/4\). The factor
\(\varrho^{-P(X)}\) then pays a fixed exponential reserve
in its fine size and the anchored count of its shapes.
For small marks, ordinary decorations and tilt objects
the larger weight costs a fixed constant, since their
covers have bounded coarse size. The constants may
depend on \(L,\varrho\), which are fixed before taking
the bare parameters small.

Thus, after retaining one activity or extraction factor,
the mixed norm is \(O(\tau)\), and the difference from
the corresponding un-tilted connected terms with the same
active factors and fillers is
\(O(\sigma\tau)\). The same statements hold in an
activity direction by substituting one varied mark or by
Cauchy's estimate on the activity ball. For pure terms
of perturbative order at least three retain three powers
of \(g\) in the small-radius estimate and three powers of
\(g^{1/4}\) at the large radius, where the output factor
\(g^{9/4}\) gives the same \(O(g^3)\) estimate. A second-order
local factor is counted with two powers at either radius.
In the tilt difference retain in addition one factor of
\(\sigma\), giving \(O(\sigma g^3)\). All statements use fixed
kinetic discs and involve the same coefficient weights as
the un-tilted estimates.
In particular, long covariance links by themselves do not
force an estimate on unboundedly high derivatives of an activity.

For pure terms of order one and two the one- and two-mark
contributions are the dressed mean and covariance on the
covers in \eqref{eq:kinetic-normalized-cumulant}. They must
be subtracted at those covers and simultaneously included
in the retained polynomial reference. The estimates above
bound their supported expansions with the corresponding
first- and second-order costs. A bound of order \(g^3\) for the
full pure remainder additionally requires this subtraction
and a uniformly bounded outgoing reference realizing it.
The substitution of the actual background
\(\psi=(I+CH)^{-1}\phi\) also requires a compatible support
reassignment, since an incoming factor on one cover is
then a function of background values outside that cover.

\paragraph{The long second-order terms in the dressed reference.}
The connected estimate at fixed kinetic coefficient requires a
second-order subtraction on connected supports of unbounded
length. This requirement already appears in the Gaussian heat of
the stable quartic in an independent background.

\begin{proposition}[Dressed quartic pairs and their supports]
\label{prop:kinetic-quartic-pairs}
Use the hypotheses of Lemma~\ref{lem:kinetic-dressed-slice}.
In the independent real background \(\psi\), let
\(V(\psi)=\sum_x\psi_x^4/4\) and use the dressed covariance
\(\widetilde C\). The coefficient of \(\psi_x^3\psi_y^3\),
for \(x\ne y\), in the second-order term of
\(\log\E_{\widetilde C}\exp[-gV(\psi+X)]\) is
\(g^2\widetilde C(x,y)\).
For any prescribed exponential size weight \(e^{A|U|}\),
the part using \(\widetilde C-C\) admits a connected-cover
allocation with anchored sum of normalized sextic coefficients
at most \(C_A g^2\lambda^2(|z|+|u|)\) at radius \(\ell_b\),
and first derivative at most \(C_Ag^2\lambda^2\), for
small \(\lambda(|z|+|u|)\). A fixed finite filter on the two
external vertices has the same upper bounds.

If the bulk \(C\) is nonzero and the kinetic perturbation
is nonzero, of finite range and of zero row sum, a reference
whose logarithm contains only fixed-diameter stencils cannot
have remainder \(O(g^3)\) for this example uniformly as
\(g\downarrow0\). A separately factored quadratic is also
permitted in this statement.
\end{proposition}

\begin{proof}
Taylor expansion for real \(g\ge0\) gives the first and second
terms \(-g\E V(\psi+X)\) and
\(g^2\Var(V(\psi+X))/2\). The expansion is valid in
zero-field jets through degree six on each fixed finite
cube: the third right derivative and the required field
jets are bounded on bounded real backgrounds by Gaussian
polynomial moments, using \(V\ge0\). In the variance the
product of the two linear fluctuation terms at \(x,y\) is
\(\psi_x^3\psi_y^3\widetilde C(x,y)\); no term with more
fluctuation fields at either vertex has this monomial.
The two orientations cancel the factor \(1/2\).

For each \(m\ge1\) use the term \(C^{(m)}\) of
Lemma~\ref{lem:kinetic-dressed-slice}. For each endpoint
pair record the endpoint tiles and a shortest grid path
between them, contained in the cube. Include the filter
stencils by a fixed enlargement at the two ends, if
there is a filter. The path has at most \(C(m+1)\) labels.
A tested sextic coefficient has a factor at most
\(C\ell_b^6=C b^{-6}\). For each pair of tiles there are
at most \(C b^8\) site pairs, and the kernel factor
is bounded by \(C\lambda b^{-2}(C\lambda (|z|+|u|))^m\).
Thus the normalized contribution for one tile pair
is bounded by
\(C g^2\lambda(C\lambda(|z|+|u|))^m\).
For a prescribed root tile the anchored count of
endpoint tile pairs is at most \(C(m+1)^8\), since
both endpoints are within \(C(m+1)\) tiles of that root.
The size weight contributes at most \(e^{C A(m+1)}\).
The geometric sum gives the claimed bound. Marking one
Neumann factor gives the derivative bound. The finite
filter consumes only its bounded absolute row norms and
bounded enlargement of the covers.

For the last claim take the unfiltered example in a
large periodic cube. The Fourier argument following
Lemma~\ref{lem:kinetic-exact-compensation} shows
that the dressed covariance has entries at arbitrarily
large separations in infinite volume. The finite torus
coefficients approach these by Fourier Riemann sums
(or by periodizing the absolutely summable Neumann series).
Fix a torus large enough to have a nonzero entry at sites
farther apart than the reference stencils. A sum of
logarithms of factors on those stencils has zero
\(\psi_x^3\psi_y^3\) coefficient, as has any quadratic.
On this fixed torus a circle-product error of order
\(g^3\) in the mixed norm gives order \(g^3\) in its
small-radius jets near zero, also after taking the local
logarithm. Here the reference has normalized vacuum
bounded away from zero and bounded small-field jets,
so the finite-jet inverse is bounded for this fixed torus.
This would force
\(g^2\widetilde C(x,y)=O(g^3)\), contradicting the
chosen nonzero entry.
\end{proof}

This computation bounds the long sextic coefficients in the
size norm in which they would have to be retained. A local
reference that includes only a finite collection of momentum
jets has no such coefficients. Treating the long second-order
terms as part of the reference still requires their connected
transport together with the nonlocal background shift.

\paragraph{The bubble under a fixed change of kinetic coefficient.}
For a fixed kinetic coefficient, the filtered bubble changes by
a uniformly small relative amount. The comparison uses the same
coefficient at both ends of a scale increment. This qualification
matters when the kinetic coefficient itself varies along a
trajectory.

\begin{lemma}[Bulk bubble comparison under kinetic dressing]
\label{lem:kinetic-dressed-bubble}
Let \(d=4\), \(q(k)=R^{-2}(1-\widehat D_R(k))\), and
\(W_j\) be the unfiltered cumulative bulk covariance of the
finite-propagation slices, excluding the terminal covariance.
For small real \(z\), define
\begin{equation}
  W_j^{[z]}(k)=\frac{\widehat W_j(k)}{1+zq(k)\widehat W_j(k)} .
\end{equation}
Let \(\mathfrak b_j^{[z]}\) be the filtered bubble increment
with \(\widehat W_j,\widehat W_{j+1}\) replaced by
\(W_j^{[z]},W_{j+1}^{[z]}\) in its Fourier formula.
Then, with the undressed bubble \(\mathfrak b_j\),
\begin{equation}
 \mathfrak b_j^{[z]}=(1+O(\lambda |z|))\mathfrak b_j,
 \qquad
 |\partial_z\mathfrak b_j^{[z]}|\le C\lambda\mathfrak b_j .
\end{equation}
The constants are uniform in the ordinary and merged scales.
This comparison uses the same \(z\) in both terms of an increment.
\end{lemma}

\begin{proof}
The local Schur and polynomial bounds of Section~\ref{sec:bubble}
give for each individual bulk slice at side \(b\)
\begin{equation}
  0\le\widehat C_b(k)\le
       C_m\lambda b^2(1+b^2q(k))^{-m}
       \qquad(m>1).
\end{equation}
Indeed in the local split the auxiliary bulk precision has a
positive lower bound of order \(r+R^2 q(k)\) after shifting by
\(\kappa r\), where \(r=(R/b)^2\). This follows from the Schur
complement with its upper block bounded and uniformly positive.
The observation has uniformly bounded norm. Applying the
polynomial majorant, then multiplying by \(a_R=\lambda R^2\),
gives the estimate. The constants can include the bounded
number of late scales with a larger displacement relative to
\(r\). For \(q(k)>0\), split the sum over geometric scales at
\(b^2q(k)=1\). The two parts of the sum for
\(q(k)\widehat W_j(k)\) are bounded by geometric sums with terms
\(C\lambda b^2q(k)\) and \(C\lambda(b^2q(k))^{1-m}\).
Thus \(0\le q(k)\widehat W_j(k)\le C\lambda\) for all \(j,k\),
including the merged scales. At \(q=0\) the same bound holds
identically.

Write \(v=\widehat W_j(k)\), \(u=\widehat W_{j+1}(k)\). Both
are nonnegative and \(v\le u\). Pointwise,
\begin{equation}
 \left(\frac{u}{1+zq u}\right)^2
     -\left(\frac{v}{1+zq v}\right)^2
 =\int_v^u\frac{2t}{(1+zq t)^3}\,dt.
\end{equation}
The integrand differs from \(2t\) by at most
\(C\lambda |z|\,2t\), and its \(z\)-derivative has magnitude
at most \(C\lambda\,2t\). Multiplying by the nonnegative
filter weight \(|\widehat{\overline T}(k)|^4\) and integrating
proves both estimates.
\end{proof}

\paragraph{Exact first-order quartic packet in the bulk.}
The change in covariance also alters the filtered square produced
by the first heat of the quartic. We identify the entire
first-order polynomial before projecting onto its local mass and
kinetic parts.

\begin{lemma}[Transport of the centered quartic at fixed kinetic coefficient]
\label{lem:kinetic-wick}
Work in a periodic bulk calculation (with period larger than the
original covariance ranges) in dimension four. Let \(W_j\) be the
cumulative covariance, let \(J_R\) have symbol \(q\) as in
Lemma~\ref{lem:kinetic-dressed-bubble}, and fix a sufficiently
small real \(z\). Put
\begin{equation}
 H_*=zJ_R,\qquad
 T_j=T(I+W_jH_*)^{-1},\qquad
 v_j=(T\Sigma_j T^*)(0),\qquad
 \Sigma_j=W_j(I+H_*W_j)^{-1},
\end{equation}
where \(T\) is the normalized even bulk filter. Fix a horizon \(M\)
ending before the terminal slice and put \(v_*=v_M\).
For \(\mathcal H_4(t;a)=t^4-6at^2+3a^2\) one has
\begin{equation}
 \frac{P_{W_j}\left[e^{-\langle\phi,H_*\phi\rangle/2}
               \mathcal H_4((T\phi)_x;v_*)\right]}
      {P_{W_j}\left[e^{-\langle\phi,H_*\phi\rangle/2}\right]}
     =\mathcal H_4((T_j\phi)_x;v_*-v_j).
\end{equation}
After the merged step,
\(0\le v_*-v_j\le C\lambda b_j^{-2}\); the merged input
allows the bound \(C\lambda R^{-2}\). Thus the first-order
tadpole can be transported with its actual filtered square
under a fixed kinetic coefficient, with no additional
first-order source if this entire polynomial is retained.
\end{lemma}

\begin{proof}
Apply \eqref{eq:kinetic-heat} with \(W_j\) in place of \(C\). Its use is justified by
the bound
\(\sup_k |zq(k)\widehat W_j(k)|\le C\lambda|z|\) from
Lemma~\ref{lem:kinetic-dressed-bubble}. Square completion sends the observed mean
to \((T_j\phi)_x\) and gives added Gaussian variance \(v_j\).
The coefficient of \(t^4\) in
\(\exp(tu-at^2/2)\), multiplied by \(4!\), is
\(\mathcal H_4(u;a)\). Gaussian integration of this generating
function changes \(a\) to \(a-v_j\), proving the identity.
For \(k\) fixed, the derivative of
\(w\mapsto w/(1+zq(k)w)\) is positive and bounded by a uniform
constant on the interval of cumulative covariances. The filtered
diagonal of the difference is therefore bounded by the sum of
the filtered slice diagonals between \(j\) and \(M\).
The ordinary scale bound sums to \(C\lambda b_j^{-2}\).
The merged bound follows by summing the individual ultraviolet
scales from \(R\). The semigroup identity
\(P_C P_{W_j}=P_{W_j+C}\) gives the transport assertion.
\end{proof}

The formula identifies the whole first-order reference before
localizing any quadratic jet. Its field filter depends on the
cumulative covariance and loses finite range. Using this
formula in the nonlinear step would require a decomposition of
that reference into local factors with controlled activities
and coefficient variations.

\subsection{Terminal absorption of a marginal quadratic form in dimension four}
\label{sec:kinetic-terminal-absorption}

A terminal form comparable to a small multiple of the Dirichlet
kinetic energy need not be small at the quartic radius. It can
instead be included in the Gaussian law, provided the scalar
mass is centered in the modified law. In this subsection use
the normalized Dirichlet form
\begin{equation}
 J_R=R^{-2}P_\Lambda L_RP_\Lambda^*,
 \qquad |f|_{J_R}^2=\langle f,J_Rf\rangle.
\end{equation}
Multiplying \(J_R\) by a bounded positive constant gives bulk
quadratic Taylor coefficient \(|k|^2\). In particular this
representative can be used for the local kinetic jet.

To integrate a small quadratic perturbation into the terminal
Gaussian law, we need its action on the conditioned covariance,
on the principal profile and on the integrated uniform source. The
following energy estimates give all three controls in the same
Dirichlet form.

\begin{lemma}[Terminal energy and hard covariance]
\label{lem:kinetic-terminal-energy}
Take \(d=4\), \(1/2<p<1\), and \(n\ge C R\). For the
terminal covariance of Section~\ref{sec:ext-cov}, let
\(C_h=\mathcal S_\ell(C_T)\) be the covariance conditioned on
\(\ell^{\mathsf T}\phi=0\). Then
\begin{equation}
 C_h\le C\lambda J_R^{-1},\qquad
 |f|_{J_R}\le C n\|f\|_{F_{n,p}} .
\end{equation}
In particular \(|u_T|_{J_R}\le C n^{-1}\) and
\(|C_<\one|_{J_R}\le C_R n^3\).
\end{lemma}

\begin{proof}
Put \(L=P_\Lambda L_RP_\Lambda^*\). The
one-dimensional zero-extension walk with uniform positive and
negative steps of size at most \(R\) has symbol bounded by
\(C\min\{R^2|k|^2,1\}\) when considered in one coordinate.
The Fourier lower bound for \(L_R\) therefore dominates the
normalized form of this walk on the cube. For each
\(\lfloor R/4\rfloor\le j\le\lfloor R/2\rfloor\) follow
each step-\(j\) chain in that coordinate to its first exterior
point. Telescoping and Cauchy--Schwarz on a chain with
\(O(n/j)\) sites bound its squared \(\ell^2\) norm by
\(C(n/j)^2\) times its squared differences. Averaging over
the chosen steps proves
\(J_R\ge c n^{-2}I\).

The kernel of the split \(A\) in
Lemma~\ref{lem:local-split} is nonnegative, is supported
in the spread-out cube, and satisfies \(A(x)\le C D_R(x)\)
for \(x\ne0\). If \(p_A=\sum_x A(x)\), then
\(G=p_A I-A_\Lambda\) is the corresponding zero-extension
jump form, so \(0\le G\le C L\). Since both forms have
bounded operator norm,
\begin{equation}
 \|L^{1/2}A_\Lambda f\|_2
 \le p_A\|L^{1/2}f\|_2+
 C\langle f,G f\rangle^{1/2}
 \le C\|L^{1/2}f\|_2 .
\end{equation}
By duality \(A_\Lambda L^{-1}A_\Lambda\le C L^{-1}\).
The Schur identity and \(B_\Lambda-A_\Lambda=\zeta L\)
give
\begin{equation}
 V_1=A_\Lambda+A_\Lambda(B_\Lambda-A_\Lambda)^{-1}A_\Lambda
 \le C L^{-1}.
\end{equation}
For positive covariances, conditioning is monotone in the
quadratic-form order, by the formula
\(
 \langle f,\mathcal S_\ell(C)f\rangle=
 \inf_{t\in\mathbb R}\langle f+t\ell,C(f+t\ell)\rangle
\).
Use \(C_T\le a_R V_c\),
\(\mathcal S_\ell(V_c)\le C V_1\), and \(a_R=\lambda R^2\)
to prove the asserted hard bound.

For the energy estimate set \(M=\|f\|_{F_{n,p}}\).
If \(x\) is within distance \(C R\) of the exterior and
\(|y-x|_\infty\le R\), the difference of the zero-extended
values is at most \(C M(R/n)^p\).
There are at most \(C n^3 R\) points in this collar, and
the jump row sums are one. Division by \(R^2\) gives an
energy at most \(C M^2 n^2 (R/n)^{2p-1}\).
Elsewhere, with depth \(r\ge C R\), telescoping along
nearest-neighbor paths gives a difference at most
\(C R M n^{-p} r^{p-1}\). At each depth there are at most
\(C n^3\) points, so the energy outside the collar is at most
\begin{equation}
 C M^2 n^{3-2p}\sum_{r=1}^n r^{2p-2}\le C M^2 n^2 .
\end{equation}
The two stated applications now use the profile and
integrated-source bounds of Section~\ref{sec:ext-cov}.
\end{proof}

Absorbing a kinetic form changes the variance in the principal
direction. A scalar mass shift restores this variance. The same
shift also changes the linear source that must be used to recover
the initial magnetic moment.

\begin{proposition}[Centering and source restoration with a terminal kinetic form]
\label{prop:kinetic-terminal-absorption}
For real symmetric \(Q\) satisfying \(-qJ_R\le Q\le qJ_R\),
let \(Q_m=Q+m n^{-2}I\). For sufficiently small \(q\), there
is a unique \(m=m_c(Q)\) in a fixed neighbourhood of zero for
which multiplication by
\(e^{-\langle\phi,Q_m\phi\rangle/2}\) preserves the scalar
Gaussian variance \(\sigma_T^2\) in the decomposition along
\(a=\ell^{\mathsf T}\phi\). Moreover, \(|m_c(Q)|\le Cq\),
and the resulting hard Gaussian and soft profile satisfy the
bounds of Assumption~\ref{ass:terminal} uniformly for \(Q\)
in this form ball. The scalar recentering is real analytic.

Suppose the exact terminal interaction factors as
\(F_T=e^{-\langle\phi,Q_{m_c}\phi\rangle/2}\overline F_T\),
where \(\overline F_T\) has the interaction and error
bounds of Assumption~\ref{ass:terminal}. Then both the
terminal and initial auxiliary magnetic second moments have
order \(V^{3/2}g^{-1/2}\).
\end{proposition}

\begin{proof}
Write \(u=u_T\), \(J=J_R\), and \(M=Q_m\).
Poincar\'e gives
\(-q'J\le M\le q'J\) with \(q'=C(q+|m|)\).
In particular
\(\|C_h^{1/2} M C_h^{1/2}\|\le C\lambda q'\).
Apply \eqref{eq:kinetic-heat} with covariance \(C_h\)
to the interaction at each value of \(a\). The hard
covariance and profile become
\begin{equation}
 C_{h,M}=(I+C_h M)^{-1}C_h,\qquad
 u_M=(I+C_h M)^{-1}u=u-C_{h,M} M u .
\end{equation}
The added soft precision is
\(r(M)=\langle u,M(I+C_h M)^{-1}u\rangle\).
The tilt determinant is independent of \(a\), and
\(\ell^{\mathsf T}u_M=1\) because \(C_h\ell=0\).
For \(q'\) small, \(C_{h,M}\le 2 C_h\); consequently its
Gaussian exponential norm moment follows by adding an
independent centered Gaussian to attain covariance \(2C_h\)
and applying conditional Jensen to
\(f\mapsto\exp(c n^2\|f\|_{F_{n,p}}^2)\).

For each weighted unit test stencil \(D\) entering
\(F_{n,p}\), the same hard moment gives
\(\|C_{h,M}^{1/2}D^*\|_2\le C n^{-1}\).
Using the form norm and Lemma~\ref{lem:kinetic-terminal-energy},
\begin{equation}
 |D(C_{h,M}Mu)|\le
 \|C_{h,M}^{1/2}D^*\|_2
 \|C_{h,M}^{1/2}Mu\|_2
 \le C n^{-1}\sqrt{\lambda}\,q' |u|_J
 \le C q'n^{-2}.
\end{equation}
Thus \(\|u_M-u\|_{F_{n,p}}\le C q'n^{-2}\), and the
magnetic overlap remains of order \(n^2\). These inequalities
also give \(|u_M|_J\le C n^{-1}\). In particular
\(|r(Q)|\le C q n^{-2}\). Differentiating the resolvent
identity, for any symmetric perturbation \(D_M\), gives
\begin{equation}
 D r(M)[D_M]=\langle u_M,D_M u_M\rangle .
\end{equation}
Hence \(\partial_m r(Q_m)=n^{-2}\|u_M\|_2^2\asymp n^{-2}\).
The upper bound follows from \(\|u_M\|_{F_{n,p}}\le C n^{-2}\)
and the lower bound from \(\ell^{\mathsf T}u_M=1\).
The intermediate value theorem yields the unique zero and
the stated bound on \(m_c\). Analyticity follows from the
invertible hard resolvent and the implicit function theorem.
In particular, for a form direction \(D_Q\) of relative
\(J\)-bound \(q_D\), implicit differentiation gives
\begin{equation}
 Dm_c(Q)[D_Q]
 =-n^2\frac{\langle u_M,D_Q u_M\rangle}{\|u_M\|_2^2},
 \qquad |Dm_c(Q)[D_Q]|\le C q_D.
\end{equation}
The profile variation is
\(-C_{h,M}(D_Q+n^{-2}Dm_c(Q)[D_Q]I)u_M\), whose
\(F_{n,p}\) norm is at most \(C q_D n^{-2}\) by the same
stencil estimate.

At this zero the full tilted law is Gaussian with
decomposition \(a u_M+Y_M\) and the original soft variance.
The proof of Proposition~\ref{prop:terminal} applied with
\(\overline F_T\) now gives the terminal second moment. To
restore the earlier source put \(v=C_<\one\). Gaussian
translation in the earlier slices gives, including the shifted
quadratic factor, the exact generating function up to a positive
constant independent of \(t\):
\begin{equation}
 e^{t^2(\one^{\mathsf T}C_<\one-\langle v,Mv\rangle)/2}
 \E_{M}\!\left[e^{t\langle\one-Mv,\phi\rangle}
                 \overline F_T(\phi+t v)\right].
\end{equation}
Here \(\E_M\) is Gaussian expectation in the tilted law.
Lemma~\ref{lem:kinetic-terminal-energy} gives
\begin{equation}
 \begin{aligned}
 |\langle v,M v\rangle|&\le C_R q' n^6,\
 |\langle Mv,u_M\rangle|&\le C_R q' n^2,\
 \Var(\langle Mv,Y_M\rangle)&\le C_R(q')^2n^6 .
 \end{aligned}
\end{equation}
Choose the form ball at fixed \(R\) so that the source overlap
\(|\langle\one-Mv,u_M\rangle|\) remains bounded below by
\(c n^2\). The hard part of the source has standard deviation
\(O_R(n^3)\). In the terminal envelope argument all insertions
of this hard linear Gaussian are bounded by Cauchy--Schwarz
with the same exponential moment reserve. Thus the terminal
second moment of the corrected linear form is
\(\asymp V^{3/2}g^{-1/2}\). Differentiating twice at zero,
the terms with a derivative of \(\overline F_T\) are controlled
exactly by its directional bounds, with
\(\|v\|_{F_{n,p}}/h\le C_R V^{3/4}g^{1/4}\). As in
Proposition~\ref{prop:terminal}, these terms and the
two constant prefactor contributions are \(O_R(V^{3/2})\).
This proves the initial bound.
\end{proof}

The mass shift is of order \(q\) and can exceed the terminal thermal
window. Applying this comparison to a physical trajectory would
require an outgoing chart that isolates the form and controls
\(\overline F_T\) and its first variation near the recentered mass.
The four-dimensional proof uses backward kinetic shooting on the
shrinking domain of Section~\ref{sec:backward-kinetic}.

% \section{Direct covariance comparison at the earlier cut}
% \label{app:direct-comparison}

% Use the notation of Section~\ref{sec:d4-prescribed}.

% One can also quantify what the same released slice gives
% without changing the microscopic field. Keep \(k=J-1\) and
% put \(r=A_cw_0\). We first record the sharper microscopic bound
% \begin{equation}
%  w_0^{\mathsf T}A_c w_0\le C a_R .
% \end{equation}
% In the notation of Lemma~\ref{lem:aligned-curvature-margin}
% introduce the row operator on the two auxiliary blocks
% \begin{equation}
%  U_1=(A^{-1/2}\mathcal L_\Lambda^*,\ c'A^{1/2}B^{-1}).
% \end{equation}
% It has bounded norm because
% \(\mathcal L_\Lambda^*\mathcal L_\Lambda=A\) and \(B^{-1}\)
% is bounded. Block multiplication gives
% \(U_1\mathcal Q=a_R A^{1/2}G^{-1}\mathcal O\), including
% cancellation of the entire second column. Hence
% \(A^{1/2}w_0=U_1\mathcal Q f(\mathcal Q)y\).
% The inequality follows from \(0\le s f(s)\le1\).

% Take \(\underline q_J=\frac12\min_{[0,s_{\max}]}q_{J-1}(s)\), where
% \([0,s_{\max}]\) is the fixed auxiliary spectral interval.
% Appendix~\ref{app:gaussian-scalar} proves that each fixed slice
% polynomial has a strictly positive minimum on its spectral interval.
% In particular \(W_k\succeq \underline q_J A_c\), so
% \(r^{\mathsf T}W_k^{-1}r\le C a_R/\underline q_J\).

% For the direct comparison the column is
% \(v_{0,\Sigma}=tT\ell+\varepsilon r\) and
% \(V_{0,\Sigma}=t\sigma^2+O(|\varepsilon|\lambda n^2)\).
% Cauchy--Schwarz in the inverse quadratic form and
% Lemma~\ref{lem:aligned-curvature-margin} give
% \begin{equation}
%  v_{0,\Sigma}^{\mathsf T} C_k^{-1}v_{0,\Sigma}
%  \le t(\sigma^2-c\lambda n^2)
%       +C|\varepsilon|\sigma\sqrt{a_R/\underline q_J}
%       +C\varepsilon^2 a_R/\underline q_J .
% \end{equation}
% Thus the direct subtraction in \eqref{eq:original-curvature-hard}
% has variance margin at least \(c'\lambda n^2\) whenever
% \begin{equation}
%  |\varepsilon|\le c_1\min\{1,\lambda(n/R)\sqrt{x_k \underline q_J}\},
% \end{equation}
% with \(c_1\) sufficiently small. This includes the unperturbed
% point. The arbitrary-power upper bound on \(q_{J-1}\) at any
% fixed positive spectral argument shows that \(\underline q_J\) decays
% faster than any power of \(R/n\). The estimate therefore gives
% a shrinking interval for this direct comparison. The full
% matching rectangle is covered by \eqref{eq:prescribed-positive}
% in the transformed field.

\section*{AI Statement}

This document is generated with a collaboration with ChatGPT-Astra 6.0.

\begin{thebibliography}{99}
\bibitem{obs-extension}
J.~Park, \emph{Gaussian and non-Gaussian finite-size scalings of the
spread-out Ising model at the free-boundary pseudocritical point},
GitHub preprint,
\href{https://github.com/jiwoonGit/phi4/blob/main/FSS/FSS_FBC-obs-extension.pdf}{DOI:10.5281/zenodo.24167641} (2026).

\bibitem{Aizenman1982}
M.~Aizenman, \emph{Geometric analysis of $\varphi^4$ fields and Ising
models. Parts I and II}, Comm. Math. Phys. \textbf{86} (1982), 1--48.
\href{https://doi.org/10.1007/BF01205659}{doi:10.1007/BF01205659}.

\bibitem{AizenmanSpanning1997}
M.~Aizenman, \emph{On the number of incipient spanning clusters},
Nucl. Phys. B \textbf{485} (1997), 551--582.
\href{https://doi.org/10.1016/S0550-3213(96)00626-8}{doi:10.1016/S0550-3213(96)00626-8}.

\bibitem{AizenmanDuminilCopin2021}
M.~Aizenman and H.~Duminil-Copin, \emph{Marginal triviality of the
scaling limits of critical 4D Ising and $\phi_4^4$ models},
Ann. of Math. \textbf{194} (2021), 163--235.
\href{https://doi.org/10.4007/annals.2021.194.1.3}{doi:10.4007/annals.2021.194.1.3}.

\bibitem{AizenmanFernandez1988}
M.~Aizenman and R.~Fern\'andez, \emph{Critical exponents for long-range
interactions}, Lett. Math. Phys. \textbf{16} (1988), 39--49.
\href{https://doi.org/10.1007/BF00398169}{doi:10.1007/BF00398169}.

\bibitem{BauerschmidtBrydgesSlade2014}
R.~Bauerschmidt, D.~C.~Brydges, and G.~Slade,
\emph{Scaling limits and critical behaviour of the 4-dimensional
$n$-component $|\varphi|^4$ spin model}, J. Stat. Phys.
\textbf{157} (2014), 692--742.
\href{https://doi.org/10.1007/s10955-014-1060-5}{doi:10.1007/s10955-014-1060-5}.

\bibitem{BlancRenaudieNachmias}
A.~Blanc-Renaudie and A.~Nachmias,
\emph{Critical percolation on the discrete torus in high dimensions},
\href{https://arxiv.org/abs/2512.19672}{arXiv:2512.19672} (2025).

\bibitem{BorgsEtAlI}
C.~Borgs, J.~T.~Chayes, R.~van der Hofstad, G.~Slade, and J.~Spencer,
\emph{Random subgraphs of finite graphs. I. The scaling window under the
triangle condition}, Random Structures Algorithms \textbf{27} (2005),
137--184.
\href{https://doi.org/10.1002/rsa.20051}{doi:10.1002/rsa.20051}.

\bibitem{BorgsEtAlII}
C.~Borgs, J.~T.~Chayes, R.~van der Hofstad, G.~Slade, and J.~Spencer,
\emph{Random subgraphs of finite graphs. II. The lace expansion and the
triangle condition}, Ann. Probab. \textbf{33} (2005), 1886--1944.
\href{https://doi.org/10.1214/009117905000000260}{doi:10.1214/009117905000000260}.

\bibitem{BrezinZinnJustin1985}
E.~Br\'ezin and J.~Zinn-Justin, \emph{Finite size effects in phase transitions},
Nucl. Phys. B \textbf{257} (1985), 867--893.
\href{https://doi.org/10.1016/0550-3213(85)90379-7}{doi:10.1016/0550-3213(85)90379-7}.

\bibitem{BrydgesSpencer1985}
D.~C.~Brydges and T.~Spencer, \emph{Self-avoiding walk in 5 or more
dimensions}, Comm. Math. Phys. \textbf{97} (1985), 125--148.
\href{https://doi.org/10.1007/BF01206182}{doi:10.1007/BF01206182}.

\bibitem{CamiaJiangNewman2021}
F.~Camia, J.~Jiang, and C.~M.~Newman,
\emph{The effect of free boundary conditions on the Ising model in high
dimensions}, Probab. Theory Related Fields \textbf{181} (2021), 311--328.
\href{https://doi.org/10.1007/s00440-021-01041-9}{doi:10.1007/s00440-021-01041-9}.

\bibitem{ChatterjeeHanson2020}
S.~Chatterjee and J.~Hanson, \emph{Restricted percolation critical
exponents in high dimensions},
Comm. Pure Appl. Math. \textbf{73} (2020), 2370--2429.
\href{https://doi.org/10.1002/cpa.21938}{doi:10.1002/cpa.21938}.

\bibitem{Dyson1969}
F.~J.~Dyson, \emph{Existence of a phase-transition in a one-dimensional
Ising ferromagnet}, Comm. Math. Phys. \textbf{12} (1969), 91--107.
\href{https://doi.org/10.1007/BF01645907}{doi:10.1007/BF01645907}.

\bibitem{FisherBarber1972}
M.~E.~Fisher and M.~N.~Barber, \emph{Scaling theory for finite-size effects
in the critical region}, Phys. Rev. Lett. \textbf{28} (1972), 1516--1519.
\href{https://doi.org/10.1103/PhysRevLett.28.1516}{doi:10.1103/PhysRevLett.28.1516}.

\bibitem{FisherMaNickel1972}
M.~E.~Fisher, S.-k.~Ma, and B.~G.~Nickel, \emph{Critical exponents for
long-range interactions}, Phys. Rev. Lett. \textbf{29} (1972), 917--920.
\href{https://doi.org/10.1103/PhysRevLett.29.917}{doi:10.1103/PhysRevLett.29.917}.

\bibitem{flores2016}
E.~Flores-Sola, B.~Berche, R.~Kenna and M.~Weigel,
\href{https://doi.org/10.1103/PhysRevLett.116.115701}
{\emph{Role of Fourier modes in finite-size scaling above the upper
critical dimension}}, Phys. Rev. Lett. \textbf{116} (2016), 115701.

\bibitem{Frohlich1982}
J.~Fr\"ohlich, \emph{On the triviality of $\lambda\phi_d^4$ theories and
the approach to the critical point in $d\geq4$ dimensions},
Nucl. Phys. B \textbf{200} (1982), 281--296.
\href{https://doi.org/10.1016/0550-3213(82)90088-8}{doi:10.1016/0550-3213(82)90088-8}.

\bibitem{GrimmEtAl2017}
J.~Grimm, E.~M.~El\c{c}i, Z.~Zhou, T.~M.~Garoni, and Y.~Deng,
\emph{Geometric explanation of anomalous finite-size scaling in high
dimensions}, Phys. Rev. Lett. \textbf{118} (2017), 115701.
\href{https://doi.org/10.1103/PhysRevLett.118.115701}{doi:10.1103/PhysRevLett.118.115701}.

\bibitem{HaraHofstadSlade2003}
T.~Hara, R.~van der Hofstad, and G.~Slade,
\emph{Critical two-point functions and the lace expansion for spread-out
high-dimensional percolation and related models}, Ann. Probab.
\textbf{31} (2003), 349--408.
\href{https://doi.org/10.1214/aop/1046294314}{doi:10.1214/aop/1046294314}.

\bibitem{HaraSladePercolation}
T.~Hara and G.~Slade, \emph{Mean-field critical behaviour for percolation
in high dimensions}, Comm. Math. Phys. \textbf{128} (1990), 333--391.
\href{https://doi.org/10.1007/BF02108785}{doi:10.1007/BF02108785}.

\bibitem{HaraSlade1991}
T.~Hara and G.~Slade, \emph{Critical behaviour of self-avoiding walk in
five or more dimensions}, Bull. Amer. Math. Soc. (N.S.) \textbf{25}
(1991), 417--423.
\href{https://doi.org/10.1090/S0273-0979-1991-16085-4}{doi:10.1090/S0273-0979-1991-16085-4}.

\bibitem{HutchcroftMichtaSlade2023}
T.~Hutchcroft, E.~Michta, and G.~Slade,
\emph{High-dimensional near-critical percolation and the torus plateau},
Ann. Probab. \textbf{51} (2023), 580--625.
\href{https://doi.org/10.1214/22-AOP1608}{doi:10.1214/22-AOP1608}.

\bibitem{LiuSAWPlateau}
Y.~Liu, \emph{A general approach to massive upper bound for two-point
function with application to self-avoiding walk torus plateau},
Electron. J. Probab. \textbf{30} (2025), article 110.
\href{https://doi.org/10.1214/25-EJP1371}{doi:10.1214/25-EJP1371}.

\bibitem{LiuParkSlade2026}
Y.~Liu, J.~Park, and G.~Slade,
\emph{Universal finite-size scaling in high-dimensional critical phenomena},
Phys. Rev. E \textbf{113} (2026), 024143.
\href{https://doi.org/10.1103/xq7j-zll9}{doi:10.1103/xq7j-zll9}.

\bibitem{LuijtenBlote1997}
E.~Luijten and H.~W.~J.~Bl\"ote, \emph{Classical critical behavior of spin
models with long-range interactions}, Phys. Rev. B \textbf{56} (1997),
8945--8958.
\href{https://doi.org/10.1103/PhysRevB.56.8945}{doi:10.1103/PhysRevB.56.8945}.

\bibitem{MichtaTorWalk}
E.~Michta, \emph{The scaling limit of the weakly self-avoiding walk on a
high-dimensional torus}, Electron. Commun. Probab. \textbf{28} (2023), 1--13.
\href{https://doi.org/10.1214/23-ECP531}{doi:10.1214/23-ECP531}.

\bibitem{HierarchicalFBC}
E.~Michta, J.~Park, and G.~Slade, \emph{Boundary conditions and universal
finite-size scaling for the hierarchical $|\varphi|^4$ model in dimensions
4 and higher}, Comm. Pure Appl. Math. \textbf{78} (2025), 2001--2118.
\href{https://doi.org/10.1002/cpa.22256}{doi:10.1002/cpa.22256}.

\bibitem{Panis2026}
R.~Panis, \emph{Triviality of the scaling limits of critical Ising
and $\varphi^4$ models with effective dimension at least four},
Ann. Probab. \textbf{54} (2026), 892--972.
\href{https://doi.org/10.1214/25-AOP1782}{doi:10.1214/25-AOP1782}.

\bibitem{RGMap}
J.~Park, \emph{A renormalisation group map for short- and long-ranged
weakly coupled \(\lvert\varphi\rvert^4\) models in \(d\ge4\) at and
above the critical point},
\href{https://arxiv.org/abs/2511.03495}{arXiv:2511.03495} (2025).

\bibitem{TorusFSS}
J.~Park, \emph{Torus scaling limits and the plateau of the critical weakly
coupled $|\varphi|^4$ model in $d\ge4$},
\href{https://arxiv.org/abs/2511.06321}{arXiv:2511.06321} (2025).

\bibitem{IsingCorrelation}
J.~Park, \emph{Correlation function of the fractional spread-out Ising
models at their upper critical dimension and critical temperature},
GitHub preprint,
\href{https://github.com/jiwoonGit/phi4/blob/main/fractional_spread_out_ising_rg/spread_out_ising_rg.pdf}{DOI:10.5281/zenodo.24167641} (2026).

\bibitem{RudnickGaspariPrivman1985}
J.~Rudnick, G.~Gaspari, and V.~Privman,
\emph{Effect of boundary conditions on the critical behavior of a finite
high-dimensional Ising model}, Phys. Rev. B \textbf{32} (1985), 7594--7597.
\href{https://doi.org/10.1103/PhysRevB.32.7594}{doi:10.1103/PhysRevB.32.7594}.

\bibitem{SladeLR}
G.~Slade, \emph{Critical exponents for long-range \(O(n)\) models below
the upper critical dimension}, Comm. Math. Phys. \textbf{358} (2018), 343--436.
\href{https://doi.org/10.1007/s00220-017-3024-5}{doi:10.1007/s00220-017-3024-5}.

\bibitem{SladeWSAWPlateau}
G.~Slade, \emph{The near-critical two-point function and the torus plateau
for weakly self-avoiding walk in high dimensions},
Math. Phys. Anal. Geom. \textbf{26} (2023), article~6.
\href{https://doi.org/10.1007/s11040-023-09447-8}{doi:10.1007/s11040-023-09447-8}.

\bibitem{WittmannYoung2014}
M.~Wittmann and A.~P.~Young, \emph{Finite-size scaling above the upper
critical dimension}, Phys. Rev. E \textbf{90} (2014), 062137.
\href{https://doi.org/10.1103/PhysRevE.90.062137}{doi:10.1103/PhysRevE.90.062137}.

\end{thebibliography}

\end{document}
